\documentclass[11pt]{article}

\usepackage[margin=1in]{geometry}
\usepackage[T1]{fontenc}
\usepackage[utf8]{inputenc}
\usepackage{lmodern}
\usepackage{microtype}
\usepackage{amsmath,amssymb,amsthm}
\usepackage{mathtools}
\usepackage{hyperref}
\usepackage{booktabs}
\usepackage{enumitem}
\usepackage{listings}
\usepackage{xcolor}

\hypersetup{
  colorlinks=true,
  linkcolor=blue!50!black,
  citecolor=blue!50!black,
  urlcolor=blue!50!black
}

\lstdefinelanguage{Lean4}{
  morekeywords={
    import,namespace,end,open,section,noncomputable,abbrev,
    def,theorem,lemma,where,by,fun,let,in,have,show,exact,
    intro,apply,simpa,calc
  },
  sensitive=true,
  morecomment=[l]{--},
  morecomment=[s]{/-}{-/},
  morestring=[b]"
}

\lstset{
  language=Lean4,
  basicstyle=\ttfamily\small,
  keywordstyle=\bfseries\color{blue!50!black},
  commentstyle=\itshape\color{green!40!black},
  stringstyle=\color{red!50!black},
  showstringspaces=false,
  breaklines=true,
  frame=single,
  framerule=0.3pt,
  xleftmargin=0.5em,
  xrightmargin=0.5em,
  keepspaces=true,
  columns=fullflexible,
  literate=
    {¬}{{$\neg$}}1
    {∃}{{$\exists$}}1
    {→}{{$\to$}}1
    {≤}{{$\leq$}}1
}

\newtheorem{theorem}{Theorem}[section]
\newtheorem{lemma}[theorem]{Lemma}
\theoremstyle{remark}
\newtheorem{remark}[theorem]{Remark}

\title{P vs NP Proof Attempts:\\
Crux-First and Steelman Lean Notes on a Weakness-Based Route}
\author{}
\date{\today}

\begin{document}
\maketitle

\begin{abstract}
This note records the current Lean work around a claimed proof of
$P \neq NP$ via ``quantale weakness and geometric complexity.''  The project now
proceeds on two tracks in parallel.

The \emph{blocker} track asks where the argument must fail unless an additional theorem
is supplied.  The \emph{optimistic} track assumes the high-level strategy may work and
builds the missing background theory from the ground up so that any genuine repair can
plug into a machine-checked scaffold.

The current Lean work proves concrete obstructions.  First, a decoder that is merely
``local on a bounded-degree neighborhood of radius $\Theta(\log m)$'' is \emph{not} thereby
polylogarithmic in size: the visible neighborhood already has polynomial size in $m$.
Second, if one does not prove an additional compression theorem, then the class of all local
rules on that neighborhood is doubly exponential in the neighborhood size, so it cannot be
uniformly represented by a polylogarithmic code budget.

Before the switching story even starts, the sampler already carries another burden.  The
manuscript fixes the VV width at \(k=c_1\log m\), but the cited \(\Omega(1/m)\) uniqueness
probability comes from averaging over many scales of \(k\).  The new Lean obstruction packages the
fixed-width variance bound: if the retained-solution count has mean \(\mu>1\) and variance at most
\(\mu\), then exact isolation occurs on at most a \(\mu/(\mu-1)^2\) fraction of hash choices.
So fixed logarithmic width can only isolate often when the pre-hash solution set is already
essentially polynomial in size.

And even the manuscript's repeated \emph{pairwise-independent columns} language cannot shoulder the
hashing argument by itself.  A separate Lean countermodel shows that pairwise-independent columns do
not imply universal hashing: one can have pairwise-uniform columns together with a deterministic
three-column linear dependency, so a nonzero difference hashes to zero with probability \(1\).

Nor can the manuscript's \(\delta\)-biased right-hand side rescue the hash moments by itself.
Another Lean module shows that once a candidate's label map is genuinely uniform over the hash-label
space, its average hit mass against \emph{any} finite probability weight on labels is exactly the
uniform value \(1/|L|\); and if a pair of candidates has jointly uniform labels, then the weighted
average joint-hit mass is exactly \(1/|L|^2\).  So any useful effect of the \(\delta\)-biased
right-hand side must come from failure of the underlying label maps to be genuinely uniform, not
from the bias itself.

And even exact pre-conditioning uniformity is not enough by itself.  A further Lean countermodel
shows that conditioning on a label-dependent event can collapse a perfectly uniform label map to a
point mass.  So any appeal to pre-conditioning \(2\)-universality or uniform labels must be
supplemented by a separate theorem about what survives the manuscript's later conditioning on
uniqueness.
The newest Lean strengthening records the corresponding independence failure in reusable form:
inside any conditioning event, deterministic coupling \(E\leftrightarrow F\) is conditionally
independent exactly in the degenerate cases where \(E\) is empty or full on that conditioning
event.  The same exact boundary now also covers complement coding \(E\leftrightarrow\neg F\).
The four-point diagonal model is now an instance, not merely a one-off computation: two Boolean
coordinates can be exactly independent before conditioning, while conditioning on the diagonal
event makes them coupled and nondegenerate.  Thus the v11 uniqueness-conditioned route needs a real
local--global decorrelation theorem, not just pre-conditioning independence of the VV layer.

Third, even one natural repair to the calibration step has a sharp failure mode: if an
involution flips the target bit while preserving the retained local features, then every
classifier using only those invariant features has exactly \(1/2\) average accuracy.  So one
cannot repair the sign-flip calibration bug merely by dropping the involution-sensitive part of
the local input and hoping marginalization will preserve a domination claim.

Even before that charitable repair, the manuscript's exact local input formula already has a hard
consistency problem: it defines the post-switch input as \(u=(z,a_i,b)\) while the involution
\(T_i\) toggles \(b\) by \(a_i\).  So full-input invariance can hold only on the zero-column case
\(a_i=0\).

The new bridge module sharpens that fork: on the exact same \(T_i\)-action, the reduced projection
\((z,a_i)\) is invariant and therefore carries zero symmetry-respecting signed signal, while any
advantage obtained by also retaining \(b\) is bounded by the mass on nonzero VV columns, since
those are exactly the points where \(b\) resolves an involution pair.

Sharper still, the same involution hypotheses force every retained feature fiber to contain
exactly half target-\(1\) and half target-\(0\) points.  So even soft scores or conditional-mean
estimators on those invariant features have zero margin on every feature value.

And this is not just a uniform-counting artifact: any involution-invariant finite weighting still
assigns equal total mass to the target-\(1\) and target-\(0\) parts of every retained feature
fiber.  So non-uniform but symmetry-respecting sampling does not restore a usable margin.

In fact, any soft score computed only from those involution-invariant retained features has zero
weighted signed correlation with the target.  So even a richer invariant scoring statistic carries
no usable directional signal before genuine symmetry breaking is introduced.

More sharply, even when the weighting is not exactly involution-invariant, the entire invariant-score
signal is still carried only by the antisymmetric part of that weighting.  So this rescue route now
needs a real theorem that substantial useful mass lives in the weight asymmetry itself.  The
support-level version is formal too: if the weights are symmetric at every nonzero-score point, the
signal is zero; any nonzero signal must expose a point where the score is nonzero and the weight
changes across the involution.

Likewise, if one augments the invariant local view by a side channel, then any advantage gained by
classifying on the pair \((u,v)\) is still bounded by the total mass where the side channel actually
separates an orbit point from its involution partner.

Equivalently, any claimed domination gap over chance now forces a matching amount of orbit-resolving
mass.  So one cannot even state a serious advantage theorem without simultaneously proving a serious
resolution theorem.

More generally, keeping a controlled amount of non-invariant information only helps on the slice
where it really breaks the symmetry.  Any residual slice on which the retained features are still
involution-invariant remains exactly at chance under symmetry-respecting weights.

Quantitatively, any excess over chance is bounded by the weight of the
symmetry-broken region.  So a proof that keeps only a little non-invariant information must also
show that almost all relevant mass actually lies outside the unresolved symmetric slice.  The
support-level version is now formal too: if all positive-weight points remain in the unresolved
slice, strict advantage is impossible; if strict advantage is claimed, Lean produces a
positive-weight point outside that slice.

Fourth, the natural ``short global program'' rescue also fails on its own: if the switched input
really has only \(n = O(\log m)\) visible bits, then an arbitrary local rule on those bits can be
hard-wired by a truth table of size \(2^n = \mathrm{poly}(m)\).  So polynomial-size decoders do
not by themselves force any small subclass inside the full local-rule space.

Fifth, the most creative metagraph-style rescue faces its own exactness barrier: if a quotient
state must capture the full interface behavior of a \(k\)-bit local gadget, then there are
\(2^{k2^k}\) such boundary transducers.  At binary-log width this already forces at least linear
code size in \(m\), so a small exact metagraph/message quotient cannot appear by default.

Sixth, even the tree-message route does not follow from tree-likeness alone: a fixed perfect
binary tree architecture of depth \(n\) already realizes every Boolean rule on \(n\) visible bits,
just by changing its leaf labels.

The v13 crux ledger adds two smaller but sharper obligations.  A \(P=NP\) upper-bound decoder that
uses an arbitrary SAT-search output cannot recover a hidden message unless every satisfying witness
has the same message readout.  The current Lean interface packages a SAT-search output as an
arbitrary satisfying witness and proves that correctness for every valid SAT-search output is
exactly the fixed-message readout theorem; in the satisfiable case, this is equivalent to
single-message readout constancy, even after existentially quantifying over the claimed output
message.  The generic blocker is now theorem-level: any two satisfying witnesses with different
readouts produce two valid SAT-search outputs with different readouts and rule out every
fixed-message decoder correct for all valid outputs.  Bare satisfiability plus a readout map is not
enough; the two-point Boolean model is just the smallest instance of this general obstruction.
The newest Lean strengthening separates this from the weaker ``selected witness'' contract:
\[
  \texttt{CorrectForSomeSatSearchOutput}
\]
is only existence of one satisfying witness with the desired readout.  Distinct satisfying
readouts certify two incompatible selected-output messages, while still ruling out any message
correct for all valid SAT-search outputs:
\[
  \texttt{exists\_two\_messages\_correctForSomeSatSearchOutput\_of\_distinct\_satisfying\_readouts},
\]
\[
  \texttt{exists\_correctForSomeSatSearchOutput\_and\_not\_exists\_correctForAll\_of\_distinct\_satisfying\_readouts}.
\]
Thus a repair that silently chooses a favorable satisfying witness is not the same theorem as a
decoder robust to arbitrary SAT-search output.
\texttt{PNPv13ConcreteSATReadoutObstruction.lean} now instantiates the same point with ordinary
CNF syntax rather than an abstract witness type.  The concrete formula is the one-variable
tautological clause \(x\lor\neg x\).  Lean proves that both the false and true assignments satisfy
the CNF, that reading out the value of \(x\) gives distinct messages, and therefore
\[
  \texttt{oneVarTautologyCNF\_not\_exists\_correctForAllSatSearchOutputs}.
\]
It also proves selected-output correctness for both messages.  Thus the manuscript-facing missing
assumption is not ``the realization is satisfiable'' but a genuine theorem that all satisfying
assignments of the actual CNF realization have the same message readout, or that SAT search is
constrained by a stronger output contract.  The positive side is now formal at the same concrete CNF
interface:
\[
  \texttt{cnf\_exists\_correctForAllSatSearchOutputs\_iff\_singleMessageReadout\_of\_satisfying\_assignment}.
\]
For any satisfiable concrete CNF, an arbitrary-output SAT-search decoder message exists exactly when
the satisfying assignments have single-message readout constancy.
\texttt{PNPv13ManuscriptCNFReadout.lean} then formalizes the manuscript's own Appendix~I.23
bridge rather than only the abstract or one-variable surfaces.  Its interface records the supported
public instance \(Y\), the CNF \(F_Y\), the extracted witness \(\mathrm{Wit}_Y(\alpha)\), the
semantic witness relation \(R(Y,W)\), the witness message \(M(W)\), and the fixed projection
\(\pi_M(\alpha)\).  Lean proves
\[
  \texttt{ManuscriptCNFReadoutData.singleMessageReadout\_of\_singleMessagePromise},
\]
namely: if satisfying CNF assignments extract to valid witnesses, if the projection agrees with the
extracted witness message, and if valid witnesses for the supported \(Y\) satisfy the v13
single-message promise, then all satisfying CNF assignments have the same projected message.  With
one satisfying assignment this yields
\[
  \texttt{ManuscriptCNFReadoutData.exists\_correctForAllSatSearchOutputs\_of\_singleMessagePromise}.
\]
The same module now proves the converse under the natural completeness condition for the manuscript
extraction layer.  \texttt{ManuscriptCNFReadoutData.CNFExtractionComplete} says every supported
valid witness is represented by some satisfying CNF assignment with the same message projection, and
\texttt{ManuscriptCNFReadoutData.SupportedCNFSatisfiable} records that the supported CNF has at
least one satisfying assignment.  Under these two hypotheses Lean proves
\[
  \texttt{ManuscriptCNFReadoutData.supportedArbitraryOutputSATSearchCorrect\_iff\_singleMessagePromise}.
\]
Thus replacing the fixed-message theorem by an arbitrary-output SAT-search procedure does not weaken
the semantic obligation: for complete supported CNF realizations, SAT-search correctness is exactly
the single-message promise.  The one-variable manuscript countermodel is also marked complete and
satisfiable, via
\[
  \texttt{ambiguousOneVarManuscriptCNFReadoutData\_cnfExtractionComplete}
\]
and
\[
  \texttt{ambiguousOneVarManuscriptCNFReadoutData\_not\_supportedArbitraryOutputSATSearchCorrect}.
\]
A complementary boundary now proves that the completeness hypothesis is not cosmetic.  The
\texttt{unrepresentedOneVarManuscriptCNFReadoutData} model has a constant CNF projection, so
\[
  \texttt{unrepresentedOneVarManuscriptCNFReadoutData\_supportedArbitraryOutputSATSearchCorrect}
\]
holds, but it also has an unrepresented valid semantic witness with the opposite message:
\[
  \texttt{unrepresentedOneVarManuscriptCNFReadoutData\_satSearchCorrect\_and\_not\_singleMessagePromise},
  \qquad
  \texttt{unrepresentedOneVarManuscriptCNFReadoutData\_not\_cnfExtractionComplete}.
\]
Thus CNF-side arbitrary-output correctness can only be promoted back to the semantic
single-message promise after proving that every supported valid witness is represented by a
satisfying CNF assignment.
The same file also proves the boundary is real: a one-variable CNF countermodel satisfies the
compiler/extraction/projection shape while failing the single-message promise and while admitting no
message correct for all SAT-search outputs.  Thus the current minimal manuscript obligation is not
the local CNF compiler, but Proposition~4.10 / Lemma~I.23's premise: locked-layer rigidity must be
instantiated for the actual gauge-buffered locked construction so that every valid extracted witness
has the same message.
The same Lean module now pushes one step further into Proposition~4.10 itself.  The
\texttt{LockedLayerCNFReadoutData} refinement exposes a locked-layer view of each semantic witness,
the locked-layer acceptance predicate, and the message carried by that locked view.  Lean proves
\[
  \texttt{LockedLayerCNFReadoutData.singleMessagePromise\_of\_lockedLayerRigidity},
\]
which is the formal Proposition~4.10 argument: if every valid witness satisfies the locked layer and
its witness message is the locked-layer message, then Hypothesis~4.9's rigidity of accepted locked
views implies the witness-level single-message promise.  Composed with Appendix~I.23, this gives
\[
  \texttt{LockedLayerCNFReadoutData.exists\_correctForAllSatSearchOutputs\_of\_lockedLayerRigidity}.
\]
The same surface can be instantiated either with the raw Section~4 predicate
\(\mathrm{Lock}_Y\), matching Hypothesis~4.9, or with the Appendix~D locked-completion predicate
\(\mathrm{Lock}_Y \wedge \mathrm{Read}_Y\), matching Hypothesis~D.8.  Lean now records the exact
contrapositive certificate:
\[
  \texttt{LockedLayerCNFReadoutData.not\_lockedLayerRigidity\_of\_distinct\_accepted\_lockMessages}.
\]
Thus an attempted concrete formalization of the v13 locked core has a crisp pass/fail target: prove
that no two accepted lock completions over the same supported public syntax carry different messages,
or exhibit such a pair and the readout-constancy route fails at Proposition~4.10.  At the witness
level the analogous obstruction is
\[
  \texttt{ManuscriptCNFReadoutData.not\_singleMessagePromise\_of\_distinct\_valid\_witnessMessages}.
\]
The remaining unproved manuscript burden has therefore been localized to the actual locked-layer
construction: prove Hypothesis~4.9 for the concrete \(\mathrm{Lock}_Y\) predicates without smuggling
the message as public data.  A matching one-variable locked-layer countermodel shows that the
locked-layer-shaped interface without rigidity still permits two accepted locked views with different
messages and no arbitrary-output SAT-search decoder message.
\texttt{PNPv13AppendixDLockedCore.lean} now formalizes the Appendix~D version directly.  A locked
completion is the manuscript's triple \((z,\xi_{\mathrm{lock}},M)\) satisfying both
\(\mathrm{Lock}_Y(z,\xi_{\mathrm{lock}},M)\) and
\(\mathrm{Read}_Y(z,\xi_{\mathrm{lock}},M)\).  Appendix~D.8 is represented as
\texttt{AppendixDLockedCore.LockedMessageRigidity}; D.7 plus D.8 gives an existential, choice-free
locked-message specification via
\[
  \texttt{AppendixDLockedCore.exists\_lockedMessageSpec\_of\_lockSatisfiable\_of\_lockedMessageRigidity}.
\]
The Appendix~D.41/D.42 witness layer and Appendix~I semantic CNF layer then prove, respectively,
\[
  \texttt{AppendixDWitnessData.singleMessagePromise\_of\_lockedMessageRigidity},
\]
\[
  \texttt{AppendixICNFReadoutData.singleMessageReadout\_of\_lockedMessageRigidity},
  \qquad
  \texttt{AppendixICNFReadoutData.semantic\_i26\_items\_of\_lockedMessageRigidity}.
\]
Thus the actual Appendix~D/I interface guarantees fixed CNF readout only because D.8 is supplied.
It is not a consequence of the finite locked-core syntax itself.  Lean proves this with
\texttt{ambiguousAppendixDLockedCore}: one public lock syntax, two quotient states, trivial
\(\mathrm{Lock}\), and deterministic \(\mathrm{Read}(z,M)\equiv(M=z)\).  This core satisfies D.7
and deterministic readout, but
\[
  \texttt{ambiguousAppendixDLockedCore\_not\_lockedMessageRigidity}
\]
and its Appendix~I realization still has no message correct for arbitrary satisfying SAT-search
outputs:
\[
  \texttt{ambiguousAppendixICNFReadoutData\_not\_exists\_correctForAll}.
\]
\texttt{PNPv13AppendixDLocalLockedCore.lean} then strengthens this no-go to
the bounded-local predicate language of Appendix~D.5.  It defines
\texttt{BoolLocalConstraint} as a finite-scope Boolean predicate, packages
finite conjunctions with a shared arity bound as \texttt{BoundedBoolLocalPredicate}, and forgets the
local presentation back to the Appendix~D semantic core through
\texttt{AppendixDLocalBoolLockedCore.toAppendixDLockedCore}.  The one-bit model
\texttt{ambiguousAppendixDLocalBoolLockedCore} has empty local \(\mathrm{Lock}\)
clauses and a single arity-two local \(\mathrm{Read}\) clause enforcing
\(M=z\).  Lean proves that this bounded-local core satisfies D.7 and fixed-state
readout determinism, but still refutes D.8:
\[
  \texttt{ambiguousAppendixDLocalBoolLockedCore\_not\_lockedMessageRigidity}.
\]
Thus bounded arity and locality in D.5 do not generate cross-completion message
constancy.  Appendix~D.8 remains the load-bearing global rigidity theorem that
must be proved for the manuscript's concrete locked construction.
\texttt{PNPv13AppendixDPublicSyntaxDiscipline.lean} adds the public-syntax
discipline from Appendix~D.29--D.31 on top of this local core.  A public syntax
record has template and surface atoms, both normalizing through a neutral
alphabet; Lean proves the D.31 core statement that every listed public template
or surface atom has a neutral normal form.  The same one-bit local core can be
equipped with singleton neutral public syntax:
\[
  \texttt{ambiguousPublicSyntaxDisciplined\_not\_lockedMessageRigidity}.
\]
It still satisfies D.7 and deterministic readout.  Therefore the no-public-tag
discipline is necessary hygiene, but it is not a substitute for the semantic
locked-message rigidity theorem.
\texttt{PNPv13AppendixDProductConstruction.lean} then starts on the concrete
route suggested by Appendix~D.36--D.41.  It formalizes the product witness
shape \((x,g,z,\xi_{\mathrm{lock}},\xi_{\mathrm{buf}},\xi_{\mathrm{aux}},M)\)
and the D.38 product relation as locked acceptance, readout consistency, gauge,
buffer, and auxiliary clauses.  Lean proves the D.41 bridge
\[
  \texttt{AppendixDProductConstruction.singleMessagePromise\_of\_lockedMessageRigidity},
\]
extracting a locked completion from every valid product witness and applying
locked-core rigidity.  It also proves the exact contrapositive:
\[
  \texttt{AppendixDProductConstruction.not\_lockedMessageRigidity\_of\_distinct\_valid\_witnessMessages}.
\]
The neutral-public one-bit core lifts through a product construction with
trivial gauge, buffer, and auxiliary clauses:
\[
  \texttt{ambiguousAppendixDProductConstruction\_not\_lockedMessageRigidity}.
\]
This matches Appendix~D.56's suggested record split: the product construction
uses a \texttt{LockCore} component carrying a rigidity proof.  The remaining
non-interface burden is still a proof of that component theorem for the actual
locked predicates.
\texttt{PNPv13AppendixDProductCNFBridge.lean} now connects this product
relation to Appendix~I's CNF extraction/projection layer.  The positive theorem
\[
  \texttt{AppendixDProductCNFReadoutData.singleMessageReadout\_of\_lockedMessageRigidity}
\]
shows that the product-CNF route closes exactly when the underlying locked core
has D.8 rigidity.  The matched contrapositive is sharper for crux work:
\[
  \texttt{AppendixDProductCNFReadoutData.not\_lockedMessageRigidity\_of\_distinct\_satisfying\_projections}.
\]
Thus two satisfying assignments of the realized CNF with different fixed
product-message projections refute the locked-core theorem needed by the whole
product route.  The same file instantiates this with the one-variable
tautological CNF: both satisfying assignments extract to valid Appendix~D
product witnesses, their singleton-block messages differ, and therefore
\[
  \texttt{ambiguousAppendixDProductCNFReadoutData\_not\_lockedMessageRigidity},
  \qquad
  \texttt{ambiguousAppendixDProductCNFReadoutData\_not\_exists\_correctForAll}.
\]
This is the current steelman boundary for the Appendix~D product repair: the
CNF compiler and product witness clauses are not enough; the author must prove
that the concrete locked predicates exclude every such pair of satisfying
assignments with distinct product messages.
\texttt{PNPv13AppendixDProductExactCNFBridge.lean} strengthens the same test
against a more charitable repair: suppose the CNF realization is exact enough
that every valid product witness has a satisfying CNF encoding with the same
fixed projected message.  Lean packages this as
\texttt{AppendixDProductExactCNFReadoutData}.  Even then, exactness only
transfers ambiguity back to the CNF:
\[
  \texttt{AppendixDProductExactCNFReadoutData.not\_singleMessageReadout\_of\_distinct\_valid\_witnessMessages},
\]
\[
  \texttt{AppendixDProductExactCNFReadoutData.not\_exists\_correctForAllSatSearchOutputs\_of\_distinct\_valid\_witnessMessages}.
\]
The exact one-variable product countermodel proves that every valid product
witness encodes into the tautological one-variable CNF, yet the two explicit
valid witnesses still have different messages:
\[
  \texttt{ambiguousAppendixDProductExactCNFReadoutData\_not\_singleMessageReadout},
  \qquad
  \texttt{ambiguousAppendixDProductExactCNFReadoutData\_not\_lockedMessageRigidity}.
\]
So exact CNF realization is not a repair theorem either.  It only moves the
burden to the same semantic locked-core claim: prove the actual product
relation has no two valid witnesses over the same public syntax with different
messages.
\texttt{PNPv13AppendixDProductBijectionCNFBridge.lean} tests the strongest
natural version of this repair: require satisfying CNF assignments and valid
product witnesses to round-trip through extraction and encoding.  Lean proves
that this stronger bijective interface still gives exactly the same boundary:
\[
  \texttt{AppendixDProductBijectionCNFReadoutData.singleMessageReadout\_iff\_validProductWitnessSingleMessage}.
\]
Thus bijectivity does not manufacture readout constancy; it merely identifies
CNF readout constancy with single-message constancy of valid product witnesses.
The one-variable product model is made bijective by encoding a witness as its
Boolean message and extracting the corresponding canonical witness.  It still
proves
\[
  \texttt{bijectiveAmbiguousAppendixDProduct\_not\_singleMessageReadout},
  \qquad
  \texttt{bijectiveAmbiguousAppendixDProduct\_not\_lockedMessageRigidity}.
\]
So the remaining manuscript burden is now insensitive to this whole family of
CNF-realization repairs: unrestricted, exact, and bijective CNF interfaces all
reduce to the same missing theorem, namely global locked-core message rigidity
for the actual Appendix~D construction.
The deterministic readout predicate in Definition~4.4 is now separated as well.  The refinement
\texttt{DeterministicReadoutLockedLayerData} records a readout state, a predicate
\(\mathrm{Read}_Y(\mathrm{state},M)\), and determinism of that predicate for each fixed state.  Lean
proves
\[
  \texttt{DeterministicReadoutLockedLayerData.lockedLayerRigidity\_of\_acceptedReadStateConstancy},
\]
so deterministic readout plus constancy of the accepted readout state implies Hypothesis~4.9 and
therefore the arbitrary-output SAT-search theorem.  But Lean also proves a one-variable
countermodel:
\[
  \texttt{ambiguousOneVarDeterministicReadoutLockedLayerData\_not\_acceptedReadStateConstancy}.
\]
In that model every readout state has a unique message, but there are two accepted readout states
with different messages.  Thus deterministic readout is not the missing theorem; the manuscript must
prove a global cross-completion rigidity statement for the actual locked layer.  The reusable
negative criterion is now explicit as
\[
  \texttt{DeterministicReadoutLockedLayerData.not\_acceptedReadStateConstancy\_of\_distinct\_accepted\_readStates}.
\]
Separately,
unconditional one-coordinate half-success estimates do not imply a product small-success bound.
Two identical half-mass events have half-mass intersection, not quarter-mass intersection.  Thus the
product theorem needs a genuinely sequential conditional hypothesis over admissible histories.  The
same Lean file now records this repair target positively on the four-point Boolean square: a
first-step half-success bound plus a second-step conditional half-success bound implies the
two-step product bound.  The correlated model fails exactly the second, conditional premise.
The separate finite switching module generalizes this to arbitrary finite histories: if every
successive event halves the current admissible history class, then the whole tower has the expected
product bound.  What remains is not product arithmetic, but an instantiation theorem for the actual
switched histories.  The same generic layer also proves the corresponding failure principle: one
failed prefix half-step blocks sequential admissibility for the entire remaining history.  The
newest monotonicity lemmas also rule out a spurious edge case: if the current history class is
already empty, the next prefix half-step is automatic.  Hence real blockers must occur at
positive-mass conditioning prefixes.  The generic layer now also records the global contrapositive:
if the final intersection is too large for the claimed tower, then no sequential half-admissibility
proof can exist at all.  The newest localization theorem strengthens this from a global audit to a
specific missing premise: any such product-bound violation yields a decomposition
\texttt{pre ++ E :: suffix} whose accumulated prefix fails the conditional half-step.  The v13
layer lifts the same statement from event lists to switched-coordinate lists, so a failed product
count identifies an exact switched coordinate whose induced success event does not halve the
current success history.  The v13 layer also specializes the product check to abstract atom
ledgers, semantic concrete-cell certificates, and ordered fixed-field certificates.  Thus a proposed
switched list must first pass the tower-product count check before any finer cell structure can
rescue it.  The v13 instantiation layer now sharpens this further.
Lean first proves a useful negative interface fact: an abstract atom ledger with arbitrary
per-cell counts is equivalent to the generic prefix half-step, because a successful half-step can
always be represented by one summary cell.  It then proves the same collapse for an existential
``concrete'' cell-map certificate: the constant cell map is enough whenever the prefix half-step
already holds.  The latest Lean theorem lifts this prefix collapse through the whole ordered
switching list: both the abstract atom-ledger switching interface and the existential concrete-cell
switching interface are equivalent to generic sequential half-admissibility of the induced success
events.  So a real manuscript instantiation cannot be an existential over cell maps.  It must
instead name the actual finite history field before each switched coordinate and prove cellwise
half-success for that fixed field.  The newest atom-ledger structure theorem clarifies
the role of the concrete cell map: once a finite \texttt{cellOf} is supplied, the prefix and
next-success decomposition equations are automatic finite partition identities.  Thus decomposition
is not an independent escape hatch; the remaining mathematical burden is the cellwise half-success
inequality for the actual cells.  The fixed-field interface records this obligation:
the supplied field's cells must decompose both the current success prefix and the prefix plus the
next success event, and the next success event must occupy at most half of each cell.  This
has now been rewritten in Lean in the more diagnostic form
\texttt{v13FieldPrefixInstantiation\_iff\_success\_le\_failure}: after the finite
cell is fixed, a certificate exists exactly when every cell has at least as many
next-failure prefix points as next-success prefix points.  The supporting theorem
\texttt{v13ConcretePrefixCount\_eq\_success\_add\_failure} proves that this is not
a new assumption but the exact split of each cell's prefix mass into success and
failure parts.  This fixed-field condition is strictly stronger than the global
prefix half-step: on the Boolean square,
first-coordinate success is globally half-admissible from the empty history, but the fixed field
that reveals that same first coordinate has a `true' cell where success has conditional mass \(1\),
so no fixed-field certificate exists.  Lean now also lifts this fixed-field premise to an ordered
fielded switching list: each step carries its own supplied field, and the full fixed-field
instantiation is equivalent to the recursive list of these cellwise half-success inequalities.
Thus the fixed-field decomposition equalities are no longer open proof obligations; the remaining
fielded burden is exactly the half-success ledger over the named cells.  A full list of fixed-field
certificates implies the generic sequential half-admissibility and hence the finite tower product
bound.  Conversely, a single failed fixed-field prefix certificate blocks the whole remaining
fielded suffix, and the same Boolean first-coordinate one-step model instantiates that blocker.
The newest fixed-field theorem abstracts the reason: any positive-mass cell on which the next
success event is total violates the half-success inequality.  In particular, if the supplied field
reveals the next success event itself, then a fixed-field certificate can exist only when that next
success prefix has zero mass.  The latest necessary-condition theorem gives the same fact in
positive form: under any valid fixed-field certificate, every cell containing a successful point
must also contain an unsuccessful prefix point.  A singleton successful atom is therefore already
fatal.  The newest strengthening generalizes the success-revealing-field blocker:
\texttt{V13HistoryFieldDeterminesSuccessOn} says that every field cell is homogeneous for the next
success event on the current history support, and the corresponding positive-success blocker
proves that any such field is certifiable only when the next-success prefix has zero mass.  The
older global predicate \texttt{V13HistoryFieldDeterminesSuccess} is now just a special case.  Thus
a proposed v13 history field cannot already decide the next success bit after conditioning on any
positive-mass prefix; it must leave enough same-cell counterexamples to satisfy the
success-versus-failure balance.  Lean also checks the repair-resistant distinction: on the Boolean
square, the second-coordinate field does \emph{not} globally determine first-coordinate success,
but after conditioning on the diagonal history event it does determine that success bit, and the
positive diagonal success prefix has no fixed-field certificate or same-cell failure matching.  The
Boolean first-coordinate field now fails by the global special case, not only by an arithmetic
calculation on one named cell.  The blocker has also been lifted from the head of a fielded suffix
to an arbitrary
later position: a failed fixed-field prefix certificate after any preceding fielded list blocks the
whole original list, with the current history equal to the old history extended by the preceding
success events.  In the Boolean square, the one-cell field certifies the first half-success cut, but
the first-coordinate field is still impossible as the second repeated-coordinate field because it
determines the repeated success event on a positive-mass prefix.  The newest interface exposes the
same burden constructively: \texttt{V13FieldFailureMatching} requires, for every supplied field
cell, an injection from successful prefix points into unsuccessful prefix points in that same cell.
Lean proves this matching interface equivalent to the fixed-field certificate, and the recursive
fielded version is equivalent to the whole fixed-field switching instantiation.  Thus a proposed
repair can now be asked for an explicit same-cell matching at each switched position, not merely a
global counting slogan.  The matching interface also gives the contrapositive in the form needed
for audit: if the accumulated next-success prefix has positive mass and the recursive matching
proof exists, then the supplied field at that position cannot determine the next success bit on the
accumulated history support.  The newest pointwise extraction theorem makes the same obligation
more operational: for every successful prefix point, a recursive matching proof yields an actual
failed prefix point in the same supplied field cell.  Consequently, one successful point whose
accumulated-history cell contains only successes already blocks both the matching formulation and
the fixed-field certificate.  The diagonal Boolean-square example now also fails by this pointwise
criterion, not only by count arithmetic.
The same v13-specific module also checks the smallest
positive-mass failure mode for sequential repetition, and this is now a general finite theorem
rather than only a Boolean-square calculation: if an event has positive mass after the current
history and the next switched step repeats that same exact-success event, then the second prefix
half-step is false.  In the Boolean-square instance the first success prefix has count \(2\), the
repeated-success prefix still has count \(2\), and the second prefix half-step is false.  Thus no
v13 switching certificate can exist for any positive-mass repeated two-step switching claim.  The
Boolean-square example also fails the global product-bound audit: two identical half-mass success
events still have half-mass intersection, so the required quarter-mass tower bound is false.  The
product-bound contrapositive now reaches the
operational fixed-field matching layer as well: a tower-product violation rules out the recursive
same-cell matching proof directly.  In particular, even using the non-revealing one-cell field for
both repeated Boolean-square cuts cannot instantiate the fixed-field route; the fielded product
bound and the recursive matching obligation both fail.  The newest localization theorem also
turns either a failed recursive matching proof or a tower-product violation into a concrete failed
fielded position \(pre ++ item :: suffix\), with the accumulated history computed from \(pre\).
For the repeated one-cell Boolean-square model, Lean identifies the exact obstruction: after the
first success event, the one-cell field's only accumulated-history cell has next-failure count
\(0\) and next-success count \(2\), so it is the concrete bad cell for the second repeated success
test.

The route synthesis module now separates this blocker from the stronger claim that every possible
v13 repair is dead.  It names the exact route-class assumptions:
\texttt{V13RepeatedPositivePivotFrom} for an abstract switched list,
\texttt{V13RepeatedPositiveFieldedPivotFrom} for a fixed-field list, and
\texttt{V13FieldedBadCellFrom} for a localized same-cell balance failure.  Under those hypotheses
Lean proves that the corresponding abstract, fixed-field, and recursive matching formulations are
blocked.  Conversely, if a proposed manuscript instantiation avoids the repeated-positive-pivot
pattern, Lean now names the repair theorem it owes: prove the recursive
\texttt{V13FieldSwitchingFailureMatchingFrom} obligation for the actual ordered fielded list.  That
theorem entails both \texttt{V13NoRepeatedPositiveFieldedPivotFrom} and
\texttt{V13NoFieldedBadCellFrom}, and it also supplies the finite tower product bound.  Thus the
current conclusion is a decision tree: either the construction satisfies the repeated-pivot or
bad-cell hypotheses and is blocked, or it must provide the explicit cellwise same-cell matching
repair theorem.

Lean now also closes the most direct renaming repair for repeated pivots.  The module
\texttt{PNPv13ForcedStepObstruction.lean} proves the finite-count fact behind it: if an
accumulated history has positive mass and pointwise implies the next success event, then adding
that event leaves the history count unchanged, so the next event cannot be a conditional half-step.
This is not limited to literal event repetition.  The fielded v13 theorem packages the condition
as \texttt{V13ForcedPositiveFieldedStepFrom} and proves that it blocks both
\texttt{V13FieldSwitchingInstantiatedFrom} and the equivalent recursive same-cell matching
formulation.  The Boolean-square anchor makes the distinction concrete: first-coordinate success
and second-coordinate success are both valid one-cell half-cuts, but after both are imposed the
diagonal success event is forced.  The diagonal event is globally half-mass, and it is not the same
success event as either coordinate event, yet it cannot serve as the third conditional cut.

One natural refinement repair is now tested separately.  The module
\texttt{PNPv13FieldRefinement.lean} defines a genuine finite field refinement as a projection from
fine cells to coarse cells commuting with the pointwise cell maps.  It proves that a coarse bad cell
cannot be hidden by such a refinement: if the coarse cell has fewer next-failure points than
next-success points, then some fine cell over it has the same strict imbalance.  Consequently,
subdividing a bad supplied cell cannot produce a fixed-field prefix certificate for the same
history and switched step.  This does not rule out changing to a non-refining adaptive field, but
that move is now exactly the stronger repair theorem already named above: prove the recursive
same-cell matching obligation for the actual new fielded list.

There is now a second refinement-side obstruction, aimed at a different escape
hatch.  \texttt{PNPv13FieldRefinementDeterminesSuccessObstruction.lean} proves that
support-relative success determination is monotone under refinement: if a coarse
field already determines the next success event on the current history support,
then every finer field does too, because each fine cell lies inside one coarse
cell.  So refinement cannot rescue a field that was already revealing too much.
The Boolean-square regression makes this concrete by starting from the diagonal
history event: the second-coordinate field already determines repeated
first-coordinate success on that conditioned support, and refining it all the
way to the identity field still blocks both the fixed-field certificate route
and the equivalent recursive same-cell matching route.

The next refinement-side file removes the remaining ambiguity about what a
refinement repair would have to accomplish.  The module
\texttt{PNPv13FieldRefinementFailedPrefixObstruction.lean} proves the abstract
failure-monotonicity theorem: once the coarse field already fails the
fixed-field prefix obligation at a given accumulated history, every genuine
refinement of that field fails there as well.  So refinement does not merely
preserve the two previously isolated causes of failure; it preserves failed
prefixes as such.  The Boolean-square regression again uses the diagonal
history followed by repeated first-coordinate success: the second-coordinate
field already fails at that conditioned prefix, and refining it to the
identity field still blocks both fixed-field route formulations.

Assessment after the v13 switching packets: the formal work does not prove that no conceivable
v13-style switching strategy can work, because one residual positive target remains--name the
actual manuscript history fields and prove their recursive cellwise half-success bounds.  But it
does close the main bookkeeping escape hatches.  Abstract ledgers and existential concrete cell
maps add no strength beyond generic sequential half-admissibility; decomposition by finite cells is
automatic; and fixed fields reduce exactly to the cellwise half-success inequalities, with
positive-mass total-success cells, repeated success events, and forced later success events blocked
by small finite models.

The switched-history burden is now also packaged at the actual-family level.
\texttt{ActualSwitchedHistoryWrapper.lean} defines the manuscript-facing
wrapper as a one-family implication from one concrete fielded switching proof
\[
  \texttt{V13FieldSwitchingInstantiatedFrom hist items}
\]
to either the exact-visible compression target or the bundled clocked payload
for the actual switched predictor family.  Lean proves the positive route and
the blocker in matching sharp form.  On the positive side, any actual switched
family carrying \texttt{ZABDecisionListData} gets this wrapper immediately, so
the route closes once the manuscript identifies its actual predictors with the
shared exact \((zfeat(z),a,b)\) decision-list class.  On the negative side, if
the fielded switching premise is actually true and the actual switched family
is still surjective onto all Boolean rules on the exact visible surface below
the visible threshold, then the wrapper itself is impossible:
\[
  \texttt{not\_switchedHistoryExactVisibleCompressionWrapper\_of\_true\_fieldedSwitching\_of\_surjective\_predict\_of\_lt\_surfaceCard},
\]
\[
  \texttt{not\_switchedHistoryClockedKpolyFiniteLearningWrapper\_of\_true\_fieldedSwitching\_of\_surjective\_predict\_of\_lt\_predictorCard}.
\]
So the remaining switched-history route is no longer a generic bridge theorem:
it must either factor through the already formalized bounded exact ZAB class or
prove a genuine strict-subclass theorem for the actual switched predictor
family itself.
The new regression theorem
\texttt{actualSwitchedLocalInterface\_not\_forall\_switchedHistoryExactVisibleCompressionWrapper\_of\_lt\_surfaceCard}
pushes this one step further.  Bare
\texttt{ActualSwitchedLocalInterface} assumptions do not even support a
universal switched-history wrapper schema across all actual-local families: the
full-rule actual switched family is already a counterexample below the same
visible threshold.  So the surviving burden is sharply one-family and
manuscript-specific.
And once one concrete switched history is fixed and its fielded switching
premise is actually true, Lean now shows that the remaining wrapper is not a
new theorem-schema at all: it is exactly the finite-image theorem for the
actual switched family.  In
\texttt{ActualSwitchedHistoryRouteReduction.lean}, the manuscript-facing wrapper
is proved equivalent to the existing exact-visible compression target, to the
existing clocked finite-learning payload, and therefore to
\texttt{predictorFamily.FinitePredictorCover (2\^{}s)}.  The same file now also
pushes the reduction through the two other standard finite-image packagings:
finite selected-index representative covers and finite quotient-code
presentations.  So any surviving repair must prove one and therefore all of
those equivalent finite-image theorems directly; changing the wrapper packaging
no longer changes the mathematical burden.
Moreover, in the surjective regime the same reduction now exposes the exact
numerical obligation directly: any surviving exact-visible or clocked wrapper
forces
\[
  2^{|\mathrm{ExactVisiblePostSwitchSurface}(Z,k)|} \le 2^s.
\]
Applying \texttt{Nat.pow\_le\_pow\_iff\_right} immediately sharpens this to the
direct width bound
\[
  |\mathrm{ExactVisiblePostSwitchSurface}(Z,k)| \le s.
\]
Lean records this in
\texttt{switchedHistoryExactVisibleCompressionWrapper\_surfaceCard\_le\_of\_true\_fieldedSwitching\_of\_surjective\_predict}
and its clocked analogue, together with the direct refuters
\texttt{not\_switchedHistoryExactVisibleCompressionWrapper\_of\_true\_fieldedSwitching\_of\_surjective\_predict\_of\_s\_lt\_surfaceCard}
and
\texttt{not\_switchedHistoryClockedKpolyFiniteLearningWrapper\_of\_true\_fieldedSwitching\_of\_surjective\_predict\_of\_s\_lt\_surfaceCard}.
So below the visible threshold the switched-history route does not merely fail
qualitatively; it now contradicts the exact small-class width statement the
manuscript would have to establish.
The same route now has a concrete bitvector specialization.  If
\(Z = \mathrm{BitVec}\,n\), the wrapper is claimed at the manuscript-shaped
budget \(r + 2k + 1\), the extractor width satisfies \(r \le n\), and the
visible width is nontrivial \((2 \le n + 2k)\), Lean proves
\[
  r + 2k + 1 < |\mathrm{ExactVisiblePostSwitchSurface}(\mathrm{BitVec}\,n,k)|
    = 2^{n+2k},
\]
hence any true switched-history wrapper already contradicts surjectivity of the
actual switched family:
\texttt{not\_switchedHistoryExactVisibleCompressionWrapper\_of\_true\_fieldedSwitching\_of\_surjective\_predict\_of\_le\_of\_two\_le}
and its clocked analogue.  So even before identifying the actual family with
the shared exact ZAB class, the active switched-history route already inherits
the manuscript-budget obstruction on bit-valued latent data.
This packet is now generalized one step further.  On \(Z = \mathrm{BitVec}\,n\),
any true switched-history wrapper for a surjective actual switched family
directly implies
\[
  2^{n+2k} \le s,
\]
captured by
\texttt{switchedHistoryExactVisibleCompressionWrapper\_bitVecSurfaceCard\_le\_of\_true\_fieldedSwitching\_of\_surjective\_predict}
and its clocked analogue.  Consequently every explicit budget gap
\(s < 2^{n+2k}\) already refutes the route-specific wrapper itself:
\texttt{not\_switchedHistoryExactVisibleCompressionWrapper\_of\_true\_fieldedSwitching\_of\_surjective\_predict\_of\_budget\_lt}
and the matching clocked statement.  The manuscript-shaped budget
\(r + 2k + 1\) is now only one corollary of this more direct truth-table
threshold blocker.  Lean now also packages the full-width manuscript case
directly: once \(2 \le n + 2k\), the switched-history wrapper is ruled out at
budget \(n + 2k + 1\) itself by
\texttt{not\_switchedHistoryExactVisibleCompressionWrapper\_of\_true\_fieldedSwitching\_of\_surjective\_predict\_of\_two\_le\_visibleWidth}
and its clocked analogue, so the route no longer gets to hide behind an
auxiliary extractor-width parameter when it is really using the full visible
surface.  The Lean packet now goes one step further: once \(2 \le n + 2k\),
every switched-history wrapper with budget
\[
  s \le n + 2k + 1
\]
is impossible on \(Z = \mathrm{BitVec}\,n\), via
\texttt{not\_switchedHistoryExactVisibleCompressionWrapper\_of\_true\_fieldedSwitching\_of\_surjective\_predict\_of\_le\_fullWidthBudget}
and the matching clocked statement.  So the bitvector route no longer has any
room left anywhere in the entire linear budget interval up to the manuscript's
full-width threshold.

At the same time, the optimistic Lean work now isolates the precise positive interface one
would need for the switching lemma to go through: a \emph{small explicitly encoded
hypothesis family} whose realized classifier class is bounded by its code budget, together
with a perspective-relative notion of accessible input surface inspired by the existing
Hyperseed formalization.

These results do not by themselves disprove the whole paper.  They do isolate an essential gap:
the switching/locality step must be strengthened by a real theorem showing that the specific
switched predictors lie in an exceptionally small subclass, not merely in the class of all local
functions.
\end{abstract}

\tableofcontents

\section{Context}

The proof strategy under review is the October 2025 manuscript
\emph{``P != NP: A Non-Relativizing Proof via Quantale Weakness and Geometric Complexity.''}
At a high level, the paper proposes the following chain:
\begin{enumerate}[leftmargin=1.5em]
  \item Construct a masked random Unique-SAT ensemble.
  \item Show that every short polynomial-time decoder can be wrapped into a predictor that is
  local on a constant fraction of blocks.
  \item Use neutrality and sparsification to show local predictors succeed only with exponentially
  small probability across many independent blocks.
  \item Convert that small-success statement into a lower bound on time-bounded Kolmogorov
  complexity.
  \item Contradict the $P=NP$ upper bound obtained from Unique-SAT self-reduction.
\end{enumerate}

The current crux-first effort does not attempt the full chain at once.  Instead it attacks the
most fragile inference in the middle:

\begin{quote}
``Switched predictor is local at logarithmic radius'' $\Longrightarrow$
``switched predictor belongs to a polylogarithmic hypothesis class.''
\end{quote}

That implication is not automatic.  It needs either a structural theorem about the switched
predictor or a direct circuit/description-length bound.  The Lean work below shows why.

\section{Current Lean Files}

The current Lean modules under
\texttt{Mettapedia/Computability/PNP/}:
\begin{center}
\begin{tabular}{ll}
\toprule
\textbf{File} & \textbf{Role} \\
\midrule
\texttt{HypothesisClass.lean} & optimistic code-indexed hypothesis families and Hyperseed-visible inputs \\
\texttt{InfinitaryHMLObstruction.lean} & HML over infinitely many minimal contexts is not finite-code presentable \\
\texttt{InfinitaryHMLObstructionRegression.lean} & regression checks for finite-decoder and finite-encoder HML obstructions \\
\texttt{LocalityObstruction.lean} & Log-radius locality is larger than polylog size \\
\texttt{CompressionObstruction.lean} & Arbitrary local rules do not admit short uniform encodings \\
\texttt{FixedWidthIsolationObstruction.lean} & fixed-width hashing isolates only on a small fraction once the expected survivor count exceeds \(1\) \\
\texttt{PairwiseSurvivorMoments.lean} & pairwise first and factorial-second moment bounds imply the fixed-width variance budget \\
\texttt{PairwiseCandidateBridge.lean} & candidate-wise \(0/1\) hits plus pairwise off-diagonal bounds instantiate the survivor-moment model \\
\texttt{PairwiseColumnsObstruction.lean} & pairwise-independent columns still need not form a universal hash family \\
\texttt{RhsBiasIrrelevance.lean} & once label maps are uniform, any finite RHS bias averages to the same single-hit and pair-hit moments \\
\texttt{TwoUniversalRhsIrrelevance.lean} & exact finite \(2\)-universality of the hash family makes RHS bias irrelevant to the first two hit moments \\
\texttt{ConditioningObstruction.lean} & conditioning can collapse uniformity and destroy exact independence \\
\texttt{ApproximateDecorrelationObstruction.lean} & distinguishing conditioned recoding fibers leave an explicit approximate-independence defect \\
\texttt{ApproximateDecorrelationObstructionRegression.lean} & regression checks for coupled, anti-coupled, and arbitrary-fiber product-defect bounds \\
\texttt{SymmetrizationObstruction.lean} & unbiased majority does not beat Bayes without a margin \\
\texttt{PostSwitchInputObstruction.lean} & exact \(u=(z,a_i,b)\) is \(T_i\)-invariant iff \(a_i=0\) \\
\texttt{PostSwitchForkObstruction.lean} & exact \(T_i\)-fork: invariant \((z,a_i)\) has zero signal, \(b\)-side-channel gain needs a nonzero-column support witness \\
\texttt{PostSwitchForkObstructionRegression.lean} & regression checks for the exact \(b\)-side-channel resolution and nonzero-column witness demands \\
\texttt{VisiblePostSwitchSurface.lean} & isolates the exact visible surface, invariant projection, and fork-visible projection \\
\texttt{ExactSwitchedFamily.lean} & isolates the exact switched-family compression target and ERM reuse theorem \\
\texttt{ExactVisibleCompressionObstruction.lean} & if the exact family spans all exact visible rules, no sub-surface-cardinality code budget can cover it \\
\texttt{OrbitNeutralityObstruction.lean} & involution-invariant feature collapse forces exact half accuracy \\
\texttt{FiberNeutralityObstruction.lean} & invariant feature fibers have exact zero conditional margin \\
\texttt{WeightedFiberNeutralityObstruction.lean} & symmetry-respecting weights still give zero conditional margin \\
\texttt{InvariantScoreObstruction.lean} & invariant soft scores have zero weighted signed signal \\
\texttt{WeightAsymmetryObstruction.lean} & invariant-score signal comes only from antisymmetric weight mass, with a nonzero-score asymmetric-weight witness \\
\texttt{WeightAsymmetryObstructionRegression.lean} & regression checks for invariant-score asymmetric-weight witnesses \\
\texttt{SideChannelResolutionObstruction.lean} & side channels help only on orbit-resolving mass \\
\texttt{ResolutionDemandObstruction.lean} & any doubled or strict half-accuracy domination gap forces matching orbit-resolving mass and a positive-weight resolving point \\
\texttt{ResolutionDemandObstructionRegression.lean} & regression checks for positive-advantage, strict-half-accuracy, positive-resolution, and resolving-point consequences \\
\texttt{ResidualSideInformationResolutionCost.lean} & source-determined residual side information over an invariant base has zero resolving mass and cannot supply advantage \\
\texttt{RandomizedResidualObstruction.lean} & randomized residuals reduce to ordinary side channels on sample/coin products, with support-sensitive and deterministic coin-slice witnesses on both positive-advantage and strict-half surfaces \\
\texttt{RandomizedResidualObstructionRegression.lean} & regression checks for the randomized residual product-space reduction, supportwise blocker, and deterministic coin-slice witnesses \\
\texttt{PostSwitchRandomizedResidualObstruction.lean} & exact post-switch support blocks finite-coin randomized residual advantage, including the strict half-accuracy surface, and forces resolving deterministic coin slices \\
\texttt{PostSwitchRandomizedResidualObstructionRegression.lean} & regression checks for the post-switch randomized residual blocker and deterministic coin-slice witness theorems \\
\texttt{PostSwitchResidualWitness.lean} & concrete active-orbit witness on the exact post-switch surface: positive resolved mass, pure residual obstruction package, and supported prediction separation \\
\texttt{PostSwitchResidualWitnessRegression.lean} & regression checks for the exact post-switch active-orbit resolved-mass and pure residual obstruction package endpoints \\
\texttt{ResidualSymmetryObstruction.lean} & any unresolved symmetric slice still stays at chance \\
\texttt{AsymmetryBudgetObstruction.lean} & advantage is bounded by the symmetry-broken mass, and strict advantage exposes outside support \\
\texttt{AsymmetryBudgetObstructionRegression.lean} & regression checks for the asymmetry-budget support witness \\
\texttt{GlobalWeaknessObstruction.lean} & a family decoded only from one global weakness scalar cannot cover all Boolean classifiers on a nontrivial state space \\
\texttt{GlobalWeaknessObstructionRegression.lean} & regression checks for distinction and non-distinction weakness-family non-surjectivity \\
\texttt{ShortProgramObstruction.lean} & polynomial-size programs already realize all log-width local rules \\
\texttt{MetagraphObstruction.lean} & exact metagraph/interface quotients have state explosion \\
\texttt{MetagraphObstructionRegression.lean} & regression checks for finite quotient-state and quotient-code metagraph obstructions \\
\texttt{TreeMessageObstruction.lean} & fixed tree architectures already realize all local rules \\
\texttt{TreeMessageObstructionRegression.lean} & regression checks for tree evaluation bijectivity and finite message-state lower bounds \\
\texttt{ClockedKpolyBridge.lean} & finite-counting kernel for clocked \(s\)-bit decoder families \\
\texttt{ClockedKpolyBridgeRegression.lean} & regression checks for the clocked bridge kernel \\
\texttt{ClockedKpolyTargetInterface.lean} & exact-surface bit-budget targets imply the clocked recovery kernel \\
\texttt{ClockedKpolyTargetInterfaceRegression.lean} & regression checks for the clocked target interface \\
\texttt{ClockedKpolyCompressionObstruction.lean} & direct obstruction to low-bit clocked realization of a fully expressive exact switched family \\
\texttt{ClockedKpolyCompressionObstructionRegression.lean} & regression checks for the clocked realization obstruction \\
\texttt{ClockedKpolyBridgeEquivalence.lean} & clocked realization is equivalent to the exact bit-budget compression target \\
\texttt{ClockedKpolyBridgeEquivalenceRegression.lean} & regression checks that the clock parameter adds no expressivity \\
\texttt{ExactVisibleImageBudget.lean} & exact and clocked bit budgets are equivalent to finite predictor-image covers of size at most \(2^s\) \\
\texttt{ExactVisibleImageBudgetRegression.lean} & regression checks for the predictor-image budget equivalences \\
\texttt{ClockedKpolyGapAssessment.lean} & bundles the clocked finite-learning payload and proves it is exactly the finite predictor-image obligation \\
\texttt{ClockedKpolyActualGapClosure.lean} & closes the payload for the concrete exact ZAB ERM family and refutes deriving it from the bare exact-visible interface \\
\texttt{ProductSuccessKpolyBridgeObstruction.lean} & product-success and fielded-switching facts alone do not imply exact-visible \(Kpoly\) compression or the clocked finite-learning payload, reduce positive uses to finite-image obligations, and refute any already-surjective exact-visible family below threshold \\
\texttt{ProductSuccessKpolyBridgeObstructionRegression.lean} & regression checks for product/fielded bridge counterexamples, clocked-payload semantic-boundary equivalences, and surjective-family refuters \\
\texttt{ActualSwitchedLocalSupportObstruction.lean} & finite actual-local observed support can be small while the full exact-visible local-rule image remains uncompressed; uniform-section, reachable-support quotient, exact-decoder, downstream-property, point-query, query-family, and adaptive-query repairs force support or quotient invariance \\
\texttt{ActualSwitchedLocalSupportObstructionRegression.lean} & regression checks for finite-support smallness, uniform-section transfer, reachable-support quotient loss, exact-decoder failure, observable-property loss, point-query/query-family/adaptive-query reachability, and exact-visible finite-image lower bounds \\
\texttt{ActualSwitchedLocalPluginLookupObstruction.lean} & unrestricted plug-in lookup on the exact visible alphabet is the full Boolean rule family, hence not a small-image bridge \\
\texttt{ActualSwitchedLocalPluginLookupSparseThresholdERMObstruction.lean} & the same literal plug-in lookup endpoint also cannot be identified with the shared sparse-threshold ERM or recovery route below the unconditional point-block visible-budget threshold \\
\texttt{ActualSwitchedLocalPluginMajorityObstruction.lean} & unrestricted per-input plug-in majority counts still realize every exact-visible Boolean rule \\
\texttt{ActualSwitchedLocalPluginSampleMajorityObstruction.lean} & finite-sample plug-in majority still realizes every exact-visible Boolean rule via graph samples \\
\texttt{ActualSwitchedLocalBoundedSampleMajorityObstruction.lean} & bounded sample-level plug-in majority is full-rule expressive exactly once the sample bound reaches the exact visible alphabet size; any fixed tie-break can differ from its default only on sampled inputs; an input-dependent fallback is already a full lookup side channel; a structured fallback family can only add sparse sampled overrides around one selected fallback rule, with explicit bit-code fallback budget lower bounds \\
\texttt{ActualSwitchedLocalBitFallbackWrapperObstruction.lean} & standalone zero-sample and one-sample wrapper exactness: zero sampled overrides add no hidden expressivity beyond the fallback decoder, and one sampled override still pays an explicit additive bit budget \\
\texttt{ActualSwitchedLocalFallbackZeroSampleImage.lean} & generic structural zero-sample image equivalence: for an arbitrary fallback family, the radius-zero wrapper has exactly the fallback predictor image, so finite covers, representative covers, quotient codes, exact-visible compression targets, clocked exact-visible realizations, and clocked payload targets are equivalent on both sides \\
\texttt{ActualSwitchedLocalBitFallbackZeroSampleImage.lean} & bit-coded structural zero-sample image equivalence: the radius-zero wrapper has exactly the fallback decoder predictor image, so finite-image budgets and clocked exact-visible realizations are equivalent on both sides \\
\texttt{ActualSwitchedLocalBitFallbackZeroSampleProbeObstruction.lean} & background zero-sample probe transport: any injective finite probe family already realized by the fallback decoder transfers unchanged to radius-zero covers, quotient budgets, exact-visible compression, and clocked realizations, but this does not by itself close a named route \\
\texttt{FiniteIIDAgreement.lean} & finite PMF agreement mass, finite agreement-region lower bounds, and IID sample mass \\
\texttt{FiniteIIDFamilyBound.lean} & finite-family deceptive sample mass, union bounds, and bad-code region lower bounds \\
\texttt{FiniteIIDAgreementRegression.lean} & regression checks for support, region-mass, large-region positive-mass disagreement, and positive-mass agreement boundaries \\
\texttt{BadCodeAgreementObstruction.lean} & region-mass, large-region disagreement, and atomic-measure obstructions to strict bad-code agreement bounds in the exact ZAB ERM route \\
\texttt{BadCodeAgreementObstructionRegression.lean} & regression checks for region-mass, large-region disagreement, and pure-atom bad-code obstructions \\
\texttt{ABDecisionListObstruction.lean} & raw \((a,b)\) fixed-order decision lists are a strict subclass of all raw Boolean rules for \(k>0\) \\
\texttt{ABDecisionListObstructionRegression.lean} & regression checks for full raw-rule and exact-surface pullback refuters \\
\texttt{ABDecisionListRouteCrux.lean} & route-facing contradiction for the named \((a,b)\) decision-list repair: factor-through plus decision-list realization is inconsistent with positive-width full expressivity on the reduced \((a,b)\) surface \\
\texttt{SharedABFeatureRouteCrux.lean} & route-facing contradictions for the named shared raw \((a,b)\) affine-feature, sparse-threshold, and decision-list repairs: full reduced expressivity makes each positive route surface inconsistent below its own combiner budget \\
\texttt{SharedABFeatureRouteCruxRegression.lean} & regression checks for the invariant-lift shared-feature route contradictions \\
\texttt{SharedABFeatureObstruction.lean} & shared raw \((a,b)\) affine, sparse-threshold, and decision-list classes cannot realize full raw rules, or oversized finite probe images, below their combiner budgets \\
\texttt{SharedABFeatureObstructionRegression.lean} & regression checks for reduced and exact-pullback shared-feature refuters \\
\texttt{ABVisibleInvariantSurjectivityObstruction.lean} & raw \((a,b)\)-invariant exact families cannot realize, supportwise agree with, or perfectly distributionally recover exact-surface targets that distinguish hidden \(z\)-fibers \\
\texttt{ABVisibleInvariantSurjectivityObstructionRegression.lean} & regression checks for hidden-\(z\) projection, positive-mass disagreement, perfect-agreement blockers, and canonical AB-code refuters \\
\texttt{PNPv13CruxLedger.lean} & v13 decoder readout constancy and product-success finite obstructions \\
\texttt{PNPv13CruxLedgerRegression.lean} & regression checks for the v13 crux ledger \\
\texttt{PNPv13EvidenceNormalization.lean} & exact residual-atom boundary for v13 CD evidence normalization, with positive-score residual obstruction and three-atom countermodel \\
\texttt{PNPv13EvidenceNormalizationRegression.lean} & regression checks for evidence-normalization equivalences and the positive residual toy blocker \\
\texttt{PNPv13TraceCapture.lean} & exact trace-fiber factorization boundary for v13 observer capture, with collapsed-trace Boolean-observer blocker \\
\texttt{PNPv13TraceCaptureRegression.lean} & regression checks for trace decodability, same-trace conflicts, and universal Boolean observer capture \\
\texttt{PNPv13TraceCaptureCardinalityObstruction.lean} & universal Boolean trace capture iff injectivity on the target-relevant support, forcing any finite trace alphabet to be at least as large as the relevant state set \\
\texttt{PNPv13TraceCaptureCardinalityObstructionRegression.lean} & regression checks for the finite trace-cardinality lower bound and its toy finite-alphabet blockers \\
\texttt{PNPv13TraceCaptureImageBudgetObstruction.lean} & the same trace-capture boundary restated in the standard finite predictor-image language, forcing the full Boolean classifier-space inequality on the target-relevant support \\
\texttt{PNPv13TraceCaptureImageBudgetObstructionRegression.lean} & regression checks for the trace-decoder-family image-budget bridge and its toy finite-alphabet blockers \\
\texttt{PNPv13RandomCoinTraceCapture.lean} & random-coin trace capture reduced to deterministic coin-slice trace capture, with collapsed random-observer blocker \\
\texttt{PNPv13RandomCoinTraceCaptureRegression.lean} & regression checks for coinwise trace decodability, coinwise conflicts, and Boolean random-observer capture \\
\texttt{PNPv13AtomicEvidenceBudget.lean} & exact multiplicity/no-double-counting boundary for v13 coordinate-summed evidence budgets, with a one-atom/two-coordinate blocker \\
\texttt{PNPv13AtomicEvidenceBudgetRegression.lean} & regression checks for no-double-counting, over-budget exposure, and the double-counted toy surface \\
\texttt{PNPv13AtomicEvidenceBudgetExactnessCharacterization.lean} & the missing dual atomic-budget obstruction: under-budget coordinate sums expose omitted positive-cost atoms, and exact reuse of the atom budget needs both no double counting and no positive-cost omission \\
\texttt{PNPv13AtomicEvidenceBudgetExactnessCharacterizationRegression.lean} & regression checks for the atomic-budget exactness theorem, a one-atom exact canary, and a missed-positive-atom toy blocker \\
\texttt{PNPv13AtomicEvidenceBudgetCancellationObstruction.lean} & exact atomic-budget equality still does not certify correctness: a double-counted positive-cost atom can cancel an omitted positive-cost atom while the total budget matches exactly \\
\texttt{PNPv13AtomicEvidenceBudgetCancellationObstructionRegression.lean} & regression checks for the atomic-budget cancellation countermodel and the failure of equality-only repair theorems \\
\texttt{PNPv13AtomicEvidenceBudgetEqualitySharpness.lean} & exact budget equality becomes sharp once one structural side condition is known: no double counting forces no omission, and charging every positive-cost atom forces no double counting \\
\texttt{PNPv13AtomicEvidenceBudgetEqualitySharpnessRegression.lean} & regression checks for the equality-sharpness theorem packet, the missed-positive under-budget surface, the double-counted over-budget surface, and equality-surface equivalence canaries \\
\texttt{PNPv13AtomicEvidenceBudgetExactnessReduction.lean} & packages the actual exactness-side conjunction and proves that exact equality plus side-condition equivalence still does not imply it \\
\texttt{PNPv13AtomicEvidenceBudgetExactnessReductionRegression.lean} & regression checks for exactness-side reduction on the exact surface and the cancellation witness where equality and equivalence hold but exactness still fails \\
\texttt{PNPv13ACCEIEnvelope.lean} & finite ACCEI/PNLD skeleton: Markov-style bad-coordinate pruning, division-form good-coordinate lower bounds, a sharp two-coordinate Markov-loss witness, arbitrary-rate sequential product bounds, product-bound contrapositive, and failed-prefix localization \\
\texttt{PNPv13ACCEIEnvelopeRegression.lean} & regression checks for ACCEI pruning, good-coordinate lower bounds, Markov-loss sharpness, sequential rate bounds, product-bound obstruction, failed-prefix localization, and half-step specialization \\
\texttt{PNPv13ConcreteSATReadoutObstruction.lean} & concrete one-variable CNF instantiation of the SAT-search readout ambiguity \\
\texttt{PNPv13ConcreteSATReadoutObstructionRegression.lean} & regression checks for the concrete CNF readout blocker \\
\texttt{PNPv13ManuscriptCNFReadout.lean} & manuscript-shaped bridges from locked-layer rigidity and witness single-message promise to CNF readout constancy, complete-CNF SAT-search equivalence, and the incomplete-extraction boundary \\
\texttt{PNPv13ManuscriptCNFReadoutRegression.lean} & regression checks for the manuscript CNF readout bridges, complete-CNF SAT-search equivalence, incomplete-extraction boundary, and one-variable countermodels \\
\texttt{PNPv13AppendixDLockedCore.lean} & Appendix D locked-completion and Appendix I semantic CNF readout bridge, plus a D.5--D.7 deterministic-readout countermodel to D.8 \\
\texttt{PNPv13AppendixDLockedCoreRegression.lean} & regression checks for Appendix D.8 closure and the finite locked-core countermodel \\
\texttt{PNPv13AppendixDLocalLockedCore.lean} & bounded-local Boolean Appendix D.5 locked cores and an arity-two local countermodel to D.8 \\
\texttt{PNPv13AppendixDLocalLockedCoreRegression.lean} & regression checks for the bounded-local locked-core countermodel \\
\texttt{PNPv13AppendixDPublicSyntaxDiscipline.lean} & Appendix D.29--D.31 public syntax normalization and a neutral-public-syntax countermodel to D.8 \\
\texttt{PNPv13AppendixDPublicSyntaxDisciplineRegression.lean} & regression checks for public-syntax normalization and D.8 failure \\
\texttt{PNPv13AppendixDProductConstruction.lean} & Appendix D.36--D.41 product witness relation, D.41 bridge from locked-core rigidity, and lifted product countermodel \\
\texttt{PNPv13AppendixDProductConstructionRegression.lean} & regression checks for product single-message closure and product-level D.8 obstruction \\
\texttt{PNPv13AppendixDProductCNFBridge.lean} & Appendix D product witness relation connected to Appendix I CNF extraction/projection, plus product-CNF countermodel \\
\texttt{PNPv13AppendixDProductCNFBridgeRegression.lean} & regression checks for product-CNF readout closure and distinct-projection D.8 blockers \\
\texttt{PNPv13AppendixDProductExactCNFBridge.lean} & exact reverse-encoding CNF realization for product witnesses and exact-CNF ambiguity blocker \\
\texttt{PNPv13AppendixDProductExactCNFBridgeRegression.lean} & regression checks for exact-CNF readout, SAT-search, and locked-rigidity blockers \\
\texttt{PNPv13AppendixDProductBijectionCNFBridge.lean} & bijective exact CNF/product witness round-trip interface and proof that readout constancy is equivalent to valid-witness message constancy \\
\texttt{PNPv13AppendixDProductBijectionCNFBridgeRegression.lean} & regression checks for bijective CNF round trips and the remaining ambiguity blockers \\
\texttt{FiniteSwitchingAdmissibility.lean} & generic finite history admissibility, monotone history counts, and tower-product theorem \\
\texttt{FiniteSwitchingAdmissibilityRegression.lean} & regression checks for the generic switching interface \\
\texttt{PNPv13SwitchingInstantiation.lean} & v13 summary/concrete equivalences, fixed history-field and ordered fielded switching interfaces, and product bridge \\
\texttt{PNPv13SwitchingInstantiationRegression.lean} & regression checks for the v13 switching-instantiation interfaces \\
\texttt{PNPv13RouteSynthesis.lean} & precise v13 fixed-field route blockers and named repair obligations \\
\texttt{PNPv13RouteSynthesisRegression.lean} & regression checks for the v13 route synthesis layer \\
\texttt{PNPv13ForcedStepObstruction.lean} & non-literal forced-success-step obstruction for v13 switching, with Boolean-square forced-diagonal anchor \\
\texttt{PNPv13ForcedStepObstructionRegression.lean} & regression checks for forced-step blockers and the forced-diagonal Boolean-square model \\
\texttt{PNPv13FieldRefinement.lean} & field refinements preserve bad-cell fixed-field obstructions \\
\texttt{PNPv13FieldRefinementRegression.lean} & regression checks for the v13 field-refinement blocker \\
\texttt{PNPv13FieldRefinementRouteObstruction.lean} & field refinement lifts a bad coarse cell to a whole fielded-route blocker \\
\texttt{PNPv13FieldRefinementRouteObstructionRegression.lean} & regression checks for the route-level field-refinement blocker \\
\texttt{PNPv13FieldRefinementDeterminesSuccessObstruction.lean} & refinement preserves support-relative success determination, so finer fields still block the fixed-field route \\
\texttt{PNPv13FieldRefinementDeterminesSuccessObstructionRegression.lean} & regression checks for the refined success-determining field obstruction \\
\texttt{PNPv13FieldRefinementFailedPrefixObstruction.lean} & refinement preserves failure of the coarse fixed-field prefix itself, so failed coarse prefixes cannot be repaired by subdivision \\
\texttt{PNPv13FieldRefinementFailedPrefixObstructionRegression.lean} & regression checks for refinement monotonicity of failed coarse prefixes \\
\texttt{CruxSynthesis.lean} & meta-synthesis ledger with theorem-backed v13 route, forced-step, evidence-normalization, full trace-capture factorization/cardinality/image-budget entries, full atomic-budget exactness entries, strengthened residual-side-information and randomized-residual subrepair entries, including the cancellation-based equivalence blocker, ACCEI good-pruning/sharpness/envelope-product entries, a concrete SAT readout anchor, finite-coin route coverage, and live escape classes \\
\texttt{CruxSynthesisRegression.lean} & regression checks that the current PNP packet is not stop-grade and that v13/finite-coin coverage, including the strengthened trace-capture, atomic-budget, residual-side-information, and randomized-residual narrow-boundary surfaces and ACCEI good pruning and sharpness, remains scoped \\
\bottomrule
\end{tabular}
\end{center}

These compile in the current \texttt{mettapedia} workspace via \texttt{lake env lean}
(with selected files also built to \texttt{.olean} artifacts).

\section{Two Parallel Programs}

\subsection{The Blocker Program}

The blocker program is adversarial.  It asks:
\begin{enumerate}[leftmargin=1.5em]
  \item Is the visible radius already too large for any naive \texttt{poly(log m)} claim?
  \item Is the full class of local rules too large for a short description?
  \item Can the manuscript's fixed \(k=c_1\log m\) VV layer really achieve \(\Omega(1/m)\)
  uniqueness probability on large solution sets, or does fixed-width isolation already collapse
  once the expected survivor count is bigger than \(1\)?
  \item Is the manuscript's weaker \emph{pairwise-independent columns} language alone enough to
  justify universal hashing, or can higher-order linear dependencies still make a nonzero
  difference hash to zero deterministically?
  \item If the hash labels are already uniform or pairwise-uniform, can the \(\delta\)-biased
  right-hand side still change the first and pair hit moments at all?
  \item Even if the pre-conditioning label map is uniform, can later conditioning on uniqueness
  still collapse the local label distribution completely?
  \item Even if two coordinates are independent before conditioning, can a uniqueness-style
  conditioning event make them dependent unless a separate decorrelation theorem is supplied?
  \item Does majority symmetrization fail whenever the local conditional mean has zero margin?
  \item Is the exact post-switch input \(u=(z,a_i,b)\) really preserved by the paper's involution,
  or only in the zero-column case \(a_i=0\)?
  \item Once that exact fork is made explicit, does the charitable repair reduce either to zero
  invariant-feature signal on \((z,a_i)\) or to advantage bounded by the mass on nonzero VV columns?
  \item Does the obvious ``drop the bad coordinate'' calibration patch collapse any feature-only
  classifier back to exact neutrality?
  \item Does the softer ``estimate the conditional mean on invariant features'' repair still have
  exact zero margin on every retained feature fiber?
  \item Does non-uniform but involution-symmetric weighting avoid that zero-margin collapse?
  \item Can an arbitrary invariant soft score recover directional signal, or does it also have zero
  signed correlation with the target?
  \item If the weighting is not exactly symmetric, does invariant-score signal come only from the
  antisymmetric part of that weighting?
  \item If one keeps a non-invariant side channel, is total advantage still bounded by the mass
  where that side channel actually resolves involution pairs?
  \item Conversely, does any claimed domination advantage force a comparably large orbit-resolving
  mass for that side channel?
  \item If one keeps some controlled non-invariant information, do residual unresolved slices still
  remain exact-chance unless their mass is proved negligible?
  \item Can a small symmetry-breaking slice still yield a large global advantage, or is the excess
  over chance quantitatively bounded by the symmetry-broken mass?
  \item Does ``short global decoder'' alone still allow the full local-rule class at log-width?
  \item Does an exact finite-state metagraph/message quotient on a log-width interface already
  require too many states?
  \item Does tree-likeness alone force a small exact message/interpreter class?
  \item Does arbitrary SAT-search output recover a unique hidden message without a theorem that all
  satisfying witnesses have the same readout?
  \item Do unconditional half-success marginals imply a product small-success bound without
  sequential conditional admissibility?
\end{enumerate}

The current Lean answer to all listed questions is ``yes.''  Therefore the published proof cannot be
accepted \emph{as written} unless it supplies additional structural theorems.
The short-global-decoder item is now sharpened in a weakness-specific form:
\texttt{GlobalWeaknessObstruction.lean} proves that if every member of an indexed
predictor family factors through the single global
\texttt{distinctionWeakness} scalar, then the family is not surjective onto
all Boolean classifiers on any nontrivial state space:
\[
  \texttt{not\_surjective\_familyFactorsThroughDistinctionWeakness\_of\_nontrivial}.
\]
The same statement is proved for \texttt{nonDistinctionWeakness}.  Thus the
formalized GSLT weakness bridge cannot be used as a full encoded-family
readout unless the route supplies a genuinely per-instance or per-code
observable rather than only a global evidence summary.

\subsection{The Optimistic Program}

The optimistic program asks what would be enough to save the route.
The current working answer is:
\begin{quote}
locality alone is not enough, but \emph{locality plus explicit compressibility} may be.
\end{quote}

Concretely, the switching theorem would need to produce not merely a local rule on a
radius-$O(\log m)$ neighborhood, but a family of predictors realized by a code space of
size at most \(\mathrm{polylog}(m)\) bits.  In that case the realized hypothesis class is
small, and standard finite-class learning bounds become plausible again.

The new module \texttt{HypothesisClass.lean} formalizes exactly this interface:
\begin{enumerate}[leftmargin=1.5em]
  \item an \emph{encoded family} is a decoder from a finite code space into classifiers;
  \item the realized class has cardinality at most the size of the code space;
  \item therefore an $s$-bit family realizes at most $2^s$ classifiers;
  \item no such family can surject onto the full local-rule class once $s < 2^n$.
\end{enumerate}

This is the correct positive abstraction: if this switching argument is salvageable, it must
land in one of these small encoded subclasses rather than in the full local-rule space.

\subsection{Hyperseed Perspective Surface}

The Hyperseed formalization in \texttt{Mettapedia/Hyperseed/} already provides a useful
observer-relative interface:
\begin{itemize}[leftmargin=1.5em]
  \item a \emph{perspective} determines what is observable,
  \item a budget determines what is affordable,
  \item a guard determines what is admitted.
\end{itemize}

The optimistic PNP background layer reuses this by treating a decoder's \emph{available
input surface} as a subtype of the Hyperseed \texttt{availableRegion}.  This does not solve
switching.  It does provide a clean ground-up semantics for what a bounded observer can even
use as input, without baking in a bespoke notion of locality for every later argument.

\subsection{Meredith/HML Observable Surface}

The Meredith/HML side of the route has a separate finite-presentability obligation.  The module
\texttt{InfinitaryHMLObstruction.lean} shows that if the type of certified minimal contexts is
infinite, then even the depth-\(1\) HML fragment is infinite.  The rho-calculus instance already
has an infinite ladder of such minimal contexts, so its depth-\(1\) fragment and full HML formula
type are both infinite.

The original formal obstruction was one-sided: no finite code type can decode surjectively onto
those formula surfaces.  The newest strengthening records the dual finite-encoding boundary:
\[
  \texttt{not\_injective\_depthFragment\_one\_to\_finite\_code},
  \qquad
  \texttt{not\_injective\_hmlFormula\_to\_finite\_code}.
\]
The concrete rho versions,
\[
  \texttt{not\_injective\_rhoDepthFragment\_one\_to\_finite\_code},
  \qquad
  \texttt{not\_injective\_rhoHMLFormula\_to\_finite\_code},
\]
block the repair move of replacing an infinite observable surface by a finite injective code
without first proving a genuine finite quotient theorem for the actual contexts used by the route.

\section{Optimistic Background Theorem}

The basic positive theorem is elementary but indispensable:

\begin{theorem}[Code budget bounds realized hypothesis class]
Let \(H\) be a family of Boolean classifiers realized by decoding an \(s\)-bit code:
\[
  \mathrm{decode} : \{0,1\}^s \to \mathcal{F}.
\]
Then the realized class satisfies
\[
  |\mathrm{range}(\mathrm{decode})| \le 2^s.
\]
\end{theorem}

\noindent
This is formalized in \texttt{HypothesisClass.lean} as
\texttt{BitEncodedClassifierFamily.card\_realized\_le}.

The contrapositive is equally important:

\begin{theorem}[Short codes cannot cover the full local-rule class]
If \(s < 2^n\), then no \(s\)-bit decoder can surject onto the full set of Boolean rules on
\(n\) visible bits.
\end{theorem}

\noindent
This is formalized as
\texttt{not\_surjective\_decode\_to\_fullLocalRule\_of\_lt}.

These statements are not the final theorem.  They are the right \emph{background theory} for
the optimistic route: they tell us exactly what shape of compression theorem the switching
step must deliver.

\section{Obstruction I: Log-Radius Locality Is Already Polynomial}

Suppose a factor graph has bounded degree $d \ge 2$, and a predictor is allowed to inspect a
radius-$r$ neighborhood around one target bit.  A rough but unavoidable upper/lower size model is
\[
  \text{visible nodes} \asymp d^r.
\]
If $r = \Theta(\log m)$, then this is
\[
  d^{c \log m} = m^{c \log d},
\]
which is polynomial in $m$, not polylogarithmic in $m$.

The Lean module \texttt{LocalityObstruction.lean} records this in several exact forms.

\begin{theorem}[Binary-log lower bound]
For every $d,m \in \mathbb{N}$ with $d \ge 2$,
\[
  m \le d^{\log_2 m + 1}.
\]
\end{theorem}

\noindent
This is formalized as the theorem
\texttt{inputSize\_le\_neighborhoodSize\_at\_binaryLogRadius}.

\begin{theorem}[Power-of-two specialization]
For every $d,n \in \mathbb{N}$ with $d \ge 2$,
\[
  2^n \le d^{\log_2(2^n)} = d^n.
\]
\end{theorem}

\noindent
This appears as
\texttt{powerOfTwo\_le\_neighborhoodSize\_at\_exactBinaryLogRadius}.

\begin{theorem}[Asymptotic dominance]
Fix $k \in \mathbb{N}$ and real constants $d>1$, $c>0$.  Then
\[
  (\log x)^k = o\!\left(d^{c \log x}\right)
\]
as $x \to \infty$.
\end{theorem}

\noindent
This is formalized as
\texttt{polylog\_isLittleO\_logRadiusNeighborhood}, and the corresponding negative
Big-O statement appears as
\texttt{logRadiusNeighborhood\_not\_isBigO\_polylog}.

\subsection{Interpretation}

This does \emph{not} say that every local rule requires polynomial-size circuits.
It says something weaker but still decisive: locality at logarithmic radius, by itself, does not
justify a \texttt{poly(log m)} complexity claim.  The radius already exposes a polynomial-sized
input window.

So if the paper wants a tiny post-switch hypothesis class, it must prove a \emph{compression
theorem} about the specific local rules that arise after switching.  That is a separate and
nontrivial burden.

\section{Obstruction Ib: Fixed-Width VV Isolation Does Not Follow from the Cited Lemma}

The manuscript's main ensemble fixes
\[
  k=c_1\log m
\]
in the VV layer.  But the cited isolation lemma proving uniqueness probability \(\Omega(1/m)\)
chooses \(k\) \emph{uniformly across all scales} \(0,\dots,m-1\).  Those are not the same claim.

The new Lean module \texttt{FixedWidthIsolationObstruction.lean} packages the exact finite bound
one needs at fixed width.

\subsection{Lean theorem}

\begin{theorem}[Variance-budget bound for exact isolation]
Let \(Y\) be the number of candidates retained by a random hash choice, on a finite hash space.
If
\[
  \mathbb E[Y]=\mu>1
  \qquad\text{and}\qquad
  \mathrm{Var}(Y)\le \mu,
\]
then
\[
  \Pr[Y=1]\ \le\ \frac{\mu}{(\mu-1)^2}.
\]
\end{theorem}

\noindent
This is formalized as
\texttt{exactOneProb\_le\_of\_variance\_budget}.

The new companion module \texttt{PairwiseSurvivorMoments.lean} formalizes the abstract bridge from
pairwise-survivor moments to this variance budget.  In that file, a finite model satisfying
\[
  \mathbb E[Y]=\mu
  \qquad\text{and}\qquad
  \mathbb E[Y(Y-1)]\le \mu^2
\]
is shown to obey \(\mathrm{Var}(Y)\le \mu\), and therefore the same exact-isolation upper bound.
This separates the remaining manuscript obligation cleanly: one still has to prove that the actual
fixed-width VV sampler instantiates those pairwise moments.

The newer bridge module \texttt{PairwiseCandidateBridge.lean} takes one more step toward the actual
VV picture.  It proves that if the retained-candidate indicators are \(0/1\)-valued, all have the
same marginal hit rate \(p\), and distinct candidates satisfy the pairwise upper bound
\[
  \mathbb E[X_a X_b] \le p^2,
  \qquad a\neq b,
\]
then the total survivor count satisfies the exact abstract interface of
\texttt{PairwiseSurvivorMomentModel}.  Concretely:
\begin{itemize}[leftmargin=1.5em]
  \item \texttt{candidateHitCount\_sum} proves the first moment formula;
  \item \texttt{candidateHitCount\_factorialSecondMoment\_le} proves the factorial second-moment
  bound;
  \item \texttt{pairwiseCandidateMomentModel} packages these into the survivor-moment model used by
  the fixed-width isolation theorem.
\end{itemize}

So the remaining gap has narrowed again.  To recover the manuscript's fixed-width uniqueness claim,
one now needs an end-to-end theorem that the actual VV hash family satisfies these candidate-wise
marginal and pairwise bounds at the fixed logarithmic width used later in the paper.

\subsection{Why this hits the manuscript}

In the source text, the supporting VV lemma states the classical bound
\[
  \Pr[|S\cap h^{-1}(0^k)|=1]\ge c/m
\]
only after choosing \(k\) uniformly from \(0,\dots,m-1\).  But the actual ensemble later fixes
\(k=c_1\log m\) and then still says isolation succeeds with \(\Omega(1/m)\) probability.

That gap is not cosmetic.  Under the standard pairwise-independent VV heuristic, if a formula has
\(N\) satisfying assignments and the fixed hash width is \(k\), then the retained-count random
variable has mean
\[
  \mu = N/2^k
\]
and variance at most \(\mu\).  The Lean theorem above therefore yields
\[
  \Pr[\text{exact isolation}]
  \ \lesssim\
  \frac{1}{\mu}
  \ =\
  \frac{2^k}{N},
\]
up to a universal constant once \(\mu\) is bounded away from \(1\).

So with fixed \(k=O(\log m)\), the sampler can isolate with noticeable probability only when the
pre-hash solution set is already \(N=O(2^k)=\mathrm{poly}(m)\).  The manuscript, however, only
cites a much weaker upper bound of the form \(N\le 2^{\alpha m}\), which still permits an
exponential solution set.  On such instances, fixed logarithmic width gives exponentially small
exact-isolation probability rather than \(\Omega(1/m)\).

\subsection{Consequence}

This means the paper now needs a real theorem of one of the following forms:
\begin{enumerate}[leftmargin=1.5em]
  \item on the relevant random-3CNF mass, the satisfiable formulas entering the VV stage already
  have only \(\mathrm{poly}(m)\) satisfying assignments; or
  \item the ensemble does not actually use fixed \(k=c_1\log m\), and the later local-input and
  finite-alphabet claims must be rewritten accordingly.
\end{enumerate}

Without one of those repairs, the promised block distribution \(D_m\) is not supported by the
cited VV argument.

\section{Obstruction Ic: Pairwise-Independent Columns Do Not Imply Universal Hashing}

The ensemble definition samples \(A\) from ``any \(2\)-universal distribution with
pairwise-independent columns,'' and the later robustness discussion repeatedly foregrounds the
\emph{pairwise-independent columns} part.  The new Lean module
\texttt{PairwiseColumnsObstruction.lean} shows that this weaker condition is not enough on its own.

\subsection{Lean theorem}

\begin{theorem}[Pairwise columns need not be universal]
There exists a one-bit linear hash family with three columns such that:
\begin{enumerate}[leftmargin=1.5em]
  \item each pair of columns is jointly uniform over \(\mathbb{F}_2^2\); but
  \item some nonzero input difference hashes to \(0\) on \emph{every} seed.
\end{enumerate}
Hence pairwise-independent columns do not by themselves imply the universal-hashing bound required
for VV isolation.
\end{theorem}

\noindent
This is formalized by the bijectivity theorems
\texttt{col12\_bijective}, \texttt{col13\_bijective}, and \texttt{col23\_bijective}, together with
\texttt{threeColumnHash\_badDifference} and
\texttt{threeColumnHash\_not\_universal}.

\subsection{Why this matters}

This does \emph{not} refute a hash family that is genuinely \(2\)-universal.  It does refute a
common informal shortcut: one cannot replace real universal hashing by mere pairwise independence
of the columns.  Higher-order linear dependencies can survive that weaker condition and completely
destroy the needed collision bound.

So the manuscript now owes a sharper statement here too:
\begin{enumerate}[leftmargin=1.5em]
  \item either prove that the actual fixed-width VV family is genuinely \(2\)-universal in the
  sense needed for the candidate-wise moment bridge; or
  \item stop appealing to pairwise-independent columns as if that weaker property alone carried the
  isolation argument.
\end{enumerate}

\section{Obstruction Id: RHS Bias Is Irrelevant Once The Label Maps Are Uniform}

The manuscript samples the VV right-hand side \(b\) from a \(\delta\)-biased source and suggests
that this robustness can be folded into the collision/isolation analysis.  The new Lean module
\texttt{RhsBiasIrrelevance.lean} isolates the exact point where that can and cannot matter.

\subsection{Lean theorems}

\begin{theorem}[Uniform labels erase single-hit bias]
Let \(L\) be a finite label space with any probability weight \(w\) on \(L\), and suppose a finite
seed space \(S\) splits as \(S \cong \{1,\dots,d\} \times L\) with \(d>0\), so that a candidate's
label map is exactly the projection onto \(L\).  Then
\[
  \frac{1}{|S|}\sum_{s \in S} w(f(s)) = \frac{1}{|L|}.
\]
\end{theorem}

\noindent
This is formalized by \texttt{average\_mass\_of\_uniform\_labels}.

\begin{theorem}[Jointly uniform labels erase pair-hit bias]
Under the same hypotheses, if a pair of candidates has jointly uniform labels
\(S \cong \{1,\dots,d\} \times L \times L\) with \(d>0\), then
\[
  \frac{1}{|S|}
  \sum_{s \in S} \mathbf{1}_{f(s)=g(s)}\, w(f(s))
  = \frac{1}{|L|^2}.
\]
\end{theorem}

\noindent
This is formalized by \texttt{average\_joint\_mass\_of\_uniform\_label\_pairs}.

\subsection{Why this matters}

This does not prove that the manuscript's actual fixed-width VV family has perfectly uniform label
maps.  It proves the sharper conditional statement: \emph{if} the candidate-wise and pairwise label
maps are genuinely uniform, then the choice of a \(\delta\)-biased right-hand-side distribution is
irrelevant to the first and pair-hit moments.  The averaged values are exactly the same as under a
uniform right-hand side.

So the manuscript cannot get extra collision-control power from \(\delta\)-bias alone.  Any real
effect of the right-hand side must come from a failure of the underlying label maps to be uniform in
the first place.  This narrows the burden again:
\begin{enumerate}[leftmargin=1.5em]
  \item either prove that the actual fixed-width family's label maps are \emph{not} uniform in a
  way that helps the needed moment bounds; or
  \item stop treating \(\delta\)-bias itself as if it repaired the missing universal-hashing step.
\end{enumerate}

The new bridge module \texttt{TwoUniversalRhsIrrelevance.lean} sharpens this one step further in
the manuscript's own language: if the actual \(A\)-family is genuinely finite \(2\)-universal, then
the weighted single-hit and pair-hit averages are exactly the uniform values for \emph{every}
finite right-hand-side distribution \(w\).  So on the first two VV moments, the entire burden sits
on the \(A\)-family, not on the \(\delta\)-biased choice of \(b\).

But even this sharper bridge is only a \emph{pre-conditioning} statement.  The new module
\texttt{ConditioningObstruction.lean} shows that this is the end of what one can get for free:
conditioning a perfectly uniform label map on a single label fiber makes the conditional label
distribution a point mass.  So no argument may pass from pre-conditioning uniformity to the
distribution conditioned on uniqueness without an additional theorem controlling that conditioning
step.

The same file now formalizes the sharper v11-style warning using only four points.  Under uniform
counting on \(\{0,1\}^2\), the first-coordinate and second-coordinate truth events satisfy the
cross-multiplied independence equation:
\[
  \texttt{conditioningBoolPair\_unconditioned\_first\_second\_independent}.
\]
But after conditioning on the diagonal event \(x_1=x_2\), the conditional independence equation
fails:
\[
  \texttt{not\_conditioningBoolPair\_conditioned\_on\_diagonal\_first\_second\_independent}.
\]
The reason is now isolated as a general finite-count theorem:
\[
  \texttt{countIndependentWithin\_coupled\_iff\_degenerate\_on\_condition}.
\]
If \(E\) and \(F\) are pointwise equivalent on the conditioning event \(C\), then the
cross-multiplied conditional independence equation holds if and only if either \(C\cap E\) is empty
or \(C\cap \neg E\) is empty.  The older nonconstant blocker is now just the positive/negative
corollary:
\[
  \texttt{not\_countIndependentWithin\_of\_coupled\_nonconstant\_on\_condition}.
\]
If \(E\) and \(F\) are pointwise equivalent on the conditioning event \(C\), and both \(C\cap E\)
and \(C\cap \neg E\) have positive finite count, then the cross-multiplied conditional
independence equation cannot hold.  The diagonal Boolean square supplies the smallest witness to
the hypotheses, but the theorem applies to any uniqueness-style conditioning event that makes two
formerly separate predicates determine one another while both truth values remain present.
The complement-coded variant is now formalized by
\[
  \texttt{countIndependentWithin\_anticoupled\_iff\_degenerate\_on\_condition}
\]
and its positive/negative corollary
\[
  \texttt{not\_countIndependentWithin\_of\_anticoupled\_nonconstant\_on\_condition}.
\]
Thus replacing a copied conditioned bit by its negation does not evade the obstruction: if
\(E\leftrightarrow\neg F\) throughout \(C\), conditional independence again forces \(E\) to be
empty or full on \(C\).  Any repair must show genuine residual randomness after conditioning, not
just change the encoding of the determined bit.
The newest Lean refinement closes the corresponding arbitrary Boolean recoding escape.  For any
Boolean map \(r:\{0,1\}\to\{0,1\}\), the theorem
\[
  \texttt{countIndependentWithin\_boolRecoding\_iff\_degenerate\_on\_condition}
\]
states that if \(F\) is exactly the conditioned recoding \(r(E)\) on \(C\), and \(r\) is
nonconstant, then conditional independence again holds only in the same degenerate empty/full
cases.  Its contrapositive,
\[
  \texttt{boolRecoding\_constant\_of\_countIndependentWithin\_of\_nonconstant\_on\_condition},
\]
says that if \(E\) has both truth values inside \(C\) and \(F=r(E)\) remains conditionally
independent of \(E\), then \(r\) must be constant.  So the proof route cannot repair the
conditioning gap by applying a different deterministic Boolean post-processing to the copied
conditioned bit; it needs a real theorem that the post-conditioned variable still contains
residual randomness not determined by \(E\).
The classification is now exact rather than only one-sided:
\[
  \texttt{countIndependentWithin\_boolRecoding\_iff\_constant\_or\_degenerate\_on\_condition}.
\]
For deterministic Boolean recodings of a conditioned source bit, conditional independence holds
precisely when the recoding is constant or the source bit is empty/full on the conditioning event.
Equivalently, under positive mass on both \(C\cap E\) and \(C\cap\neg E\),
\[
  \texttt{countIndependentWithin\_boolRecoding\_iff\_constant\_of\_nonconstant\_on\_condition}.
\]
Thus the only deterministic recoding repair that survives the finite-count audit is the
information-erasing constant one, which is useless for a route that needs a discriminating
post-conditioned predictor.
This does not refute a future proof of the required v11 decorrelation estimate.  It does block the
naive inference from pre-conditioning independence of \(F\), \(A\), or \(b\) to independence after
conditioning on the USAT/uniqueness promise.

The latest Lean strengthening closes the immediate ``richer deterministic label''
variant of the same repair.  For any decidable codomain \(\alpha\) and any deterministic
recoding \(r:\{0,1\}\to\alpha\), every output fiber
\[
  \{\omega : r(\mathbf{1}_{E(\omega)}) = y\}
\]
is conditionally independent of the source bit \(E\), on a nondegenerate
conditioning event, if and only if that fiber does not distinguish the two
source-bit values:
\[
  r(1)=y \quad\Longleftrightarrow\quad r(0)=y.
\]
This is formalized as
\[
  \texttt{countIndependentWithin\_boolToAnyRecoding\_fiber\_iff\_fiber\_constant\_of\_nonconstant\_on\_condition}.
\]
Thus changing the output type of a deterministic recoding of the same conditioned bit
does not evade the audit: every output value that separates \(E\)-true states from
\(E\)-false states yields a fiber still coupled or anti-coupled to the source.

The quantitative extension, \texttt{ApproximateDecorrelationObstruction.lean}, closes
one immediate ``make it approximate'' variant on the same finite-count surface.  It
defines the two oriented cross-multiplied defects
\[
  \Delta_+(C,E,F)=|C|\,|C\cap E\cap F|-|C\cap E|\,|C\cap F|,
\]
and
\[
  \Delta_-(C,E,F)=|C\cap E|\,|C\cap F|-|C|\,|C\cap E\cap F|.
\]
A symmetric tolerance-\(\tau\) approximate independence claim must bound both
defects by \(\tau\).  Lean proves that under coupling \(E\leftrightarrow F\) inside
\(C\),
\[
  \texttt{countIndependentWithinForwardGap\_coupled}
  \quad\Longrightarrow\quad
  \Delta_+(C,E,F)=|C\cap E|\,|C\cap\neg E|,
\]
and under anti-coupling \(E\leftrightarrow\neg F\) inside \(C\),
\[
  \texttt{countIndependentWithinBackwardGap\_anticoupled}
  \quad\Longrightarrow\quad
  \Delta_-(C,E,F)=|C\cap E|\,|C\cap\neg E|.
\]
Consequently, the theorem
\[
  \texttt{not\_countIndependentWithinTolerance\_of\_coupled\_product\_gt}
\]
blocks any approximate repair whose tolerance is smaller than this explicit
source-fiber product, with the analogous anti-coupled theorem for
\(\Delta_-\).  The same module now lifts this quantitative result to every
decidable output fiber of an arbitrary deterministic recoding \(r:\{0,1\}\to\alpha\).
If the fiber over \(y\) distinguishes the two source-bit values,
\[
  \neg\bigl(r(1)=y \Longleftrightarrow r(0)=y\bigr),
\]
then
\[
  \texttt{not\_countIndependentWithinTolerance\_of\_boolToAnyRecoding\_fiber\_distinguishes\_product\_gt}
\]
blocks any symmetric tolerance below the same product
\(|C\cap E|\,|C\cap\neg E|\).  The obligation form is now also formalized:
\[
  \texttt{conditionedSourceFiberProduct\_le\_of\_CountIndependentWithinTolerance\_boolToAnyRecoding\_fiber\_distinguishes}.
\]
It says that any successful approximate-independence certificate for an
informative deterministic fiber forces
\[
  |C\cap E|\,|C\cap\neg E| \le \tau.
\]
The newest regression-pinned form removes one more escape hatch: if the two
conditioned source-bit fibers are known only through lower bounds \(L_1\) and
\(L_0\), then
\[
  \texttt{lowerBounds\_mul\_le\_of\_CountIndependentWithinTolerance\_boolToAnyRecoding\_fiber\_distinguishes}
\]
forces \(L_1L_0\le\tau\).  The direct contradiction form,
\[
  \texttt{not\_countIndependentWithinTolerance\_of\_boolToAnyRecoding\_fiber\_distinguishes\_lowerBounds\_mul\_gt},
\]
rules out the same approximate-independence certificate whenever
\(\tau<L_1L_0\).  Thus a repair cannot separately claim substantial mass on
both conditioned sides and a small approximate-independence tolerance for an
informative deterministic fiber; those two claims are now a direct finite-count
contradiction.
The latest Lean strengthening packages the same audit at the whole-output
level.  The predicate
\[
  \texttt{CountIndependentWithinToleranceOutput}
\]
requires every output fiber of a deterministic variable \(Y:\Omega\to\alpha\)
to satisfy the same tolerance bound.  For a nonconstant recoding
\(Y(\omega)=r(\mathbf{1}_{E(\omega)})\), Lean proves
\[
  \texttt{conditionedSourceFiberProduct\_le\_of\_CountIndependentWithinToleranceOutput\_boolToAnyRecoding\_nonconstant}
\]
and its lower-bound and contradiction forms.  Thus a repair cannot evade the
fiber theorem by claiming approximate independence only for the recoded output
variable as a whole: the \(r(1)\)-fiber already forces the same
\(|C\cap E|\,|C\cap\neg E|\le\tau\) obligation.
The newest finite-coin refinement covers the first randomized recoding subcase.
For \(Y(\omega,c)=r(\mathbf{1}_{E(\omega)},c)\), if a single output value \(y\)
uniformly selects the true source side for every coin,
\[
  \forall c,\ r(1,c)=y \quad\text{and}\quad \forall c,\ r(0,c)\ne y,
\]
or uniformly selects the false side with the inequalities reversed, then
\[
  \texttt{conditionedSourceFiberProduct\_le\_of\_CountIndependentWithinToleranceOutput\_coinwise\_trueFiber}
\]
and
\[
  \texttt{conditionedSourceFiberProduct\_le\_of\_CountIndependentWithinToleranceOutput\_coinwise\_falseFiber}
\]
force the same product/tolerance obligation on the lifted source/coin product
space.  The lower-bound, strict-tolerance contradiction, and zero-tolerance
forms are regression-pinned as well.  The follow-up source-count form uses
\[
  \texttt{finiteEventCount\_fst\_lift}
\]
to factor source-only events on \(\Omega\times\mathsf{Coin}\):
\[
  |(C\cap E)\times\mathsf{Coin}|=|C\cap E|\,|\mathsf{Coin}|.
\]
Consequently the coinwise true/false-fiber certificate must satisfy
\[
  (|C\cap E|\,|\mathsf{Coin}|)(|C\cap\neg E|\,|\mathsf{Coin}|)\le\tau,
\]
with matching lower-bound and strict-tolerance contradiction forms
\texttt{sourceLowerBounds\_coinProduct\_le\_of\_CountIndependentWithinToleranceOutput\_*}
and
\texttt{not\_countIndependentWithinToleranceOutput\_of\_coinwise\_*\_sourceLowerBounds\_mul\_gt}.
The zero-tolerance boundary is now source-level as well.  The helper
\texttt{finiteEventCount\_fst\_lift\_pos\_of\_pos} proves that any positive
source-only event remains positive on \(\Omega\times\mathsf{Coin}\) when the
finite coin space is nonempty.  Hence
\[
\begin{gathered}
  \texttt{not\_countIndependentWithinToleranceOutput\_zero\_of\_coinwise\_trueFiber\_sourceNonconstant},\\
  \texttt{not\_countIndependentWithinToleranceOutput\_zero\_of\_coinwise\_falseFiber\_sourceNonconstant}
\end{gathered}
\]
block uniform true-side or false-side distinguishing fibers from satisfying
zero-tolerance output independence under the source-side hypotheses
\(|C\cap E|>0\), \(|C\cap\neg E|>0\), and
\(\mathsf{Coin}\ne\emptyset\).  This removes the need to state the
nondegeneracy premise directly on the lifted product space.
The arbitrary finite-coin zero-tolerance boundary is now exact under the same
source-side nondegeneracy hypotheses.  Lean proves
\[
  \texttt{countIndependentWithinToleranceOutput\_zero\_iff\_finiteCoinRecoding\_fiberBalanced\_of\_sourceNonconstant},
\]
which says that zero-tolerance output independence on
\(\Omega\times\mathsf{Coin}\) holds if and only if every output fiber is
balanced between the true and false source sides.  The proof uses the one-sided
zero blockers
\[
\begin{gathered}
  \texttt{not\_countIndependentWithinToleranceOutput\_zero\_of\_finiteCoinRecoding\_trueMinusFalse\_sourceNonconstant},\\
  \texttt{not\_countIndependentWithinToleranceOutput\_zero\_of\_finiteCoinRecoding\_falseMinusTrue\_sourceNonconstant}
\end{gathered}
\]
to rule out any positive imbalance in either direction.
The next Lean refinement removes the uniform-fiber assumption for arbitrary
finite-coin recodings of the same source bit.  For each output value \(y\), let
\[
  T_y=\#\{c:r(1,c)=y\},\qquad F_y=\#\{c:r(0,c)=y\}.
\]
Lean proves the exact mixed-coin identities
\[
\begin{aligned}
  \Delta^+_y &=
    |C\cap E|\,|C\cap\neg E|\,|\mathsf{Coin}|\,(T_y-F_y),\\
  \Delta^-_y &=
    |C\cap E|\,|C\cap\neg E|\,|\mathsf{Coin}|\,(F_y-T_y),
\end{aligned}
\]
as
\[
  \texttt{countIndependentWithinForwardGap\_finiteCoinRecoding\_fiber}
\]
and
\[
  \texttt{countIndependentWithinBackwardGap\_finiteCoinRecoding\_fiber}.
\]
Thus a finite-coin randomized recoding can repair the same-source fiber only by
balancing the true-side and false-side coin counts for every output fiber, or
by making the source-fiber product small enough for the chosen tolerance.  The
strict blocker theorems
\[
  \texttt{not\_countIndependentWithinToleranceOutput\_of\_finiteCoinRecoding\_trueMinusFalse\_gt}
\]
and
\[
  \texttt{not\_countIndependentWithinToleranceOutput\_of\_finiteCoinRecoding\_falseMinusTrue\_gt}
\]
turn any unbalanced fiber into a direct contradiction when the claimed
tolerance is below the corresponding scaled imbalance.  The source-lower-bound
forms replace the exact source counts by \(L_1\le |C\cap E|\) and
\(L_0\le |C\cap\neg E|\), yielding the manuscript-facing obligations
\[
  L_1L_0|\mathsf{Coin}|(T_y-F_y)\le\tau
  \quad\text{and}\quad
  L_1L_0|\mathsf{Coin}|(F_y-T_y)\le\tau,
\]
with matching strict contradiction forms.
The current Lean packet packages this as a threshold quantization theorem:
below one lower-bounded source/coin imbalance block, approximate output
independence is already equivalent to exact balanced fibers.  The main form is
\[
  \texttt{countIndependentWithinToleranceOutput\_sourceLowerBounds\_lt\_coinProduct\_iff\_finiteCoinRecoding\_fiberBalanced},
\]
assuming \(L_1\le |C\cap E|\), \(L_0\le |C\cap\neg E|\), and
\(\tau<L_1L_0|\mathsf{Coin}|\).  The exact source-count specialization is
\[
  \texttt{countIndependentWithinToleranceOutput\_lt\_sourceCoinProduct\_iff\_finiteCoinRecoding\_fiberBalanced}.
\]
Thus small additive slack cannot hide a single finite-coin fiber imbalance;
it must either reach the first source/coin defect block or force exact
balance.
This does not close randomized residuals that use information beyond a
finite-coin recoding of the conditioned source bit, but it precisely closes the
finite-coin cancellation loophole for the same-source recoding surface.
Thus richer deterministic labels do not evade the quantitative audit: every
informative deterministic output fiber must either pay this finite-count defect
or prove that the relevant product is small after the manuscript's
normalization.  In the finite-coin case the analogous obligation is the scaled
coin-fiber imbalance, so exact coin-count balance is now the formal escape
condition rather than an informal repair slogan.
The balanced case is now formalized directly.  Lean defines
\[
  \texttt{FiniteCoinRecodingFiberBalanced}\ r
  \quad\equiv\quad
  \forall y,\ T_y=F_y,
\]
and proves
\[
  \texttt{countIndependentWithinToleranceOutput\_zero\_of\_finiteCoinRecoding\_fiberBalanced}.
\]
Thus a perfectly balanced finite-coin recoding really does decorrelate the
marginal output from the source bit at zero tolerance.  This is not a blocker;
it is the exact cancellation case.  The companion retained-coin theorem records
the price of using that cancellation while still preserving source information
through the random seed.  If the certified output is \((r(E(\omega),c),c)\)
and for some retained coin \(c_0\) one has \(r(1,c_0)\ne r(0,c_0)\), then Lean
proves
\[
  \texttt{countIndependentWithinForwardGap\_coinRetained\_trueFiber}
\]
with defect
\[
  |C\cap E|\,|C\cap\neg E|\,|\mathsf{Coin}|.
\]
The corresponding obligation and strict-blocker forms
\[
  \texttt{sourceCoinProduct\_le\_of\_CountIndependentWithinToleranceOutput\_coinRetained\_trueFiber}
\]
and
\[
  \texttt{not\_countIndependentWithinToleranceOutput\_of\_coinRetained\_trueFiber\_sourceProduct\_gt}
\]
show that retaining the coin revives the finite-count obstruction.  The newest
Lean strengthening turns this into an exact below-threshold boundary:
\[
  \texttt{finiteCoinRecodingFiberBalanced\_coinRetained\_iff\_pointwise\_eq}
\]
states that balanced fibers for the retained output \((r(b,c),c)\) are
equivalent to pointwise source-blindness, \(\forall c,\ r(1,c)=r(0,c)\).  Hence
below one lower-bounded source/coin block,
\[
  \texttt{countIndependentWithinToleranceOutput\_coinRetained\_sourceLowerBounds\_lt\_coinProduct\_iff\_pointwise\_eq}
\]
says approximate output independence of the retained-coin output is possible
exactly when every retained coin already carries no source-dependent output.
The retained-side counterpart now makes the same boundary fiberwise rather
than coinwise.  Lean proves
\[
  \texttt{finiteCoinRecodingFiberBalanced\_sideRetained\_iff\_sideFiber\_count\_eq},
\]
which says that balanced fibers for \((r(b,c),s(c))\) are equivalent to
equality of true and false counts in every retained output/side fiber:
\[
  \forall y,t,\quad
  \#\{c:r(1,c)=y\wedge s(c)=t\}
  =
  \#\{c:r(0,c)=y\wedge s(c)=t\}.
\]
Consequently, below one lower-bounded source/coin block,
\[
  \texttt{countIndependentWithinToleranceOutput\_sideRetained\_sourceLowerBounds\_lt\_coinProduct\_iff\_sideFiber\_count\_eq}
\]
turns approximate independence of the retained-side output into this exact
side-fiber multiset-balance condition.  A side channel can still aggregate
coins, but it cannot hide an imbalanced true/false multiset inside any
observed side/output fiber below the first source/coin threshold.
Therefore
the balanced-fiber escape has a sharp fork: the marginal output can be exactly
decorrelated only by erasing same-source evidence from output fibers.  This is
not limited to literally retaining the whole coin.  Lean now proves the
two-sided side-channel generalization
\[
  \texttt{countIndependentWithinForwardGap\_sideRetained\_trueFiber}
  \quad\text{and}\quad
  \texttt{countIndependentWithinBackwardGap\_sideRetained\_falseFiber}.
\]
For any retained side channel \(s:\mathsf{Coin}\to\mathsf{Side}\), if the
side fiber \(s(c)=s(c_0)\) contains no false-side output equal to
\(r(1,c_0)\), the output \((r(E,c),s(c))\) has forward defect
\[
  |C\cap E|\,|C\cap\neg E|\,|\mathsf{Coin}|\,
  \#\{c:r(1,c)=r(1,c_0)\wedge s(c)=s(c_0)\}.
\]
Symmetrically, if the same kind of side fiber contains no true-side output
equal to \(r(0,c_0)\), the selected false-side fiber has the same scaled
formula as a backward defect, with \(r(0,c_0)\) in place of \(r(1,c_0)\).
The obligation and strict-blocker forms
\[
  \texttt{sideChannelProduct\_le\_of\_CountIndependentWithinToleranceOutput\_sideRetained\_trueFiber}
  \quad\text{and}\quad
  \texttt{sideChannelProduct\_le\_of\_CountIndependentWithinToleranceOutput\_sideRetained\_falseFiber}
\]
and
\[
  \texttt{not\_countIndependentWithinToleranceOutput\_of\_sideRetained\_trueFiber\_product\_gt}
  \quad\text{and}\quad
  \texttt{not\_countIndependentWithinToleranceOutput\_of\_sideRetained\_falseFiber\_product\_gt}
\]
show that any side information capable of isolating such a source-bearing
fiber must pay the corresponding scaled product in the tolerance.  Lean now
removes the literal-pair assumption as well.  For any postprocessing
\(p:\alpha\times\mathsf{Side}\to\beta\), if the observed value
\(p(r(1,c_0),s(c_0))\) is distinct from every false-side observed value, then
\[
  \texttt{countIndependentWithinForwardGap\_postprocessedSide\_trueFiber}
\]
gives the same forward-defect formula with the selected observed-fiber count
\[
  \#\{c:p(r(1,c),s(c))=p(r(1,c_0),s(c_0))\}.
\]
The theorem
\[
  \texttt{countIndependentWithinBackwardGap\_postprocessedSide\_falseFiber}
\]
is the symmetric false-side version.  Thus hashing, compressing, or otherwise
observing the retained side channel does not evade the audit whenever some
observed fiber still separates the source side; the corresponding strict
blockers
\[
  \texttt{not\_countIndependentWithinToleranceOutput\_of\_postprocessedSide\_trueFiber\_product\_gt}
  \quad\text{and}\quad
  \texttt{not\_countIndependentWithinToleranceOutput\_of\_postprocessedSide\_falseFiber\_product\_gt}
\]
force the same scaled product into the tolerance.  Lean now
also records the exact below-threshold boundary for this hashed or otherwise
postprocessed side-channel case.  It proves
\[
  \texttt{finiteCoinRecodingFiberBalanced\_postprocessedSide\_iff\_observedFiber\_count\_eq},
\]
which says that balanced fibers for \(p(r(b,c),s(c))\) are exactly equality of
true-side and false-side counts for every observed postprocessed value:
\[
  \forall z,\quad
  \#\{c:p(r(1,c),s(c))=z\}
  =
  \#\{c:p(r(0,c),s(c))=z\}.
\]
Consequently,
\[
  \texttt{countIndependentWithinToleranceOutput\_postprocessedSide\_sourceLowerBounds\_lt\_coinProduct\_iff\_observedFiber\_count\_eq}
\]
says that below one lower-bounded source/coin block, approximate independence
of the postprocessed retained-side output is possible exactly under this
observed-fiber balance condition.  Hashing can merge side/output fibers, but
it must merge them into perfectly balanced observed fibers; otherwise the
finite-count obstruction reappears.
Lean also records the quantitative form behind this exact boundary:
\[
  \texttt{sourceLowerBounds\_observedFiberImbalanceProducts\_le\_of\_CountIndependentWithinToleranceOutput\_postprocessedSide}.
\]
Writing
\[
  T_z=\#\{c:p(r(1,c),s(c))=z\},\qquad
  F_z=\#\{c:p(r(0,c),s(c))=z\},
\]
the theorem says that any approximate-independence certificate with source
lower bounds \(\ell_1,\ell_0\) must satisfy, for every observed value \(z\),
\[
  \ell_1\ell_0|\mathsf{Coin}|(T_z-F_z)\le \tau
  \quad\text{and}\quad
  \ell_1\ell_0|\mathsf{Coin}|(F_z-T_z)\le \tau.
\]
The strict blockers
\[
  \texttt{not\_countIndependentWithinToleranceOutput\_of\_postprocessedSide\_observedFiber\_trueMinusFalse\_sourceLowerBounds\_gt}
\]
and
\[
  \texttt{not\_countIndependentWithinToleranceOutput\_of\_postprocessedSide\_observedFiber\_falseMinusTrue\_sourceLowerBounds\_gt}
\]
therefore block any hashed retained-side repair whose merged observed fibers
are merely less isolated but still imbalanced.  Postprocessing has to buy
actual cancellation, not just coarser observation.
Lean now also exposes the predicate-level downstream consequence:
\[
  \texttt{postprocessedSide\_outputPredicate\_count\_eq\_of\_countIndependentWithinToleranceOutput\_sourceLowerBounds\_lt\_coinProduct}.
\]
For every decidable property \(P\) of the postprocessed retained-side
observation, below the same source/coin block,
\[
  \#\{c:P(p(r(1,c),s(c)))\}
  =
  \#\{c:P(p(r(0,c),s(c)))\}.
\]
The strict blockers
\[
  \texttt{not\_countIndependentWithinToleranceOutput\_of\_postprocessedSide\_outputPredicate\_trueCount\_gt\_falseCount\_sourceLowerBounds\_lt\_coinProduct}
\]
and
\[
  \texttt{not\_countIndependentWithinToleranceOutput\_of\_postprocessedSide\_outputPredicate\_falseCount\_gt\_trueCount\_sourceLowerBounds\_lt\_coinProduct}
\]
rule out either direction of downstream predicate advantage.  The
downstream-decoder consequence is now explicit as well:
\[
  \texttt{postprocessedSide\_outputDecoder\_correctCount\_eq\_card\_of\_countIndependentWithinToleranceOutput\_sourceLowerBounds\_lt\_coinProduct}.
\]
For every Boolean decoder \(d\) of the postprocessed retained-side observation,
below the same source/coin block Lean proves
\[
  \#\{c:d(p(r(1,c),s(c)))=1\}
  +
  \#\{c:d(p(r(0,c),s(c)))=0\}
  =
  |\mathsf{Coin}|.
\]
Thus any decoder that does better than half on the two coin fibers refutes the
approximate-independence certificate.  The strict blocker is
\[
  \texttt{not\_countIndependentWithinToleranceOutput\_of\_postprocessedSide\_outputDecoder\_correctCount\_gt\_card\_sourceLowerBounds\_lt\_coinProduct}.
\]
Lean now records the uniform-recovery special cases directly.  A downstream
predicate that contains every true-side observed value and no false-side
observed value is blocked by
\[
  \texttt{not\_countIndependentWithinToleranceOutput\_of\_postprocessedSide\_outputPredicate\_separates\_sourceLowerBounds\_lt\_coinProduct}.
\]
A Boolean decoder that uniformly recovers the source bit after postprocessing
is blocked by
\[
  \texttt{not\_countIndependentWithinToleranceOutput\_of\_postprocessedSide\_outputDecoder\_recovers\_sourceLowerBounds\_lt\_coinProduct}.
\]
Finally, if the postprocessed true-side and false-side observed supports are
disjoint, then membership in the true-side support gives such a decoder, and
Lean proves the corresponding support-separation blocker
\[
  \texttt{not\_countIndependentWithinToleranceOutput\_of\_postprocessedSide\_sourceRangesSeparated\_sourceLowerBounds\_lt\_coinProduct}.
\]
The same result is also recorded in positive form as a required observed
collision:
\[
  \texttt{exists\_postprocessedSide\_true\_false\_collision\_of\_countIndependentWithinToleranceOutput\_sourceLowerBounds\_lt\_coinProduct}.
\]
Thus a deterministic postprocessing repair that still claims below-threshold
approximate independence must identify at least one true-side coin and one
false-side coin whose observed postprocessed values coincide.  It is not enough
to hide the retained pair syntactically; the repair must create actual
true/false support overlap in the finite coin observation.
Lean then strengthens this from a single collision to pointwise support
coverage:
\[
  \texttt{postprocessedSide\_pointwise\_true\_false\_collisions\_of\_countIndependentWithinToleranceOutput\_sourceLowerBounds\_lt\_coinProduct}.
\]
Equivalently, every true-side coin has a false-side witness with the same
postprocessed observation, and every false-side coin has a true-side witness.
The two one-sided witness forms are
\[
\begin{gathered}
  \texttt{exists\_false\_postprocessedSide\_collision\_of\_countIndependentWithinToleranceOutput\_sourceLowerBounds\_lt\_coinProduct},\\
  \texttt{exists\_true\_postprocessedSide\_collision\_of\_countIndependentWithinToleranceOutput\_sourceLowerBounds\_lt\_coinProduct}.
\end{gathered}
\]
Thus a repair cannot satisfy the below-threshold certificate by arranging only
one accidental overlap while leaving most finite-coin observations separated.
Lean now also records why even this pointwise collision target is not enough:
\[
  \texttt{postprocessedSide\_pointwise\_collisions\_do\_not\_imply\_coinEquiv\_preserving}.
\]
The witness is a three-coin, two-observation example with true-side
multiplicities \(2,1\) and false-side multiplicities \(1,2\).  Every observed
value appears on both source sides, so the pointwise collision obligations are
satisfied in both directions, but no bijective coin matching can preserve the
postprocessed observation.  This blocks the parallel repair route that accepts
support overlap as a replacement for the exact finite multiplicity matching.
Lean further isolates the remaining ``choose a witness'' repair.  A one-sided
true-to-false coin map that preserves the postprocessed observation is enough
only if it is injective:
\[
  \texttt{exists\_postprocessedSide\_coinEquiv\_preserving\_of\_injective\_coinMap\_preserving}.
\]
Equivalently, when no preserving coin equivalence exists, every preserving
one-sided witness map must be non-injective:
\[
  \texttt{not\_injective\_postprocessedSide\_coinMap\_preserving\_of\_no\_coinEquiv}.
\]
The same three-coin skew example realizes this obstruction in the strongest
finite form:
\[
  \texttt{postprocessedSide\_oneSided\_coinMap\_preserving\_does\_not\_imply\_coinEquiv\_preserving}.
\]
It has preserving one-sided maps, but every such map collapses two true-side
coins onto one false-side witness.  Thus witness choice alone cannot replace
the multiplicity-preserving bijection demanded by the postprocessed-side route.
The finite example is not the main theorem here.  Lean now records the abstract
repair equivalence:
\[
  \texttt{exists\_injective\_postprocessedSide\_coinMap\_preserving\_iff\_exists\_coinEquiv\_preserving}.
\]
For any finite coin type, an injective one-sided witness map is exactly the
original preserving coin-equivalence obligation.  With decidable observed
outputs this becomes the observed-fiber count condition:
\[
  \texttt{exists\_injective\_postprocessedSide\_coinMap\_preserving\_iff\_observedFiber\_count\_eq}.
\]
Consequently, any observed-fiber count mismatch forces every preserving
one-sided witness map to be non-injective:
\[
  \texttt{not\_injective\_postprocessedSide\_coinMap\_preserving\_of\_observedFiber\_count\_ne}.
\]
This generic pigeonhole theorem is the robust blocker for the injective-witness
repair route; the three-coin skew instance merely exhibits one small mismatch.
The next Lean strengthening records the multiplicity-level obligation behind
that support statement.  Exact observed-fiber equality gives a fiberwise
equivalence, and Lean assembles these fiberwise equivalences through
\texttt{Equiv.sigmaFiberEquiv} into a global matching of coin choices:
\[
  \texttt{exists\_postprocessedSide\_coinEquiv\_preserving\_of\_countIndependentWithinToleranceOutput\_sourceLowerBounds\_lt\_coinProduct}.
\]
The matched equivalence \(e : \mathsf{Coin}\simeq\mathsf{Coin}\) satisfies
\[
  \forall c,\quad
  \mathsf{post}(r(\mathsf{true},c),\mathsf{side}(c))
  =
  \mathsf{post}(r(\mathsf{false},e(c)),\mathsf{side}(e(c))).
\]
The direct contrapositive blocker is
\[
  \texttt{not\_countIndependentWithinToleranceOutput\_of\_no\_postprocessedSide\_coinEquiv\_preserving\_sourceLowerBounds\_lt\_coinProduct}.
\]
Lean also proves the local cardinality obstruction to this matching repair.
Any proposed preserving coin equivalence forces equality of the true-side and
false-side finite counts in every observed postprocessed fiber:
\[
  \texttt{postprocessedSide\_observedFiber\_count\_eq\_of\_coinEquiv\_preserving}.
\]
Consequently, a single observed value \(z\) with unequal true/false coin-fiber
multiplicities rules out every preserving coin equivalence,
\[
  \texttt{not\_exists\_postprocessedSide\_coinEquiv\_preserving\_of\_observedFiber\_count\_ne},
\]
and below the source/coin threshold it directly refutes approximate output
independence:
\[
  \texttt{not\_countIndependentWithinToleranceOutput\_of\_postprocessedSide\_observedFiber\_count\_ne\_sourceLowerBounds\_lt\_coinProduct}.
\]
Thus the repair is not allowed to merely produce the same set of observed
values on both source sides: every observed value must occur with the same
finite multiplicity, enough to support a true/false coin bijection preserving
the postprocessed observation.
The same matching obstruction is now exposed at the downstream-predicate
interface.  Any preserving coin equivalence forces equality of the true-side
and false-side counts for every decidable predicate of the postprocessed
observation,
\[
  \texttt{postprocessedSide\_outputPredicate\_count\_eq\_of\_coinEquiv\_preserving}.
\]
Hence a single predicate-count mismatch rules out every preserving coin
equivalence,
\[
  \texttt{not\_exists\_postprocessedSide\_coinEquiv\_preserving\_of\_outputPredicate\_count\_ne},
\]
and below the source/coin threshold blocks approximate output independence:
\[
  \texttt{not\_countIndependentWithinToleranceOutput\_of\_postprocessedSide\_outputPredicate\_count\_ne\_sourceLowerBounds\_lt\_coinProduct}.
\]
For decoder-only repair attempts, Lean records the corresponding half-accuracy
coin-matching obligation:
\[
  \texttt{postprocessedSide\_outputDecoder\_correctCount\_eq\_card\_of\_coinEquiv\_preserving},
\]
with contradiction forms
\[
  \texttt{not\_exists\_postprocessedSide\_coinEquiv\_preserving\_of\_outputDecoder\_correctCount\_ne\_card}
\]
and
\[
  \texttt{not\_countIndependentWithinToleranceOutput\_of\_postprocessedSide\_outputDecoder\_correctCount\_ne\_card\_sourceLowerBounds\_lt\_coinProduct}.
\]
Thus a repair cannot hide behind a downstream test or decoder instead of
exhibiting the full observed-fiber matching.
The predicate-only route is now fully collapsed, not just blocked in one
direction.  Lean proves that downstream predicate neutrality is equivalent to
observed-fiber balance,
\[
  \texttt{postprocessedSide\_outputPredicate\_count\_eq\_iff\_observedFiber\_count\_eq},
\]
and therefore equivalent to the existence of the preserving coin equivalence,
\[
  \texttt{exists\_postprocessedSide\_coinEquiv\_preserving\_iff\_outputPredicate\_count\_eq}.
\]
Combining this with the source/coin threshold gives the exact predicate-only
boundary
\[
  \texttt{countIndependentWithinToleranceOutput\_postprocessedSide\_sourceLowerBounds\_lt\_coinProduct\_iff\_outputPredicate\_count\_eq}.
\]
Thus a proposed certificate that all downstream tests are source-neutral is not
a weaker substitute for coin matching; it is precisely the same finite
multiplicity condition in predicate form.
The same equivalence has now been recorded directly at the independence layer:
\[
  \texttt{countIndependentWithinToleranceOutput\_postprocessedSide\_sourceLowerBounds\_lt\_coinProduct\_iff\_coinEquiv\_preserving}.
\]
Below the source/coin threshold, a postprocessed-side approximate independence
certificate exists if and only if there is a true/false coin equivalence
preserving every observed postprocessed value.
This closes another natural repair attempt: hashing or postprocessing the
retained side cannot leave a recoverable Boolean trace, or even disjoint
observed supports, for the source bit unless the independence budget has
already crossed the source/coin threshold.
The finite-output postprocessing foundation also separates lossless recoding
from genuine coarsening.  It now proves
\[
  \texttt{finiteCoinRecodingFiberBalanced\_postcompose\_iff\_of\_injective}
\]
and the corresponding predicate-erasure equivalence
\[
  \texttt{finiteCoinOutputPredicateErasing\_postcompose\_iff\_of\_injective}.
\]
Thus an injective deterministic postprocessor cannot create the balanced
finite-coin erasure needed by the repair; if postprocessing repairs the
obstruction, the map must be many-to-one.  The stronger observed-support
version is
\[
  \texttt{finiteCoinRecodingFiberBalanced\_postcompose\_iff\_of\_observedLeftInverse},
\]
with the corresponding predicate-erasure form
\[
  \texttt{finiteCoinOutputPredicateErasing\_postcompose\_iff\_of\_observedLeftInverse}.
\]
Even a many-to-one map cannot create erasure if the original output can be
recovered from the postprocessed value on every observed true/false coin
sample.  Therefore a coarsening repair must actually lose original-output
information on the finite support used by the argument.  Lean now sharpens this
to the collision level:
\[
  \texttt{finiteCoinRecodingFiberBalanced\_postcompose\_iff\_of\_observedInjective}
\]
and
\[
  \texttt{exists\_observedOutput\_collision\_of\_postcompose\_fiberBalanced\_of\_not\_fiberBalanced}.
\]
If the original finite-coin recoding is not balanced but the postprocessed
recoding is balanced, then two distinct observed original outputs must have the
same postprocessed value.  Thus the many-to-one repair owes an actual observed
collision, not merely a vague lossy postprocessing step.  Lean now sharpens
the same point quantitatively:
\[
  \texttt{countIndependentWithinToleranceOutput\_finiteCoinOutput\_postcompose\_iff\_of\_observedInjective}
\]
and
\[
  \texttt{countIndependentWithinToleranceOutput\_finiteCoinOutput\_postcompose\_iff\_of\_observedLeftInverse}.
\]
These prove that support-injective postprocessing, or more generally any
postprocessing with a left inverse on the actually observed outputs, preserves
the entire finite-coin output tolerance certificate exactly.  The defect-level
forms
\[
  \texttt{countIndependentWithinOutputDefect\_finiteCoinOutput\_postcompose\_eq\_of\_observedInjective}
\]
and
\[
  \texttt{countIndependentWithinOutputDefect\_finiteCoinOutput\_postcompose\_eq\_of\_observedLeftInverse}
\]
show that the maximum output defect itself is unchanged.  Therefore a
deterministic postprocessing repair cannot lower the approximate
decorrelation budget unless it genuinely loses original-output information on
the observed support.  Lossless relabeling, or any supportwise recoverable
coarsening, is quantitatively useless as a repair.  The new follow-up file
\texttt{FiniteCoinPostprocessingImprovementObstruction.lean} now records the
strict contrapositive:
\[
  \texttt{exists\_observedOutput\_collision\_of\_postcompose\_tolerance\_of\_not\_tolerance}
\]
and
\[
  \texttt{exists\_observedOutput\_collision\_of\_postcompose\_outputDefect\_lt}.
\]
If postprocessing makes a fixed tolerance-\(\tau\) certificate newly true, or
strictly lowers the finite-coin output defect, then some two distinct observed
original outputs must map to the same postprocessed value.  The paired
cardinality forms
\[
  \texttt{exists\_postprocessedFiber\_card\_gt\_one\_of\_postcompose\_tolerance\_of\_not\_tolerance}
\]
and
\[
  \texttt{exists\_postprocessedFiber\_card\_gt\_one\_of\_postcompose\_outputDefect\_lt}
\]
state the same repair obligation as a genuinely nontrivial merged fiber.  Thus
even approximate same-source improvement still owes a concrete observed
many-to-one step; it cannot come from a supportwise lossless re-encoding.  The next Lean
strengthening is source-bias sensitive:
\[
  \texttt{exists\_oppositeBias\_observedOutput\_collision\_of\_postcompose\_fiberBalanced\_of\_not\_fiberBalanced}.
\]
It proves that a balancing postprocessor must merge, inside one postprocessed
fiber, an original output with strictly more true-side coin mass and an
original output with strictly more false-side coin mass.  Same-bias collisions
therefore cannot be used as the repair; the coarsening must perform a real
opposite-bias cancellation on the observed finite support.  The conservation
form of the same obstruction is now explicit:
\[
  \texttt{finiteCoinSignedFiberBias\_postcompose\_eq\_sum\_preimage}
\]
and
\[
  \texttt{finiteCoinRecodingFiberBalanced\_postcompose\_iff\_forall\_sum\_signedFiberBias\_preimage\_eq\_zero}.
\]
Here \(\texttt{finiteCoinSignedFiberBias}\) is the integer true-minus-false
fiber count.  A successful deterministic coarsening is therefore equivalent to
zero total signed original bias in every postprocessed fiber.  This blocks a
parallel repair route in which the author supplies collisions but not exact
cancellation certificates: the Lean statement asks for the full per-fiber
integer identity.  The degenerate endpoint is explicit
as
\[
  \texttt{finiteCoinRecodingFiberBalanced\_constant},
\]
which says a constant output is always balanced.  That is source erasure, not
evidence that the original finite-coin signal survived postprocessing.
The new file \texttt{FiniteCoinPostprocessingSameSignObstruction.lean} now
pushes the approximate story to the same cancellation boundary.  It defines
\texttt{finiteOutputFiberSignedBiasOneSided} and proves
\[
  \texttt{countIndependentWithinOutputDefect\_finiteCoinOutput\_le\_of\_signedBiasOneSided}.
\]
If every postprocessed fiber merges only nonnegative signed biases, or only
nonpositive signed biases, then deterministic postprocessing cannot strictly
lower the finite-coin output defect at all.  The contrapositive forms
\[
  \texttt{exists\_oppositeSign\_observedOutput\_collision\_of\_postcompose\_outputDefect\_lt}
\]
and
\[
  \texttt{exists\_oppositeSign\_observedOutput\_collision\_of\_postcompose\_tolerance\_of\_not\_tolerance}
\]
therefore show that strict approximate improvement requires a merged fiber that
contains both a positive-bias original output and a negative-bias original
output.  Mere same-sign coarsening is not enough: the repair must perform
genuine opposite-sign cancellation on the finite observed support.
The new file \texttt{FiniteCoinSignedBiasQuantizerObstruction.lean} turns this
into the manuscript-facing finite-quantization statement.  It defines
\texttt{finiteCoinSignedBiasQuantizerOneSided} for quantizers
\(q : \mathbb{Z} \to \beta\) of the signed true-minus-false boundary bias, and
the induced recoding \texttt{finiteCoinSignedBiasQuantizedRecoding}.  The core
theorem
\[
  \texttt{countIndependentWithinOutputDefect\_finiteCoinOutput\_le\_of\_finiteCoinSignedBiasQuantizerOneSided}
\]
says that if every quantizer cell lies entirely in the nonnegative biases or
entirely in the nonpositive biases, then quantizing the boundary bias cannot
strictly lower the finite-coin output defect.  The contrapositive theorems
\[
  \texttt{not\_finiteCoinSignedBiasQuantizerOneSided\_of\_quantizedOutput\_outputDefect\_lt}
\]
and
\[
  \texttt{not\_finiteCoinSignedBiasQuantizerOneSided\_of\_quantizedOutput\_tolerance\_of\_not\_tolerance}
\]
therefore say that any successful finite quantization repair must use some
bucket that crosses zero.  The explicit collision forms
\[
  \texttt{exists\_oppositeSign\_quantizerCollision\_of\_quantizedOutput\_outputDefect\_lt}
\]
and
\[
  \texttt{exists\_oppositeSign\_quantizerCollision\_of\_quantizedOutput\_tolerance\_of\_not\_tolerance}
\]
state that one quantizer value must identify a positive-bias and
negative-bias original output.  So a side-preserving bucketization of the
boundary field cannot be the repair; any helpful finite quantizer must perform
genuine cross-zero cancellation.
The new file
\texttt{FiniteCoinSignedBiasDitheredQuantizerObstruction.lean}
lifts the same obstruction to the manuscript's explicit finite-public-dither
variant.  Lean models the dithered quantizer by adjoining the seed to the
observed output,
\[
  \texttt{finiteCoinOutputWithDither},
\]
and then applying a deterministic map of the pair
\((\text{signed bias}, \text{seed})\).  The theorem
\[
  \texttt{countIndependentWithinOutputDefect\_finiteCoinOutputWithDither\_le\_of\_finiteCoinSignedBiasDitheredQuantizerOneSided}
\]
says that if every dithered quantizer bucket contains only nonnegative signed
biases or only nonpositive signed biases, uniformly across all seeds, then the
dithered quantization still cannot strictly lower the finite-coin output
defect.  The tolerance form
\[
  \texttt{not\_countIndependentWithinToleranceOutput\_finiteCoinSignedBiasDitheredQuantizedRecoding\_of\_not\_tolerance\_of\_oneSided}
\]
shows that such a one-sided dithered quantizer also cannot make a failing
tolerance certificate newly true.  Consequently, any successful finite
dithered quantizer must have a bucket that merges a positive-bias
output/seed pair with a negative-bias output/seed pair, as recorded in
\[
  \texttt{exists\_oppositeSign\_ditheredQuantizerCollision\_of\_quantizedOutput\_outputDefect\_lt}
\]
and
\[
  \texttt{exists\_oppositeSign\_ditheredQuantizerCollision\_of\_quantizedOutput\_tolerance\_of\_not\_tolerance}.
\]
So “independent dither” does not create a new escape hatch here: it is still
just deterministic postprocessing on a larger finite observed support, and any
helpful bucket must perform genuine cross-zero cancellation.
The follow-up file
\texttt{FiniteCoinSignedBiasDitheredQuantizerExactCancellation.lean}
pushes the same dither lane to the exact zero-tolerance boundary.  The theorem
\[
  \texttt{finiteCoinSignedFiberBias\_finiteCoinSignedBiasDitheredQuantizedRecoding\_eq\_sum\_preimage}
\]
states that the signed bias of one dithered quantizer output is exactly the sum
of the original signed biases over the corresponding \((\text{output},
\text{seed})\) preimage fiber.  Therefore
\[
  \texttt{countIndependentWithinToleranceOutput\_zero\_finiteCoinSignedBiasDitheredQuantizedRecoding\_iff\_forall\_sum\_signedFiberBias\_preimage\_eq\_zero}
\]
shows that exact output balance after dithered quantization is equivalent to a
literal integer cancellation law in every dithered bucket.  The contrapositive
\[
  \texttt{not\_countIndependentWithinToleranceOutput\_zero\_finiteCoinSignedBiasDitheredQuantizedRecoding\_of\_exists\_nonzero\_sum\_signedFiberBias\_preimage}
\]
blocks exact repair as soon as one dithered quantizer bucket has nonzero total
signed bias on the enlarged support.  So public dither does not merely require
an opposite-sign collision somewhere; every surviving bucket must satisfy exact
integer cancellation.  Lean now also identifies this law with the older
retained-side postprocessing boundary rather than treating it as a separate
escape route: the theorem
\[
  \texttt{finiteCoinSignedBiasDitheredQuantizedRecoding\_exactCancellation\_iff\_observedBucketCount\_eq}
\]
shows that exact signed-bias cancellation is exactly the statement that every
observed dithered quantizer bucket contains equally many true-side and
false-side witnesses, and the route-facing theorem
\[
  \texttt{countIndependentWithinToleranceOutput\_zero\_finiteCoinSignedBiasDitheredQuantizedRecoding\_iff\_observedBucketCount\_eq}
\]
packages zero-tolerance success in precisely that bucketwise count-equality
form.
Lean now
also proves this erasure claim first in predicate form,
\[
  \texttt{not\_finiteCoinRecodingFiberBalanced\_of\_outputPredicate\_separates},
\]
and then in decoder form:
\[
  \texttt{not\_finiteCoinRecodingFiberBalanced\_of\_outputDecoder\_recovers}.
\]
For any nonempty coin space, no output-only property can hold on every
true-side marginal output and fail on every false-side marginal output while
the recoding remains fiber-balanced; the Boolean decoder result is just the
special case where the property is ``the decoder returns \(1\).''  If a repair
instead weakens this to a partial output-predicate advantage, the aggregate
count theorem
\[
  \texttt{finiteCoinRecodingFiberBalanced\_outputPredicate\_count\_eq}
\]
still blocks it: for every decidable output predicate \(P\),
\[
  \#\{c:P(r(1,c))\}=\#\{c:P(r(0,c))\}.
\]
The strict forms
\[
  \texttt{not\_finiteCoinRecodingFiberBalanced\_of\_outputPredicate\_trueCount\_gt\_falseCount}
\]
and
\[
  \texttt{not\_finiteCoinRecodingFiberBalanced\_of\_outputPredicate\_falseCount\_gt\_trueCount}
\]
turn either directional predicate-count advantage into a contradiction to
fiber balance.  Specializing to a Boolean decoder gives the sharper
half-accuracy statement
\[
  \texttt{finiteCoinRecodingFiberBalanced\_outputDecoder\_correctCount\_eq\_card}.
\]
If \(D:\alpha\to\{0,1\}\), then balanced fibers force
\[
  \#\{c:D(r(1,c))=1\}+\#\{c:D(r(0,c))=0\}=|\mathsf{Coin}|.
\]
Thus a marginal output decoder is correct on exactly half of the two
source-side coin trials; Lean's strict form
\[
  \texttt{not\_finiteCoinRecodingFiberBalanced\_of\_outputDecoder\_correctCount\_gt\_card}
\]
blocks any better-than-half aggregate decoder.  If a repair instead preserves
the source by carrying the coin as side information, the source-product
obstruction returns, scaled by \(|\mathsf{Coin}|\).
This still does not refute every possible approximate
decorrelation strategy; it does prove that the manuscript cannot simply weaken
exact conditional independence to approximate conditional independence on a
nondegenerate same-source surface, including these whole-output deterministic
and arbitrary finite-coin randomized recoding subcases, without supplying a real
quantitative small-product, coin-fiber balance, or normalization estimate.

\paragraph{Current blocker status.}
This is now a high-confidence local blocker for the deterministic Boolean
conditioning repair pattern and its immediate arbitrary-codomain deterministic
extension: copied bits, complemented bits, all nonconstant deterministic Boolean
recodings, and every deterministic recoding of the same bit into a decidable
codomain have exactly the same empty/full degeneracy boundary under the
finite-count fiber audit.  It is not yet a global \(>90\%\) blocker for the
whole PNP route.  A future repair could still avoid the theorem by proving that
the post-conditioned variable uses residual information not determined by the
conditioned source bit, by working with a randomized residual object, or by
supplying a quantitatively justified approximate theorem whose allowed error
actually dominates the finite-count product defect on the relevant conditioned
surface.
To upgrade this from a local crux to a route-killing crux, the remaining task is
to connect the manuscript's uniqueness-conditioned variables to this exact
deterministic same-source recoding surface, or to extend the theorem to the
richer post-conditioned object actually used by the proof attempt.

It should therefore not be presented as a reason for the author to abandon the
PNP project.  Its proper role is diagnostic: it says that a proof continuing
along this nearby line must supply a real post-conditioning decorrelation
bridge, not merely a different deterministic Boolean recoding of the same
conditioned bit.  A stop-grade crux would need a broader route-class theorem
showing that the manuscript's natural repairs still fall inside a blocked
formal family, or an equally broad counterexample to the required conditioning
bridge itself.

\section{Obstruction II: The Full Local Rule Class Is Enormous}

Let $n$ be the number of visible Boolean features after switching.  Then the unconstrained class of
local predictors is simply
\[
  \{0,1\}^{\{0,1\}^n},
\]
the set of all Boolean functions on $n$ bits.  Its size is
\[
  2^{2^n}.
\]
Therefore any \emph{uniform} code space that represents all such rules injectively must have at
least $2^n$ bits.

This is formalized in \texttt{CompressionObstruction.lean}.

\begin{theorem}[Counting bound for local rules]
Let \texttt{LocalRule n} denote the type of Boolean functions on \texttt{n} visible bits, and
\texttt{BitCode s} the type of \texttt{s}-bit codewords.  If
\[
  s < 2^n,
\]
then there is no injective map
\[
  \texttt{LocalRule n} \to \texttt{BitCode s}.
\]
\end{theorem}

\noindent
This is the theorem \texttt{no\_injective\_bitCode\_of\_lt}.

\begin{theorem}[Binary-log neighborhood corollary]
If $d \ge 2$ and
\[
  s < 2^m,
\]
then there is no injective map
\[
  \texttt{LocalRule}\!\left(d^{\log_2 m + 1}\right) \to \texttt{BitCode s}.
\]
\end{theorem}

\noindent
This is the theorem
\texttt{no\_uniform\_code\_below\_expInput\_for\_binaryLogNeighborhood}.

\begin{theorem}[Size lower bound on the full class]
For $d \ge 2$,
\[
  2^{2^m} \le
  \left|\texttt{LocalRule}\!\left(d^{\log_2 m + 1}\right)\right|.
\]
\end{theorem}

\noindent
This is the theorem
\texttt{expExpInput\_le\_card\_localRule\_binaryLogNeighborhood}.

\subsection{Interpretation}

Again, this does not say that the switched rule class \emph{is} the full local rule class.  It says
that without a theorem proving otherwise, one cannot pass from
``local on $n$ visible bits'' to ``describable by only $\mathrm{polylog}(m)$ bits.''

That gap is exactly where the current program remains vulnerable.

\section{What These Results Break}

The present Lean work targets the following inference schema:
\[
  \text{short decoder}
  \Longrightarrow
  \text{switched local predictor}
  \Longrightarrow
  \text{tiny ACC$^0$/ERM hypothesis}.
\]

The first arrow is the paper's switching claim.  The second arrow is where the present obstruction
lands.

\begin{remark}
If the switched predictor is only known to be some arbitrary function of a log-radius neighborhood,
then:
\begin{enumerate}[leftmargin=1.5em]
  \item the neighborhood already has polynomial size in the ambient parameter, and
  \item the class of all such local rules is far too large to admit a polylogarithmic uniform code.
\end{enumerate}
\end{remark}

Therefore a successful proof attempt must supply one of the following:
\begin{enumerate}[leftmargin=1.5em]
  \item a theorem that the switched predictors belong to a very special subclass
  (for example, a family with explicit circuit/description bound), or
  \item a direct compression algorithm showing that the relevant local rules can be encoded in
  \(\mathrm{polylog}(m)\) bits for reasons stronger than locality.
\end{enumerate}

Without such a theorem, the proof's route from locality to tiny hypothesis class is unsupported.

\section{Obstruction III: Zero-Margin Symmetrization Cannot Concentrate}

The symmetrization route in the paper makes a stronger claim than the two counting obstructions
above.  It says, roughly, that after taking a polylogarithmic number of promise-preserving sign
flips and then taking a majority vote, the resulting surrogate label agrees with the Bayes rule on
the local $\sigma$-field for all but a negligible fraction of blocks.

There is a direct obstruction here when the conditional surrogate mean lands at exactly
\(\tfrac12\).

\subsection{Lean counting lemma}

The Lean module \texttt{SymmetrizationObstruction.lean} formalizes the following fact.

\begin{theorem}[Odd-majority is exactly balanced on unbiased bits]
Let \(s\) be odd.  Among the \(2^s\) bit-vectors of length \(s\), exactly half have strict majority
\texttt{true} and half have strict majority \texttt{false}.
\end{theorem}

\noindent
The key formal statements are:
\begin{itemize}[leftmargin=1.5em]
  \item \texttt{majority\_compl}: complementing all bits flips the majority bit when \(s\) is odd;
  \item \texttt{card\_majority\_true\_eq\_card\_majority\_false}: the majority-\texttt{true} and
  majority-\texttt{false} classes have equal cardinality;
  \item \texttt{majority\_tie\_break\_not\_high\_probability}: if a Bayes rule breaks a
  \(p=\tfrac12\) tie toward \texttt{true}, then majority on \(s\) unbiased samples matches that rule
  on only half the sample space.
\end{itemize}

\subsection{Application to the Current Symmetrization Wrapper}

Now consider the paper's symmetrized surrogate label for a fixed block \(j\) and bit \(i\).  Take
the simplest possible decoder \(P\): the decoder that always outputs \(0\).

Under a sign flip \(\sigma\), the paper's own back-mapping rule gives
\[
  Y^{(r)}_{j,i} = \langle a_{j,i}, \sigma^{(r)} \rangle
\]
over \(\mathbb{F}_2\), because the decoder output is \(0\) and the back-map xors in the VV shift.

If \(a_{j,i} \neq 0\), then \(\langle a_{j,i}, \sigma \rangle\) is an unbiased bit under uniform random
\(\sigma\).  Hence the sequence
\[
  Y^{(1)}_{j,i}, \dots, Y^{(s)}_{j,i}
\]
is a fair sample of surrogate bits, and for odd \(s\) the majority vote is exactly balanced.

But the paper's Bayes rule is written with threshold \(1[p \ge \tfrac12]\), which at \(p=\tfrac12\)
chooses the value \(1\).  Therefore, in this zero-margin case, the majority vote matches the stated
Bayes rule on only half the seed choices, not with probability \(1 - m^{-\Omega(1)}\).

\subsection{Consequence}

This is materially stronger than the earlier ``local does not imply simple'' objection.  It says:

\begin{quote}
The concentration-to-Bayes step in the symmetrization argument is false unless an additional
margin theorem is proved.
\end{quote}

In other words, one must show that for the relevant blocks and bits,
\[
  \left| p_{j,i} - \tfrac12 \right|
\]
is bounded below by more than the sampling noise scale.  Without such a margin, majority voting
cannot recover the Bayes rule uniformly, even for the constant-zero decoder.

\section{Obstruction IIIb: The Exact Post-Switch Input Is Not \(T_i\)-Invariant}

Before moving to a charitable repair, the manuscript's own exact definitions already create a
direct problem.

Section 3.2 defines the promise-preserving involution
\[
  T_i(F^h, A, b) = (F^{\tau_i h}, A, b \oplus a_i),
\]
while Section 3.3 defines the post-switch local input
\[
  u_i(\Phi) = (z(F^h), a_i, b).
\]
Appendix A.5 then says this same involution preserves \(u=(z,a_i,b)\).  The new Lean module
\texttt{PostSwitchInputObstruction.lean} formalizes the exact contradiction point.

\subsection{Lean theorem}

\begin{theorem}[Full-input invariance holds only on zero VV columns]
Let \(u = (z,a,b)\) be the exact post-switch local input, and let \(T\) act on it by
\[
  T(z,a,b) = (z,a,b \oplus a).
\]
Then
\[
  T(u) = u
  \qquad\Longleftrightarrow\qquad
  a = 0.
\]
Equivalently, if the VV column \(a\) is nonzero, then the involution necessarily changes the
right-hand side coordinate \(b\).
\end{theorem}

\noindent
The key formal statements are:
\begin{itemize}[leftmargin=1.5em]
  \item \texttt{vvToggle\_eq\_self\_iff\_zero}: \(b \oplus a = b\) iff \(a=0\);
  \item \texttt{tiInputMap\_eq\_self\_iff\_zeroColumn}: the exact local input is \(T_i\)-invariant
  iff the VV column is zero;
  \item \texttt{invariantProjection\_tiInputMap}: the reduced projection \((z,a)\) is always
  preserved;
  \item \texttt{tiInputMap\_changes\_b\_of\_nonzeroColumn}: if \(a \neq 0\), then \(b\) really
  changes.
\end{itemize}

\subsection{Why this matters}

This is not a soft complaint about omitted detail.  It is a direct incompatibility between the
paper's own formulas:
\begin{enumerate}[leftmargin=1.5em]
  \item the exact local input keeps \(b\),
  \item the involution toggles \(b\) by \(a_i\),
  \item yet the calibration appendix says the involution preserves \(u=(z,a_i,b)\).
\end{enumerate}

That can only be true on the special case \(a_i=0\).  So the calibration argument cannot be right
as written on the full local input.

\subsection{Consequence}

This forces the proof into a fork:
\begin{enumerate}[leftmargin=1.5em]
  \item either keep the full input \((z,a_i,b)\), in which case the involution is not feature
  preserving and a new calibration theorem is needed;
  \item or drop the changing coordinate \(b\), in which case one lands in the invariant-feature
  regime attacked next.
\end{enumerate}

\section{Obstruction IIIc: The Exact Manuscript Fork Already Lands in the Symmetry Budget}

The previous section exposed the formula-level inconsistency in the manuscript as written.  The new
Lean module \texttt{PostSwitchForkObstruction.lean} shows that the most charitable exact repair is
already constrained by the later symmetry-budget files.

\subsection{Lean theorems}

\begin{theorem}[The exact \(T_i\)-action is an involution and \(b\) resolves exactly the nonzero columns]
On the exact local input \(u=(z,a,b)\) with
\[
  T(z,a,b) = (z,a,b \oplus a),
\]
the map \(T\) is an involution.  Moreover,
\[
  b(Tu) \neq b(u)
  \qquad\Longleftrightarrow\qquad
  a \neq 0.
\]
Equivalently, the \(b\)-coordinate distinguishes an involution pair exactly on the nonzero-column
slice.
\end{theorem}

\begin{theorem}[The exact fork]
Let the target bit flip under the exact involution \(T\), and let the weighting be \(T\)-invariant.
Then:
\begin{enumerate}[leftmargin=1.5em]
  \item every soft score depending only on the reduced invariant projection \((z,a)\) has zero
  weighted signed correlation with the target;
  \item every classifier using \(((z,a),b)\) has success advantage over chance bounded by the
  total weight of the nonzero-column slice.
\end{enumerate}
\end{theorem}

\noindent
The key formal statements are:
\begin{itemize}[leftmargin=1.5em]
  \item \texttt{tiInputMap\_involutive}: the manuscript's exact local update really is an
  involution;
  \item \texttt{resolvedBySideChannel\_tiInputMap\_b\_iff\_nonzeroColumn}: the \(b\)-side channel
  resolves an orbit pair exactly when \(a \neq 0\);
  \item \texttt{weighted\_signedScore\_sum\_eq\_zero\_on\_invariantProjection}: the reduced exact
  input \((z,a)\) has zero signed signal under symmetry-respecting weights;
  \item \texttt{doubledAdvantage\_invariantProjection\_with\_b\_le\_nonzeroColumnMass}: any
  classifier on \(((z,a),b)\) pays for its domination gap with nonzero-column mass;
  \item \texttt{pos\_nonzeroColumnMass\_of\_pos\_doubledAdvantage\_invariantProjection\_with\_b}:
  any positive \(b\)-side-channel advantage forces positive nonzero-column mass;
  \item \texttt{exists\_b\_resolving\_input\_of\_pos\_doubledAdvantage\_invariantProjection\_with\_b}:
  any positive \(b\)-side-channel advantage exposes a positive-weight input where
  \(b(Tu)\ne b(u)\);
  \item \texttt{exists\_nonzeroColumn\_of\_pos\_doubledAdvantage\_invariantProjection\_with\_b}:
  the same advantage exposes a positive-weight input with \(a\ne 0\);
  \item \texttt{not\_exactInputInvariantWeightedSupport\_of\_pos\_doubledAdvantage\_invariantProjection\_with\_b}:
  such positive advantage refutes the manuscript-style exact-input-invariant support premise;
  \item \texttt{pos\_nonzeroColumnMass\_of\_total\_lt\_two\_mul\_weightedCorrectMass\_invariantProjection\_with\_b}:
  the same nonzero-column obligation now also holds on the strict half-accuracy surface
  \(\mathrm{total}<2\cdot\mathrm{correct}\);
  \item \texttt{not\_total\_lt\_two\_mul\_weightedCorrectMass\_invariantProjection\_with\_b\_of\_exactInputInvariantWeightedSupport}
  and
  \texttt{not\_exactInputInvariantWeightedSupport\_of\_total\_lt\_two\_mul\_weightedCorrectMass\_invariantProjection\_with\_b}:
  exact-input-invariant support blocks the strict half-accuracy route just as decisively as it
  blocks positive doubled advantage.
\end{itemize}

\subsection{Why this matters}

This turns the manuscript fork into a direct quantitative obligation.

If one drops \(b\), then the surviving exact input is already inside the invariant-feature regime:
no symmetry-respecting signed signal remains.  If one keeps \(b\), then the proof is no longer
allowed to say merely that ``a small side channel helps.''  It must prove that the resulting
domination gap is supported by a sufficiently large nonzero-column slice in the actual weighted
domain.  More sharply, the strengthened Lean API now extracts an actual positive-weight witness
on that slice, so an aggregate domination claim must identify where the local input really changes
under \(b\mapsto b\oplus a_i\).  The same API closes the immediate manuscript escape: if the same weighted
support is still required to be exact-input invariant, then it has zero nonzero-column mass, so any
positive \(b\)-side-channel advantage contradicts that invariance premise outright.

So this exact bridge does not merely identify a local inconsistency.  It shows that the natural
repair lands immediately in the later neutrality and resolution-budget obstructions.

\section{Obstruction IV: Involution-Invariant Feature Collapse Forces Exact Neutrality}

Once the exact full-input inconsistency is noticed, the most natural charitable repair is to drop
the involution-sensitive coordinate \(b\) and keep only the involution-fixed part of the local view,
say a feature map of the form
\[
  u' = (z, a_i).
\]

This repair looks attractive because it restores feature invariance.  The new Lean module
\texttt{OrbitNeutralityObstruction.lean} shows why this move cannot, by itself, support a
domination theorem.

\subsection{Lean theorem}

\begin{theorem}[Feature-preserving involutions force exact half accuracy]
Let \(\alpha\) be a finite sample space, let \(\tau : \alpha \to \alpha\) be an involution, let
\(u : \alpha \to U\) be a feature map, and let \(y : \alpha \to \{0,1\}\) be the target bit.
Assume
\[
  u(\tau x) = u(x)
  \qquad\text{and}\qquad
  y(\tau x) = 1 - y(x)
\]
for every \(x \in \alpha\).  Then for every classifier \(h : U \to \{0,1\}\),
\[
  2 \cdot \#\{x : h(u(x)) = y(x)\} = |\alpha|.
\]
Equivalently, the average accuracy of \(h \circ u\) is exactly \(1/2\).
\end{theorem}

\noindent
The key formal statements are:
\begin{itemize}[leftmargin=1.5em]
  \item \texttt{correct\_under\_pair\_iff\_incorrect}: correctness at \(x\) is equivalent to
  incorrectness at \(\tau x\);
  \item \texttt{card\_correct\_eq\_card\_incorrect}: the correct and incorrect classes have equal
  cardinality;
  \item \texttt{two\_mul\_card\_correct\_eq\_card}: exactly half the finite sample space is
  classified correctly.
\end{itemize}

\subsection{Why this hits the charitable calibration patch}

Suppose the repair to the calibration bug is:
\begin{enumerate}[leftmargin=1.5em]
  \item remove the involution-sensitive coordinate \(b\) from the local input;
  \item prove neutrality by pairing each instance with its \(T_i\)-flipped partner;
  \item then learn or evaluate a classifier using only the reduced feature map \(u'\).
\end{enumerate}

If the paired instances really preserve \(u'\) and flip the target bit, then the theorem above
applies immediately.  Every feature-only classifier on \(u'\) has exact \(1/2\) average accuracy.
So this route cannot yield a strict domination statement, a high-probability Bayes recovery
statement, or any per-block success advantage above chance.

\subsection{Consequence}

This is a stronger objection than ``the paper needs more detail.''  It says:

\begin{quote}
If the calibration fix works by dropping the coordinate changed by the involution and then
classifying from the invariant remainder alone, that fix destroys any chance of beating
\(1/2\) accuracy on the paired finite domain.
\end{quote}

So a viable rescue must do something more subtle than mere feature collapse.  For example, it
would have to:
\begin{enumerate}[leftmargin=1.5em]
  \item keep some controlled non-invariant information,
  \item switch from hard classification to a softer statistic whose expectation matters,
  \item or show that the actual learning/evaluation domain is not the balanced finite orbit space
  modeled by the theorem above.
\end{enumerate}

\section{Obstruction IVb: Invariant Feature Fibers Have Zero Margin}

The most natural reply to Obstruction IV is:
\begin{quote}
Fine, hard classification on the invariant features fails.  But one can still estimate a
\emph{soft score} or conditional mean on those features and only use high-margin feature values.
\end{quote}

This is the right stronger target.  The new Lean module
\texttt{FiberNeutralityObstruction.lean} blocks that repair under the same involution
hypotheses.

\subsection{Lean theorem}

\begin{theorem}[Every invariant feature fiber is exactly balanced]
Let \(\alpha\) be a finite sample space, let \(\tau : \alpha \to \alpha\) be an involution, let
\(u : \alpha \to U\) be a feature map, and let \(y : \alpha \to \{0,1\}\) be the target bit.
Assume
\[
  u(\tau x) = u(x)
  \qquad\text{and}\qquad
  y(\tau x) = 1 - y(x)
\]
for every \(x \in \alpha\).  Then for every retained feature value \(v \in U\),
\[
  2 \cdot \#\{x : u(x)=v \text{ and } y(x)=1\}
  =
  \#\{x : u(x)=v\}.
\]
Equivalently, conditioned on any retained feature value \(v\), the target bit has exact mean
\(1/2\).
\end{theorem}

\noindent
The key formal statements are:
\begin{itemize}[leftmargin=1.5em]
  \item \texttt{featureFiberTrueEquivFalse}: on each fiber \(u^{-1}(v)\), the involution pairs
  the \texttt{true} points with the \texttt{false} points;
  \item \texttt{card\_featureFiberTrue\_eq\_card\_featureFiberFalse}: the two parts of each fiber
  have equal cardinality;
  \item \texttt{two\_mul\_card\_featureFiberTrue\_eq\_card\_featureFiber}: the \texttt{true}
  points occupy exactly half of each fiber.
\end{itemize}

\subsection{Why this hits the soft-score repair}

This is stronger than the hard-classifier obstruction.  It says the retained feature map does not
merely fail to support a good classifier; it fails to support \emph{any nonzero conditional
margin}.  On every retained feature value \(v\), the exact local average is already
\[
  \mathbb{E}[\,y \mid u=v\,] = \tfrac12.
\]

So if the repair route is:
\begin{enumerate}[leftmargin=1.5em]
  \item drop the involution-sensitive coordinate,
  \item keep only the invariant summary \(u'\),
  \item estimate a soft score or conditional mean from \(u'\),
  \item then use only ``high-margin'' \(u'\)-values,
\end{enumerate}
it fails immediately on the balanced finite orbit model.  There are no high-margin retained
feature values at all.

\subsection{Consequence}

This narrows the remaining rescue space again.  A successful patch now has to do more than
``replace hard labels by soft scores.''  It must either:
\begin{enumerate}[leftmargin=1.5em]
  \item retain some controlled non-invariant information,
  \item show that the actual averaging domain is not fiberwise balanced under the involution,
  \item or use a structure theorem stronger than feature-invariant orbit pairing.
\end{enumerate}

\section{Obstruction IVc: Symmetry-Respecting Weights Still Have Zero Margin}

Another natural reply is:
\begin{quote}
The real evaluation or training average is not uniform on the finite orbit space.  Once one uses
the correct non-uniform weighting, the invariant features may regain a useful soft margin.
\end{quote}

The new Lean module \texttt{WeightedFiberNeutralityObstruction.lean} blocks that repair whenever
the weighting itself respects the same involution symmetry.

\subsection{Lean theorem}

\begin{theorem}[Orbit-symmetric weights preserve zero fiber margin]
Let \(\alpha\) be a finite sample space, let \(\tau : \alpha \to \alpha\) be an involution, let
\(u : \alpha \to U\) be a retained feature map, let \(y : \alpha \to \{0,1\}\) be the target bit,
and let \(w : \alpha \to M\) be a weight into any additive commutative monoid.  Assume
\[
  u(\tau x) = u(x), \qquad y(\tau x) = 1-y(x), \qquad w(\tau x) = w(x)
\]
for every \(x \in \alpha\).  Then for every retained feature value \(v\),
\[
  \sum_{x : u(x)=v,\ y(x)=1} w(x)
  =
  \sum_{x : u(x)=v,\ y(x)=0} w(x).
\]
\end{theorem}

\noindent
The key formal statements are:
\begin{itemize}[leftmargin=1.5em]
  \item \texttt{weightedFeatureFiberTrueMass}: the weighted mass of the \texttt{true} part of one
  feature fiber;
  \item \texttt{weightedFeatureFiberFalseMass}: the weighted mass of the \texttt{false} part;
  \item \texttt{weightedFeatureFiberTrueMass\_eq\_weightedFeatureFiberFalseMass}: the exact balance
  theorem under involution-invariant weights.
\end{itemize}

\subsection{Why this closes the non-uniform averaging reply}

This matters because it removes a very common escape hatch.  One can no longer say:
\begin{quote}
The balanced finite orbit space was just a counting toy model.  The real proof uses a more
careful empirical or distributional weighting, so the invariant local statistic may still have a
nonzero expectation gap.
\end{quote}

Not if that weighting respects the same sign-flip involution.  The theorem shows that every such
weighting still gives exact zero conditional margin on every retained feature fiber.  In
particular, normalizing an involution-invariant finite weight into a probability distribution
still leaves
\[
  \mathbb{E}[\,y \mid u=v\,] = \tfrac12
\]
for every retained feature value \(v\) with nonzero mass.

\subsection{Consequence}

So the remaining escape routes are narrower again.  To revive a margin on the retained local
input, one must now do something stronger than:
\begin{enumerate}[leftmargin=1.5em]
  \item pass from hard labels to soft scores,
  \item pass from uniform counts to non-uniform but symmetry-respecting weights,
  \item or restrict attention to invariant-feature fibers alone.
\end{enumerate}

What survives must either break the involution symmetry in the weighting/domain, retain some
controlled non-invariant information, or prove a much stronger ensemble-specific structural
theorem.

\section{Obstruction IVd: Residual Symmetric Slices Still Stay At Chance}

The next natural reply is:
\begin{quote}
Fine.  We keep a little non-invariant information, so the exact sign-flip symmetry is broken.  We
do not need to resolve every orbit, only enough of them.
\end{quote}

This is the right next theorem to ask for.  The new Lean module
\texttt{ResidualSymmetryObstruction.lean} makes the burden precise: every residual slice on which
the retained information still fails to distinguish the involution partners remains exact-chance
under symmetry-respecting weights.

\subsection{Lean theorem}

\begin{theorem}[Any unresolved symmetric slice remains balanced]
Let \(\alpha\) be a finite sample space, let \(\tau : \alpha \to \alpha\) be an involution, let
\(p \subseteq \alpha\) be a \(\tau\)-stable slice, let \(u : \alpha \to U\) be the retained local
features, let \(y : \alpha \to \{0,1\}\) be the target bit, let \(w : \alpha \to M\) be a weight
into an additive commutative monoid, and let \(h : U \to \{0,1\}\) be any classifier.  Assume on
the slice \(p\) that
\[
  u(\tau x) = u(x), \qquad y(\tau x) = 1-y(x), \qquad w(\tau x) = w(x).
\]
Then the total weighted mass of correctly classified points in \(p\) equals the total weighted
mass of incorrectly classified points in \(p\).
\end{theorem}

\noindent
The key formal statements are:
\begin{itemize}[leftmargin=1.5em]
  \item \texttt{weightedCorrectMass}: total weighted correct mass for a feature-only classifier;
  \item \texttt{weightedIncorrectMass}: total weighted incorrect mass;
  \item \texttt{weightedCorrectMass\_eq\_weightedIncorrectMass}: the global weighted half-accuracy
  theorem under exact involution symmetry;
  \item \texttt{weightedCorrectMass\_eq\_weightedIncorrectMass\_on\_slice}: the residual-slice
  version obtained by restricting to a \(\tau\)-stable subtype.
\end{itemize}

\subsection{Why this hits the ``controlled non-invariant info'' reply}

This is the right quantitative obstruction.  A small amount of non-invariant information can only
help on the part of the domain where it actually breaks the orbit symmetry.  Every remaining
slice on which the feature map is still involution-invariant contributes only chance performance.

So the paper can no longer get by with a vague statement of the form:
\begin{quote}
We keep a little symmetry-breaking local information, and that should be enough to restore a
useful predictor.
\end{quote}

To make that route work, one now needs a real theorem that the unresolved symmetric slice carries
negligible weight, or else a theorem that the retained information breaks the involution on almost
all of the relevant mass.

\subsection{Consequence}

This does not yet show that \emph{every} controlled symmetry-breaking repair fails.  It does show
that such a repair only helps to the extent that it resolves most of the mass.  What remains must
therefore prove one of:
\begin{enumerate}[leftmargin=1.5em]
  \item the residual symmetric slice has negligible weight,
  \item the retained non-invariant information separates almost all involution pairs,
  \item or the actual evaluation domain is not well-modeled by a symmetry-respecting weighted
  slice at all.
\end{enumerate}

\section{Obstruction IVe: Advantage Is Bounded By Symmetry-Broken Mass}

The next creative rescue move is the quantitative version of the previous one:
\begin{quote}
Maybe we only break the symmetry on a relatively small region, but that is enough because the
classifier can exploit that region very efficiently and still gain a substantial global advantage.
\end{quote}

The new Lean module \texttt{AsymmetryBudgetObstruction.lean} blocks that move in the natural
finite weighted model.

\subsection{Lean theorem}

\begin{theorem}[Asymmetry budget bound]
Let \(\alpha\) be a finite sample space with weight \(w : \alpha \to \mathbb{N}\).  Let
\(p \subseteq \alpha\) be a \(\tau\)-stable slice on which the retained feature map is still
invariant under an involution \(\tau\), the target bit is flipped by \(\tau\), and the weights are
\(\tau\)-invariant.  Then for every classifier \(h\),
\[
  2 \cdot \mathrm{CorrectMass}(h)
  \;\le\;
  \mathrm{TotalMass} + \mathrm{OutsideMass}(p).
\]
Equivalently, the global excess over chance is bounded by the total mass outside the unresolved
symmetric slice.
\end{theorem}

\begin{theorem}[Strict advantage exposes outside support]
Under the same hypotheses, if
\[
  \mathrm{TotalMass}
  \;<\;
  2 \cdot \mathrm{CorrectMass}(h),
\]
then there is a concrete point \(x\) with \(w(x)>0\) and \(x\notin p\).  Conversely, if every
positive-weight point lies in \(p\), then strict advantage over half accuracy is impossible.
\end{theorem}

\noindent
The key formal statements are:
\begin{itemize}[leftmargin=1.5em]
  \item \texttt{two\_mul\_correctSliceMass\_eq\_sliceMass}: on the unresolved symmetric slice, the
  correct mass is exactly half of the slice mass;
  \item \texttt{weightedCorrectMass\_eq\_correctSlice\_add\_correctOutside}: total correct mass
  splits into the unresolved slice contribution plus the symmetry-broken contribution;
  \item \texttt{correctOutsideMass\_le\_outsideMass}: outside the unresolved slice, even perfect
  classification cannot use more mass than is actually there;
  \item \texttt{two\_mul\_weightedCorrectMass\_le\_total\_plus\_outside}: the final asymmetry-budget
  inequality;
  \item \texttt{not\_total\_lt\_two\_mul\_weightedCorrectMass\_of\_support\_subset}: if all
  positive-weight support remains inside the unresolved slice, strict advantage is blocked;
  \item \texttt{pos\_outsideMass\_of\_total\_lt\_two\_mul\_weightedCorrectMass}: strict advantage
  forces positive outside mass;
  \item \texttt{exists\_support\_outside\_of\_total\_lt\_two\_mul\_weightedCorrectMass}: strict
  advantage exposes a positive-weight sample outside the unresolved slice.
\end{itemize}

\subsection{Why this hits the last soft escape}

This is stronger than saying ``the residual slice is still at chance.''  It says the \emph{entire
global advantage} is quantitatively budgeted by the part of the domain where symmetry is actually
broken.

So one can no longer say:
\begin{quote}
We only need a little symmetry-breaking local information.  Perhaps it affects a small region, but
that small region could still be enough to drive the whole lower-bound argument.
\end{quote}

Not without a theorem that the symmetry-broken region itself carries substantial mass.  The Lean
bound shows that if the unresolved symmetric slice has mass close to the total, then the global
accuracy remains close to \(1/2\).  The support witness sharpens the audit: a strict global
advantage claim cannot remain purely aggregate, since it must identify at least one positive-mass
point where the repair has actually left the unresolved symmetric slice.

\subsection{Consequence}

At this point, the remaining route is no longer merely ``add some asymmetry.''  It must prove a
real mass-transfer statement:
\begin{enumerate}[leftmargin=1.5em]
  \item either almost all relevant mass leaves the involution-symmetric slice after the repair,
  \item or the weighted domain relevant to the proof is not symmetry-respecting in the sense used
  above,
  \item or the real predictor uses stronger ensemble-specific information than the current local
  symmetry analysis captures.
\end{enumerate}

\section{Obstruction IVf: Invariant Soft Scores Have Zero Signed Signal}

Another creative move is:
\begin{quote}
Perhaps the proof should not be phrased as conditional means on feature fibers at all.  Maybe one
can compute a richer soft score from the invariant local summary and then feed that score into a
more delicate downstream argument.
\end{quote}

This is a real strengthening of the repair attempt.  The new Lean module
\texttt{InvariantScoreObstruction.lean} blocks the core signal claim under the same involution
symmetry.

\subsection{Lean theorem}

\begin{theorem}[Invariant soft scores have zero signed correlation]
Let \(\alpha\) be a finite sample space, let \(\tau : \alpha \to \alpha\) be an involution, let
\(u : \alpha \to U\) be a retained feature map, let \(y : \alpha \to \{0,1\}\) be the target bit,
let \(w : \alpha \to \mathbb{N}\) be a \(\tau\)-invariant weight, and let
\(\mathrm{score} : U \to \mathbb{Z}\) be any soft score computed only from the retained features.
Assume
\[
  u(\tau x) = u(x), \qquad y(\tau x) = 1-y(x), \qquad w(\tau x) = w(x)
\]
for every \(x \in \alpha\).  Writing
\[
  \mathrm{sgnTarget}(y) =
  \begin{cases}
    1 & y=1, \\
    -1 & y=0,
  \end{cases}
\]
one has
\[
  \sum_{x \in \alpha} w(x)\,\mathrm{score}(u(x))\,\mathrm{sgnTarget}(y(x)) = 0.
\]
\end{theorem}

\noindent
The key formal statements are:
\begin{itemize}[leftmargin=1.5em]
  \item \texttt{targetSign}: the \(\{\pm 1\}\)-encoding of a Boolean target;
  \item \texttt{targetSign\_flip}: flipping the Boolean target negates that signed encoding;
  \item \texttt{weighted\_signedScore\_sum\_eq\_zero}: the weighted zero-correlation theorem for
  any invariant soft score.
\end{itemize}

\subsection{Why this hits the richer invariant-scoring repair}

This closes a subtler loophole than the earlier fiber-balance statements.  Those showed that the
conditional mean on each invariant feature value is exactly \(1/2\), and therefore that hard
classifiers or direct mean estimators on those features cannot beat chance.  The new theorem says
more: even before one thresholds or calibrates anything, \emph{the score itself} has zero signed
signal against the target.

So one can no longer say:
\begin{quote}
We are not really using a classifier on the invariant features.  We only need a good enough
numerical score from those features, and then a later nonlinear step can extract the useful
signal.
\end{quote}

Not under the same involution symmetry.  Any score depending only on the invariant retained
features is paired with its involution partner at equal weight and opposite signed target.  The
positive and negative contributions cancel exactly.

\subsection{Consequence}

Taken together with the previous sections, this leaves very little room for an invariant-feature
rescue.  The remaining route must now prove something outside the invariant scoring model:
\begin{enumerate}[leftmargin=1.5em]
  \item either the actual score retains genuinely non-invariant information,
  \item or the relevant domain/weighting is not symmetry-respecting,
  \item or the ensemble has a stronger special structure that invalidates the involution pairing
  hypotheses on the real predictor family.
\end{enumerate}

\section{Obstruction IVg: Invariant-Score Signal Comes Only From Weight Asymmetry}

The next and now very natural reply is:
\begin{quote}
Fine.  Perhaps the retained score inputs stay involution-invariant, but the real training or
evaluation weighting is not exactly symmetric.  Then an invariant soft score may still pick up a
useful signed signal.
\end{quote}

This is the last serious invariant-score escape hatch.  The new Lean module
\texttt{WeightAsymmetryObstruction.lean} makes that route exact: any invariant-score signal comes
entirely from the antisymmetric part of the weighting.

\subsection{Lean theorem}

\begin{theorem}[Only antisymmetric weight mass contributes to invariant-score signal]
Let \(\alpha\) be a finite sample space, let \(\tau : \alpha \to \alpha\) be an involution, let
\(u : \alpha \to U\) be a retained feature map, let \(y : \alpha \to \{0,1\}\) be the target bit,
let \(w : \alpha \to \mathbb{N}\) be an arbitrary weight, and let
\(\mathrm{score} : U \to \mathbb{Z}\) be any soft score on the retained features.  Assume
\[
  u(\tau x) = u(x), \qquad y(\tau x) = 1-y(x)
\]
for every \(x \in \alpha\).  Then
\[
  2 \sum_{x \in \alpha} w(x)\,\mathrm{score}(u(x))\,\mathrm{sgnTarget}(y(x))
  =
  \sum_{x \in \alpha} (w(x)-w(\tau x))\,\mathrm{score}(u(x))\,\mathrm{sgnTarget}(y(x)).
\]
\end{theorem}

\begin{theorem}[Nonzero invariant-score signal exposes asymmetric scored support]
Under the same hypotheses, suppose the invariant signed-score sum is nonzero:
\[
  \sum_{x \in \alpha} w(x)\,\mathrm{score}(u(x))\,\mathrm{sgnTarget}(y(x)) \neq 0.
\]
Then there is a concrete point \(x\) such that
\[
  w(x) \neq w(\tau x)
  \qquad\text{and}\qquad
  \mathrm{score}(u(x)) \neq 0.
\]
Conversely, if \(w(\tau x)=w(x)\) at every point with nonzero score, the signed-score sum is zero.
\end{theorem}

\noindent
The key formal statements are:
\begin{itemize}[leftmargin=1.5em]
  \item \texttt{signedScoreContribution}: one point's contribution to the invariant signed-score
  sum;
  \item \texttt{antisymmetricWeightContribution}: the corresponding contribution of the
  weight-difference term \(w(x)-w(\tau x)\);
  \item \texttt{two\_mul\_signedScore\_sum\_eq\_antisymmetricWeight\_sum}: the exact decomposition
  identity;
  \item \texttt{signedScore\_sum\_eq\_zero\_of\_score\_support\_weight\_symmetric}: zero signal
  when all nonzero-score points have symmetric weights;
  \item \texttt{exists\_weight\_asymmetric\_score\_point\_of\_signedScore\_sum\_ne\_zero}: nonzero
  signal exposes a nonzero-score point with asymmetric weight.
\end{itemize}

\subsection{Why this narrows the weighting rescue}

This is stronger than the earlier exact-symmetry obstruction.  It says not just that symmetric
weights give zero signal, but that \emph{all} invariant-score signal is sourced only by the
antisymmetric part of the weighting.  The involution-invariant part cancels exactly.

So one can no longer say:
\begin{quote}
The true weighting is probably a little asymmetric, and that may be enough.
\end{quote}

Not without quantifying that asymmetry.  The witness theorem rules out an even weaker repair where
the weights are asymmetric only on score-zero points: such asymmetry contributes nothing.  The proof
now needs a real bound of the form:
\begin{quote}
the antisymmetric weight budget is large enough on nonzero-score mass in the relevant domain to
support the claimed advantage.
\end{quote}

\subsection{Consequence}

This pushes the remaining route into a much narrower quantitative corner.  To save invariant local
scoring, one must now prove one of:
\begin{enumerate}[leftmargin=1.5em]
  \item the actual proof-relevant weighting has a substantial antisymmetric component on
  nonzero-score support,
  \item the score retains genuinely non-invariant information, or
  \item the real ensemble/predictor pair is not governed by the involution model above.
\end{enumerate}

\section{Obstruction IVh: Side Channels Help Only On Orbit-Resolving Mass}

At this point the natural pivot is:
\begin{quote}
Fine.  We stop asking invariant features alone to do the whole job.  We keep a controlled
involution-sensitive side channel \(v\), and classify on the pair \((u,v)\).
\end{quote}

This is the right next charitable target.  The new Lean module
\texttt{SideChannelResolutionObstruction.lean} packages the earlier symmetry budget results into the
exact theorem that this move needs to face.

\subsection{Lean theorem}

\begin{theorem}[Side-channel advantage is bounded by orbit-resolving mass]
Let \(\alpha\) be a finite sample space, let \(\tau : \alpha \to \alpha\) be an involution, let
\(u : \alpha \to U\) be a retained invariant feature map, let \(v : \alpha \to V\) be an arbitrary
side channel, let \(y : \alpha \to \{0,1\}\) be the target bit, let \(w : \alpha \to \mathbb{N}\)
be a \(\tau\)-invariant weight, and let \(h : U \times V \to \{0,1\}\) be any classifier.  Assume
\[
  u(\tau x) = u(x), \qquad y(\tau x) = 1-y(x), \qquad w(\tau x) = w(x)
\]
for every \(x \in \alpha\).  Define the \emph{resolved mass} to be the total weight of the points
for which the side channel separates the involution pair:
\[
  \mathrm{ResolvedMass}(v)
  =
  \sum_{x : v(\tau x) \neq v(x)} w(x).
\]
Then
\[
  2 \cdot \mathrm{CorrectMass}(h \circ (u,v))
  \;\le\;
  \mathrm{TotalMass} + \mathrm{ResolvedMass}(v).
\]
Equivalently, the entire global excess over chance is bounded by the weight of the domain where the
side channel actually resolves the orbit pairing.
\end{theorem}

\noindent
The key formal statements are:
\begin{itemize}[leftmargin=1.5em]
  \item \texttt{unresolvedBySideChannel}: the slice on which \(v(\tau x)=v(x)\);
  \item \texttt{resolvedMass}: the weight outside that unresolved slice;
  \item \texttt{two\_mul\_weightedCorrectMass\_pair\_le\_total\_plus\_resolvedMass}: the final
  resolution-budget inequality for classifiers on \((u,v)\).
\end{itemize}

\subsection{Why this hits the controlled-side-channel rescue}

This is stronger than the earlier residual-symmetry prose, because it says exactly what the side
channel must buy.  Keeping some involution-sensitive information does \emph{not} by itself create a
large advantage.  The benefit is quantitatively capped by the mass where that side channel really
separates involution partners.

So one can no longer say:
\begin{quote}
We only need a small side channel.  Once it is present, a clever classifier on \((u,v)\) may still
recover a substantial domination advantage.
\end{quote}

Not without a theorem that the side channel resolves involution pairs on a substantial fraction of
the relevant mass.  The unresolved slice remains chance, and the total excess over chance is budgeted
by what lies outside it.

\subsection{Consequence}

This leaves the controlled-side-channel route in a much sharper form.  To save it, one must now
prove one of:
\begin{enumerate}[leftmargin=1.5em]
  \item the retained side channel separates involution pairs on most of the proof-relevant mass,
  \item the weighting itself has a large antisymmetric component that can be exploited,
  \item or the real predictor is not well-modeled by this involution-paired finite setup.
\end{enumerate}

\section{Obstruction IVi: Any Domination Gap Forces Matching Resolution}

The previous section bounded success by resolved mass.  The converse reading is now just as
important:
\begin{quote}
if the proof claims a real domination gap over chance from a retained side channel, that claim
already commits it to a quantitatively comparable orbit-resolution theorem.
\end{quote}

The new Lean module \texttt{ResolutionDemandObstruction.lean} packages that commitment as a direct
corollary.

\subsection{Lean theorem}

\begin{theorem}[Doubled advantage is bounded by resolved mass]
For a weighted classifier \(h\), define the \emph{doubled advantage}
\[
  \mathrm{DoubledAdv}(h)
  =
  2 \cdot \mathrm{CorrectMass}(h) - \mathrm{TotalMass}.
\]
Under the same involution hypotheses as in Obstruction IVh, one has
\[
  \mathrm{DoubledAdv}(h \circ (u,v))
  \le
  \mathrm{ResolvedMass}(v).
\]
\end{theorem}

\noindent
The key formal statements are:
\begin{itemize}[leftmargin=1.5em]
  \item \texttt{doubledAdvantage}: twice the weighted success gap over chance;
  \item \texttt{doubledAdvantage\_pair\_le\_resolvedMass}: the final converse bound;
  \item \texttt{resolvedMass\_ge\_of\_le\_doubledAdvantage\_pair}: a claimed doubled gap
  \(\delta\) forces at least \(\delta\) units of orbit-resolving mass;
  \item \texttt{resolvedMass\_ge\_of\_total\_plus\_delta\_le\_two\_mul\_weightedCorrectMass\_pair}:
  the same burden stated directly at the half-accuracy surface;
  \item \texttt{not\_pos\_doubledAdvantage\_pair\_of\_resolvedMass\_eq\_zero}: zero
  orbit-resolving side-channel mass blocks every positive doubled advantage;
  \item \texttt{resolvedMass\_pos\_of\_pos\_doubledAdvantage\_pair}: any positive
  doubled advantage witnesses positive orbit-resolving mass;
  \item \texttt{resolvedMass\_pos\_of\_total\_lt\_two\_mul\_weightedCorrectMass\_pair}:
  strict success above half accuracy witnesses positive orbit-resolving mass
  without passing through the subtraction-defined gap;
  \item \texttt{resolvedMass\_eq\_zero\_of\_supportwise\_unresolved}: if \(v(\tau x)=v(x)\)
  at every positive-weight point, the resolved mass is zero;
  \item \texttt{exists\_resolving\_point\_of\_pos\_resolvedMass}: positive resolved mass
  exposes a positive-weight point \(x\) with \(v(\tau x)\ne v(x)\);
  \item \texttt{exists\_resolving\_point\_of\_pos\_doubledAdvantage\_pair}: any positive
  doubled advantage exposes such a point directly;
  \item \texttt{exists\_resolving\_point\_of\_total\_lt\_two\_mul\_weightedCorrectMass\_pair}:
  any strict half-accuracy advantage exposes such a point directly.
\end{itemize}

The newest deterministic witness form is important for the residual-side-information repair.
It is no longer enough to prove only an aggregate side-channel bound.  A claimed positive
advantage must identify a proof-relevant support point where the residual information really changes
under the involution.  If every positive-weight point is supportwise unresolved, Lean now derives the
zero-resolved-mass blocker directly.

The same demand is now specialized to the manuscript's exact post-switch fork.  The theorem
\[
  \texttt{nonzeroColumnMass\_ge\_of\_le\_doubledAdvantage\_invariantProjection\_with\_b}
\]
says that any claimed doubled advantage from classifying on \(((z,a_i),b)\) forces matching mass on
nonzero \(a_i\)-columns.  Its zero-mass corollary
\[
  \texttt{not\_pos\_doubledAdvantage\_invariantProjection\_with\_b\_of\_zero\_nonzeroColumnMass}
\]
blocks the residual \(b\)-side-channel bridge whenever the proof-relevant weighted support contains
no nonzero \(a_i\)-column mass.
The support-witness form is now explicit:
\[
  \texttt{exists\_b\_resolving\_input\_of\_pos\_doubledAdvantage\_invariantProjection\_with\_b},
\]
\[
  \texttt{exists\_nonzeroColumn\_of\_pos\_doubledAdvantage\_invariantProjection\_with\_b}.
\]
Thus a positive \(b\)-side-channel claim must produce a concrete positive-weight post-switch input
whose \(b\)-coordinate changes under \(T_i\), equivalently a positive-weight input on a nonzero
\(a_i\)-column.

The latest refinement closes the immediate escape hatch under the manuscript's exact local-input
invariance premise.  Lean defines
\[
  \texttt{exactInputInvariantWeightedSupport}\ w
\]
to mean that every positive-weight post-switch input is fixed by the full-input involution
\(T_i\), not merely by the reduced projection.  Then
\[
  \texttt{nonzeroColumnMass\_eq\_zero\_of\_exactInputInvariantWeightedSupport}
\]
proves that this premise forces zero mass on nonzero \(a_i\)-columns, because every nonzero
\(a_i\)-column changes the full input \((z,a_i,b)\) under \(b \mapsto b \oplus a_i\).  Conversely,
\[
  \texttt{not\_exactInputInvariantWeightedSupport\_of\_pos\_nonzeroColumnMass}
\]
proves that supplying positive nonzero-column mass contradicts exact-input-invariant weighted
support.  Combining this with the advantage bound gives
\[
  \texttt{not\_pos\_doubledAdvantage\_and\_exactInputInvariantWeightedSupport}.
\]
Thus this repair is no longer left at ``prove there is enough nonzero-column mass.''  Under the
same exact-input invariance premise needed by the local \(T_i\)-symmetry argument, positive
\(b\)-side-channel advantage and the required support condition are formally inconsistent.

The randomized-residual escape now has the same quantitative boundary.  The module
\texttt{RandomizedResidualObstruction.lean} models a randomized residual as a deterministic side
channel on the product of the original finite sample space with a finite coin space.  The theorem
\[
  \texttt{doubledAdvantage\_randomizedResidual\_le\_resolvedMass}
\]
proves that randomized-residual advantage is bounded by the joint mass where the coin-indexed
residual separates involution partners, and
\[
  \texttt{randomizedResidualResolvedMass\_ge\_of\_le\_doubledAdvantage}
\]
turns any claimed doubled gap \(\delta\) into at least \(\delta\) units of joint resolving mass.
Two corollaries pin the non-escape cases:
\[
  \texttt{not\_pos\_doubledAdvantage\_randomizedResidual\_of\_coinwise\_unresolved}
\]
blocks coinwise-unresolved randomized residuals, while
\[
  \texttt{not\_pos\_doubledAdvantage\_randomizedResidual\_of\_fixed\_source\_support}
\]
blocks arbitrary randomized residuals when all positive source weight is already fixed by the
involution.  The support-sensitive strengthening is
\[
  \texttt{not\_pos\_doubledAdvantage\_randomizedResidual\_of\_supportwise\_unresolved}.
\]
It only requires unresolvedness on source/coin pairs with positive joint weight.  Its contrapositive
form is now also theorem-level:
\[
  \texttt{exists\_resolving\_coin\_of\_pos\_doubledAdvantage\_randomizedResidual}.
\]
Thus randomization is not a free repair of the residual side-information gap; positive finite-coin
advantage must exhibit a real positive-weight source/coin pair where the residual changes across
the involution.
The newest strengthening removes the last possible ``only in the mixture''
reading of that witness.  From positive randomized resolving mass Lean now
extracts a deterministic positive-weight coin slice with positive ordinary
side-channel resolving mass:
\[
  \texttt{exists\_positive\_coin\_resolvedMass\_of\_pos\_randomizedResidualResolvedMass}.
\]
From positive randomized advantage it gets the same conclusion directly:
\[
  \texttt{exists\_positive\_coin\_resolvedMass\_of\_pos\_doubledAdvantage\_randomizedResidual}.
\]
So a finite-coin randomized residual repair cannot defer all symmetry breaking
to an aggregate average over coins.  One actual coin value of positive
coin-weight already carries a standard deterministic side-channel obstruction.
The new strengthening shows that Lean can package that deterministic slice all
the way up to the ordinary residual-side-information obstruction packet:
\[
  \texttt{exists\_positive\_coin\_obstruction\_package\_of\_pos\_doubledAdvantage\_randomizedResidual}
\]
and
\[
  \texttt{exists\_positive\_coin\_obstruction\_package\_of\_total\_lt\_two\_mul\_weightedCorrectMass\_randomizedResidual}.
\]
So even before specializing to the manuscript's exact post-switch surface, the
finite-coin randomized-residual route no longer stands apart from the ordinary
residual-side lane: strict half-accuracy advantage already forces a
positive-weight deterministic coin slice carrying the same non-determination,
collision, resolving-point, and source-only-failure witnesses.

The post-switch specialization now removes that remaining randomized repair under the exact
support premise already used above.  The module
\texttt{PostSwitchRandomizedResidualObstruction.lean} proves
\[
  \texttt{randomizedResidualResolvedMass\_tiInputMap\_eq\_zero\_of\_exactInputInvariantWeightedSupport}.
\]
So exact-input-invariant weighted support on \((z,a_i,b)\) leaves zero joint resolving mass for
every finite-coin randomized residual.  Equivalently,
\[
  \texttt{not\_exactInputInvariantWeightedSupport\_of\_pos\_randomizedResidualResolvedMass\_tiInputMap}
\]
says that supplying positive randomized joint resolution contradicts the support premise.  The
combined theorem
\[
  \texttt{not\_pos\_doubledAdvantage\_randomizedResidual\_and\_exactInputInvariantWeightedSupport}
\]
blocks positive finite-coin randomized residual advantage on the exact post-switch surface whenever
the route keeps exact local-input invariance of its positive-weight support.
The supportwise post-switch version,
\[
  \texttt{not\_pos\_doubledAdvantage\_randomizedResidual\_invariantProjection\_of\_supportwise\_unresolved},
\]
removes the need to assume fixed support globally: it is enough that every positive-weight
input/coin pair has the same residual value after the local-input switch.  Conversely,
\[
  \texttt{exists\_resolving\_coin\_of\_pos\_doubledAdvantage\_randomizedResidual\_invariantProjection}
\]
says that any positive finite-coin randomized-residual advantage on the exact post-switch surface
must expose a positive-weight input/coin pair resolving that switch.
The deterministic-slice form is now specialized too:
\[
  \texttt{exists\_positive\_coin\_resolvedMass\_of\_pos\_doubledAdvantage\_randomizedResidual\_invariantProjection}.
\]
Thus on the exact post-switch surface, any positive finite-coin randomized
advantage must identify a positive-weight coin whose ordinary residual side
channel has positive mass resolving the local-input switch.
The same boundary now holds on the strict half-accuracy surface.  Lean proves
\[
  \texttt{not\_total\_lt\_two\_mul\_weightedCorrectMass\_randomizedResidual\_invariantProjection\_of\_exactInputInvariantWeightedSupport}
\]
and
\[
  \texttt{not\_total\_lt\_two\_mul\_weightedCorrectMass\_randomizedResidual\_and\_exactInputInvariantWeightedSupport},
\]
so exact-input-invariant support blocks not only positive doubled advantage but
also any claim that a finite-coin randomized residual beats half accuracy on
the exact post-switch surface.  Conversely,
\[
  \texttt{exists\_positive\_coin\_resolvedMass\_of\_total\_lt\_two\_mul\_weightedCorrectMass\_randomizedResidual\_invariantProjection}
\]
and
\[
  \texttt{not\_exactInputInvariantWeightedSupport\_of\_total\_lt\_two\_mul\_weightedCorrectMass\_randomizedResidual\_invariantProjection}
\]
show that even this weaker-looking route already exposes a positive-weight
deterministic coin slice with ordinary resolving mass and therefore directly
refutes the manuscript's exact-support premise.  The collapse is now stronger:
\[
  \texttt{exists\_positive\_coin\_obstruction\_package\_of\_total\_lt\_two\_mul\_weightedCorrectMass\_randomizedResidual\_invariantProjection}
\]
shows that the same strict-half post-switch witness already carries the full
ordinary invariant-projection residual obstruction package, not merely positive
resolving mass.  So the post-switch randomized-residual repair no longer
survives by retreating from positive doubled advantage to the strict
half-accuracy formulation.

\texttt{CruxSynthesis.lean} now records this as the narrow
\texttt{PNPRandomizedResidualSubrepairClass} subledger, backed by the theorem
\texttt{V13RandomizedResidualSubcoverage}.  The synthesis layer no longer
leaves these finite-coin witnesses only in the local obstruction files: it
explicitly records the ordinary source/coin product-space budget, the
supportwise-unresolved zero-mass and no-advantage blockers, the deterministic
positive-weight coin-slice witness, and the exact post-switch analogues.
This remains deliberately narrow.  The broad
\texttt{PNPRepairClass.randomizedResidual} entry is still uncovered in the
current packet; the new synthesis theorem only says that any finite-coin
randomized-residual repair must surface one of those explicit witnesses or
fail the present Lean boundary.

\subsection{Why this matters}

This is the strongest current version of the side-channel blocker.  It no longer says only that
resolved mass is \emph{sufficient} for advantage.  It says a claimed advantage is itself evidence
that substantial resolution must already be present.

So one can no longer say:
\begin{quote}
Perhaps the predictor gets a modest but important edge from a very small side channel, and only
later amplifies that edge through clever structure.
\end{quote}

Not unless that modest edge is matched by actual orbit-resolving mass.  The theorem makes the
burden explicit: any nontrivial domination statement already entails a nontrivial local resolution
statement.

\subsection{Consequence}

This is close to the form needed to challenge the route on its own terms.  To save the current
style of proof, one now needs a theorem of the form:
\begin{quote}
the switched predictor's retained side channel, deterministic or randomized, resolves involution
pairs on enough relevant mass to support the claimed lower-bound-driving advantage.
\end{quote}

Without that, the claimed domination gap has no quantitative room to exist.

\section{Obstruction V: Short Programs Still Cover the Full Log-Width Local Rule Class}

Another natural reply is:
\begin{quote}
The switched predictor comes from a short global program, so even if arbitrary local rules are
huge, only a much smaller subclass can actually arise from polynomial-time decoders.
\end{quote}

This is not true without an additional restriction beyond mere global shortness.

\subsection{Lean theorem}

The new module \texttt{ShortProgramObstruction.lean} formalizes the basic truth-table fact.

\begin{theorem}[Truth tables realize all local rules]
For every \(n\), there is a surjective decoder
\[
  \{0,1\}^{2^n} \twoheadrightarrow \{\,f : \{0,1\}^n \to \{0,1\}\,\}.
\]
Equivalently, every Boolean rule on \(n\) visible bits can be hard-wired by a truth table of
length \(2^n\).
\end{theorem}

\noindent
The key formal statements are:
\begin{itemize}[leftmargin=1.5em]
  \item \texttt{truthTableDecode\_surjective}: every local rule is realized by a
  \(2^n\)-bit truth table;
  \item \texttt{truthTableFamily\_realizes\_all}: the realized class of the truth-table family is
  the full local-rule space;
  \item \texttt{card\_realized\_truthTableFamily}: this realized class has exactly cardinality
  \(2^{2^n}\).
\end{itemize}

\subsection{Binary-log consequence}

At binary-log width the truth-table size is already linear in the ambient parameter:
\[
  2^{\lfloor \log_2 m \rfloor + 1} \le 2m \qquad (m \ne 0).
\]
This is formalized as
\texttt{truthTableBits\_binaryLogWidth\_le\_twoMul}.

So if the switched interface really has only \(n = \lfloor \log_2 m \rfloor + 1\) visible bits,
then a merely linear-size code budget already suffices to realize the \emph{entire} local-rule
class on that interface.

\subsection{Consequence}

This directly blocks the short-program rescue in its naive form:

\begin{quote}
``The global decoder is polynomial-size'' does not imply
``the induced local rule lies in a tiny subclass.''
\end{quote}

On the contrary, polynomial-size descriptions are already large enough to hard-wire every Boolean
rule on \(O(\log m)\) visible bits.  So if one wants a true small subclass theorem, it must use
something much stronger than global shortness alone: symmetry, algebraic normal form, a tree
interpreter with bounded message alphabet, or some other genuine structural restriction.

\section{Obstruction VI: Exact Metagraph Quotients Have Interface-State Explosion}

The strongest remaining charitable rescue is roughly this:
\begin{quote}
The switched neighborhood first collapses to a canonical metagraph or message-passing quotient,
and the witness bit is computed by a fixed interpreter on that quotient.
\end{quote}

This is a serious idea.  It is also constrained by a generic exactness barrier.

\subsection{Lean theorem}

The module \texttt{MetagraphObstruction.lean} models a quotient state by its exact interface
behavior on \(k\) visible boundary bits.  The semantic object is then a boundary transducer
\[
  (\{0,1\}^k \to \{0,1\}^k).
\]

\begin{theorem}[Boundary transducers are exponentially numerous]
There are exactly
\[
  2^{k 2^k}
\]
boundary transducers on \(k\) visible bits.
\end{theorem}

\noindent
The key formal statements are:
\begin{itemize}[leftmargin=1.5em]
  \item \texttt{card\_boundaryMap}: exact \(k\)-bit interface semantics have cardinality
  \(2^{k2^k}\);
  \item \texttt{no\_surjective\_bitCode\_to\_boundaryMap\_of\_lt}: no \(s\)-bit code can cover the
  full interface-semantic space when \(s < k2^k\);
  \item \texttt{boundaryMap\_card\_le\_of\_surjective\_interpret}: any finite state space whose
  interpreter covers all exact boundary transducers has at least \(2^{k2^k}\) states;
  \item \texttt{not\_surjective\_quotientCode\_boundaryMap\_of\_stateCard\_lt}: factoring an
  arbitrary code through a smaller quotient state space cannot repair this obstruction;
  \item \texttt{boundaryMap\_stateCard\_lowerBound\_binaryLogWidth}: at
  \(k = \lfloor \log_2 m \rfloor + 1\), exact interface coverage forces at least \(2^m\) quotient
  states;
  \item \texttt{no\_exact\_metagraphQuotient\_below\_inputSize\_binaryLogWidth}: at
  \(k = \lfloor \log_2 m \rfloor + 1\), there is no exact quotient with fewer than \(m\) code bits.
\end{itemize}

The newest extension closes a small but important loophole in the first version of this
obstruction.  One could try to let a large code first choose a finite metagraph state, and then use
a fixed interpreter from states to exact boundary behavior.  Surjectivity of the composite already
makes the interpreter surjective, so the state space itself must have the full boundary-map
cardinality.  At binary-log boundary width, any exact universal quotient-state repair therefore
requires at least \(2^m\) states, even before accounting for how those states are encoded.

\subsection{Why this matters}

If a metagraph/message quotient is meant to be \emph{exact}, then it must preserve enough state to
reconstruct the gadget's interface behavior under gluing.  The theorem above says that, in the
absence of an additional structural restriction, the exact interface-semantic space is already far
too large for a tiny finite-state algebra.

This does not yet show that the current ensemble realizes the full transducer class.  What it
does show is that a small exact quotient theorem cannot follow from tree-likeness or locality
alone.  To survive, the metagraph route needs a new theorem that the actual switched/VV gadgets
lie in a drastically smaller semantic subclass.

\subsection{Consequence}

So the metagraph rescue is now narrowed to something like:
\begin{quote}
not ``there is an exact small quotient of all local interfaces,''
but ``the specific random masked+isolated interfaces collapse to a much smaller algebra for
ensemble-specific reasons.''
\end{quote}

That is a legitimate possible theorem.  It is also substantially stronger than anything currently
written in the paper, and it is the next point where a decisive break may occur.

\section{Obstruction VII: Tree-Likeness Alone Does Not Yield a Small Message Class}

Another natural repair is:
\begin{quote}
On locally tree-like neighborhoods, the witness bit is computed by a small message-passing or
dynamic-program interpreter on the tree.
\end{quote}

This too needs a stronger premise than mere tree structure.

\subsection{Lean theorem}

The new module \texttt{TreeMessageObstruction.lean} defines a fixed perfect binary tree
architecture of depth \(n\), where internal nodes merely branch on the next input bit and only the
leaves carry data.

\begin{theorem}[Perfect binary trees realize all local rules]
For every Boolean rule
\[
  f : \{0,1\}^n \to \{0,1\},
\]
there is a depth-\(n\) perfect binary decision tree whose evaluation equals \(f\).
\end{theorem}

\noindent
The key formal statements are:
\begin{itemize}[leftmargin=1.5em]
  \item \texttt{eval\_encodeDecisionTree}: every local rule is exactly realized by a leaf-labeled
  depth-\(n\) perfect tree;
  \item \texttt{evalDecisionTree\_injective} and \texttt{evalDecisionTree\_bijective}: for this fixed
  perfect architecture, extensional equality of semantics is exactly equality of the leaf labels;
  \item \texttt{evalDecisionTree\_surjective}: the semantics map from such trees onto
  \texttt{LocalRule n} is surjective;
  \item \texttt{no\_surjective\_bitCode\_to\_decisionTreeSemantics\_of\_lt}: therefore no
  \(s\)-bit exact code can cover all depth-\(n\) tree semantics when \(s < 2^n\);
  \item \texttt{localRule\_card\_le\_of\_surjective\_treeMessageInterpreter}: any finite
  tree-message state space whose interpreter covers all exact depth-\(n\) semantics has at least
  \(2^{2^n}\) states;
  \item \texttt{not\_surjective\_quotientCode\_treeMessageInterpreter\_of\_stateCard\_lt}: factoring
  a code through a smaller finite tree-message state cannot repair the obstruction;
  \item \texttt{treeMessage\_stateCard\_lowerBound\_binaryLogDepth}: at depth
  \(n=\lfloor\log_2 m\rfloor+1\), exact coverage forces at least \(2^m\) message states.
\end{itemize}

\subsection{Why this matters}

This is a direct answer to the naive tree-message rescue:

\begin{quote}
tree-likeness by itself does \emph{not} imply semantic simplicity.
\end{quote}

Even with an extremely rigid tree architecture, the semantic class is already the full Boolean
function space on the visible inputs.  The newest Lean strengthening also removes the ``many trees
may collapse to few semantics'' escape for this architecture: evaluation is injective, so the
semantics are exactly the leaf truth table.  If a repair instead proposes a finite tree-message
state space plus a fixed interpreter to exact depth-\(n\) trees, surjectivity onto all exact
semantics forces that state space to have the full local-rule cardinality.  So if one wants a
genuine tree-message theorem, it must use special properties of the actual SAT+VV messages, not
just the fact that the underlying graph is a tree.

\subsection{Consequence}

The surviving charitable tree route is now much narrower:
\begin{quote}
not ``all tree-like local neighborhoods admit a tiny interpreter,''
but ``the specific masked+isolated tree messages belong to a sharply smaller algebra.''
\end{quote}

That is again a real possible theorem.  It is also much stronger than the current text, and it is
precisely the kind of statement that now has to be made explicit to keep this proof approach alive.

\section{What These Results Do Not Yet Prove}

The current Lean modules do \emph{not} prove that the full argument fails.  In particular,
they do not yet address:
\begin{enumerate}[leftmargin=1.5em]
  \item whether the masking and isolation ensemble enforces extra algebraic structure on local rules;
  \item whether the switched predictors might collapse to a subclass analogous to low-depth message
  passing or tree computations;
  \item whether the neutrality and sparsification lemmas are themselves correct in the precise form
  used in the paper;
  \item whether a different, stronger switching theorem could bypass the present obstruction.
\end{enumerate}

So the current status is best described as:
\begin{quote}
The proof attempt is not yet broken in full, but one of its central complexity-compression
transitions has not been justified, and the default combinatorics point in the opposite direction.
\end{quote}

\section{Stop-Grade Synthesis Ledger}

The newest meta-synthesis module, \texttt{CruxSynthesis.lean}, makes the
status discipline formal.  It defines a coarse type of live repair classes and
represents a crux packet by the repair classes it covers.  A packet is
\emph{stop-grade} only if it covers every named repair class in that ledger:
\[
  \texttt{PNPCruxPacket.StopGrade packet}
  \quad\equiv\quad
  \forall r:\texttt{PNPRepairClass},\ \texttt{packet.covers r}.
\]

The current accumulated PNP packet covers the deterministic same-source
conditioning repairs, the v13 repeated-pivot, forced-step, and bad-cell fixed-field repairs,
field refinement of a bad cell, invariant-score repairs, bounded side-channel
repairs, short-global-decoder repairs, CD evidence normalization as an exact
residual-atom boundary, trace capture as the full finite
trace-fiber/cardinality/image-budget boundary, the atomic evidence budget as
the full finite exactness boundary, the
same-source finite-count approximation subcase, and the abstract
clocked finite-uniform kernel behind the later \(Kpoly\) counting step.  Lean
records this exact covered list as \texttt{currentPNPCoveredRepairClasses}.

The v13 entries in this ledger are now tied directly to the v13 route modules
rather than recorded as bare Boolean placeholders.  \texttt{CruxSynthesis.lean}
imports \texttt{PNPv13RouteSynthesis.lean}, \texttt{PNPv13ForcedStepObstruction.lean}, and
\texttt{PNPv13FieldRefinement.lean}, defines the coverage propositions
\[
\begin{gathered}
  \texttt{V13RepeatedPositiveFieldedPivotCoverage},\quad
  \texttt{V13ForcedPositiveFieldedStepCoverage},\quad
  \texttt{V13FieldedBadCellCoverage},\\
  \texttt{V13FieldRefinementBadCellCoverage},\quad
  \texttt{V13EvidenceNormalizationCoverage},
  \quad \texttt{V13TraceCaptureFactorizationCoverage},\\
  \texttt{V13AtomicEvidenceBudgetCoverage},
\end{gathered}
\]
and proves
\[
  \texttt{currentPNPLocalCruxPacket\_theoremBacked\_v13\_status}.
\]
That theorem says exactly that those seven v13 ledger entries are covered while
the current packet is still not stop-grade.  This is a route-interface claim:
it does not assert that a SAT/Cook--Levin construction, or the manuscript's
actual switched family, has been formalized inside the ledger.

Inside that synthesis layer, the atomic-budget entry is no longer just the
original one-sided no-double-counting statement.  The proposition
\texttt{V13AtomicEvidenceBudgetCoverage} now packages the dual under-budget
omission witness, the exactness-side conjunction
\texttt{ExactnessSideConditions}, the equality-surface reductions from exact
budget balance to either structural side condition, and the cancellation anchor
showing that exact scalar equality plus
\(
  \texttt{NoDoubleCounting} \leftrightarrow
  \texttt{AllPositiveCostAtomsCharged}
\)
still does not repair the budget step.

The trace-capture entry in that same synthesis layer has now been promoted in
parallel.  \texttt{V13TraceCaptureFactorizationCoverage} no longer stops at
deterministic/randomized fiber-constancy equivalences.  It also packages the
injectivity criterion
\(
  \texttt{CapturesAllBoolObservers}
  \leftrightarrow
  \texttt{Function.Injective}\ \texttt{traceOnRelevant}
\),
the finite trace-alphabet lower bound
\(
  |\texttt{relevant states}| \le |\texttt{Trace}|
\),
and the Boolean decoder-space lower bound
\(
  2^{|\texttt{relevant states}|} \le 2^{|\texttt{Trace}|}
\).
So the synthesis packet now carries not only the pointwise collapsed-fiber
conflict but also the finite cardinality and finite image-budget obstructions
for the same route class.

The refinement route itself has now been strengthened one layer further.
\texttt{PNPv13FieldRefinementRouteObstruction.lean} lifts the earlier
prefix-level refinement theorem to the whole fielded switching list: if a
coarse field already has a bad cell after some explicit fielded prefix, then
replacing that step's field by any genuine refinement still leaves a localized
bad cell at the same position.  Lean then blocks both
\texttt{V13FieldSwitchingInstantiatedFrom} and
\texttt{V13FieldSwitchingFailureMatchingFrom} for the refined suffix by routing
through \texttt{V13FieldedBadCellFrom}.  The Boolean-square regression fixes a
concrete two-step example where a one-cell second cut is refined to the
second-coordinate field and the whole route still fails.

\texttt{PNPv13EvidenceNormalization.lean} supplies the fourth entry.  It
abstracts a CD evidence atom by four predicates: target relevance, neutrality,
safe-buffer absorption, and hidden-gauge absorption.  Lean proves
\[
  \texttt{normalizesNonNeutral\_iff\_not\_hasResidualAtom},
\]
so the manuscript's evidence-normalization theorem is exactly the claim that
there is no target-relevant non-neutral atom outside both safe-buffer and
hidden-gauge classes.  The positive-score weakening is also exact:
\[
  \texttt{positiveAtomsCovered\_iff\_not\_positiveResidualScore}.
\]
The three-atom toy model
\texttt{v13ToyEvidenceSurface} has one safe atom, one gauge atom, and one
positive-score residual atom.  The synthesis anchor
\texttt{v13EvidenceNormalizationCoverage\_anchor\_toyResidual\_blocks}
therefore records the current crux without overstating it: any repair must
prove that the actual CD construction has no such residual positive atom, not
merely rename residual evidence as safe or gauge.

\texttt{PNPv13TraceCapture.lean} supplies the next v13 ledger boundary.  It
abstracts a trace map on target-relevant states and proves
\[
  \texttt{observerTraceDecodable\_iff\_constantOnTraceFibers}.
\]
Thus replacing an adaptive observer by a finite CD trace is exactly a quotient
factorization theorem: the observer's output must be constant on every
target-relevant trace fiber.  The file also proves the all-Boolean-observer
endpoint
\[
  \texttt{capturesAllBoolObservers\_iff\_not\_relevantTraceCollision}.
\]
So a trace that collapses two target-relevant states cannot capture every
Boolean downstream observer.  The toy surface
\texttt{v13CollapsedTraceSurface} collapses two relevant states to one trace
while \texttt{v13ToyTraceObserver} separates them, and the synthesis anchor
\texttt{v13TraceCaptureFactorizationCoverage\_anchor\_collapsedTrace\_blocks}
records the corresponding non-decodability and universal-Boolean-capture
failure.  This does not prove that the manuscript's actual CD traces collide;
it pins the exact theorem any trace-capture repair must supply.

The next strengthening turns that qualitative boundary into a finite-alphabet
lower bound.  The module
\texttt{PNPv13TraceCaptureCardinalityObstruction.lean} defines the restricted
map
\[
  \texttt{traceOnRelevant} : \{s \mid \texttt{targetRelevant}\ s\} \to \texttt{Trace}
\]
and proves
\[
  \texttt{capturesAllBoolObservers\_iff\_traceOnRelevant\_injective}.
\]
Thus universal Boolean trace capture is not merely absence of one explicit
collision witness; it is injectivity of the proposed trace alphabet on the
entire target-relevant support.  Consequently any finite trace alphabet must
satisfy
\[
  |\{s \mid \texttt{targetRelevant}\ s\}| \le |\texttt{Trace}|,
\]
formalized as
\[
  \texttt{relevantStateCard\_le\_traceCard\_of\_capturesAllBoolObservers}
\]
and contrapositively
\[
  \texttt{not\_capturesAllBoolObservers\_of\_traceCard\_lt\_relevantStateCard}.
\]
The regression file records both a positive canary and finite blockers:
identity traces on all relevant Boolean states do capture every Boolean
observer, but one trace for two relevant states and two traces for four
relevant states are already impossible.  So a manuscript repair that proposes
finite CD traces must now prove a target-relevant injective coding theorem, not
just an abstract factorization slogan.

The next file places this same obstruction on the standard finite-cover route
used elsewhere in the project.  The module
\texttt{PNPv13TraceCaptureImageBudgetObstruction.lean} defines
\texttt{traceDecoderFamily}, whose indices are Boolean decoders
\(\texttt{Trace} \to \texttt{Bool}\) and whose inputs are the target-relevant
states themselves, represented as the subtype
\(\{s \mid \texttt{targetRelevant}\ s\}\).  Universal Boolean trace capture
then makes this indexed family surjective onto the full Boolean classifier
space on the relevant support, via
\[
  \texttt{traceDecoderFamily\_surjective\_of\_capturesAllBoolObservers}.
\]
But the family of all trace decoders has size only \(2^{|\texttt{Trace}|}\), so
it carries an explicit finite predictor-image cover of that same size.  Feeding
this into the project's standard finite-image lower bound yields
\[
  2^{|\{s \mid \texttt{targetRelevant}\ s\}|} \le 2^{|\texttt{Trace}|},
\]
formalized as
\[
  \texttt{boolClassifierCard\_le\_traceClassifierCard\_of\_capturesAllBoolObservers}
\]
and contrapositively
\[
  \texttt{not\_capturesAllBoolObservers\_of\_traceClassifierCard\_lt}.
\]
So the manuscript's trace-capture route is now blocked not only pointwise by
an injectivity obligation, but also globally in the same finite predictor-image
language that the later `Kpoly' bridge already uses.  The regression file
records the positive identity-trace canary and the same two toy blockers in the
image-budget form: one trace cannot realize all four Boolean classifiers on two
relevant states, and two traces cannot realize all sixteen Boolean classifiers
on four relevant states.

\texttt{PNPv13RandomCoinTraceCapture.lean} tightens the same trace-capture
entry at the manuscript's random-coin reduction point.  A coin-indexed observer
\(A(s,c)\) is captured by a trace/coin decoder exactly when every deterministic
coin slice \(s\mapsto A(s,c)\) is trace-decodable:
\[
  \texttt{randomObserverTraceDecodable\_iff\_forall\_coin\_observerTraceDecodable}.
\]
With an inhabited output type this is equivalently coinwise constancy on
target-relevant trace fibers, via
\[
  \texttt{randomObserverTraceDecodable\_iff\_constantOnTraceFibers}.
\]
Thus random coins do not repair a quotient collision; they only index a family
of deterministic quotient obligations.  The collapsed two-state trace now has
\texttt{v13ToyRandomTraceObserver}, and the synthesis anchor
\texttt{v13TraceCaptureFactorizationCoverage\_anchor\_collapsedRandomTrace\_blocks}
records that one conflicting coin slice already blocks randomized trace
decodability.

\texttt{PNPv13AtomicEvidenceBudget.lean} supplies the next v13 ledger boundary.
It abstracts coordinates charging evidence atoms with natural-number costs.
Lean defines the coordinate-summed cost by grouping by atom multiplicity:
an atom charged by \(m\) coordinates contributes \(m\) times its cost.  The
central theorem
\[
  \texttt{coordinateSummedCost\_le\_atomBudget\_of\_noDoubleCounting}
\]
proves that the once-per-atom budget bounds the coordinate-summed cost only
under the explicit no-double-counting premise.  The converse diagnostic
\[
  \texttt{doubleCountedAtom\_of\_atomBudget\_lt\_coordinateSummedCost}
\]
proves that any strict over-budget coordinate sum exposes a positive-cost atom
with multiplicity at least two.  The toy surface
\texttt{v13DoubleCountedBudgetSurface} has two coordinates charging one shared
unit-cost atom, so \(\texttt{atomBudget}=1\) while
\(\texttt{coordinateSummedCost}=2\).  The synthesis anchor
\texttt{v13AtomicEvidenceBudgetCoverage\_anchor\_doubleCountedAtom\_blocks}
therefore fixes the exact repair theorem: a manuscript budget bound must prove
no positive-cost evidence atom is double-counted, or else pay the multiplicity.

\texttt{PNPv13AtomicEvidenceBudgetExactnessCharacterization.lean} adds the
missing dual obstruction.  It defines
\texttt{AllPositiveCostAtomsCharged} and
\texttt{MissedPositiveCostAtom} and proves
\[
  \texttt{atomBudget\_le\_coordinateSummedCost\_of\_allPositiveCostAtomsCharged},
\]
so charging every positive-cost atom at least once forces the once-per-atom
budget to be a lower bound on the coordinate-summed cost.  The converse
diagnostic
\[
  \texttt{missedPositiveCostAtom\_of\_coordinateSummedCost\_lt\_atomBudget}
\]
shows that any strict under-budget coordinate sum exposes a positive-cost atom
with multiplicity zero.  Finally,
\[
  \texttt{coordinateSummedCost\_eq\_atomBudget\_of\_noDoubleCounting\_of\_allPositiveCostAtomsCharged}
\]
proves the exact reuse criterion: a manuscript that wants to identify the
coordinate-summed cost with the once-per-atom budget must prove both that no
positive-cost atom is double-counted and that no positive-cost atom is omitted.
The toy surface \texttt{v13MissedPositiveBudgetSurface} charges one unit-cost
atom and misses another, so \(\texttt{coordinateSummedCost}=1\) while
\(\texttt{atomBudget}=2\).  This blocks the creative repair where the author
avoids multiplicity inflation only by silently leaving some positive-cost
evidence atoms unpaid.

\texttt{PNPv13AtomicEvidenceBudgetCancellationObstruction.lean} blocks the next
escape hatch.  The toy surface \texttt{v13CancellationBudgetSurface} lets both
coordinates charge atom \(0\) while atom \(1\) is never charged, with both
atoms having unit cost.  Lean proves
\[
  \texttt{coordinateSummedCost}=\texttt{atomBudget}=2
\]
on that surface, while simultaneously proving
\texttt{DoubleCountedAtom} and \texttt{MissedPositiveCostAtom}.  The theorems
\[
  \texttt{not\_forall\_coordinateSummedCost\_eq\_atomBudget\_implies\_noDoubleCounting},
\]
\[
  \texttt{not\_forall\_coordinateSummedCost\_eq\_atomBudget\_implies\_allPositiveCostAtomsCharged},
\]
and
\[
  \texttt{not\_forall\_coordinateSummedCost\_eq\_atomBudget\_implies\_exactness\_side\_conditions}
\]
make the route consequence explicit: exact aggregate budget balance does not
entail either structural side condition, nor their conjunction.  So the
manuscript cannot replace per-atom accounting theorems by a single scalar
equality and claim the budget step is repaired by net cancellation.

\texttt{PNPv13AtomicEvidenceBudgetEqualitySharpness.lean} closes the next repair
escape that remains after the cancellation countermodel.  The cancellation toy
surface shows that exact equality alone does not imply either structural side
condition, but the new file proves that once either side condition is supplied,
the other is forced by equality.  Concretely,
\[
  \texttt{coordinateSummedCost\_lt\_atomBudget\_of\_noDoubleCounting\_of\_missedPositiveCostAtom}
\]
shows that under no-double-counting, any missed positive-cost atom makes the
coordinate-summed cost strictly too small, while
\[
  \texttt{atomBudget\_lt\_coordinateSummedCost\_of\_allPositiveCostAtomsCharged\_of\_doubleCountedAtom}
\]
shows that if every positive-cost atom is charged, then any double-counted
positive-cost atom makes the coordinate-summed cost strictly too large.
Therefore the equality-surface theorems
\[
  \texttt{allPositiveCostAtomsCharged\_of\_coordinateSummedCost\_eq\_atomBudget\_of\_noDoubleCounting}
\]
and
\[
  \texttt{noDoubleCounting\_of\_coordinateSummedCost\_eq\_atomBudget\_of\_allPositiveCostAtomsCharged}
\]
prove that on the exact-budget surface, the two structural obligations are
equivalent.  This does not rescue the manuscript's route, because the same
cancellation surface still satisfies exact equality while violating both
conditions; instead it sharpens the crux by isolating the precise missing
content of any equality-based repair.  The author now needs a real theorem for
at least one side condition, not just an aggregate scalar identity.

\texttt{PNPv13AtomicEvidenceBudgetExactnessReduction.lean} makes that last
sentence formal and route-facing.  It packages the real budget-step obligation
as
\[
  \texttt{ExactnessSideConditions}
  \;:=\;
  \texttt{NoDoubleCounting} \wedge \texttt{AllPositiveCostAtomsCharged}
\]
and proves that this conjunction is exactly what suffices for exact budget
reuse.  On the equality surface, Lean then proves
\[
  \texttt{ExactnessSideConditions}
  \;\leftrightarrow\;
  \texttt{NoDoubleCounting}
\]
and also
\[
  \texttt{ExactnessSideConditions}
  \;\leftrightarrow\;
  \texttt{AllPositiveCostAtomsCharged}.
\]
So once a manuscript proves either structural side condition, it has proved the
whole exactness packet.  But the same file also proves the sharper obstruction
\[
  \texttt{not\_forall\_coordinateSummedCost\_eq\_atomBudget\_and\_sideConditionEquivalence\_implies\_exactnessSideConditions}.
\]
The cancellation surface is again the witness: it satisfies exact scalar
equality and the equality-surface equivalence
\(
  \texttt{NoDoubleCounting} \leftrightarrow
  \texttt{AllPositiveCostAtomsCharged}
\),
yet still fails \texttt{ExactnessSideConditions}.  This blocks the creative
repair where one cites only the induced equivalence of the two side conditions
and silently treats that as if one of them had been proved.

\texttt{PNPv13ACCEIEnvelope.lean} adds the finite skeleton for the manuscript's
ACCEI/PNLD envelope step without asserting the analytic envelope hypotheses.
The Markov-pruning theorem
\[
  \texttt{acceiBadCount\_mul\_succ\_threshold\_le\_of\_sum\_le}
\]
shows that, for a finite coordinate set, a total envelope budget bounds the
number of coordinates whose envelope exceeds a threshold.  The strengthened
finite pruning packet records the complement and division forms
\[
  \texttt{acceiBadCount\_eq\_card\_sub\_acceiGoodCount},\qquad
  \texttt{acceiBadCount\_le\_budget\_div\_succ\_threshold\_of\_sum\_le},
\]
and the large-good-set lower bound
\[
  \texttt{card\_sub\_budget\_div\_succ\_threshold\_le\_acceiGoodCount\_of\_sum\_le}.
\]
Thus the finite consequence of an averaged ACCEI budget is exactly a Markov
loss of at most
\(\texttt{budget}/(\texttt{threshold}+1)\) bad coordinates.  Lean also proves
\[
  \texttt{exists\_accei\_good\_of\_budget\_div\_succ\_threshold\_lt\_card},
\]
so at least one coordinate survives only when that finite loss is strictly
smaller than the switched-coordinate set.  This pins the repair obligation:
the manuscript cannot infer an arbitrarily large pruned coordinate set from a
global average envelope without paying this explicit loss term or proving a
stronger envelope hypothesis.  The sharpness packet makes this obstruction
concrete.  The envelope
\[
  \texttt{oneGoodOneBadACCEIEnvelope}
\]
has one coordinate at \(0\) and one coordinate at
\(\texttt{threshold}+1\).  Lean proves
\[
  \texttt{oneGoodOneBadACCEIEnvelope\_markov\_badCount\_sharp},\qquad
  \texttt{oneGoodOneBadACCEIEnvelope\_good\_lower\_bound\_sharp},
\]
so both the bad-count Markov inequality and the good-coordinate lower bound
are met with equality.  It also proves
\[
  \texttt{oneGoodOneBadACCEIEnvelope\_not\_all\_good},
\]
showing that a total budget of \(\texttt{threshold}+1\) does not imply all
coordinates survive pruning.  The sequential
theorem
\[
  \texttt{finiteHistory\_product\_bound\_of\_sequentialRateFrom}
\]
generalizes the earlier half-step tower product: if every next event has
finite one-step count rate at most
\(\texttt{numerator}/\texttt{denominator}\), then the whole history satisfies
the corresponding product bound.  Its contrapositive
\[
  \texttt{not\_sequentialRateAdmissibleFrom\_of\_product\_bound\_violation}
\]
is the Lean-side obstruction: any proposed ACCEI/PNLD pruning route that
violates the finite product bound cannot have the required one-step rate
certificate.  The stronger localization theorem
\[
  \texttt{exists\_failed\_prefixRateStep\_at\_append\_cons\_of\_product\_bound\_violation}
\]
turns such a violation into an explicit preceding history, next event, and
remaining suffix where the claimed one-step rate inequality fails.  Thus a
repair cannot hide behind a global product statement; it must supply every
local conditional rate cut.  The synthesis ledger records this as
\texttt{V13ACCEIEnvelopeProductCoverage}; this is deliberately only the finite
counting skeleton, not a proof that the manuscript's stopped-observer ACCEI
envelope exists or has the claimed budget.

The same module now gives the finite-coin stack its own conservative
route-coverage map, rather than letting the surrounding prose imply more than
the Lean theorems prove.  Lean names the covered finite-coin subrepair classes
as \texttt{PNPFiniteCoinSubrepairClass} and records the current covered list as
\texttt{currentPNPFiniteCoinCoveredSubrepairs}.  The covered items are:
\[
\begin{array}{ll}
\texttt{sameSourceFiniteCoinMarginalRecoding}
  & \text{same-source finite-coin marginal recodings of the source bit,}\\
\texttt{predicateErasure}
  & \text{balanced fibers erase every decidable output predicate,}\\
\texttt{decoderHalfAccuracy}
  & \text{balanced fibers force every Boolean decoder to be half-accurate,}\\
\texttt{exactScaledFiberImbalance}
  & \text{each output-fiber defect is }|\mathsf{Coin}|\text{ times imbalance,}\\
\texttt{subCardinalityToleranceQuantization}
  & \text{tolerance below }|\mathsf{Coin}|\text{ collapses to exact balance,}\\
\texttt{sourceLowerBoundToleranceQuantization}
  & \text{tolerance below a source/coin block collapses to balance,}\\
\texttt{retainedCoinTolerancePointwiseCollapse}
  & \text{retained-coin tolerance below a source/coin block collapses pointwise,}\\
\texttt{retainedSideToleranceSideFiberBalance}
  & \text{retained-side tolerance below a source/coin block collapses fiberwise,}\\
\texttt{sourceNonconstantUniformFiberZeroBlock}
  & \text{source-side nondegeneracy blocks uniform zero-tolerance fibers,}\\
\texttt{sourceNonconstantZeroIffBalanced}
  & \text{source-side nondegenerate zero tolerance iff balanced fibers,}\\
\texttt{deterministicOutputPostprocessing}
  & \text{exact erasure is stable under deterministic output maps,}\\
\texttt{injectivePostprocessingErasureEquivalence}
  & \text{lossless output maps cannot create finite-coin erasure,}\\
\texttt{observedLeftInversePostprocessingErasureEquivalence}
  & \text{support-recoverable postprocessing cannot create erasure,}\\
\texttt{observedInjectivePostprocessingErasureEquivalence}
  & \text{support-injective postprocessing cannot create erasure,}\\
\texttt{postprocessingObservedOutputCollisionObligation}
  & \text{postprocessed repair of imbalance must collide observed outputs,}\\
\texttt{postprocessingObservedOppositeBiasCollisionObligation}
  & \text{a balancing collision must merge opposite source-bias outputs,}\\
\texttt{postprocessingSignedBiasCancellationCertificate}
  & \text{postprocessed repair requires zero signed bias in each merged fiber,}\\
\texttt{postprocessedSideToleranceObservedFiberBalance}
  & \text{postprocessed retained-side tolerance collapses to observed balance,}\\
\texttt{postprocessedSideObservedFiberImbalanceLowerBound}
  & \text{postprocessed observed-fiber imbalance forces tolerance budget,}\\
\texttt{postprocessedSidePredicateErasure}
  & \text{postprocessed downstream predicates have equal side counts,}\\
\texttt{postprocessedSideDecoderHalfAccuracy}
  & \text{postprocessed downstream decoders are half-accurate below threshold,}\\
\texttt{postprocessedSideSourceRecoveryBlock}
  & \text{postprocessed downstream source recovery is blocked below threshold,}\\
\texttt{postprocessedSideCollisionObligation}
  & \text{postprocessed independence forces true/false observed overlap,}\\
\texttt{postprocessedSidePointwiseCollisionObligation}
  & \text{every observed coin value is witnessed on both source sides,}\\
\texttt{postprocessedSidePointwiseCollisionInsufficient}
  & \text{pointwise support overlap need not provide coin matching,}\\
\texttt{postprocessedSideOneSidedCoinMapInsufficient}
  & \text{one-sided witness maps need injectivity/bijectivity,}\\
\texttt{postprocessedSideInjectiveWitnessMapEquivalence}
  & \text{injective witness maps are exactly coin matching,}\\
\texttt{postprocessedSideCoinMatchingObligation}
  & \text{observed finite multiplicities support a preserving coin bijection,}\\
\texttt{postprocessedSideCoinMatchingCardinalityBlock}
  & \text{one observed multiplicity mismatch blocks that coin matching,}\\
\texttt{postprocessedSideCoinMatchingPredicateBlock}
  & \text{one downstream predicate mismatch blocks that coin matching,}\\
\texttt{postprocessedSidePredicateMatchingEquivalence}
  & \text{all downstream predicates neutral iff coin matching exists,}\\
\texttt{postprocessedSideApproxIndependenceCoinMatchingEquivalence}
  & \text{below threshold, approximate independence iff coin matching exists,}\\
\texttt{finitePreimagePostprocessingBlowup}
  & \text{finite-output postprocessing changes }\tau\text{ to }m\tau.
\end{array}
\]
The theorem \texttt{currentPNPFiniteCoinCoveredSubrepairs\_exact} pins this
list on the Lean side.  Several anchor theorems tie the list back to actual
mathematics rather than bookkeeping alone:
\[
\begin{gathered}
\texttt{finiteCoinCoverage\_anchor\_predicateErasure},\quad
\texttt{finiteCoinCoverage\_anchor\_decoderHalfAccuracy},\\
\texttt{finiteCoinCoverage\_anchor\_exactScaledFiberImbalance},\quad
\texttt{finiteCoinCoverage\_anchor\_subCardinalityToleranceQuantization},\\
\texttt{finiteCoinCoverage\_anchor\_sourceLowerBoundToleranceQuantization},\\
\texttt{finiteCoinCoverage\_anchor\_retainedCoinTolerancePointwiseCollapse},\\
\texttt{finiteCoinCoverage\_anchor\_retainedSideToleranceSideFiberBalance},\\
\texttt{finiteCoinCoverage\_anchor\_deterministicOutputPostprocessing},\quad
\texttt{finiteCoinCoverage\_anchor\_injectivePostprocessingErasureEquivalence},\\
\texttt{finiteCoinCoverage\_anchor\_observedLeftInversePostprocessingErasureEquivalence},\\
\texttt{finiteCoinCoverage\_anchor\_observedInjectivePostprocessingErasureEquivalence},\\
\texttt{finiteCoinCoverage\_anchor\_postprocessingObservedOutputCollisionObligation},\\
\texttt{finiteCoinCoverage\_anchor\_postprocessingObservedOppositeBiasCollisionObligation},\\
\texttt{finiteCoinCoverage\_anchor\_postprocessingSignedBiasCancellationCertificate},\\
\texttt{finiteCoinCoverage\_anchor\_postprocessedSideToleranceObservedFiberBalance},\\
\texttt{finiteCoinCoverage\_anchor\_postprocessedSideObservedFiberImbalanceLowerBound},\\
\texttt{finiteCoinCoverage\_anchor\_postprocessedSidePredicateErasure},\\
\texttt{finiteCoinCoverage\_anchor\_postprocessedSideDecoderHalfAccuracy},\\
\texttt{finiteCoinCoverage\_anchor\_postprocessedSideSourceRecoveryBlock},\\
\texttt{finiteCoinCoverage\_anchor\_postprocessedSideCollisionObligation},\\
\texttt{finiteCoinCoverage\_anchor\_postprocessedSidePointwiseCollisionObligation},\\
\texttt{finiteCoinCoverage\_anchor\_postprocessedSidePointwiseCollisionInsufficient},\\
\texttt{finiteCoinCoverage\_anchor\_postprocessedSideOneSidedCoinMapInsufficient},\\
\texttt{finiteCoinCoverage\_anchor\_postprocessedSideInjectiveWitnessMapEquivalence},\\
\texttt{finiteCoinCoverage\_anchor\_postprocessedSideCoinMatchingObligation},\\
\texttt{finiteCoinCoverage\_anchor\_postprocessedSideCoinMatchingCardinalityBlock},\\
\texttt{finiteCoinCoverage\_anchor\_postprocessedSideCoinMatchingPredicateBlock},\\
\texttt{finiteCoinCoverage\_anchor\_postprocessedSidePredicateMatchingEquivalence},\\
\texttt{finiteCoinCoverage\_anchor\_postprocessedSideApproxIndependenceCoinMatchingEquivalence},\\
\texttt{finiteCoinCoverage\_anchor\_finitePreimagePostprocessingBlowup},\\
\texttt{finiteCoinCoverage\_anchor\_coinwiseTrueFiber\_sourceNonconstant\_zeroBlock},\quad
\texttt{finiteCoinCoverage\_anchor\_coinwiseFalseFiber\_sourceNonconstant\_zeroBlock},\\
\texttt{finiteCoinCoverage\_anchor\_sourceNonconstantZeroIffBalanced}.
\end{gathered}
\]

The same synthesis file now also records a narrow randomized-residual subledger
\texttt{PNPRandomizedResidualSubrepairClass}.  This is intentionally not a
claim that broad randomized residual repairs are closed.  It records the
finite-coin product-space facts now proved by
\texttt{RandomizedResidualObstruction.lean} and
\texttt{PostSwitchRandomizedResidualObstruction.lean}:
\[
\begin{array}{ll}
\texttt{productSideChannelBudget}
  & \text{finite coins are an ordinary product-space side channel,}\\
\texttt{supportwiseUnresolvedZeroMass}
  & \text{supportwise unresolved residuals have zero joint resolving mass,}\\
\texttt{supportwiseUnresolvedNoAdvantage}
  & \text{supportwise unresolved residuals have no positive advantage,}\\
\texttt{positiveAdvantageResolutionWitness}
  & \text{positive advantage yields a positive-weight resolving coin,}\\
\texttt{positiveAdvantageDeterministicCoinSliceWitness}
  & \text{positive advantage yields a deterministic resolving coin slice,}\\
\texttt{strictHalfDeterministicCoinSliceWitness}
  & \text{strict half-accuracy advantage still yields a deterministic resolving coin slice,}\\
\texttt{strictHalfDeterministicCoinSliceResidualObstructionPackage}
  & \text{strict half-accuracy advantage already yields a deterministic slice carrying the ordinary residual obstruction package,}\\
\texttt{postSwitchSupportwiseUnresolvedNoAdvantage}
  & \text{the exact post-switch surface inherits the supportwise blocker,}\\
\texttt{postSwitchExactSupportNoStrictHalfAdvantage}
  & \text{exact post-switch support blocks strict half-accuracy randomized residual advantage,}\\
\texttt{postSwitchPositiveAdvantageResolutionWitness}
  & \text{positive post-switch advantage yields a resolving input/coin pair,}\\
\texttt{postSwitchPositiveAdvantageDeterministicCoinSliceWitness}
  & \text{positive post-switch advantage yields a deterministic resolving coin slice,}\\
\texttt{postSwitchStrictHalfDeterministicCoinSliceWitness}
  & \text{strict half-accuracy post-switch advantage still yields a deterministic resolving coin slice,}\\
\texttt{postSwitchStrictHalfDeterministicCoinSliceResidualObstructionPackage}
  & \text{the exact post-switch strict-half slice already carries the invariant-projection residual obstruction package.}
\end{array}
\]
The route anchors
\[
\begin{gathered}
\texttt{randomizedResidualCoverage\_anchor\_supportwise\_unresolved\_no\_positive\_advantage},\\
\texttt{randomizedResidualCoverage\_anchor\_positive\_advantage\_witnesses\_resolving\_coin},\\
\texttt{randomizedResidualCoverage\_anchor\_positive\_advantage\_witnesses\_deterministic\_coin\_slice},\\
\texttt{randomizedResidualCoverage\_anchor\_strict\_half\_advantage\_witnesses\_deterministic\_coin\_slice},\\
\texttt{randomizedResidualCoverage\_anchor\_strict\_half\_advantage\_coin\_slice\_forces\_residual\_obstruction\_package},\\
\texttt{randomizedResidualCoverage\_anchor\_postSwitch\_exact\_support\_no\_strict\_half\_advantage},\\
\texttt{randomizedResidualCoverage\_anchor\_postSwitch\_positive\_advantage\_witnesses\_resolving\_coin},\\
\texttt{randomizedResidualCoverage\_anchor\_postSwitch\_positive\_advantage\_witnesses\_deterministic\_coin\_slice},\\
\texttt{randomizedResidualCoverage\_anchor\_postSwitch\_strict\_half\_advantage\_witnesses\_deterministic\_coin\_slice},\\
\texttt{randomizedResidualCoverage\_anchor\_postSwitch\_strict\_half\_advantage\_coin\_slice\_forces\_residual\_obstruction\_package}
\end{gathered}
\]
tie that subledger to the actual obstruction theorems.

This matters because the synthesis layer no longer records the randomized
post-switch lane only at the positive-doubled-advantage threshold.  Even after
weakening to the strict half-accuracy formulation, the same finite-coin route
is now Lean-visible as already forcing a deterministic ordinary residual slice,
and even that slice is no longer a separate escape hatch: Lean immediately
upgrades it to the full residual-side obstruction package, both generically and
on the exact post-switch support surface.

This is deliberately a subledger.  The theorem
\texttt{finiteCoinStackUncoveredBroadRepairClasses\_exact} states that the
broad classes not closed by the finite-coin stack are exactly the current
synthesis gaps.  In particular, the finite-coin stack is recorded as coverage
inside \texttt{sameSourceFiniteCountApproximation}, not as a proof of
\texttt{randomizedResidual}, \texttt{residualSideInformation},
\texttt{approximateDecorrelation}, or \texttt{kpolyCompressionBridge}.

It also records the named repair classes still outside the packet:
\[
\begin{array}{ll}
\texttt{residualSideInformation} & \text{post-conditioned data not determined by the source bit,}\\
\texttt{randomizedResidual} & \text{randomized post-conditioning residual objects,}\\
\texttt{approximateDecorrelation} & \text{broad quantitative replacement for exact finite-count independence,}\\
\texttt{kpolyCompressionBridge} & \text{the remaining manuscript-specific exact-visible / clocked realization bridge.}
\end{array}
\]

The residual-side-information gap now has a first narrow collision packet in
\texttt{ResidualSideInformationObstruction.lean}.  It defines
\texttt{SideInfoDeterminedBy} for residual data that factors through a source
summary, \texttt{SideInfoCollisionOverBase} for two worlds with the same source
summary and different residual values,
\texttt{PositiveWeightSideInfoCollisionOverBase} for the same collision with a
positive-weight supported endpoint, and
\texttt{SourceOnlyPredicateCapturesSideEq} for source-only equality tests on a
residual value.  Lean proves
\[
\begin{gathered}
\texttt{not\_sideInfoDeterminedBy\_of\_same\_base\_ne\_side},\\
\texttt{not\_sourceOnlyPredicateCapturesSideEq\_left\_of\_same\_base\_ne\_side},\\
\texttt{not\_sourceOnlyBooleanClassifier\_of\_same\_base\_ne\_target}.
\end{gathered}
\]
Thus a same-source residual collision cannot be rebranded as source-only data,
even for the single equality predicate naming one residual value.  The concrete
two-point witness
\texttt{residualSideInfoToyBase}/\texttt{residualSideInfoToySide} pins this
failure without any probabilistic or asymptotic assumptions:
\[
\begin{gathered}
\texttt{residualSideInfoToy\_collision},\qquad
\texttt{residualSideInfoToy\_not\_determinedBy\_base},\\
\texttt{residualSideInfoToy\_not\_sourceOnlyPredicate\_false},\qquad
\texttt{residualSideInfoToy\_not\_sourceOnlyBooleanClassifier}.
\end{gathered}
\]
\texttt{CruxSynthesis.lean} records this as the narrow
\texttt{PNPResidualSideInformationSubrepairClass} subledger, with route anchors
named \texttt{residualSideInformationCoverage\_anchor\_*}.  The newest
strengthening packages these anchors as the theorem-backed packet
\texttt{V13ResidualSideInformationSubcoverage}.  This remains deliberately
narrow: it does not claim broad coverage of
\texttt{PNPRepairClass.residualSideInformation}.  What it does record is the
stronger invariant-base audit boundary now available in Lean.  Positive doubled
advantage or strict half-accuracy advantage must expose a supported residual
collision, a supported residual equality test that is not source-only, and a
positive-support prediction separation witness; conversely, supported
source-only residual predicates or supportwise source-only classifier behavior
block both advantage formulations.  Any repair that keeps varying residual data
must either prove that the residual is source-determined or expose a genuine
side channel to the resolving-mass theory.

The follow-up packet
\texttt{ResidualSideInformationResolutionCost.lean} removes one tempting
interpretation of the ``source-determined'' escape.  If the base feature is
itself invariant under the manuscript involution and the residual side
information factors through that base feature, then Lean proves
\[
\begin{gathered}
\texttt{side\_eq\_under\_involution\_of\_determinedBy\_invariant\_base},\\
\texttt{resolvedMass\_eq\_zero\_of\_determinedBy\_invariant\_base},\\
\texttt{not\_pos\_doubledAdvantage\_pair\_of\_determinedBy\_invariant\_base},\\
\texttt{not\_total\_lt\_two\_mul\_weightedCorrectMass\_pair\_of\_determinedBy\_invariant\_base}.
\end{gathered}
\]
Thus such residuals are not merely source-only in a functional sense; they have
zero orbit-resolving mass and cannot beat chance in either the doubled-advantage
or strict half-accuracy formulations.  The contrapositive boundary is also
formalized:
\[
\begin{gathered}
\texttt{not\_sideInfoDeterminedBy\_of\_pos\_doubledAdvantage\_invariant\_base},\\
\texttt{not\_sideInfoDeterminedBy\_of\_total\_lt\_two\_mul\_weightedCorrectMass\_invariant\_base}.
\end{gathered}
\]
The 2026-05-05 residual-side packet strengthens this negative boundary into a
positive witness obligation.  Lean now proves
\[
\begin{gathered}
\texttt{exists\_positive\_same\_base\_ne\_side\_of\_pos\_doubledAdvantage\_invariant\_base},\\
\texttt{exists\_positive\_same\_base\_ne\_side\_of\_total\_lt\_two\_mul\_weightedCorrectMass\_invariant\_base}.
\end{gathered}
\]
The supported-collision packet then packages the same witness at the residual
subledger surface:
\[
\begin{gathered}
\texttt{positiveWeightSideInfoCollisionOverBase\_of\_pos\_doubledAdvantage\_invariant\_base},\\
\texttt{positiveWeightSideInfoCollisionOverBase\_of\_total\_lt\_two\_mul\_weightedCorrectMass\_invariant\_base},\\
\texttt{exists\_positive\_not\_sourceOnlyPredicateCapturesSideEq\_of\_pos\_doubledAdvantage\_invariant\_base},\\
\texttt{exists\_positive\_not\_sourceOnlyPredicateCapturesSideEq\_of\_total\_lt\_two\_mul\_weightedCorrectMass\_invariant\_base}.
\end{gathered}
\]
The contrapositive no-advantage form is also now formalized:
\[
\begin{gathered}
\texttt{not\_pos\_doubledAdvantage\_pair\_of\_supported\_sourceOnlyPredicateCapturesSideEq\_invariant\_base},\\
\texttt{not\_total\_lt\_two\_mul\_weightedCorrectMass\_pair\_of\_supported\_sourceOnlyPredicateCapturesSideEq\_invariant\_base}.
\end{gathered}
\]
Thus a successful residual-side-information repair over an invariant base must
not merely prove that the residual is not source-derived.  Any claimed
advantage must expose a positive-weight point \(x\) where the base still agrees
across the involution, \(\texttt{base}(\tau x)=\texttt{base}(x)\), but the
residual side information changes, \(\texttt{side}(\tau x)\ne\texttt{side}(x)\).
Equivalently, the repair must provide an actual same-source residual collision
on support, and a supported residual equality predicate that cannot be
source-only, or else move to a different, explicitly symmetry-breaking base.
If all positive-weight residual equality tests are already source-only, Lean
now blocks both positive doubled advantage and strict half-accuracy advantage.
The next repair layer is classifier-level rather than value-level: the
classifier may syntactically mention \((\texttt{base},\texttt{side})\), but its
realized positive-support predictions might still factor through the invariant
base.  This is now blocked as well.  The generic supportwise lemma
\[
\begin{gathered}
\texttt{not\_total\_lt\_two\_mul\_weightedCorrectMass\_of\_supportwise\_prediction\_invariant},\\
\texttt{not\_pos\_doubledAdvantage\_of\_supportwise\_prediction\_invariant}
\end{gathered}
\]
shows that any positive-support prediction invariant under the involution stays
at chance.  The residual specialization introduces
\texttt{SupportwiseSourceOnlyPairClassifier} and proves
\[
\begin{gathered}
\texttt{prediction\_eq\_under\_involution\_of\_supportwiseSourceOnlyPairClassifier\_invariant\_base},\\
\texttt{not\_pos\_doubledAdvantage\_pair\_of\_supportwiseSourceOnlyPairClassifier\_invariant\_base},\\
\texttt{not\_total\_lt\_two\_mul\_weightedCorrectMass\_pair\_of\_supportwiseSourceOnlyPairClassifier\_invariant\_base}.
\end{gathered}
\]
Thus a residual-side-information repair over an invariant base cannot survive
by saying that only the final classifier, not each residual equality test, is
source-only on support.  If the realized predictions factor through the base
where the weight is positive, Lean blocks both advantage formulations.  This
still does not close arbitrary proof-relevant side channels; it makes the
invariant-base residual repair boundary exact up to genuinely
involution-sensitive supported predictions.  The new 2026-05-08 packet turns
that classifier-level boundary into a positive audit obligation.  At the
generic level, Lean now proves
\[
\begin{gathered}
\texttt{exists\_positive\_prediction\_ne\_of\_pos\_doubledAdvantage},\\
\texttt{exists\_positive\_prediction\_ne\_of\_total\_lt\_two\_mul\_weightedCorrectMass},
\end{gathered}
\]
so any supportwise involution-invariant classifier prediction immediately
contradicts either advantage formulation.  The invariant-base residual
specialization packages this as
\[
\begin{gathered}
\texttt{exists\_positive\_prediction\_ne\_of\_pos\_doubledAdvantage\_pair\_invariant\_base},\\
\texttt{exists\_positive\_prediction\_ne\_of\_total\_lt\_two\_mul\_weightedCorrectMass\_pair\_invariant\_base}.
\end{gathered}
\]
Hence a successful residual-side-information repair over an involution-invariant
base must now do more than exhibit a supported residual value variation.  It
must exhibit a positive-weight point \(x\) where the final classifier actually
separates the involution pair,
\[
h(\texttt{base}(\tau x), \texttt{side}(\tau x))
  \ne h(\texttt{base}(x), \texttt{side}(x)).
\]
Residual variation without supported prediction separation is no longer enough
to claim that the side channel is genuinely proof-relevant to the classifier.
Update note, 2026-05-08: the active invariant-base residual-side-information
boundary is no longer merely ``not source-determined'' or even ``not
source-only on support''.  The exact remaining audit obligation is the
prediction-separation witness above.  The broad
\texttt{PNPRepairClass.residualSideInformation} entry in
\texttt{CruxSynthesis.lean} is still not stop-grade covered, but any repair
that stays over an involution-invariant base must now surface explicit
supported classifier-level separation.

Update note, 2026-05-14: the synthesis layer now packages these invariant-base
residual obligations into one route-facing theorem rather than leaving them as
separate repair chores.  Lean proves
\[
\begin{gathered}
\texttt{residualSideInformationCoverage\_anchor\_positive\_advantage\_forces\_supported\_obstruction\_package},\\
\texttt{residualSideInformationCoverage\_anchor\_strict\_half\_advantage\_forces\_supported\_obstruction\_package}.
\end{gathered}
\]
So any positive doubled advantage, or any strict half-accuracy advantage, from
\((\texttt{base},\texttt{side})\) over an involution-invariant base must now
simultaneously pay the full residual-side-information bill: the residual is not
source-determined, the realized positive-support classifier is not source-only,
there is a supported same-base residual variation witness, there is a supported
residual equality test that is not source-only, and there is supported
prediction separation.  These are no longer five independent places where a
manuscript repair can hope to wiggle through one at a time; in Lean they now
arrive as one coupled obstruction package.

Update note, 2026-05-14 (later): the residual-side-information lane now states
the quantitative bridge that the route summary had already been appealing to.
Lean proves
\[
\begin{gathered}
\texttt{residualSideInformationCoverage\_anchor\_positive\_advantage\_forces\_positive\_resolvedMass},\\
\texttt{residualSideInformationCoverage\_anchor\_strict\_half\_advantage\_forces\_positive\_resolvedMass},\\
\texttt{residualSideInformationCoverage\_anchor\_positive\_resolvedMass\_witnesses\_same\_base\_residual\_collision}.
\end{gathered}
\]
So an invariant-base residual repair no longer gets to claim only that some
supported residual witness exists.  Any positive doubled advantage or strict
half-accuracy advantage from \((\texttt{base},\texttt{side})\) must already pay
strictly positive orbit-resolving mass for \(\texttt{side}\), and that
quantitative obligation in turn forces a positive-weight same-base residual
collision.  This is the direct synthesis-layer reduction from residual-side
advantage claims to the generic side-channel resolution-demand obstruction.

Update note, 2026-05-14 (later still): the same quantitative condition is now
formalized as a classifier-free residual obstruction package.  In
\texttt{ResidualSideInformationResolutionCost.lean}, Lean proves
\[
\texttt{pos\_resolvedMass\_obstruction\_package\_invariant\_base}.
\]
So once an invariant-base residual side channel carries positive
orbit-resolving mass, the manuscript no longer gets to defer the residual audit
to later classifier choices.  Before any final Boolean predictor is discussed,
positive resolved mass already forces all of the following: the residual is not
source-determined, there is a positive-weight same-base residual collision, and
there is a supported residual equality predicate that is not source-only.  The
synthesis file now re-exports this as
\[
\texttt{residualSideInformationCoverage\_anchor\_positive\_resolvedMass\_forces\_supported\_obstruction\_package}.
\]
This further tightens the route: residual-side-information advantage first has
to pay positive resolved mass, and positive resolved mass already pays a pure
residual obstruction package even before prediction-separation obligations are
considered.

Update note, 2026-05-14 (exact post-switch witness): this pure residual packet
is now instantiated on the manuscript's exact visible post-switch surface
itself.  In \texttt{PostSwitchResidualWitness.lean}, Lean proves
\[
\begin{gathered}
\texttt{pos\_resolvedMass\_activeOrbit},\\
\texttt{exactPostSwitch\_pos\_resolvedMass\_forces\_pure\_residual\_obstruction\_package},\\
\texttt{exactPostSwitch\_activeOrbit\_realizes\_pure\_residual\_obstruction\_package}.
\end{gathered}
\]
So the active two-point exact post-switch orbit already carries positive
ordinary resolved mass for the retained side bit, and that quantitative fact
alone already forces the pure residual obstruction package on the manuscript's
own local visible surface.  Before any classifier-level prediction-separation
claim is used, Lean now extracts from this exact post-switch witness that the
retained side channel is not source-determined over the invariant projection,
that there is a positive-weight same-base residual collision, and that there is
a supported residual equality predicate that is not source-only.

Update note, 2026-05-14 (synthesis ledger integration): this exact post-switch
active-orbit packet is now re-exported through \texttt{CruxSynthesis.lean}
itself.  The theorem
\[
\texttt{residualSideInformationCoverage\_anchor\_exactPostSwitch\_activeOrbit\_pure\_residual\_package}
\]
packages the same concrete active-orbit witness at the synthesis layer, and
the residual-side subrepair ledger now records it as the explicit covered case
\[
\texttt{exactPostSwitchConcretePureResidualWitness}.
\]
So the theorem-backed route summary no longer stops at the abstract statement
that positive resolved mass forces a pure residual obstruction package.  It now
also includes the manuscript's own exact post-switch visible surface as a
concrete theorem-backed instance of that package.

Update note, 2026-05-14 (exact post-switch prediction witness): the same exact
post-switch active orbit is now also recorded at the synthesis layer as a
concrete classifier witness.  The local theorem
\[
\texttt{exactPostSwitch\_activeOrbit\_realizes\_residual\_prediction\_witness}
\]
is re-exported in \texttt{CruxSynthesis.lean} as
\[
\texttt{residualSideInformationCoverage\_anchor\_exactPostSwitch\_activeOrbit\_prediction\_witness}.
\]
So the theorem-backed route summary now includes not only the pure residual
package on the manuscript's exact visible surface, but also the concrete
classifier-level fact that this active orbit preserves the invariant
projection, has involution-symmetric weight, flips the target on positive
support, achieves positive doubled advantage, and separates a positive-weight
involution pair at the final classifier output.

Update note, 2026-05-15 (active-orbit collapse to the exact fork blocker): the
same concrete witness no longer sits only at the level of ``there exists a
positive residual prediction route.''  In
\texttt{PostSwitchResidualWitness.lean} the new theorem
\[
\texttt{exactPostSwitch\_activeOrbit\_forces\_fork\_obstruction}
\]
packages three facts at once:
\begin{enumerate}[leftmargin=1.5em]
  \item the active orbit beats half accuracy in the strict form
  \(\texttt{weightedTotalMass}<2\cdot\texttt{weightedCorrectMass}\);
  \item its support already carries positive mass on the nonzero-column slice;
  \item its weight function is therefore not exact-input invariant.
\end{enumerate}
So the manuscript's own concrete residual witness no longer leaves a separate
``residual side information'' escape route on the exact post-switch surface.
Lean now sends that witness directly into the already formalized fork
obstruction surface: any attempt to keep the same active-orbit witness while
also insisting on exact-input-invariant support is immediately contradictory,
and this remains true on the strict half-accuracy formulation, not only on the
positive doubled-advantage one.

Update note, 2026-05-17 (synthesis-layer residual/fork collapse): this exact
surface collapse is now recorded inside the residual-side synthesis ledger
itself, rather than only in the local witness file.  In
\texttt{CruxSynthesis.lean}, Lean proves
\[
\texttt{residualSideInformationCoverage\_anchor\_exactPostSwitch\_prediction\_witness\_forces\_fork\_obstruction}.
\]
The residual-side subrepair ledger also now records the new covered case
\[
\texttt{exactPostSwitchPredictionWitnessForcesForkObstruction}.
\]
So on the manuscript's concrete exact visible post-switch surface, the
classifier-level residual prediction witness is no longer tracked as a
standalone repair lane.  At the synthesis interface it is explicitly coupled to
the same fork-obstruction package: strict half-accuracy success, positive
nonzero-column mass, and failure of exact-input-invariant support arrive
together on the same active two-point orbit.

Update note, 2026-05-17 (generic exact post-switch \(b\)-route collapse): the
same fork package is now also stated at the full exact post-switch
\(((z,a_i),b)\) route, not only on the concrete active orbit.  Lean now
re-exports
\[
\texttt{postSwitchForkCoverage\_anchor\_strict\_half\_advantage\_fork\_obstruction\_package}.
\]
So for any exact post-switch classifier on \((\texttt{invariantProjection}, b)\),
strict half-accuracy advantage already forces three simultaneous obligations:
there is a positive-weight input where the retained \(b\) side channel changes
under the involution, there is a positive-weight nonzero \(a_i\) column, and
the support cannot be exact-input invariant.  This removes the last gap between
the local fork-obstruction file and the route-synthesis layer for the strict
half-accuracy exact \(b\)-route.

Update note, 2026-05-17 (finite-coin randomized residual is not an independent
strict-half route): the same synthesis layer now also re-exports
\[
\texttt{randomizedResidualCoverage\_anchor\_strict\_half\_advantage\_coin\_slice\_forces\_residual\_obstruction\_package}
\]
and
\[
\texttt{randomizedResidualCoverage\_anchor\_postSwitch\_strict\_half\_advantage\_coin\_slice\_forces\_residual\_obstruction\_package}.
\]
So on both the generic finite-coin product-space surface and the exact
post-switch invariant-projection surface, strict half-accuracy randomized
residual advantage no longer merely forces a deterministic positive-weight coin
slice.  That slice already carries the ordinary residual-side obstruction
package: failure of determination by the invariant base, positive collision
over that base, an explicit positive-weight resolving point, and failure of the
corresponding source-only side-equality capture.  Thus the finite-coin
randomized residual repair is no longer an independent synthesis lane at the
strict-half threshold; Lean collapses it directly into the residual-side route.

The broad \texttt{approximateDecorrelation} class remains outside the stop-grade
packet.  However, \texttt{ApproximateDecorrelationObstruction.lean} now blocks
the immediate same-source coupled, anti-coupled, and arbitrary deterministic
distinguishing-fiber finite-count subcases: any claimed tolerance below
\(|C\cap E|\,|C\cap\neg E|\) contradicts the proved cross-product defect
identity, and any claimed tolerance certificate for such an informative fiber
therefore entails the explicit obligation
\(|C\cap E|\,|C\cap\neg E|\le\tau\).  With lower bounds
\(L_1\le |C\cap E|\) and \(L_0\le |C\cap\neg E|\), Lean further reduces this
to \(L_1L_0\le\tau\), and directly blocks any claimed tolerance
\(\tau<L_1L_0\).  The same module now also exposes the whole-output predicate
\texttt{CountIndependentWithinToleranceOutput}; for every nonconstant
deterministic recoding \(r:\{0,1\}\to\alpha\), approximate independence of all
output fibers still forces the same product/tolerance obligation because the
\(r(1)\)-fiber distinguishes the source bit.  For finite-coin randomized
recodings \(r:\{0,1\}\to\mathsf{Coin}\to\alpha\), the same product/tolerance
obligation first holds whenever some output fiber uniformly distinguishes the
source side across all coins, and the product-space obligation factors to
\((L_1|\mathsf{Coin}|)(L_0|\mathsf{Coin}|)\le\tau\) from source-side lower
bounds \(L_1,L_0\).  The newer mixed-coin identities remove that uniformity
assumption: for every output fiber \(y\), the forward and backward defects are
exactly
\[
  |C\cap E|\,|C\cap\neg E|\,|\mathsf{Coin}|(T_y-F_y)
  \quad\text{and}\quad
  |C\cap E|\,|C\cap\neg E|\,|\mathsf{Coin}|(F_y-T_y),
\]
where \(T_y=\#\{c:r(1,c)=y\}\) and \(F_y=\#\{c:r(0,c)=y\}\).  Therefore a
finite-coin same-source randomized repair must either balance all output-fiber
coin counts or pay the scaled imbalance in the tolerance; the lower-bound
variants let this be checked from source-side estimates \(L_1,L_0\) without
knowing the exact conditioned counts.  This is now the arbitrary finite-coin
same-source recoding subcase of the randomized residual repair space.  The
balanced subcase is now also characterized: balanced fibers give exact
zero-tolerance independence for the marginal output, but if the random coin is
retained and any coin value separates \(r(1,c)\) from \(r(0,c)\), the output
\((r(E,c),c)\) has a fiber with defect
\(|C\cap E|\,|C\cap\neg E|\,|\mathsf{Coin}|\).  Hence a balanced same-source
finite-coin repair must either be a genuine erasure at the output-fiber level,
or account for retained side information under the same finite-count audit.
The retained-side theorem generalizes the retained-coin result: any side
channel whose fiber isolates a true-side output from false-side matches has a
forward defect, and any side channel whose fiber isolates a false-side output
from true-side matches has a backward defect.  In both cases the defect is
\(|C\cap E|\,|C\cap\neg E|\,|\mathsf{Coin}|\) times the size of the selected
side/output fiber.  The postprocessed-side theorem removes the remaining
syntactic escape: even after applying an arbitrary observation
\(p:\alpha\times\mathsf{Side}\to\beta\), any observed true-side fiber isolated
from all false-side observations has the same forward defect, and any observed
false-side fiber isolated from all true-side observations has the same backward
defect, scaled by the selected observed-fiber count.  Thus a repair cannot keep
source-bearing side information and evade the finite-count audit merely by
hashing or compressing the retained pair.  The output-predicate count theorem
makes ``genuine erasure'' precise:
balanced fibers force every decidable output property to have equal true-side
and false-side coin counts.  The previous output-only decoder theorem is the
perfect-recovery special case; the count theorem also rules out any strict
predicate-level advantage in either direction.  The Boolean decoder corollary
strengthens this into an exact half-accuracy law: true-side correct decodes
plus false-side correct decodes equal \(|\mathsf{Coin}|\), so no balanced
marginal decoder can beat random-guess aggregate accuracy over the two source
sides.
For the uniform true-side/false-side finite-coin fibers, the zero-tolerance
version is now phrased over the original source counts: a nonempty finite coin
space lifts positive source-side true and false conditioned fibers to the
product space, so the corresponding zero-tolerance output-independence claim is
impossible without asking the manuscript author to restate nondegeneracy on
\(\Omega\times\mathsf{Coin}\).

Update note, 2026-05-14 (postprocessed-side skew threshold counterexample):
the earlier three-coin postprocessed-side skew audit is now upgraded from a
mere support-overlap counterexample to a quantitative approximate-decorrelation
blocker.  In \texttt{ApproximateDecorrelationObstruction.lean}, Lean proves
\[
\begin{gathered}
\texttt{postprocessedSidePointwiseCollisionSkew\_true\_count\_true}=2,\\
\texttt{postprocessedSidePointwiseCollisionSkew\_false\_count\_true}=1,\\
\texttt{postprocessedSidePointwiseCollisionSkew\_not\_countIndependentWithinToleranceOutput\_of\_lt\_three}.
\end{gathered}
\]
So the same skew witness still has two-sided pointwise collisions, and it still
admits the earlier one-sided preserving true-to-false witness map, but every
approximate-independence certificate with tolerance \(\tau<3\) is formally
refuted.  The package theorem
\[
\texttt{postprocessedSidePointwiseCollisionSkew\_collision\_and\_oneSidedMap\_counterexample\_of\_lt\_three}
\]
records all three facts together.  This sharpens the route-level obstruction:
neither pointwise support overlap nor one-sided witness choice is enough to
repair postprocessed retained-side dependence below threshold; the manuscript
still has to solve the exact multiplicity-level coin-matching obligation.
The same file now proves the lower bound is \emph{sharp}.  Lean computes the
two relevant fiber defects exactly:
\[
\begin{gathered}
\texttt{postprocessedSidePointwiseCollisionSkew\_true\_forward\_gap\_eq\_three},\\
\texttt{postprocessedSidePointwiseCollisionSkew\_false\_backward\_gap\_eq\_three},
\end{gathered}
\]
with the opposite oriented defects equal to \(0\).  From those exact gap
values, Lean proves
\[
\texttt{postprocessedSidePointwiseCollisionSkew\_countIndependentWithinToleranceOutput\_iff\_three\_le}.
\]
So this three-coin skew witness is not merely a ``fails below some threshold''
example: its approximate-decorrelation threshold is formally characterized as
exactly \(\tau = 3\).  The manuscript therefore cannot salvage the route by
claiming that support-overlap plus a one-sided preserving map becomes adequate
at some smaller hidden tolerance; the obstruction already isolates the precise
minimal tolerance needed on this witness.

The theorem
\[
  \texttt{not\_stopGrade\_currentPNPLocalCruxPacket}
\]
proves that the current packet is not stop-grade, by the explicit uncovered
\texttt{residualSideInformation} class.  Companion theorems separately mark
\texttt{residualSideInformation}, \texttt{randomizedResidual},
\texttt{approximateDecorrelation}, and \texttt{kpolyCompressionBridge} as
uncovered.  The finite-coin-specific companion theorems
\[
\begin{gathered}
  \texttt{finiteCoinStack\_does\_not\_cover\_residualSideInformation},\quad
  \texttt{finiteCoinStack\_does\_not\_cover\_randomizedResidual},\\
  \texttt{finiteCoinStack\_does\_not\_cover\_approximateDecorrelation},\quad
  \texttt{finiteCoinStack\_does\_not\_cover\_kpolyCompressionBridge}
\end{gathered}
\]
make the same scoping constraint explicit for the new finite-coin work.  By
contrast,
\[
  \texttt{sameSourceFiniteCountApproximation\_covered\_currentPNPLocalCruxPacket}
\]
records the finite-count approximation subcase as covered, and
\[
  \texttt{clockedFiniteUniformKernel\_covered\_currentPNPLocalCruxPacket}
\]
records that the abstract finite-counting kernel itself is already on the
covered side of the ledger.

The current small-image fork is now integrated into the same synthesis layer.
\texttt{CruxSynthesis.lean} imports
\texttt{ClockedKpolyActualGapClosure.lean} and records the narrow
\texttt{PNPKpolySubrepairClass} entries for this fork: exact ZAB ERM payload
closure, the bit-vector specialization, equality-transfer to the bit-vector ERM
family, the bare exact-visible interface obstruction, the actual-support
quotient boundaries, the plug-in lookup/majority endpoint obstructions, and
the bounded-sample false-support, default-tie, fallback-side-channel, and
structured fallback-family sparse-override boundaries described below.
The named theorem
\[
  \texttt{currentPNPStrongestHonestStatus}
\]
assembles the present honest claim: the theorem-backed v13 entries are covered,
including the residual-atom evidence-normalization boundary and the trace-fiber
factorization boundary for observer capture, and the atomic evidence-budget
multiplicity boundary; the abstract clocked finite-uniform kernel is covered,
the concrete bit-vector exact ZAB ERM family has the clocked payload at budget
\(r+2k+1\), and the bare exact-visible interface cannot imply such payloads for
every family below the full visible surface.  The same theorem also records that
the broad
\texttt{kpolyCompressionBridge} remains uncovered and that the current packet is
not stop-grade.
The same synthesis packet now records the concrete SAT-readout repair at the
short-global-decoder entry: for Boolean readouts on the one-variable CNF
\(x\lor\neg x\), arbitrary-output SAT-search correctness is equivalent to
single-message readout constancy, while the nonconstant readout \(x\) itself
has selected-output correctness for both messages and no arbitrary-output
message.  A separate generic anchor records the manuscript-level version: once
the CNF extraction is complete and supported satisfiability is present,
supported arbitrary-output SAT-search correctness is equivalent to the semantic
single-message promise.

The ledger also records the next marginal target rather than leaving the route
choice implicit.  It defines \texttt{currentPNPNextMarginalTarget} as
\texttt{residualSideInformation} and proves that this target is still uncovered
and still belongs to the finite-coin stack's broad uncovered list.  The reason
is conservative: postprocessed-side matching and finite-coin randomized
residuals now have exact subledgers, while the clocked finite-uniform kernel is
already covered.  The new residual collision subledger blocks the first
source-only recoding escape hatch, and the residual resolution-cost subledger
shows that source-determined residuals over an invariant base have zero
resolving mass and no advantage.  But the broad bridge remains open: a useful
retained residual must still either be shown genuinely non-source-derived with
positive resolving support, or must explicitly change the invariant-base
hypotheses being used by the side-channel budget theorem.
The ledger lemma
\[
  \texttt{not\_stopGrade\_of\_not\_covers\_currentPNPNextMarginalTarget}
\]
turns that prioritization into a formal audit rule: any packet that still
misses the selected next target cannot be called stop-grade.  This is
deliberately weaker than claiming all residual repairs are broken; it records
that the current local crux families, even when useful, do not by themselves
close the residual side-information bridge.
The same file now also isolates the stronger promoted packet obtained by
adding the already formalized theorem-backed residual-side-information blocker
as a broad repair-class claim.  It defines
\texttt{currentPNPResidualPromotedPacket}, proves that this packet extends
\texttt{currentPNPLocalCruxPacket}, proves that it covers
\texttt{residualSideInformation}, and records that it still leaves exactly
\texttt{randomizedResidual}, \texttt{approximateDecorrelation}, and
\texttt{kpolyCompressionBridge} uncovered.  Accordingly
\texttt{currentPNPResidualPromotedNextMarginalTarget} is
\texttt{randomizedResidual}, and
\texttt{currentPNPResidualPromotedStatus} shows that even this promoted packet
is not stop-grade.  This does not upgrade the baseline current packet; it
records that once the existing residual obstruction layer is treated as a
broad blocker, the next remaining broad escape route is the randomized
residual lane rather than residual side information itself.
The next broad packet step is now formalized as well.  The same module defines
\texttt{RandomizedResidualCoverage} as the theorem-backed statement that
finite-coin randomized residuals stay under the ordinary resolving-mass budget,
supportwise-unresolved randomized residuals have zero joint resolving mass and
no positive doubled advantage, any surviving positive advantage already yields a
deterministic positive-weight coin slice with positive ordinary resolving mass,
and any strict half-accuracy advantage already collapses to an ordinary
residual-side-information obstruction package on such a deterministic slice.
It then defines \texttt{currentPNPRandomizedResidualPromotedPacket}, proves that
this packet extends \texttt{currentPNPResidualPromotedPacket}, and proves that
it covers both \texttt{residualSideInformation} and
\texttt{randomizedResidual}.  What remains open at that point is exactly
\texttt{approximateDecorrelation} together with the manuscript-specific
\texttt{kpolyCompressionBridge}; accordingly
\texttt{currentPNPRandomizedResidualPromotedNextMarginalTarget} is
\texttt{approximateDecorrelation}, and
\texttt{currentPNPRandomizedResidualPromotedStatus} shows that even this
stronger packet is still not stop-grade.  This is again a promoted extension
packet, not a change to the baseline current packet.
The next broad packet step is now formalized too.  The same module defines
\texttt{ApproximateDecorrelationCoverage} as the theorem-backed statement that
below the natural source/coin threshold, approximate independence of a
postprocessed retained-side observation is already equivalent to an exact
preserving true/false coin equivalence, while the concrete three-coin skew
witness still has two-sided pointwise collisions and a one-sided preserving map
but fails approximate independence for every tolerance \(\tau<3\).  It then
defines \texttt{currentPNPApproximateDecorrelationPromotedPacket}, proves that
this packet extends \texttt{currentPNPRandomizedResidualPromotedPacket}, and
proves that it covers \texttt{residualSideInformation},
\texttt{randomizedResidual}, and \texttt{approximateDecorrelation}.  What
remains open at that point is exactly the manuscript-specific
\texttt{kpolyCompressionBridge}; accordingly
\texttt{currentPNPApproximateDecorrelationPromotedNextMarginalTarget} is
\texttt{kpolyCompressionBridge}, and
\texttt{currentPNPApproximateDecorrelationPromotedStatus} shows that even this
stronger packet is still not stop-grade.  This is again a promoted extension
packet, not a change to the baseline current packet.
The final broad packet step is now formalized as well.  The same module defines
\texttt{KpolyCompressionBridgeCoverage} as the theorem-backed statement that the
manuscript's remaining product-bound-only and fielded-switching-only
exact-visible bridge schemas, and their bundled clocked finite-learning payload
variants, are exactly finite-image obligations at the same bit budget; that the
nil and Boolean-square one-step clocked bridge schemas are refuted below the
exact-surface cardinality threshold; that truly surjective exact-visible
families refute the product-bound and fielded-switching clocked bridge schemas
below the predictor-cardinality threshold; and that the bundled clocked payload
is equivalent both to finite predictor cover and to the existing exact-visible
compression target.  It then defines
\texttt{currentPNPKpolyCompressionBridgePromotedPacket}, proves that this
packet extends
\texttt{currentPNPApproximateDecorrelationPromotedPacket}, and proves that it
covers \texttt{residualSideInformation}, \texttt{randomizedResidual},
\texttt{approximateDecorrelation}, and \texttt{kpolyCompressionBridge}.  The
theorems \texttt{stopGrade\_currentPNPKpolyCompressionBridgePromotedPacket} and
\texttt{currentPNPKpolyCompressionBridgePromotedStatus} therefore record that
the current broad PNP synthesis ledger is now stop-grade.  This does not claim
a full \(P \neq NP\) proof; it says only that, at the present synthesis layer,
the currently named broad repair classes are all closed by theorem-backed route
blockers.  The same synthesis layer now also packages the final manuscript-
facing exact ERM route as \texttt{CanonicalZABERMRouteCoverage}: once a family
supplies \texttt{CanonicalZABERMRecoveryData}, Lean immediately places it on
the exact-visible compression / finite-image / clocked-payload surface already
audited by this stop-grade packet.  Thus the remaining gap is no longer a
missing small-image lemma for that named route; it is the universality bridge
showing that the manuscript's actual switched family yields
\texttt{CanonicalZABERMRecoveryData} or an equivalent code-realization
interface.

The same module now also states the exact completion criterion for this ledger.
For any packet extending the current local packet,
\[
  \texttt{stopGrade\_iff\_covers\_currentPNP\_gaps\_of\_extends\_current}
\]
proves that the packet is stop-grade if and only if it covers the four named
gap classes above.  Lean also defines the bookkeeping object
\texttt{currentPNPGapCompletedPacket}, obtained by adding exactly the current
gap list to the current local packet, and proves
\[
  \texttt{stopGrade\_currentPNPGapCompletedPacket}.
\]
It also proves
\[
  \texttt{currentPNPGapCompletedPacket\_covers\_currentPNPNextMarginalTarget},
\]
pinning the selected \(Kpoly\) target as one of the explicit obligations of any
ledger-complete packet.
This theorem is only a ledger closure: it does not supply the missing
mathematics for residual side information, randomized residuals, approximate
decorrelation, or the remaining manuscript-specific \(Kpoly\) realization
bridge.  Instead it gives a precise
target for any future route-stopping packet.
Conversely,
\[
  \texttt{not\_stopGrade\_of\_not\_covers\_currentPNP\_gaps\_of\_extends\_current}
\]
records that an extension of the current packet is still not stop-grade if it
fails to cover the current gap list.

This is intentionally not another local counterexample.  Its role is to prevent
overclaiming: the current Lean work blocks several same-route repair families,
but a fair request for the author to stop would require either covering these
remaining repair classes or proving that the manuscript route cannot use them.

\section{Next Attack Surface}

The next natural formal target is the paper's implicit compression claim.  The
first and third interface steps are now isolated in Lean:
\texttt{VisiblePostSwitchSurface.lean} defines the exact manuscript-visible
surface \(u=(z,a,b)\) together with its invariant and fork projections, and
\texttt{ExactSwitchedFamily.lean} packages the exact bit-budget theorem needed
to reuse the same-route ERM bound.  The remaining mathematical burden is now
sharper:

\begin{quote}
After switching, the witness-bit predictor is not merely local; it is simple enough to lie in a
small explicit class.
\end{quote}

To make that rigorous, one would want a theorem of the shape:
\begin{enumerate}[leftmargin=1.5em]
  \item define the exact visible post-switch data carried by the masked Unique-SAT ensemble;
  \item define the induced witness-bit function on that data;
  \item prove that this function belongs to a sharply bounded class
  (for example, bounded-depth circuits of size \(\mathrm{polylog}(m)\), or a code family of
  description length \(\mathrm{polylog}(m)\));
  \item verify that this bound is strong enough to feed the later union-bound and compression steps.
\end{enumerate}

The new module \texttt{ClockedKpolyBridge.lean} separates the last item from
the still-missing representation theorem.  It defines a clocked \(s\)-bit
decoder family and proves the finite counting kernel:
\[
  \texttt{exactRecoveryRate\_ge\_one\_sub\_uniformMissRatio\_pow}.
\]
Thus, once a target is realized by a clocked \(s\)-bit family, the existing
finite-uniform ERM bound applies with failure term
\[
  2^s \cdot \rho_X^m.
\]
On the Boolean domain it also proves the sample-margin form
\[
  \texttt{exactRecoveryRate\_ge\_one\_sub\_bool\_margin\_pow},
\]
which gives recovery rate at least \(1-(1/2)^k\) when the sample length is at
least \(s+k\).  This is not yet the manuscript's \(Kpoly\) bridge: the open
obligation is still to show that the actual switched witness-bit decoder is
realized by such a clocked bit-code family with the claimed bit and clock
budgets.

The follow-up module \texttt{ClockedKpolyTargetInterface.lean} connects this
kernel back to the exact visible post-switch surface.  Its main theorem says
that any already-proved exact visible compression target
\[
  \texttt{ExactVisibleCompressionTarget}\ G\ s
\]
can be repackaged as a clocked \(s\)-bit family and therefore yields
\[
  \texttt{exactRecoveryRate\_ge\_one\_sub\_uniformMissRatio\_pow\_of\_exactVisibleCompressionTarget}.
\]
The bridge is now exact in both directions.  The module
\texttt{ClockedKpolyBridgeEquivalence.lean} proves
\[
  \texttt{exists\_clockedExactVisibleRealization\_iff\_exactVisibleCompressionTarget}.
\]
Thus a clocked realization theorem for the exact post-switch surface is not a
separate way around the compression target.  It is precisely the same
\(s\)-bit \texttt{ExactVisibleCompressionTarget}, with the clock parameter
carried as external bookkeeping.  The companion clock-invariance theorem says
that changing the external clock bound does not alter which families are
realizable by \(s\) program bits.  Therefore the remaining \(Kpoly\) bridge has
to prove actual small-class structure for the switched predictor family; it
cannot be discharged by a new clock wrapper around an otherwise uncompressed
truth-table family.

At the synthesis-ledger level, this moves the abstract
\texttt{clockedFiniteUniformKernel} class to the covered side and leaves only
\texttt{kpolyCompressionBridge} as the manuscript-specific realization theorem
still missing.
So the remaining \(Kpoly\)-side burden is now cleaner:
\begin{enumerate}[leftmargin=1.5em]
  \item prove a manuscript-faithful exact visible compression target for the actual switched family;
  \item then feed that target through the clocked interface to obtain the finite recovery inequality.
\end{enumerate}
What is still not formalized is the first step, namely the actual
representation theorem for the manuscript decoder.

The new module \texttt{ExactVisibleImageBudget.lean} makes that first step even
more concrete.  For an indexed predictor family \(G\), Lean proves
\[
  G.\texttt{HasBitBudget}\ s
  \quad\Longleftrightarrow\quad
  \exists S,\ \bigl(\forall i,\ G_i\in S\bigr)\ \wedge\ |S|\le 2^s.
\]
The exact visible and clocked interfaces inherit the same characterization:
\[
  \texttt{exists\_clockedExactVisibleRealization\_iff\_finitePredictorCover}.
\]
Thus the broad \(Kpoly\) bridge is now reduced to a finite predictor-image
theorem for the actual switched family.  This is useful route-level progress,
not a closure claim: it does not prove the image is small, it states the exact
finite-image fact that any successful small-program repair must supply.
The same file also removes an abstraction escape hatch: a finite predictor
cover is equivalent to a finite set of representative \emph{indices} whose
actual selected predictors cover the whole family.  In Lean this is
\[
  \texttt{finitePredictorCover\_iff\_finiteIndexRepresentativeCover}.
\]
It is also equivalent to a finite quotient-code presentation:
\[
  \texttt{finitePredictorCover\_iff\_finitePredictorQuotient}.
\]
The clocked exact-visible interface now has matching negative forms:
\[
  \texttt{not\_exists\_clockedExactVisibleRealization\_iff\_not\_finiteIndexRepresentativeCover}
\]
and
\[
  \texttt{not\_exists\_clockedExactVisibleRealization\_iff\_not\_finitePredictorQuotient}.
\]
Thus a proposed finite-image proof can be read operationally either as a finite
list of actual switched predictors or as a finite code/decoder quotient, not
merely as an arbitrary external set of Boolean functions.
The file now also proves the basic closure properties needed to use that
formulation compositionally.  Finite predictor-image covers, finite
representative-index covers, and finite quotient-code presentations are
monotone in the allowed budget, pull back across an explicit factorization
\(G(x)=H(\mathrm{view}(x))\), and push back to the summary family when the view
map has a section.  These are recorded in Lean as the
\texttt{finitePredictorCover\_of\_factorsThrough},
\texttt{finitePredictorCover\_of\_factorsThrough\_of\_rightInverse},
\texttt{finiteIndexRepresentativeCover\_of\_factorsThrough},
\texttt{finiteIndexRepresentativeCover\_of\_factorsThrough\_of\_rightInverse},
\texttt{finitePredictorQuotient\_of\_factorsThrough}, and
\texttt{finitePredictorQuotient\_of\_factorsThrough\_of\_rightInverse}
lemmas.  This is still a preservation theorem, not a small-image theorem: it
keeps any proposed exact-visible compression argument honest about which finite
image bound is being transported.
Combining that transport with reduced-surface surjectivity gives a sharper
pullback obstruction.  If an exact family factors through a finite view, the
view map has a section, and the reduced family is still surjective onto all
Boolean classifiers on that view, then every finite predictor-image cover of
the exact family has size at least the full Boolean classifier-space
cardinality of the view.  The raw \((a,b)\) specialization is now explicit:
if \(G\) factors through \(\texttt{ABVisibleSurface}\ k\) and the reduced
family is still fully expressive, then no finite predictor cover of \(G\) can
have size below
\[
  2^{|\texttt{ABVisibleSurface}\ k|} = 2^{2^{2k}}.
\]
The same statement is proved for finite representative-index covers and finite
quotient-code presentations.  This blocks a concrete finite-image escape hatch:
a route cannot claim a small exact-visible image after quotienting to \((a,b)\)
unless it proves that the reduced predictor family is far smaller than the full
Boolean rule space.
The same module now proves the direct obstruction for this finite-image form:
if \(G\) is still surjective onto all Boolean rules on a finite input surface,
then every finite predictor cover, finite representative-index cover, and finite
quotient-code presentation has size at least
\[
  2^{|\text{surface}|}.
\]
For the exact visible surface this yields
\[
  \texttt{not\_exists\_clockedExactVisibleRealization\_of\_surjective\_predict\_of\_lt\_predictorCard}.
\]
The newest strengthening removes a weaker escape hatch.  Full surjectivity is
not required for the counting lower bound: if the switched family already
realizes an injectively indexed finite probe family \(P\) of classifiers, then
every finite predictor cover, representative-index cover, and quotient-code
presentation has size at least \(|P|\).  Lean records this as
\[
  \texttt{finitePredictorCover\_card\_le\_of\_injective\_realization},
\]
together with the representative and quotient variants.  The exact visible and
clocked negative forms are
\[
  \texttt{not\_exactVisibleCompressionTarget\_of\_injective\_realization\_of\_lt\_card}
\]
and
\[
  \texttt{not\_exists\_clockedExactVisibleRealization\_of\_injective\_realization\_of\_lt\_card}.
\]
Thus the route cannot replace a bit-budget proof by a finite-image cover unless
it also proves the actual predictor image is a strict subclass not merely of the
full Boolean classifier space, but of every large finite probe family it
already realizes.
The factor-map transport theorem now has the same partial-image form.  If an
exact family factors through a reduced view with a section, and the reduced
family realizes an injectively indexed finite probe family \(P\), then every
finite predictor cover, representative-index cover, and quotient-code
presentation of the original family still has size at least \(|P|\).  The
generic Lean anchors are
\[
  \texttt{finitePredictorCover\_card\_le\_of\_factorsThrough\_of\_rightInverse\_of\_injective\_realization}
\]
and the matching representative/quotient variants, with negative forms below
\(|P|\).  Thus a reduced-view repair cannot evade the finite-image audit by
avoiding full surjectivity; any large finite class already realized on the
reduced view is enough to force the same lower bound after pullback.
The concrete \((a,b)\)-reduced route now has the same partial-image
specialization.  If the exact visible family factors through
\(\texttt{abVisibleData}\) and the reduced \((a,b)\) family realizes an
injectively indexed finite probe family \(P\), then the pullback exact family
inherits the lower bounds for finite predictor covers, representative-index
covers, and quotient-code presentations.  The main Lean anchors are
\[
  \texttt{abVisible\_finitePredictorCover\_card\_le\_of\_factorsThrough\_of\_injective\_realization},
\]
\[
  \texttt{not\_exactVisibleCompressionTarget\_of\_factorsThrough\_ab\_and\_injective\_realization\_of\_lt\_card},
\]
and
\[
  \texttt{not\_exists\_clockedExactVisibleRealization\_of\_factorsThrough\_ab\_and\_injective\_realization\_of\_lt\_card}.
\]
So the \((a,b)\) reduction cannot repair the \(Kpoly\) bridge by weakening
``surjective onto all reduced rules'' to an unspecified finite image: every
finite probe family already realized on the raw \((a,b)\) surface still gives
a hard cardinal lower bound after exact-visible pullback.

The new module \texttt{ExactVisibleCompressionObstruction.lean} records the
fallback counting barrier on that exact object: if the switched family on
\(u=(z,a,b)\) is still rich enough to realize \emph{all} Boolean rules on the
exact visible surface, then no \(s\)-bit code can cover it when
\(s < |\text{surface}|\).  So any same-route rescue must now prove a genuine
strict subclass theorem for the \emph{actual} switched predictors.

The new module \texttt{ClockedKpolyCompressionObstruction.lean} states the same
point in the manuscript's remaining form.  It does not only say that a small
exact visible compression target is impossible under full expressivity.  It
also says that, under the same surjectivity hypothesis, there cannot even
exist a clocked \(s\)-bit realization family for the switched predictors when
\[
  s < |\text{ExactVisiblePostSwitchSurface}\ Z\ k|.
\]
So the still-open \(Kpoly\) bridge really is a strict-subclass obligation on
the actual switched family, not an implementation detail about packaging an
already-small family into a clocked decoder interface.

\texttt{ProductSuccessKpolyBridgeObstruction.lean} now rules out a broader
route-level shortcut to that bridge.  The v13 product-success and fixed-field
switching facts are event-history statements; by themselves they do not impose
any finite-image bound on an arbitrary exact-visible predictor family.  Lean
formalizes this with the full family
\[
  \texttt{fullExactVisibleRuleFamily}\ Z\ k,
\]
whose index is the Boolean rule itself and whose prediction map is surjective
onto all exact-visible rules.  Therefore any bridge that tried to derive
\texttt{ExactVisibleCompressionTarget} only from the finite tower product bound
is false below the surface-cardinality threshold:
\[
  \texttt{not\_productBoundOnlyExactVisibleCompressionBridge\_nil\_of\_lt\_surfaceCard}.
\]
The same is true if the bridge uses only a fixed-field v13 certificate:
\[
  \texttt{not\_fieldedSwitchingOnlyExactVisibleCompressionBridge\_nil\_of\_lt\_surfaceCard}.
\]
The obstruction is not an artifact of the empty list.  The Boolean-square
one-step model with the one-cell field has a genuine fielded certificate and
therefore the product bound, but Lean still proves both product-only and
fielded-only bridges false:
\[
  \texttt{not\_productBoundOnlyExactVisibleCompressionBridge\_boolPair\_one\_step\_of\_lt\_surfaceCard},
\]
\[
  \texttt{not\_fieldedSwitchingOnlyExactVisibleCompressionBridge\_boolPair\_one\_step\_of\_lt\_surfaceCard}.
\]
Thus the remaining \(Kpoly\) bridge cannot be a consequence of product success
alone.  It must include a semantic theorem tying the manuscript's switched
event structure to a concrete finite predictor-image cover, representative
index cover, or quotient-code presentation of the exact-visible family.
The module now makes this semantic boundary exact.  It defines
\[
  \texttt{V13FieldedProductBoundFrom}
\]
for the fielded tower product inequality and proves that, for any fixed
exact-visible family \(G\), the implication
\[
  \texttt{V13FieldedProductBoundFrom}\to
  \texttt{ExactVisibleCompressionTarget}\ G\ s
\]
is equivalent to
\[
  G.\texttt{FinitePredictorCover}\ (2^s)
\]
whenever the product premise is true:
\[
  \texttt{productBound\_true\_exactVisibleCompressionBridge\_iff\_finitePredictorCover}.
\]
The fielded-certificate version is identical:
\[
  \texttt{fieldedSwitching\_true\_exactVisibleCompressionBridge\_iff\_finitePredictorCover}.
\]
Consequently any global product-only or fielded-only bridge immediately yields
not only a finite predictor cover, but also a finite representative-index cover
and a finite quotient-code presentation for the selected exact-visible family:
\[
  \texttt{finiteIndexRepresentativeCover\_of\_productBoundOnlyExactVisibleCompressionBridge},
\]
\[
  \texttt{finitePredictorQuotient\_of\_fieldedSwitchingOnlyExactVisibleCompressionBridge}.
\]
This is the current minimal repair target.  A manuscript-side \(Kpoly\) bridge
must prove one of those finite-image objects for the actual switched predictors;
product arithmetic and fielded half-success certificates do not reduce that
semantic burden.
The latest strengthening closes the parallel presentation repair in which the
bridge target is not \texttt{ExactVisibleCompressionTarget} but the bundled
\texttt{ClockedKpolyFiniteLearningPayload}.  Lean defines the corresponding
product-bound-only and fielded-switching-only clocked-payload bridge schemas and
proves that they are still exactly demands for finite predictor-image covers:
\[
  \texttt{productBoundOnlyClockedKpolyFiniteLearningBridge\_iff\_forall\_semanticFiniteImageBridge},
\]
\[
  \texttt{fieldedSwitchingOnlyClockedKpolyFiniteLearningBridge\_iff\_forall\_semanticFiniteImageBridge}.
\]
Moreover any such clocked-payload bridge immediately implies the older
exact-visible compression bridge, so the same full-visible counterexamples
apply.  The empty and Boolean-square one-step versions are recorded as
\[
  \texttt{not\_productBoundOnlyClockedKpolyFiniteLearningBridge\_nil\_of\_lt\_surfaceCard},
\]
\[
  \texttt{not\_fieldedSwitchingOnlyClockedKpolyFiniteLearningBridge\_boolPair\_one\_step\_of\_lt\_surfaceCard}.
\]
For already-surjective exact-visible families, the direct refuter uses the
Boolean predictor-space threshold:
\[
  \texttt{not\_productBoundOnlyClockedKpolyFiniteLearningBridge\_of\_true\_productBound\_of\_surjective\_predict\_of\_lt\_predictorCard},
\]
with the analogous fielded-switching theorem.  Thus moving from compression
language to payload language does not bypass the finite-image audit.

\texttt{ActualSwitchedLocalSupportObstruction.lean} now tests the next natural
repair idea: perhaps the manuscript's actual switched-local interface only
needs to be small on the block outputs that are really observed.  Lean separates
that observed-support statement from the full exact-visible \(Kpoly\) target.
The strengthened version constructs an arbitrary finite-support actual-local
interface
\[
  \texttt{supportFullRuleActualSwitchedLocalInterface}\ \texttt{visibleOf}
\]
whose observed block-output family has truth-table budget bounded by
\(|\texttt{Block}|\), hence a predictor cover of size at most
\(2^{|\texttt{Block}|}\):
\[
  \texttt{supportFullRule\_outputFamily\_finitePredictorCover\_card}.
\]
But the selected local-rule family is still exactly the full Boolean family on
the exact visible post-switch surface.  Therefore, for every
\[
  s < |\text{ExactVisiblePostSwitchSurface}\ Z\ k|,
\]
Lean proves both
\[
  \texttt{supportFullRule\_not\_finitePredictorCover\_of\_lt\_surfaceCard}
\]
and
\[
  \texttt{supportFullRule\_not\_clockedKpolyFiniteLearningPayload\_of\_lt\_surfaceCard}.
\]
The previous one-block theorem is retained as the minimal instance.
The next possible repair is to add a uniform section from every exact-visible
input back to an actual block.  Lean proves that this is exactly the missing
bridge at this abstraction level:
\[
  \texttt{predictorFamily\_factorsThrough\_outputFamily\_of\_uniformVisibleSection},
\]
\[
  \texttt{predictorFamily\_finitePredictorCover\_of\_outputFamily\_finitePredictorCover},
\]
and
\[
  \texttt{clockedKpolyFiniteLearningPayload\_of\_outputFamily\_hasBitBudget}.
\]
So observed-output compression transfers to the full exact-visible \(Kpoly\)
payload if the actual support reaches every exact-visible input uniformly in
the output index.  But for the finite-support full-rule construction, this
section hypothesis forces the cardinal inequality
\[
  |\text{ExactVisiblePostSwitchSurface}\ Z\ k| \le |\texttt{Block}|
\]
via
\[
  \texttt{supportFullRule\_surfaceCard\_le\_blockCard\_of\_uniformVisibleSection}.
\]
Consequently
\[
  \texttt{supportFullRule\_not\_hasUniformVisibleSection\_of\_card\_lt\_surfaceCard}
\]
blocks the repair whenever the actual observed support is smaller than the
full exact-visible surface.

Lean also tests the fallback ``reachable-support quotient'' repair directly.
The observed block-output family is only the pullback of the exact-visible
family along the support map:
\[
  \texttt{supportFullRule\_outputFamily\_factorsThrough\_predictorFamily}.
\]
If \(\texttt{visibleOf}\) misses any exact-visible input, Lean constructs two
distinct full exact-visible rules that agree on every observed block:
\[
  \texttt{supportFullRule\_exists\_distinct\_rules\_same\_outputFamily\_of\_not\_surjective}.
\]
Equivalently, exact recovery from observed block outputs forces the support
map to be surjective:
\[
  \texttt{supportFullRule\_observedRuleMap\_injective\_forces\_visibleOf\_surjective}.
\]
The same obstruction is now stated in the stronger post-processing form.  If a
candidate repair supplies a decoder
\[
  \texttt{decode} : (\texttt{Block}\to\mathsf{Bool})\to
  \texttt{ExactVisibleRule}\ Z\ k
\]
that reconstructs every full exact-visible rule from its observed block trace,
then the support map must again be surjective:
\[
  \texttt{supportFullRule\_exactDecoder\_forces\_visibleOf\_surjective}.
\]
The current Lean packet also blocks the softer reply that the proof only needs
some downstream invariant of the rule, not the whole rule.  Any property
\(P : \texttt{ExactVisibleRule}\ Z\ k \to \alpha\) decoded from observed block
traces must be constant on observed-output fibers:
\[
  \texttt{supportFullRule\_propertyDecoder\_constant\_on\_observedFibers}.
\]
Lean now proves that this condition is also sufficient:
\[
  \texttt{supportFullRule\_exists\_propertyDecoder\_iff\_constant\_on\_observedFibers}.
\]
Thus a reachable-support quotient is not a free replacement for the full-rule
target.  Each later invariant must be proved to factor through the observed
quotient; otherwise the missing off-support values are semantically available
to the full exact-visible target but invisible to the actual block trace.
So any property that separates two full rules with the same observed trace
cannot be recovered:
\[
  \texttt{supportFullRule\_not\_exists\_propertyDecoder\_of\_same\_output\_ne}.
\]
In particular, if an exact-visible input \(u_0\) is outside the reachable
support, even the single bit \(r(u_0)\) is not observable from the block trace:
\[
  \texttt{supportFullRule\_not\_exists\_evalDecoder\_of\_missed\_point}.
\]
The latest point-query boundary is exact, not merely negative.  Lean proves
that a decoder for the single query \(r(u_0)\) exists precisely when \(u_0\)
is in the range of the actual support map:
\[
  \texttt{supportFullRule\_exists\_evalDecoder\_iff\_mem\_range}.
\]
Consequently, supplying decoders for all exact-visible point queries is
equivalent to full reachable-support surjectivity:
\[
  \texttt{supportFullRule\_all\_evalDecoders\_iff\_visibleOf\_surjective}.
\]
The same boundary is now available for an arbitrary family of downstream
point probes \(Q\to\texttt{ExactVisiblePostSwitchSurface}\ Z\ k\).  A decoder
for the whole query vector exists if and only if every queried point is in the
actual reachable support:
\[
  \texttt{supportFullRule\_exists\_queryDecoder\_iff\_forall\_mem\_range}.
\]
If even one requested probe is outside that support, no such query-vector
decoder exists:
\[
  \texttt{supportFullRule\_not\_exists\_queryDecoder\_of\_missed\_query}.
\]
The latest strengthening handles adaptive downstream probing.  Lean now defines
\texttt{AdaptiveBoolQueryTree}, a finite decision tree whose later
exact-visible point queries may depend on earlier rule-bit answers.  If every
point that can be queried by the tree lies in the reachable support, the whole
adaptive computation can be simulated from the observed block trace:
\[
  \texttt{supportFullRule\_exists\_adaptiveQueryDecoder\_of\_allQueriesReachable}.
\]
Conversely, adaptive decoding still factors through exactly the same quotient:
\[
  \texttt{supportFullRule\_exists\_adaptiveQueryDecoder\_iff\_constant\_on\_observedFibers}.
\]
For crux extraction this gives a compact certificate:
\[
  \texttt{supportFullRule\_not\_exists\_adaptiveQueryDecoder\_of\_same\_output\_eval\_ne}.
\]
Any two full rules that agree on every observed block but make the adaptive tree
return different outputs refute all observed-trace decoders for that tree.  This
is the form to use for deeper missed queries, because the witness rules can be
chosen to follow the same earlier branch transcript and then differ at the
off-support query.
In particular, if the root query is off support and the two root branches have
different final outputs, no observed-trace decoder exists:
\[
  \texttt{supportFullRule\_not\_exists\_rootAdaptiveQueryDecoder\_of\_missed\_point\_ne}.
\]
Thus adaptive choice of probes does not repair the support gap.  It can only
use off-support rule bits after a separate theorem proves that those queried
bits are actually reachable, or that the tree's final output is invariant on
the observed quotient.
Therefore, below surface cardinality,
\[
  \texttt{supportFullRule\_not\_observedRuleMap\_injective\_of\_card\_lt\_surfaceCard}
\]
and
\[
  \texttt{supportFullRule\_exists\_distinct\_rules\_same\_outputFamily\_of\_card\_lt\_surfaceCard}
\]
and now also
\[
  \texttt{supportFullRule\_not\_exists\_exactDecoder\_of\_card\_lt\_surfaceCard}
\]
and
\[
  \texttt{supportFullRule\_exists\_unobservable\_eval\_of\_card\_lt\_surfaceCard}
\]
show that the quotient literally forgets off-support bits of the local rule.
Thus observed actual-local support smallness is not enough.  A viable
manuscript bridge must either prove a finite-image theorem for the selected
local rules on the whole exact-visible surface, prove uniform reachability
with support at least the surface size, or explicitly replace the full-rule
recovery/ERM target by a reachable-support quotient and prove that every later
point query or invariant is supported by, or constant on, that weaker quotient
target.
The plug-in endpoint repair is now tested separately from the observed-support
quotient.  A natural response to the \(Kpoly\) gap is that the actual local
rule is not an arbitrary full rule, but the result of a finite-alphabet
plug-in lookup or plug-in majority procedure.  Lean shows that this phrasing
does not by itself reduce the image.  The literal lookup endpoint in
\texttt{ActualSwitchedLocalPluginLookupObstruction.lean} is just the full
Boolean lookup-table family:
\[
  \texttt{pluginLookupActualSwitchedLocalInterface\_surjective}.
\]
Consequently it has no clocked finite-learning payload below the full
truth-table budget, and it cannot be identified with bounded ZAB
decision-list constructor data:
\[
  \texttt{pluginLookupActualSwitchedLocalInterface\_not\_clockedKpolyFiniteLearningPayload\_of\_lt\_surfaceCard},
\]
\[
  \texttt{pluginLookupActualSwitchedLocalInterface\_not\_zabDecisionListData\_of\_lt\_surfaceCard}.
\]
The new file
\texttt{ActualSwitchedLocalPluginLookupSparseThresholdERMObstruction.lean}
pushes the same literal lookup endpoint against the later sparse-threshold
route itself.  Since unrestricted lookup is still surjective on the full
exact-visible rule family, any hypothetical shared exact sparse-threshold ERM
packet or recovery packet already inherits the unconditional visible-budget
inequality
\[
  n \le 2r + 2k + 1
\]
on \(\texttt{BitVec}\ n\).  The regression file instantiates this at
\(\texttt{BitVec}\ 2\), \(k = 0\), and extractor width \(r = 0\), so Lean
forces \(2 \le 1\) and rules out both the exact-family and recovery packets on
the manuscript's literal finite-alphabet lookup endpoint.
The counted-majority version is no better without a further structural
restriction.  \texttt{PluginMajorityCounts.ofRule} puts one vote on the desired
side at every exact-visible input, so
\[
  \texttt{pluginMajorityActualSwitchedLocalInterface\_surjective}
\]
and the same below-surface payload and ZAB-constructor refuters follow.
The follow-up file
\texttt{ActualSwitchedLocalPluginMajoritySparseThresholdERMRecoveryHeavyRegionCharacterization.lean}
closes the same later repair at the counted-majority level.  Since unrestricted
counted plug-in majority is still surjective on the full exact-visible rule
family, Lean now proves that any hypothetical sparse-threshold recovery packet
already forces
\[
  1 - \mu(\texttt{PMF.lightestPoint}\ \mu) \le q.
\]
On the skew \(\texttt{BitVec 1}\) witness \((4/5,1/5)\), the regression file
derives the concrete threshold
\[
  4/5 \le q,
\]
and therefore rules out the manuscript-style threshold \(q = 3/4\) before any
sample-budget or injective-branch argument is used.
The sample-level version removes the last abstraction: the graph sample of a
rule contains one labeled example \((u,r(u))\) for every exact-visible input
and plug-in majority recovers \(r\) pointwise:
\[
  \texttt{PluginSampleMajority.predict\_graphSample}.
\]
Thus
\[
  \texttt{pluginSampleMajorityActualSwitchedLocalInterface\_surjective}
\]
and the corresponding negative \(Kpoly\) payload and ZAB-constructor theorems
hold.
The new file
\texttt{ActualSwitchedLocalPluginSampleMajoritySparseThresholdERMObstruction.lean}
pushes that same endpoint against the later manuscript-facing shared
sparse-threshold ERM wrapper.  Since unrestricted sample-level plug-in majority
is still surjective on the full exact-visible rule family, Lean now proves that
it cannot carry
\[
  \texttt{ActualSwitchedLocalInterface.SharedExactZABSparseThresholdERMData}
\]
or the stronger recovery packet below the same point-block visible-budget
threshold.  On \(\texttt{BitVec}\ n\), any such identification would force
\[
  n \le 2r + 2k + 1.
\]
So ``finite sample plus majority'' is still not enough to enter the already
analyzed sparse-threshold ERM route: the manuscript would need a separate
theorem showing that its sample-majority endpoint really lands in that class.
The follow-up file
\texttt{ActualSwitchedLocalPluginSampleMajoritySparseThresholdERMRecoveryHeavyRegionCharacterization.lean}
closes the next obvious repair.  It no longer argues by visible-budget
shortage.  Instead it uses the same unrestricted surjectivity to transfer the
later recovery-side heavy-region theorem directly to sample-level plug-in
majority:
\[
  \texttt{pluginSampleMajorityActualSwitchedLocalInterface\_one\_sub\_apply\_lightestPoint\_le\_of\_nonempty\_sharedExactZABSparseThresholdERMRecoveryData}.
\]
So any hypothetical recovery packet on unrestricted sample-level plug-in
majority already forces
\[
  1 - \mu(\texttt{PMF.lightestPoint}\ \mu) \le q.
\]
On the same skew \(\texttt{BitVec 1}\) witness used for the later full-rule and
fallback counterexamples, the regression file proves this means
\[
  4/5 \le q,
\]
and therefore rules out the manuscript-style threshold \(q = 3/4\) even before
any bounded-code or injective-index argument is invoked.
Finally, bounded sample size gives an exact threshold rather than a free
compression theorem.  In
\texttt{ActualSwitchedLocalBoundedSampleMajorityObstruction.lean}, Lean proves
\[
  \texttt{boundedSamplePluginMajorityActualSwitchedLocalInterface\_surjective\_iff}.
\]
The endpoint is full-rule expressive exactly when the sample bound reaches
\(|\text{ExactVisiblePostSwitchSurface}\ Z\ k|\).  If the bound is below that
surface size, the all-false rule is not realizable because some input is
omitted and the tie-breaking majority returns \texttt{true}.  If the bound
reaches the surface size, the graph-sample construction restores full
expressivity, and Lean again proves the negative clocked-payload and
ZAB-constructor statements.

The bounded side has a sharper invariant than the all-false witness.  Lean now
defines
\[
  \texttt{PluginSampleMajority.falseSupport}
\]
and proves
\[
  \texttt{PluginSampleMajority.falseSupport\_predict\_card\_le\_bound},
\]
because any input omitted from the sample has zero true and false votes and
therefore returns the tie-breaking value \texttt{true}.  At the actual-local
endpoint this becomes
\[
  \texttt{boundedSamplePluginMajorityActualSwitchedLocalInterface\_falseSupport\_card\_le\_sampleBound}
\]
and
\[
  \texttt{boundedSamplePluginMajorityActualSwitchedLocalInterface\_not\_realizes\_rule\_of\_sampleBound\_lt\_falseSupport\_card}.
\]
So a bounded plug-in majority endpoint can realize only rules whose false
support has cardinality at most the sample bound.  Therefore ``finite plug-in
majority'' is not a manuscript-side repair by itself.  The missing theorem
must show either that the intended target rules have small false support under
the chosen tie-breaking convention, or that the selected samples/rules are
forced into one of the already formalized bounded ZAB/code families.

The same obstruction is now tie-break invariant.  Lean defines
\[
  \texttt{PluginSampleMajority.predictWithDefault}
\]
and
\[
  \texttt{PluginSampleMajority.nondefaultSupport}
\]
and proves that omitted inputs return the chosen default:
\[
  \texttt{PluginSampleMajority.predictWithDefault\_eq\_default\_of\_not\_mem\_map\_fst}.
\]
Consequently
\[
  \texttt{PluginSampleMajority.nondefaultSupport\_predictWithDefault\_card\_le\_bound}
\]
and, at the actual-local endpoint,
\[
  \texttt{boundedSamplePluginMajorityWithDefaultActualSwitchedLocalInterface\_nondefaultSupport\_card\_le\_sampleBound}.
\]
Any target whose nondefault support exceeds the sample bound is not realized,
and the exact surjectivity threshold is unchanged for every default:
\[
  \texttt{boundedSamplePluginMajorityWithDefaultActualSwitchedLocalInterface\_surjective\_iff}.
\]
Thus changing the plug-in majority tie-break from \texttt{true} to
\texttt{false}, or to any explicitly fixed Boolean default, is not a repair.

The remaining obvious tie-break escape is to let the fallback depend on the
visible input.  Lean now makes this case explicit rather than leaving it as a
named repair theorem.  It defines
\[
  \texttt{PluginSampleMajority.predictWithFallback}
\]
and the indexed endpoint
\[
  \texttt{boundedSamplePluginMajorityWithFallbackActualSwitchedLocalInterface}.
\]
The theorem
\[
  \texttt{boundedSamplePluginMajorityWithFallbackActualSwitchedLocalInterface\_emptySample\_predict}
\]
proves that the empty sample already returns the fallback rule exactly, and
\[
  \texttt{boundedSamplePluginMajorityWithFallbackActualSwitchedLocalInterface\_surjective}
\]
therefore proves full exact-visible rule expressivity for every sample bound,
including zero.  Consequently,
\[
  \texttt{boundedSamplePluginMajorityWithFallbackActualSwitchedLocalInterface\_not\_finitePredictorCover\_of\_lt\_surfaceCard}
\]
and
\[
  \texttt{boundedSamplePluginMajorityWithFallbackActualSwitchedLocalInterface\_not\_clockedKpolyFiniteLearningPayload\_of\_lt\_surfaceCard}
\]
hold below the exact-visible truth-table budget.  This closes the fallback
repair fork sharply: an input-dependent tie-break is not a bounded-sample
majority improvement, but a full lookup table inserted into the index.

Lean also isolates the more disciplined fallback repair where the fallback is
not arbitrary, but selected from an explicit family.  It defines
\[
  \texttt{PluginSampleMajority.disagreementSupport},
\]
\[
  \texttt{BoundedSampleFallbackFamilyIndex},
\]
and the endpoint
\[
  \texttt{boundedSamplePluginMajorityWithFallbackFamilyActualSwitchedLocalInterface}.
\]
The theorem
\[
  \texttt{boundedSamplePluginMajorityWithFallbackFamilyActualSwitchedLocalInterface\_disagreementSupport\_card\_le\_sampleBound}
\]
proves that the selected predictor can differ from its selected fallback rule
on at most \(\texttt{sampleBound}\) exact-visible inputs.  The theorem
\[
  \texttt{boundedSamplePluginMajorityWithFallbackFamilyActualSwitchedLocalInterface\_not\_realizes\_rule\_of\_forall\_sampleBound\_lt\_disagreementSupport\_card}
\]
therefore blocks realization of any target rule that lies more than
\(\texttt{sampleBound}\) points away from every fallback code.  Thus a
structured fallback repair has been reduced to a genuine finite Hamming-cover
obligation; bounded sampling alone only supplies sparse overrides around the
chosen fallback.

This reduction is now exact, not merely one-sided.  Lean defines a finite
support sample constructor
\[
  \texttt{PluginSampleMajority.supportSample}
\]
and proves
\[
  \texttt{PluginSampleMajority.predictWithFallback\_supportSample\_eq\_of\_eq\_off\_support}.
\]
The endpoint theorem
\[
  \texttt{boundedSamplePluginMajorityWithFallbackFamilyActualSwitchedLocalInterface\_realizes\_rule\_iff\_mem\_fallbackFamilyRadiusCover}
\]
then characterizes the predictor image exactly: a target rule is realized iff
it belongs to the radius-\(\texttt{sampleBound}\) Hamming cover around the
fallback family.  Thus there is no hidden extra expressive channel in the
bounded sample part of the construction; the entire repair burden is the
finite Hamming-cover burden.

Lean now records the corresponding exact repair criterion.  The theorem
\[
  \texttt{boundedSamplePluginMajorityWithFallbackFamilyActualSwitchedLocalInterface\_surjective\_iff\_fallbackFamilyRadiusCover\_eq\_univ}
\]
proves that the structured fallback endpoint is surjective exactly when the
fallback-family radius cover is the full exact-visible Boolean rule cube.  The
pointwise theorem
\[
  \texttt{boundedSamplePluginMajorityWithFallbackFamilyActualSwitchedLocalInterface\_surjective\_iff\_forall\_exists\_disagreementSupport\_card\_le\_sampleBound}
\]
unpacks the same condition: every exact-visible rule must be within Hamming
distance \(\texttt{sampleBound}\) of some fallback code.  This is the minimal
repair theorem for this fork.  A proposed structured fallback repair must
actually prove that cover; otherwise the exact image theorem supplies a target
rule outside the interface image.

The exact repair criterion now has its direct numerical lower-bound form too.
The reusable finite-cover theorem
\[
  \texttt{PluginSampleMajority.boolClassifierSpace\_card\_le\_code\_mul\_smallSubsets\_card\_of\_radiusCover\_eq\_univ}
\]
proves that any fallback-family radius cover equal to the whole Boolean rule
cube must have
\[
  2^{|\texttt{surface}|}
    \le
  |\texttt{FallbackIndex}| \cdot
    |\texttt{smallSubsets}(\texttt{surface},\texttt{sampleBound})|.
\]
The endpoint theorem
\[
  \texttt{boundedSamplePluginMajorityWithFallbackFamilyActualSwitchedLocalInterface\_boolClassifierSpace\_card\_le\_code\_mul\_smallSubsets\_card\_of\_surjective}
\]
transfers this bound to the structured fallback interface.  Thus the repair
burden is not merely qualitative coverage: full expressivity forces a concrete
cardinality lower bound on fallback codes and sparse override supports.

The Hamming-cover obligation is now counted explicitly.  Lean defines
\[
  \texttt{PluginSampleMajority.smallSubsets}
\]
for supports of size at most the sample radius and
\[
  \texttt{PluginSampleMajority.fallbackFamilyRadiusCover}
\]
for the union of Hamming balls around the fallback family.  The theorem
\[
  \texttt{PluginSampleMajority.fallbackFamilyRadiusCover\_card\_le}
\]
proves that this union has size at most
\[
  |\texttt{FallbackIndex}| \cdot
    |\texttt{smallSubsets}(\texttt{surface},\texttt{sampleBound})|.
\]
Consequently the endpoint theorem
\[
  \texttt{boundedSamplePluginMajorityWithFallbackFamilyActualSwitchedLocalInterface\_finitePredictorCover}
\]
gives an explicit predictor-image cover of at most that size, and
\[
  \texttt{boundedSamplePluginMajorityWithFallbackFamilyActualSwitchedLocalInterface\_not\_surjective\_of\_code\_mul\_smallSubsets\_card\_lt\_boolClassifierSpace}
\]
proves non-surjectivity whenever this product is below the full exact-visible
truth-table space.  This blocks the structured fallback repair unless the
fallback family is large enough, together with radius-\(\texttt{sampleBound}\)
sparse overrides, to cover the Boolean cube.

Lean now specializes this burden to fallback families indexed by a fixed bit
budget.  The endpoint
\[
  \texttt{boundedSamplePluginMajorityWithBitFallbackFamilyActualSwitchedLocalInterface}
\]
uses a universe-lifted \(\texttt{BitCode fallbackBits}\) index, so the statement
still applies to arbitrary finite hidden alphabets.  The theorem
\[
  \texttt{boundedSamplePluginMajorityWithBitFallbackFamilyActualSwitchedLocalInterface\_boolClassifierSpace\_card\_le\_bitCode\_mul\_smallSubsets\_card\_of\_surjective}
\]
proves that full-rule expressivity of a \(\texttt{fallbackBits}\)-bit fallback
family forces
\[
  2^{|\texttt{surface}|}
    \le
  2^{\texttt{fallbackBits}} \cdot
    |\texttt{smallSubsets}(\texttt{surface},\texttt{sampleBound})|.
\]
Conversely,
\[
  \texttt{boundedSamplePluginMajorityWithBitFallbackFamilyActualSwitchedLocalInterface\_not\_surjective\_of\_bitCode\_mul\_smallSubsets\_card\_lt\_boolClassifierSpace}
\]
proves non-surjectivity under the strict opposite inequality.  This moves the
repair obligation from an abstract finite fallback-family size to an explicit
code-budget condition: a proposed bit-coded fallback mechanism must pay enough
bits, plus enough sampled sparse overrides, to cover the full exact-visible
Boolean rule cube.

Lean now expands the sparse-override factor exactly.  The theorem
\[
  \texttt{PluginSampleMajority.smallSubsets\_eq\_biUnion\_powersetCard\_range}
\]
partitions the radius-\(r\) support family by support cardinality, and
\[
  \texttt{PluginSampleMajority.card\_smallSubsets\_eq\_sum\_range\_choose}
\]
proves
\[
  |\texttt{smallSubsets}(\texttt{surface},r)|
    =
  \sum_{i=0}^{r} \binom{|\texttt{surface}|}{i}.
\]
Consequently,
\[
  \texttt{boundedSamplePluginMajorityWithBitFallbackFamilyActualSwitchedLocalInterface\_boolClassifierSpace\_card\_le\_bitCode\_mul\_sum\_choose\_of\_surjective}
\]
specializes the repair burden to
\[
  2^{|\texttt{surface}|}
    \le
  2^{\texttt{fallbackBits}}
    \sum_{i=0}^{\texttt{sampleBound}}
      \binom{|\texttt{surface}|}{i}.
\]
The companion theorem
\[
  \texttt{boundedSamplePluginMajorityWithBitFallbackFamilyActualSwitchedLocalInterface\_not\_surjective\_of\_bitCode\_mul\_sum\_choose\_lt\_boolClassifierSpace}
\]
blocks full expressivity under the strict opposite inequality.  This makes the
remaining escape route quantitatively explicit: for sublinear sampled override
radius, the repair must still justify how a binomial-prefix Hamming ball times
the fallback code family covers an exponential truth-table cube.

The exact binomial-prefix count now has its own additive bit-budget
obstruction, not merely a coarse polynomial corollary.  If a proposed route
certifies
\[
  \sum_{i=0}^{\texttt{sampleBound}}
    \binom{|\texttt{surface}|}{i}
      \le 2^{\texttt{overheadBits}},
\]
then
\[
  \texttt{boundedSamplePluginMajorityWithBitFallbackFamilyActualSwitchedLocalInterface\_fallbackBits\_add\_overheadBits\_ge\_surfaceCard\_of\_sumChooseEnvelope\_le\_twoPow\_surjective}
\]
proves
\[
  |\texttt{surface}|
    \le
  \texttt{fallbackBits}+\texttt{overheadBits}.
\]
The matching theorem
\[
  \texttt{boundedSamplePluginMajorityWithBitFallbackFamilyActualSwitchedLocalInterface\_not\_surjective\_of\_sumChooseEnvelope\_le\_twoPow\_and\_fallbackBits\_add\_overheadBits\_lt\_surfaceCard}
\]
blocks full-rule expressivity under the same exact envelope assumption and a
strict total-bit deficit.  This closes the natural repair ``use the exact
Hamming-ball size instead of the polynomial envelope'': the exact count may be
smaller, but once it is represented as an overhead bit budget, the same
truth-table lower bound applies.

The exact support count is also now strict below full radius.  Lean proves
\[
  \texttt{PluginSampleMajority.smallSubsets\_ssubset\_univ\_of\_radius\_lt\_card}
\]
and
\[
  \texttt{PluginSampleMajority.card\_smallSubsets\_lt\_twoPow\_card\_of\_radius\_lt\_card},
\]
so if \(r<|\texttt{surface}|\), then the radius-\(r\) Hamming-ball support
family is a proper subfamily of all finite supports and has cardinality
strictly below \(2^{|\texttt{surface}|}\).  Equivalently,
\[
  \texttt{PluginSampleMajority.sum\_range\_choose\_lt\_twoPow\_of\_lt}
\]
proves
\[
  \sum_{i=0}^{r}\binom{n}{i}<2^n
  \qquad\text{whenever } r<n.
\]
The route-level consequence is the zero-fallback blocker
\[
  \texttt{boundedSamplePluginMajorityWithBitFallbackFamilyActualSwitchedLocalInterface\_not\_surjective\_zeroFallbackBits\_of\_sampleBound\_lt\_surfaceCard}.
\]
With only one structured fallback rule available, bounded sampled overrides
below the visible alphabet size cover a proper Hamming ball, not the full
Boolean rule cube.

The same formalization now proves that this boundary is sharp rather than an
artifact of the negative counting proof.  The generic repair-side theorem
\[
  \texttt{boundedSamplePluginMajorityWithFallbackFamilyActualSwitchedLocalInterface\_surjective\_of\_nonempty\_of\_surfaceCard\_le\_sampleBound}
\]
shows that any nonempty fallback family becomes full-rule expressive once the
sample bound reaches the exact-visible alphabet size.  Its bit-coded
specialization is
\[
  \texttt{boundedSamplePluginMajorityWithBitFallbackFamilyActualSwitchedLocalInterface\_surjective\_of\_surfaceCard\_le\_sampleBound}.
\]
Together with the strict-radius blocker, Lean derives the exact zero-fallback
dichotomy
\[
  \texttt{boundedSamplePluginMajorityWithBitFallbackFamilyActualSwitchedLocalInterface\_zeroFallbackBits\_surjective\_iff\_surfaceCard\_le\_sampleBound}.
\]
Thus the zero-fallback repair cannot be rescued by a clever tie-breaking
choice: with no fallback bits, full expressivity is equivalent to allowing
enough sampled overrides to touch every exact-visible input.

The next file,
\texttt{ActualSwitchedLocalBitFallbackWrapperSparseThresholdERMObstruction.lean},
packages the corresponding manuscript-facing consequence for the later shared
sparse-threshold ERM wrapper.  If the fallback decoder is already surjective
onto all exact-visible Boolean rules, then the empty sample leaves that
fallback rule unchanged, so the wrapped endpoint is surjective for every sample
bound.  Lean therefore transfers the same sparse-threshold visible-budget
obstruction directly to this repair path: such a wrapped endpoint cannot carry
\[
  \texttt{ActualSwitchedLocalInterface.SharedExactZABSparseThresholdERMData}
\]
or the stronger recovery packet below the same point-block threshold, and on
\(\texttt{BitVec}\ n\) any such identification would still force
\[
  n \le 2r + 2k + 1.
\]
The regression file
\texttt{ActualSwitchedLocalBitFallbackWrapperSparseThresholdERMObstructionRegression.lean}
checks this on the one-sample wrapped endpoint with the canonical
truth-table fallback decoder.  Even after adding that full-rule fallback side
channel, \(\texttt{BitVec}\ 2, k=0, r=0\) still refutes the sparse-threshold
ERM identification.  So ``majority plus a full-rule fallback decoder'' is not
a new sparse-threshold repair; it is just another surjective actual-local
endpoint.

The follow-up file
\texttt{ActualSwitchedLocalBitFallbackWrapperSparseThresholdERMRecoveryHeavyRegionCharacterization.lean}
shows that the same fallback side channel also fails at the later
\(q\)-threshold recovery stage.  Once the fallback decoder is surjective, the
wrapped endpoint is still surjective, so Lean transfers the recent surjective
heavy-region characterization directly through the wrapper.  Any hypothetical
recovery packet on that endpoint already forces
\[
  1 - \mu(\texttt{lightestPoint}) \le q.
\]
This is not another visible-width or truth-table-budget wrapper: it is the
same intrinsic lightest-point threshold now stated on the fallback-side repair
path itself.

The synthesis layer now cashes those two axes out together on one manuscript
endpoint instead of leaving them as separate local facts.  In
\texttt{CruxSynthesis.lean}, the theorem
\[
  \texttt{kpolyCoverage\_anchor\_boundedSampleMajorityWithBitFallbackFamily\_visibleWidth\_and\_lightestPoint\_threshold\_of\_nonempty\_recovery\_of\_exactVisibleRuleDecode}
\]
states that any hypothetical recovery packet on the bounded-sample
bit-fallback endpoint with the canonical truth-table decoder must satisfy both
\[
  n \le 2r + 2k + 1
  \qquad\text{and}\qquad
  1 - \mu(\texttt{lightestPoint}) \le q.
\]
Its companion theorem
\[
  \texttt{kpolyCoverage\_anchor\_boundedSampleMajorityWithBitFallbackFamily\_not\_nonempty\_recovery\_of\_exactVisibleRuleDecode\_of\_visibleWidth\_gap\_or\_lt\_one\_sub\_apply\_lightestPoint}
\]
packages the contrapositive: violating either inequality already rules out the
same recovery packet.  This is the first explicit two-axis composition here on
one exact-visible bounded-sample endpoint, rather than another one-axis
refinement.

The regression file
\texttt{ActualSwitchedLocalBitFallbackWrapperSparseThresholdERMRecoveryHeavyRegionCharacterizationRegression.lean}
uses the one-sample wrapped endpoint with the canonical truth-table fallback
decoder on the skew \(\texttt{BitVec}\ 1\) measure \((4/5,1/5)\).  Lean proves
the theorem-shape canary \(1-\mu(\texttt{lightestPoint}) \le q\), derives the
concrete threshold \(4/5 \le q\), and therefore rules out \(q = 3/4\) even
after granting both sampled overrides and a full-rule fallback decoder.  So
that fallback side channel does not escape the later recovery obstruction
either.

This repair path is now also promoted into the synthesis ledger rather than
left only in the local obstruction file.  \texttt{CruxSynthesis.lean} now
records a large family of separate narrow covered entries for this cluster: the sample-level
majority visible-budget boundary, the bit-fallback visible-budget boundary, the
fallback-side \(q\)-threshold
\texttt{PNPKpolySubrepairClass.actualLocalBoundedSampleMajorityBitFallbackRecoveryLightestPointBoundary},
the canonical exact-decoder full-rule boundary
\texttt{PNPKpolySubrepairClass.actualLocalBoundedSampleMajorityBitFallbackExactDecoderFullRuleBoundary},
the canonical exact-decoder sparse-threshold ERM visible-budget boundary
\texttt{PNPKpolySubrepairClass.actualLocalBoundedSampleMajorityBitFallbackExactDecoderSparseThresholdERMVisibleBudgetBoundary},
the canonical exact-decoder no-extractor sparse-threshold ERM visible-budget boundary
\texttt{PNPKpolySubrepairClass.actualLocalBoundedSampleMajorityBitFallbackExactDecoderNoExtractorSparseThresholdERMVisibleBudgetBoundary},
the radius-zero exact-decoder finite-cover obstruction
\texttt{PNPKpolySubrepairClass.actualLocalBoundedSampleMajorityBitFallbackZeroSampleExactDecoderFinitePredictorCoverObstruction},
the radius-zero exact-decoder representative-cover obstruction
\texttt{PNPKpolySubrepairClass.actualLocalBoundedSampleMajorityBitFallbackZeroSampleExactDecoderFiniteIndexRepresentativeCoverObstruction},
the radius-zero exact-decoder quotient obstruction
\texttt{PNPKpolySubrepairClass.actualLocalBoundedSampleMajorityBitFallbackZeroSampleExactDecoderFinitePredictorQuotientObstruction},
the radius-zero exact-decoder exact-visible compression obstruction
\texttt{PNPKpolySubrepairClass.actualLocalBoundedSampleMajorityBitFallbackZeroSampleExactDecoderExactVisibleCompressionObstruction},
the radius-zero exact-decoder clocked-realization obstruction
\texttt{PNPKpolySubrepairClass.actualLocalBoundedSampleMajorityBitFallbackZeroSampleExactDecoderClockedRealizationObstruction},
the radius-zero exact-decoder clocked-payload obstruction
\texttt{PNPKpolySubrepairClass.actualLocalBoundedSampleMajorityBitFallbackZeroSampleExactDecoderClockedPayloadObstruction},
the canonical exact-decoder visible-width boundary
\texttt{PNPKpolySubrepairClass.actualLocalBoundedSampleMajorityBitFallbackExactDecoderVisibleWidthBoundary},
the canonical exact-decoder lightest-point boundary
\texttt{PNPKpolySubrepairClass.actualLocalBoundedSampleMajorityBitFallbackExactDecoderLightestPointBoundary},
the canonical exact-decoder joint recovery boundary
\texttt{PNPKpolySubrepairClass.actualLocalBoundedSampleMajorityBitFallbackExactDecoderJointRecoveryBoundary},
the canonical exact-decoder no-extractor visible-width boundary
\texttt{PNPKpolySubrepairClass.actualLocalBoundedSampleMajorityBitFallbackExactDecoderNoExtractorVisibleWidthBoundary},
the canonical exact-decoder no-extractor lightest-point boundary
\texttt{PNPKpolySubrepairClass.actualLocalBoundedSampleMajorityBitFallbackExactDecoderNoExtractorLightestPointBoundary},
the canonical exact-decoder no-extractor joint recovery boundary
\texttt{PNPKpolySubrepairClass.actualLocalBoundedSampleMajorityBitFallbackExactDecoderNoExtractorJointRecoveryBoundary},
the canonical exact-decoder clocked-payload obstruction
\texttt{PNPKpolySubrepairClass.actualLocalBoundedSampleMajorityBitFallbackExactDecoderClockedPayloadObstruction},
the polynomial-envelope additive budget boundary
\texttt{PNPKpolySubrepairClass.actualLocalBoundedSampleMajorityBitFallbackPolynomialEnvelopeBudgetBoundary},
the polynomial-envelope surjectivity obstruction
\texttt{PNPKpolySubrepairClass.actualLocalBoundedSampleMajorityBitFallbackPolynomialEnvelopeSurjectivityObstruction},
the factored-overhead additive budget boundary
\texttt{PNPKpolySubrepairClass.actualLocalBoundedSampleMajorityBitFallbackFactorBudgetBoundary},
the factored-overhead surjectivity obstruction
\texttt{PNPKpolySubrepairClass.actualLocalBoundedSampleMajorityBitFallbackFactorSurjectivityObstruction},
the off-by-one successor-factor additive budget boundary
\texttt{PNPKpolySubrepairClass.actualLocalBoundedSampleMajorityBitFallbackSuccessorFactorBudgetBoundary},
the off-by-one successor-factor surjectivity obstruction
\texttt{PNPKpolySubrepairClass.actualLocalBoundedSampleMajorityBitFallbackSuccessorFactorSurjectivityObstruction},
the double-successor-factor additive budget boundary
\texttt{PNPKpolySubrepairClass.actualLocalBoundedSampleMajorityBitFallbackDoubleSuccessorFactorBudgetBoundary},
the double-successor-factor surjectivity obstruction
\texttt{PNPKpolySubrepairClass.actualLocalBoundedSampleMajorityBitFallbackDoubleSuccessorFactorSurjectivityObstruction},
the exact one-sample additive budget boundary
\texttt{PNPKpolySubrepairClass.actualLocalBoundedSampleMajorityBitFallbackOneSampleBudgetBoundary},
the exact one-sample surjectivity obstruction
\texttt{PNPKpolySubrepairClass.actualLocalBoundedSampleMajorityBitFallbackOneSampleSurjectivityObstruction},
the raw one-sample successor-budget boundary
\texttt{PNPKpolySubrepairClass.actualLocalBoundedSampleMajorityBitFallbackOneSampleSuccessorBudgetBoundary},
the raw one-sample successor-budget surjectivity obstruction
\texttt{PNPKpolySubrepairClass.actualLocalBoundedSampleMajorityBitFallbackOneSampleSuccessorSurjectivityObstruction},
the zero-sample additive budget boundary
\texttt{PNPKpolySubrepairClass.actualLocalBoundedSampleMajorityBitFallbackZeroSampleBudgetBoundary},
the zero-sample surjectivity obstruction
\texttt{PNPKpolySubrepairClass.actualLocalBoundedSampleMajorityBitFallbackZeroSampleSurjectivityObstruction},
the zero-sample exact-decoder boundary
\texttt{PNPKpolySubrepairClass.actualLocalBoundedSampleMajorityBitFallbackZeroSampleExactDecoderBoundary},
the zero-sample fallback-equivalence boundary
\texttt{PNPKpolySubrepairClass.actualLocalBoundedSampleMajorityBitFallbackZeroSampleFallbackEquivalence},
the zero-sample finite predictor-cover boundary
\texttt{PNPKpolySubrepairClass.actualLocalBoundedSampleMajorityBitFallbackZeroSampleFinitePredictorCoverBoundary},
the zero-sample finite representative-cover boundary
\texttt{PNPKpolySubrepairClass.actualLocalBoundedSampleMajorityBitFallbackZeroSampleFiniteIndexRepresentativeCoverBoundary},
the zero-sample finite quotient-code boundary
\texttt{PNPKpolySubrepairClass.actualLocalBoundedSampleMajorityBitFallbackZeroSampleFinitePredictorQuotientBoundary},
the zero-sample exact visible compression boundary
\texttt{PNPKpolySubrepairClass.actualLocalBoundedSampleMajorityBitFallbackZeroSampleExactVisibleCompressionBoundary},
the zero-sample clocked realization boundary
\texttt{PNPKpolySubrepairClass.actualLocalBoundedSampleMajorityBitFallbackZeroSampleClockedRealizationBoundary},
the zero-sample clocked-payload boundary
\texttt{PNPKpolySubrepairClass.actualLocalBoundedSampleMajorityBitFallbackZeroSampleClockedPayloadBoundary},
the exact one-sample product boundary
\texttt{PNPKpolySubrepairClass.actualLocalBoundedSampleMajorityBitFallbackOneSampleProductBoundary},
the exact one-sample product surjectivity obstruction
\texttt{PNPKpolySubrepairClass.actualLocalBoundedSampleMajorityBitFallbackOneSampleProductSurjectivityObstruction},
the exact two-sample quadratic product boundary
\texttt{PNPKpolySubrepairClass.actualLocalBoundedSampleMajorityBitFallbackTwoSampleQuadraticProductBoundary},
the exact two-sample quadratic product surjectivity obstruction
\texttt{PNPKpolySubrepairClass.actualLocalBoundedSampleMajorityBitFallbackTwoSampleQuadraticProductSurjectivityObstruction},
the exact two-sample quadratic-envelope budget boundary
\texttt{PNPKpolySubrepairClass.actualLocalBoundedSampleMajorityBitFallbackTwoSampleQuadraticEnvelopeBudgetBoundary},
the exact two-sample quadratic-envelope surjectivity obstruction
\texttt{PNPKpolySubrepairClass.actualLocalBoundedSampleMajorityBitFallbackTwoSampleQuadraticEnvelopeSurjectivityObstruction},
the fallback-side finite predictor-cover obstruction
\texttt{PNPKpolySubrepairClass.actualLocalBoundedSampleMajorityFallbackFinitePredictorCoverObstruction},
the fallback-side finite representative-cover obstruction
\texttt{PNPKpolySubrepairClass.actualLocalBoundedSampleMajorityFallbackFiniteIndexRepresentativeCoverObstruction},
the fallback-side finite quotient-code obstruction
\texttt{PNPKpolySubrepairClass.actualLocalBoundedSampleMajorityFallbackFinitePredictorQuotientObstruction},
the fallback-side exact visible compression obstruction
\texttt{PNPKpolySubrepairClass.actualLocalBoundedSampleMajorityFallbackExactVisibleCompressionObstruction},
the fallback-side clocked realization obstruction
\texttt{PNPKpolySubrepairClass.actualLocalBoundedSampleMajorityFallbackClockedRealizationObstruction},
the fallback-side clocked finite-learning payload obstruction
\texttt{PNPKpolySubrepairClass.actualLocalBoundedSampleMajorityFallbackClockedPayloadObstruction},
the exact fallback-family radius-cover boundary
\texttt{PNPKpolySubrepairClass.actualLocalBoundedSampleMajorityFallbackFamilyRadiusCoverBoundary},
the fallback-family finite predictor-cover boundary
\texttt{PNPKpolySubrepairClass.actualLocalBoundedSampleMajorityFallbackFamilyFinitePredictorCoverBoundary},
the fallback-family radius-cover surjectivity boundary
\texttt{PNPKpolySubrepairClass.actualLocalBoundedSampleMajorityFallbackFamilyRadiusCoverSurjectivityBoundary},
the fallback-family pointwise radius-cover surjectivity boundary
\texttt{PNPKpolySubrepairClass.actualLocalBoundedSampleMajorityFallbackFamilyPointwiseRadiusCoverSurjectivityBoundary},
the fallback-family empty-sample realization boundary
\texttt{PNPKpolySubrepairClass.actualLocalBoundedSampleMajorityFallbackFamilyEmptySampleRealizationBoundary},
the fallback-family surjective-transport boundary
\texttt{PNPKpolySubrepairClass.actualLocalBoundedSampleMajorityFallbackFamilyFallbackSurjectiveBoundary},
the fallback-family large-disagreement-support obstruction
\texttt{PNPKpolySubrepairClass.actualLocalBoundedSampleMajorityFallbackFamilyLargeDisagreementSupportObstruction},
the abstract fallback-family support-product boundary
\texttt{PNPKpolySubrepairClass.actualLocalBoundedSampleMajorityFallbackFamilySmallSubsetsProductBoundary},
the abstract fallback-family support-product surjectivity obstruction
\texttt{PNPKpolySubrepairClass.actualLocalBoundedSampleMajorityFallbackFamilySmallSubsetsProductSurjectivityObstruction},
the generic fallback-family full-radius boundary
\texttt{PNPKpolySubrepairClass.actualLocalBoundedSampleMajorityFallbackFamilyFullRadiusBoundary},
the generic fallback-family zero-sample fallback-equivalence boundary
\texttt{PNPKpolySubrepairClass.actualLocalBoundedSampleMajorityFallbackFamilyZeroSampleFallbackEquivalence},
the generic fallback-family zero-sample finite predictor-cover boundary
\texttt{PNPKpolySubrepairClass.actualLocalBoundedSampleMajorityFallbackFamilyZeroSampleFinitePredictorCoverBoundary},
the generic fallback-family zero-sample finite representative-cover boundary
\texttt{PNPKpolySubrepairClass.actualLocalBoundedSampleMajorityFallbackFamilyZeroSampleFiniteIndexRepresentativeCoverBoundary},
the generic fallback-family zero-sample finite quotient-code boundary
\texttt{PNPKpolySubrepairClass.actualLocalBoundedSampleMajorityFallbackFamilyZeroSampleFinitePredictorQuotientBoundary},
the generic fallback-family zero-sample exact visible compression boundary
\texttt{PNPKpolySubrepairClass.actualLocalBoundedSampleMajorityFallbackFamilyZeroSampleExactVisibleCompressionBoundary},
the generic fallback-family zero-sample clocked realization boundary
\texttt{PNPKpolySubrepairClass.actualLocalBoundedSampleMajorityFallbackFamilyZeroSampleClockedRealizationBoundary},
the generic fallback-family zero-sample clocked-payload boundary
\texttt{PNPKpolySubrepairClass.actualLocalBoundedSampleMajorityFallbackFamilyZeroSampleClockedPayloadBoundary},
the bit-fallback exact radius-cover boundary
\texttt{PNPKpolySubrepairClass.actualLocalBoundedSampleMajorityBitFallbackRadiusCoverBoundary},
the bit-fallback finite predictor-cover boundary
\texttt{PNPKpolySubrepairClass.actualLocalBoundedSampleMajorityBitFallbackFinitePredictorCoverBoundary},
the bit-fallback radius-cover surjectivity boundary
\texttt{PNPKpolySubrepairClass.actualLocalBoundedSampleMajorityBitFallbackRadiusCoverSurjectivityBoundary},
the bit-fallback pointwise radius-cover surjectivity boundary
\texttt{PNPKpolySubrepairClass.actualLocalBoundedSampleMajorityBitFallbackPointwiseRadiusCoverSurjectivityBoundary},
the abstract sparse-support product boundary
\texttt{PNPKpolySubrepairClass.actualLocalBoundedSampleMajorityBitFallbackSmallSubsetsProductBoundary},
the abstract sparse-support product surjectivity obstruction
\texttt{PNPKpolySubrepairClass.actualLocalBoundedSampleMajorityBitFallbackSmallSubsetsProductSurjectivityObstruction},
the generic bit-fallback full-radius boundary
\texttt{PNPKpolySubrepairClass.actualLocalBoundedSampleMajorityBitFallbackFullRadiusBoundary},
the zero-fallback exact-visible deficit obstruction
\texttt{PNPKpolySubrepairClass.actualLocalBoundedSampleMajorityBitFallbackZeroFallbackBitsDeficitObstruction},
the zero-fallback exact-radius boundary
\texttt{PNPKpolySubrepairClass.actualLocalBoundedSampleMajorityBitFallbackZeroFallbackBitsExactRadiusBoundary},
the generic sum-of-binomial product boundary
\texttt{PNPKpolySubrepairClass.actualLocalBoundedSampleMajorityBitFallbackSumChooseProductBoundary},
the generic sum-of-binomial product surjectivity obstruction
\texttt{PNPKpolySubrepairClass.actualLocalBoundedSampleMajorityBitFallbackSumChooseProductSurjectivityObstruction},
the generic sum-of-binomial-envelope budget boundary
\texttt{PNPKpolySubrepairClass.actualLocalBoundedSampleMajorityBitFallbackSumChooseEnvelopeBudgetBoundary},
the generic sum-of-binomial-envelope surjectivity obstruction
\texttt{PNPKpolySubrepairClass.actualLocalBoundedSampleMajorityBitFallbackSumChooseEnvelopeSurjectivityObstruction},
the generic surjective-actual-local joint recovery barrier
\texttt{PNPKpolySubrepairClass.surjectiveActualLocalJointRecoveryBoundary},
the generic surjective-actual-local visible-width threshold
\texttt{PNPKpolySubrepairClass.surjectiveActualLocalRecoveryVisibleWidthBoundary},
the generic surjective-actual-local lightest-point threshold
\texttt{PNPKpolySubrepairClass.surjectiveActualLocalRecoveryLightestPointBoundary},
and its no-extractor corollary
\texttt{PNPKpolySubrepairClass.surjectiveActualLocalNoExtractorLightestPointBoundary},
plus the broader unconditional half-width no-extractor boundary
\texttt{PNPKpolySubrepairClass.surjectiveActualLocalNoExtractorVisibleWidthBoundary},
the observed-support noninjective observed-rule-map bridge obstruction
\texttt{PNPKpolySubrepairClass.supportFullRuleObservedRuleMapNoninjectiveBelowSurfaceCard},
the observed-support distinct-rules-same-output bridge obstruction
\texttt{PNPKpolySubrepairClass.supportFullRuleDistinctRulesSameOutputBelowSurfaceCard},
the observed-support adaptive-query collision obstruction
\texttt{PNPKpolySubrepairClass.supportFullRuleNoAdaptiveQueryDecoderSameOutputEvalBoundary},
the observed-support missed-root adaptive-query obstruction
\texttt{PNPKpolySubrepairClass.supportFullRuleNoRootAdaptiveQueryDecoderMissedPointBoundary},
the observed-support no-exact-decoder bridge obstruction
\texttt{PNPKpolySubrepairClass.supportFullRuleNoExactDecoderBelowSurfaceCard},
and the observed-support unobservable-evaluation obstruction
\texttt{PNPKpolySubrepairClass.supportFullRuleUnobservableEvalBelowSurfaceCard}.
So the covered \texttt{Kpoly} subledger now does not merely remember the
fallback-surjective specializations; it also records both the direct
fallback-side finite-cover and clocked-payload failures, the radius-zero
exact-decoder no-compression packet, the exact-decoder visible-width
threshold, the exact-decoder lightest-point threshold, the exact-decoder
visible-width/lightest-point/joint no-extractor obstruction forms, and the
underlying generic barrier that any surjective actual-local endpoint must obey.
It now also records the generic finite-support/full-rule bridge packet:
\texttt{PNPKpolySubrepairClass.supportFullRuleObservedSmallNotExactVisibleCoverBoundary}
and
\texttt{PNPKpolySubrepairClass.supportFullRuleNotClockedPayloadBelowSurfaceCard}.
It also records the uniform-section surface-cardinality boundary
\texttt{PNPKpolySubrepairClass.supportFullRuleUniformSectionSurfaceCardBoundary}
and the below-surface no-uniform-section obstruction
\texttt{PNPKpolySubrepairClass.supportFullRuleNoUniformSectionBelowSurfaceCard}.
It now also records the positive uniform-section lift packet:
\texttt{PNPKpolySubrepairClass.supportFullRuleUniformSectionFinitePredictorCoverBoundary}
and
\texttt{PNPKpolySubrepairClass.supportFullRuleUniformSectionClockedPayloadBoundary}.
So a finite observed block-output support can still coexist with a full
exact-visible family that has no below-surface predictor cover, and the same
small-support regime still fails every clocked finite-learning payload below
surface cardinality. Any would-be uniform visible section already forces the
observed support cardinality up to the full exact-visible surface size, so a
below-surface support deficit admits no such section at all. And if such a
uniform section did exist, the observed block truth-table cover and the
resulting block-cardinality clocked payload would really transfer back to the
full exact-visible family.
It now also records the generic property-decoder rigidity packet:
\texttt{PNPKpolySubrepairClass.supportFullRulePropertyDecoderFiberConstancyBoundary},
\texttt{PNPKpolySubrepairClass.supportFullRulePropertyDecoderIffFiberConstancyBoundary},
and
\texttt{PNPKpolySubrepairClass.supportFullRuleNoPropertyDecoderSameOutputBoundary}.
So any downstream property recovered from observed block-output traces must be
constant on observed fibers, a property is decodable exactly when that
fiber-constancy holds, and any property separating two full rules with the
same observed trace is therefore undecodable.
It now also records the reachable-support query-decoder packet:
\texttt{PNPKpolySubrepairClass.supportFullRuleEvalDecoderRangeBoundary},
\texttt{PNPKpolySubrepairClass.supportFullRuleAllEvalDecodersSurjectiveBoundary},
\texttt{PNPKpolySubrepairClass.supportFullRuleQueryDecoderRangeBoundary},
and
\texttt{PNPKpolySubrepairClass.supportFullRuleNoQueryDecoderMissedQueryBoundary}.
So a single exact-visible point query is decodable exactly when that point lies
in reachable support, decoding every point query is equivalent to surjectivity
of the reachable-support map, a whole query family is decodable exactly when
all queried points are reachable, and any family containing one missed query is
therefore blocked outright.
It now also records the generic observed-quotient factorization and adaptive
reachability packet:
\texttt{PNPKpolySubrepairClass.supportFullRuleOutputFamilyFactorBoundary},
\texttt{PNPKpolySubrepairClass.supportFullRuleObservedRuleInjectiveSurjectiveBoundary},
\texttt{PNPKpolySubrepairClass.supportFullRuleExactDecoderSurjectiveBoundary},
\texttt{PNPKpolySubrepairClass.supportFullRuleAdaptiveQueryDecoderReachableBoundary},
and
\texttt{PNPKpolySubrepairClass.supportFullRuleAdaptiveQueryDecoderIffFiberConstancyBoundary}.
So the observed block-output quotient is only the pullback of the full
exact-visible family along reachable support, injectivity of that quotient or
the existence of an exact full-rule decoder already forces reachable-support
surjectivity, any adaptive exact-visible query tree with all queried points
reachable is decodable from observed traces, and such a tree is decodable
exactly when its output is constant on observed fibers.
It now also records the exact-visible image-lower-bound packet:
\texttt{PNPKpolySubrepairClass.exactVisibleRepresentativeCoverSurjectivityObstruction},
\texttt{PNPKpolySubrepairClass.exactVisiblePredictorQuotientSurjectivityObstruction},
\texttt{PNPKpolySubrepairClass.exactVisibleClockedRealizationSurjectivityObstruction},
\texttt{PNPKpolySubrepairClass.injectiveFiniteProbeExactVisibleCompressionObstruction},
and
\texttt{PNPKpolySubrepairClass.injectiveFiniteProbeClockedRealizationObstruction}.
So a fully expressive exact-visible family already blocks not just finite
predictor-image covers but also finite representative-index covers, finite
quotient-code presentations, and the clocked exact-visible realization target
below the full Boolean image size.  The payload-side obstruction is now also
split out explicitly as
\texttt{PNPKpolySubrepairClass.exactVisibleClockedFiniteLearningPayloadSurjectivityObstruction}
and
\texttt{PNPKpolySubrepairClass.injectiveFiniteProbeClockedFiniteLearningPayloadObstruction},
so a fully expressive exact-visible family already blocks the bundled clocked
finite-learning payload below the full Boolean image size, and any injectively
realized finite probe family already transfers the same payload obstruction
below the realized probe cardinality.  Likewise, any injectively realized
finite probe family already transfers the same obstruction to exact-visible
compression and clocked realization below the realized probe cardinality.
The payload-side lower-bound packet is now also explicit as
\texttt{PNPKpolySubrepairClass.exactVisibleClockedFiniteLearningPayloadSurjectivityLowerBound},
\texttt{PNPKpolySubrepairClass.injectiveFiniteProbeClockedFiniteLearningPayloadLowerBound},
\texttt{PNPKpolySubrepairClass.sectionBackedClockedFiniteLearningPayloadSurjectiveLowerBound},
\texttt{PNPKpolySubrepairClass.sectionBackedClockedFiniteLearningPayloadSurjectiveObstruction},
\texttt{PNPKpolySubrepairClass.sectionBackedInjectiveFiniteProbeClockedFiniteLearningPayloadLowerBound},
and
\texttt{PNPKpolySubrepairClass.sectionBackedInjectiveFiniteProbeClockedFiniteLearningPayloadObstruction}.
So the bundled clocked finite-learning payload itself now inherits the same
positive lower bounds under full exact-visible surjectivity, injectively
realized finite probe families, and section-backed pullbacks of both kinds,
with the corresponding below-cardinality obstructions recorded explicitly for
the two section-backed variants.
It now also records the finite-probe representative/quotient lower-bound and
section-backed pullback packet:
\texttt{PNPKpolySubrepairClass.injectiveFiniteRepresentativeIndexCoverLowerBound},
\texttt{PNPKpolySubrepairClass.injectiveFinitePredictorQuotientLowerBound},
\texttt{PNPKpolySubrepairClass.sectionBackedInjectiveFiniteProbePullbackLowerBound},
\texttt{PNPKpolySubrepairClass.sectionBackedInjectiveFiniteProbePullbackObstruction},
\texttt{PNPKpolySubrepairClass.sectionBackedInjectiveFiniteRepresentativeIndexCoverLowerBound},
\texttt{PNPKpolySubrepairClass.sectionBackedInjectiveFiniteRepresentativeIndexCoverObstruction},
\texttt{PNPKpolySubrepairClass.sectionBackedInjectiveFinitePredictorQuotientLowerBound},
and
\texttt{PNPKpolySubrepairClass.sectionBackedInjectiveFinitePredictorQuotientObstruction}.
So the same injectively realized finite probe family already forces the same
cardinality lower bound for representative-index covers and quotient-code
presentations, and every section-backed pullback of such a reduced view
inherits both the positive lower bounds and the corresponding below-cardinality
obstructions for predictor covers, representative-index covers, and quotient
codes.
It now also records the finite-image equivalence / transport / reduced-\((a,b)\)
pullback packet:
\texttt{PNPKpolySubrepairClass.finiteRepresentativeIndexCoverEquivalence},
\texttt{PNPKpolySubrepairClass.finiteImageCoverBudgetWeakening},
\texttt{PNPKpolySubrepairClass.finitePredictorCoverBudgetWeakeningBoundary},
\texttt{PNPKpolySubrepairClass.finiteRepresentativeIndexCoverBudgetWeakeningBoundary},
\texttt{PNPKpolySubrepairClass.finitePredictorQuotientBudgetWeakeningBoundary},
\texttt{PNPKpolySubrepairClass.finiteImageCoverFactorTransport},
\texttt{PNPKpolySubrepairClass.finitePredictorCoverFactorTransportBoundary},
\texttt{PNPKpolySubrepairClass.finiteRepresentativeIndexCoverFactorTransportBoundary},
\texttt{PNPKpolySubrepairClass.finitePredictorQuotientFactorTransportBoundary},
\texttt{PNPKpolySubrepairClass.sectionBackedFinitePredictorCoverSurjectiveObstruction},
\texttt{PNPKpolySubrepairClass.sectionBackedFiniteRepresentativeIndexCoverSurjectiveObstruction},
\texttt{PNPKpolySubrepairClass.sectionBackedFinitePredictorQuotientSurjectiveObstruction},
\texttt{PNPKpolySubrepairClass.sectionBackedFinitePredictorCoverSurjectiveLowerBound},
\texttt{PNPKpolySubrepairClass.sectionBackedFiniteRepresentativeIndexCoverSurjectiveLowerBound},
\texttt{PNPKpolySubrepairClass.sectionBackedFinitePredictorQuotientSurjectiveLowerBound},
\texttt{PNPKpolySubrepairClass.surjectiveFinitePredictorCoverLowerBound},
\texttt{PNPKpolySubrepairClass.surjectiveFiniteRepresentativeIndexCoverLowerBound},
\texttt{PNPKpolySubrepairClass.surjectiveFinitePredictorQuotientLowerBound},
\texttt{PNPKpolySubrepairClass.finiteImageReducedABPullbackObstruction},
\texttt{PNPKpolySubrepairClass.finiteImageReducedABInjectiveProbePullbackObstruction},
and
\texttt{PNPKpolySubrepairClass.finiteQuotientReducedABPullbackObstruction}.
So finite predictor-image covers, representative-index covers, and
quotient-code presentations are already interchangeable at this layer; the
predictor-cover / representative-cover / quotient-code budget-weakening moves
and factor-transport moves are already internalized separately; section-backed
finite summaries already force the same generic below-cardinality obstruction
for all three finite-image formulations; the same section-backed surjective
summary hypotheses now also yield the matching positive full-Boolean lower
bounds in all three formulations; a fully expressive finite family now also
carries the same generic full-Boolean lower bound in all three formulations;
and the reduced-\((a,b)\) pullback variants are already recognized as the same
below-cardinality obstruction whether stated in predictor-cover,
injective-probe, or quotient-code form.
It now also records the exact-visible compression / realization equivalence
packet:
\texttt{PNPKpolySubrepairClass.exactVisibleCompressionTargetPredictorCoverEquivalence},
\texttt{PNPKpolySubrepairClass.clockedExactVisibleRealizationPredictorCoverEquivalence},
\texttt{PNPKpolySubrepairClass.clockedFiniteLearningPayloadPredictorCoverEquivalence},
\texttt{PNPKpolySubrepairClass.clockedFiniteLearningPayloadExactVisibleCompressionEquivalence},
\texttt{PNPKpolySubrepairClass.exactVisibleCompressionTargetPredictorQuotientEquivalence},
\texttt{PNPKpolySubrepairClass.clockedExactVisibleRealizationPredictorQuotientEquivalence},
\texttt{PNPKpolySubrepairClass.exactVisibleCompressionTargetRepresentativeIndexCoverEquivalence},
\texttt{PNPKpolySubrepairClass.clockedExactVisibleRealizationRepresentativeIndexCoverEquivalence},
\texttt{PNPKpolySubrepairClass.notClockedExactVisibleRealizationRepresentativeIndexCoverEquivalence},
and
\texttt{PNPKpolySubrepairClass.notClockedExactVisibleRealizationPredictorQuotientEquivalence}.
So at the exact-visible layer, compression targets, clocked realization
targets, and the bundled clocked finite-learning payload are already nothing
more than finite predictor-image, quotient-code, or representative-index cover
claims at the same bit budget, and the failure of clocked realization is
already equivalently the failure of those finite-image schemes.
It also records the product-bound / fielded-switching finite-image bridge
packet:
\texttt{PNPKpolySubrepairClass.productBoundBridgeFiniteImageBoundary},
\texttt{PNPKpolySubrepairClass.fieldedSwitchingBridgeFiniteImageBoundary},
\texttt{PNPKpolySubrepairClass.productBoundOnlyFullVisibleObstruction},
\texttt{PNPKpolySubrepairClass.fieldedSwitchingOnlyFullVisibleObstruction},
\texttt{PNPKpolySubrepairClass.productBoundOnlySurjectiveVisibleObstruction},
and
\texttt{PNPKpolySubrepairClass.fieldedSwitchingOnlySurjectiveVisibleObstruction}.
So the manuscript's global product-bound-only and fixed-field-switching-only
exact-visible bridge schemas are already nothing more than finite-image-cover
demands, and those demands are already refuted below surface cardinality both
by the full exact-visible Boolean family and by any already-surjective
exact-visible family.
It now also records that a cardinally small observed block-output quotient
cannot injectively identify the full exact-visible family, already merges two
distinct full exact-visible rules to the same observed output trace, cannot
exactly decode the full exact-visible family, and must leave at least one
exact-visible bit unobservable.
It also records the adaptive-query consequence that any observed-output
collision separating a tree kills every decoder for that tree, and in
particular that missing a queried root point already kills every nontrivial
root adaptive-query decoder.
And it now records the one-block minimal-instance packet:
\texttt{PNPKpolySubrepairClass.oneBlockFullRuleObservedSmallNotExactVisibleCoverBoundary},
\texttt{PNPKpolySubrepairClass.oneBlockFullRuleNotClockedPayloadBelowSurfaceCard},
\texttt{PNPKpolySubrepairClass.oneBlockFullRuleNoExactDecoderOneLtSurfaceCardBoundary},
and
\texttt{PNPKpolySubrepairClass.oneBlockFullRuleUnobservableEvalOneLtSurfaceCardBoundary}.
So the one-block observed quotient still has a trivial finite predictor cover
while the full exact-visible family stays above every \(2^s\) cover below
surface cardinality, already refutes every clocked finite-learning payload
there, and once the exact-visible surface has more than one point admits
neither an exact full-rule decoder nor total observability of every
exact-visible bit.
In particular, the
manuscript's most permissive bounded-sample bit-fallback endpoint with
\texttt{exactVisibleRuleDecode} is already full-rule expressive, and on that
same endpoint the zero-sample wrapper already carries the full exact-visible
no-compression burden, while any hypothetical recovery packet must satisfy
both the visible-width ceiling and the lightest-point threshold because both
of those obstructions are already enforced at the surjective actual-local
layer itself.
It also records the exact fallback-family Hamming-cover image theorem, the
finite predictor-cover certificate, the exact surjectivity-via-full-cover
criterion, the equivalent generic fallback-family pointwise repair criterion,
the generic fallback-family empty-sample realization boundary, the generic
fallback-surjective transport boundary, the generic fallback-family
large-disagreement-support obstruction, the generic zero-sample
fallback-equivalence criterion, and now also the four fallback-side
no-compression consequences themselves (finite representative-cover, finite
quotient-code, exact visible compression, and clocked realization
obstructions), together with the equivalent manuscript-facing pointwise repair
criterion, before descending to
the abstract support-product
obstruction, the generic fallback-family full-radius boundary, the abstract
sparse-support product obstruction, the generic bit-fallback full-radius and
zero-fallback radius boundaries, the radius-zero finite-image and
clocked-realization equivalences, and now also the exact pointwise
radius-cover criteria themselves as covered ledger entries together with the
radius-zero representative-cover and
quotient-code transport forms themselves as covered ledger entries, before the
generic exact support-count product obstruction and the additive-envelope
corollary.  So the one-sample and
two-sample counting packets are now recognized inside Lean as special cases of
an already-ledgered broader Hamming-cover obstruction route, and radius zero
is no longer only a surjectivity canary.

Lean also records the coarse polynomial envelope that will usually be easier to
compose with asymptotic route statements.  The theorem
\[
  \texttt{PluginSampleMajority.sum\_range\_choose\_le\_succ\_mul\_succ\_pow}
\]
proves, for all natural \(n,r\),
\[
  \sum_{i=0}^{r} \binom{n}{i}
    \le
  (r+1)(n+1)^r.
\]
The corresponding support-count theorem
\[
  \texttt{PluginSampleMajority.card\_smallSubsets\_le\_succ\_mul\_succ\_card\_pow}
\]
is valid even when the finite input type is empty, because the base is
\(n+1\).  Therefore full expressivity of a bit-coded fallback endpoint implies
\[
  2^{|\texttt{surface}|}
    \le
  2^{\texttt{fallbackBits}}
    (\texttt{sampleBound}+1)
    (|\texttt{surface}|+1)^{\texttt{sampleBound}},
\]
as formalized by
\[
  \texttt{boundedSamplePluginMajorityWithBitFallbackFamilyActualSwitchedLocalInterface\_boolClassifierSpace\_card\_le\_bitCode\_mul\_sampleBound\_succ\_mul\_surfaceCard\_succ\_pow\_of\_surjective}.
\]
The strict reverse polynomial inequality gives non-surjectivity via
\[
  \texttt{boundedSamplePluginMajorityWithBitFallbackFamilyActualSwitchedLocalInterface\_not\_surjective\_of\_bitCode\_mul\_sampleBound\_succ\_mul\_surfaceCard\_succ\_pow\_lt\_boolClassifierSpace}.
\]
This is not a new analytic lower bound; it is a reusable corollary of the exact
finite Hamming-ball count, stated in a shape that makes any ``bounded sample plus
short fallback code'' repair compare directly against the exponential rule cube.

The same corollary is now packaged as an additive bit-budget obstruction.  If a
proposed route proves that the sparse-override polynomial envelope is bounded by
\(2^{\texttt{overheadBits}}\), then full expressivity forces
\[
  |\texttt{surface}|
    \le
  \texttt{fallbackBits}+\texttt{overheadBits}.
\]
This is formalized by
\[
  \texttt{boundedSamplePluginMajorityWithBitFallbackFamilyActualSwitchedLocalInterface\_fallbackBits\_add\_overheadBits\_ge\_surfaceCard\_of\_polynomialEnvelope\_le\_twoPow\_surjective}.
\]
The strict-deficit form
\[
  \texttt{boundedSamplePluginMajorityWithBitFallbackFamilyActualSwitchedLocalInterface\_not\_surjective\_of\_polynomialEnvelope\_le\_twoPow\_and\_fallbackBits\_add\_overheadBits\_lt\_surfaceCard}
\]
blocks the endpoint when
\[
  (\texttt{sampleBound}+1)
    (|\texttt{surface}|+1)^{\texttt{sampleBound}}
      \le 2^{\texttt{overheadBits}}
  \quad\text{and}\quad
  \texttt{fallbackBits}+\texttt{overheadBits}<|\texttt{surface}|.
\]
Thus a repair cannot hide the sampled override radius in an unspecified
``polynomial factor'': it must exhibit enough overhead bits to pay for that
factor, and the sum of fallback bits plus overhead bits must still reach the
exact-visible alphabet size.

Lean now also closes the common factored-overhead certificate.  The arithmetic
helper
\[
  \texttt{PluginSampleMajority.succ\_mul\_succ\_pow\_le\_twoPow\_add\_mul}
\]
proves that if \(r+1 \le 2^a\) and \(n+1 \le 2^b\), then
\[
  (r+1)(n+1)^r \le 2^{a+br}.
\]
Instantiating \(r=\texttt{sampleBound}\) and
\(n=|\texttt{surface}|\), the endpoint theorem
\[
  \texttt{boundedSamplePluginMajorityWithBitFallbackFamilyActualSwitchedLocalInterface\_fallbackBits\_add\_radiusBits\_add\_baseBits\_mul\_sampleBound\_ge\_surfaceCard\_of\_envelopeFactors\_le\_twoPow\_surjective}
\]
proves that full expressivity forces
\[
  |\texttt{surface}|
    \le
  \texttt{fallbackBits}
    + \texttt{radiusBits}
    + \texttt{baseBits}\cdot\texttt{sampleBound}.
\]
The strict-deficit theorem
\[
  \texttt{boundedSamplePluginMajorityWithBitFallbackFamilyActualSwitchedLocalInterface\_not\_surjective\_of\_envelopeFactors\_le\_twoPow\_and\_fallbackBits\_add\_radiusBits\_add\_baseBits\_mul\_sampleBound\_lt\_surfaceCard}
\]
blocks the same endpoint whenever \(\texttt{sampleBound}+1\le
2^{\texttt{radiusBits}}\), \(|\texttt{surface}|+1\le
2^{\texttt{baseBits}}\), and the displayed bit sum is still below
\(|\texttt{surface}|\).  This turns a common informal repair move, budgeting
the sample radius and the surface alphabet separately, into an explicit bit
inequality that can be checked without leaving the finite interface.
These two factor-certificate routes are now also recorded directly in
\texttt{CruxSynthesis.lean}, so the ledger does not stop at the generic
polynomial-factor objection: it now remembers the common separate
radius/base-bit budgeting variant too.

Lean now also records the off-by-one variant that a repair is most likely to
state directly.  The helper
\[
  \texttt{PluginSampleMajority.succ\_le\_twoPow\_succ\_of\_le\_twoPow}
\]
proves that \(\texttt{sampleBound}\le 2^{\texttt{radiusBits}}\) implies
\(\texttt{sampleBound}+1\le 2^{\texttt{radiusBits}+1}\).  Consequently
\[
  \texttt{boundedSamplePluginMajorityWithBitFallbackFamilyActualSwitchedLocalInterface\_fallbackBits\_add\_radiusBits\_succ\_add\_baseBits\_mul\_sampleBound\_ge\_surfaceCard\_of\_sampleBound\_le\_twoPow\_and\_base\_le\_twoPow\_surjective}
\]
proves the lower bound
\[
  |\texttt{surface}|
    \le
  \texttt{fallbackBits}
    +(\texttt{radiusBits}+1)
    +\texttt{baseBits}\cdot\texttt{sampleBound}.
\]
The matching theorem
\[
  \texttt{boundedSamplePluginMajorityWithBitFallbackFamilyActualSwitchedLocalInterface\_not\_surjective\_of\_sampleBound\_le\_twoPow\_and\_base\_le\_twoPow\_and\_fallbackBits\_add\_radiusBits\_succ\_add\_baseBits\_mul\_sampleBound\_lt\_surfaceCard}
\]
gives non-surjectivity under the strict reverse bit inequality.  This prevents
the endpoint from treating a bit bound on \(\texttt{sampleBound}\) as though it
were already a bit bound on \(\texttt{sampleBound}+1\).  These successor-factor
routes are now also recorded directly in \texttt{CruxSynthesis.lean}, so the
ledger remembers the raw-value off-by-one variant a manuscript is especially
likely to state verbatim.
were already a bit bound on the Hamming-ball multiplier
\(\texttt{sampleBound}+1\).

The same successor-bit accounting is now also available when the surface
alphabet itself is only bounded as \(|\texttt{surface}|\le
2^{\texttt{baseBits}}\).  The theorem
\[
  \texttt{boundedSamplePluginMajorityWithBitFallbackFamilyActualSwitchedLocalInterface\_fallbackBits\_add\_radiusBits\_succ\_add\_baseBits\_succ\_mul\_sampleBound\_ge\_surfaceCard\_of\_sampleBound\_le\_twoPow\_and\_surfaceCard\_le\_twoPow\_surjective}
\]
proves
\[
  |\texttt{surface}|
    \le
  \texttt{fallbackBits}
    +(\texttt{radiusBits}+1)
    +(\texttt{baseBits}+1)\cdot\texttt{sampleBound}.
\]
Its strict-deficit counterpart
\[
  \texttt{boundedSamplePluginMajorityWithBitFallbackFamilyActualSwitchedLocalInterface\_not\_surjective\_of\_sampleBound\_le\_twoPow\_and\_surfaceCard\_le\_twoPow\_and\_fallbackBits\_add\_radiusBits\_succ\_add\_baseBits\_succ\_mul\_sampleBound\_lt\_surfaceCard}
\]
blocks the endpoint below that bit budget.  Thus a repair must pay for both
successors in the support-count envelope; neither the sample-radius factor nor
the alphabet-base factor can be rounded down for free.  These double-successor
routes are now also recorded directly in \texttt{CruxSynthesis.lean}, so the
ledger keeps the even weaker raw-value surface-cardinality certificate too.

The radius-zero boundary is now formalized as a sharp sanity check.  Lean proves
\[
  \texttt{PluginSampleMajority.smallSubsets\_zero}
  \qquad\text{and}\qquad
  \texttt{PluginSampleMajority.card\_smallSubsets\_zero},
\]
so there is exactly one zero-sample sparse override support, namely the empty
support.  The endpoint theorem
\[
  \texttt{boundedSamplePluginMajorityWithBitFallbackFamilyActualSwitchedLocalInterface\_fallbackBits\_ge\_surfaceCard\_of\_zeroSample\_surjective}
\]
then proves that zero-sample full expressivity forces
\[
  |\texttt{surface}| \le \texttt{fallbackBits}.
\]
Equivalently,
\[
  \texttt{boundedSamplePluginMajorityWithBitFallbackFamilyActualSwitchedLocalInterface\_not\_surjective\_zeroSample\_of\_fallbackBits\_lt\_surfaceCard}
\]
blocks any zero-sample bit-coded fallback route whose fallback code has fewer
bits than exact-visible inputs.  This is the degenerate case of the Hamming-cover
product theorem where the repair has no sampled override budget at all; it
recovers the ordinary truth-table code-size barrier inside the structured
fallback interface.

The lower bound is also sharp, not merely necessary.  Lean now constructs the
exact endpoint
\[
  \texttt{boundedSamplePluginMajorityWithBitFallbackFamilyActualSwitchedLocalInterface\_surjective\_zeroSample\_exactVisibleRuleDecode}
\]
using the canonical full truth-table decoder
\(\texttt{exactVisibleRuleDecode}\).  Given a target exact-visible rule, the
surjectivity proof chooses its
\(\texttt{exactVisibleRuleEncode}\) code and the empty bounded sample; the
plug-in fallback predictor is then definitionally the decoded truth table.  The
route-level anchor
\[
  \texttt{kpolyCoverage\_anchor\_boundedSampleMajorityWithBitFallbackFamily\_surjective\_zeroSample\_exactVisibleRuleDecode}
\]
records that the zero-sample structured fallback repair has an exact boundary:
below \(|\texttt{surface}|\) fallback bits it cannot be full-rule expressive,
and at \(|\texttt{surface}|\) bits it is exactly the unrestricted full
truth-table route.

The same observation is now formalized without the radius-zero restriction.
Lean proves
\[
  \texttt{boundedSamplePluginMajorityWithFallbackFamilyActualSwitchedLocalInterface\_emptySample\_predict}
\]
and
\[
  \texttt{boundedSamplePluginMajorityWithFallbackFamilyActualSwitchedLocalInterface\_surjective\_of\_fallback\_surjective}.
\]
Thus a structured fallback endpoint is full-rule expressive for every sample
bound whenever the fallback family itself is full-rule expressive; the empty
sample simply leaves the selected fallback rule unchanged.  The bit-coded
specialization
\[
  \texttt{boundedSamplePluginMajorityWithBitFallbackFamilyActualSwitchedLocalInterface\_surjective\_exactVisibleRuleDecode}
\]
states this for the canonical full truth-table decoder at
\(|\texttt{surface}|\) bits and arbitrary \(\texttt{sampleBound}\).  Consequently,
a bounded-sample fallback repair must prove genuine smallness of the fallback
family itself.  Merely bounding the sampled sparse overrides does not matter if
the fallback decoder is already the full truth-table family.

Lean now also proves the exact zero-sample converse:
\[
  \texttt{boundedSamplePluginMajorityWithFallbackFamilyActualSwitchedLocalInterface\_zeroSample\_surjective\_iff\_fallback\_surjective}
\]
and its bit-coded specialization
\[
  \texttt{boundedSamplePluginMajorityWithBitFallbackFamilyActualSwitchedLocalInterface\_zeroSample\_surjective\_iff\_fallback\_surjective}.
\]
At radius zero the wrapper contributes no hidden expressivity at all: full-rule
expressivity of the endpoint is equivalent to full-rule expressivity of the
fallback family or decoder.  The route anchors
\[
  \texttt{kpolyCoverage\_anchor\_boundedSampleMajorityWithFallbackFamily\_zeroSample\_surjective\_iff\_fallback\_surjective}
\]
and
\[
  \texttt{kpolyCoverage\_anchor\_boundedSampleMajorityWithBitFallbackFamily\_zeroSample\_surjective\_iff\_fallback\_surjective}
\]
make this available to later synthesis arguments.
Lean now packages the full generic structural form in
\texttt{ActualSwitchedLocalFallbackZeroSampleImage.lean}.  For an arbitrary
fallback family, the theorems
\[
  \texttt{FallbackZeroSampleImage.finitePredictorCover\_iff\_fallbackFamily},
  \texttt{FallbackZeroSampleImage.finiteIndexRepresentativeCover\_iff\_fallbackFamily},
\]
\[
  \texttt{FallbackZeroSampleImage.finitePredictorQuotient\_iff\_fallbackFamily},
  \texttt{FallbackZeroSampleImage.exactVisibleCompressionTarget\_iff\_fallbackFamily},
\]
and
\[
  \texttt{FallbackZeroSampleImage.exists\_clockedExactVisibleRealization\_iff\_fallbackFamily}
\]
show that the radius-zero wrapper contributes no hidden image or compression
content at all: every finite-cover, representative-cover, finite-quotient,
exact-visible-compression, or clocked-realization claim is exactly the
corresponding claim for the raw fallback family itself.  Combining the same
compression equivalence with
\[
  \texttt{kpolyCoverage\_anchor\_clockedFiniteLearningPayload\_iff\_exactVisibleCompressionTarget}
\]
now yields the further route anchor
\[
  \texttt{kpolyCoverage\_anchor\_boundedSampleMajorityWithFallbackFamily\_zeroSample\_clockedPayload\_iff\_fallbackFamily},
\]
so radius zero preserves clocked finite-learning payload targets just as rigidly
as it preserves image and compression targets.  In short: at radius zero the
structured fallback wrapper is only a reindexing of the fallback family, not a
new small-image mechanism.
These zero-sample routes are now also recorded directly in
\texttt{CruxSynthesis.lean}, so the ledger keeps the degenerate no-override
boundary rather than leaving it only as theorem-level scaffolding.
The smaller extracted module
\texttt{ActualSwitchedLocalBitFallbackWrapperObstruction.lean} packages this
point as a standalone wrapper boundary: at radius zero the fallback wrapper has
no hidden compressive effect at all, so any claimed gain must come from
genuinely restricting the fallback decoder rather than from the wrapper.
Lean now packages the bit-coded specialization in
\texttt{ActualSwitchedLocalBitFallbackZeroSampleImage.lean}.  The theorem
\[
  \texttt{BitFallbackZeroSampleImage.finitePredictorCover\_iff\_fallbackDecoderFamily}
\]
proves that the radius-zero wrapper and the raw fallback decoder family have
exactly the same predictor image, not merely the same surjectivity boundary.
Consequently,
\[
  \texttt{BitFallbackZeroSampleImage.exactVisibleCompressionTarget\_iff\_fallbackDecoderFamily}
\]
and
\[
  \texttt{BitFallbackZeroSampleImage.exists\_clockedExactVisibleRealization\_iff\_fallbackDecoderFamily}
\]
reduce every zero-sample small-image or clocked-realization claim to the
corresponding claim for the fallback decoder family itself.  Thus radius zero
adds no hidden compression content at all; it is only a reindexing of the same
fallback image.  Combining the same radius-zero compression equivalence with
\[
  \texttt{kpolyCoverage\_anchor\_clockedFiniteLearningPayload\_iff\_exactVisibleCompressionTarget}
\]
now yields the further route anchor
\[
  \texttt{kpolyCoverage\_anchor\_boundedSampleMajorityWithBitFallbackFamily\_zeroSample\_clockedPayload\_iff\_fallbackDecoderFamily},
\]
so radius zero preserves clocked finite-learning payload targets just as rigidly
as it preserves finite-cover, quotient, exact-visible-compression, and
clocked-realization targets.
This structural radius-zero packet is now also promoted into
\texttt{CruxSynthesis.lean}: the synthesis ledger records the finite-cover,
representative-cover, finite-quotient, exact-visible-compression, and
clocked-realization equivalence entries, together with the clocked-payload
boundary entry, rather than keeping them only in the extracted zero-sample
image module.
Lean also now packages the complementary exact-decoder obstruction at the same
layer.  The theorem
\[
  \texttt{kpolyCoverage\_anchor\_boundedSampleMajorityWithBitFallbackFamily\_not\_clockedPayload\_of\_exactVisibleRuleDecode\_of\_lt\_surfaceCard}
\]
shows that any exact visible-rule decoder for the bit-fallback family already
inherits the standard below-surface cardinality obstruction to clocked
finite-learning payloads, and its radius-zero specialization
\[
  \texttt{kpolyCoverage\_anchor\_boundedSampleMajorityWithBitFallbackFamily\_zeroSample\_exactVisibleRuleDecode\_not\_clockedPayload\_of\_lt\_surfaceCard}
\]
records that the degenerate wrapper does not evade this obstruction either.
The follow-on background module
\texttt{ActualSwitchedLocalBitFallbackZeroSampleProbeObstruction.lean}
extends this from full-image equivalence to finite probe transport.  Its
theorems
\[
  \texttt{BitFallbackZeroSampleProbeObstruction.finitePredictorCover\_card\_le\_of\_fallback\_injective\_realization}
\]
and
\[
  \texttt{BitFallbackZeroSampleProbeObstruction.not\_exactVisibleCompressionTarget\_of\_fallback\_injective\_realization\_of\_lt\_card}
\]
(together with the representative-cover, quotient, and clocked variants)
show that any injectively indexed finite probe family already realized by the
fallback decoder survives unchanged through the radius-zero wrapper.  This is
still background rather than a decisive crux theorem: one must separately
prove that the fallback decoder family itself realizes a probe family large
enough to contradict the claimed budget, so the transport result alone does
not rule out every legal strengthening of the route.

The radius-one boundary is now formalized as the first nontrivial sparse
override count.  Lean proves
\[
  \texttt{PluginSampleMajority.smallSubsets\_one}
  \qquad\text{and}\qquad
  \texttt{PluginSampleMajority.card\_smallSubsets\_one},
\]
so the allowed supports are exactly the empty support and one singleton support
for each exact-visible input, hence
\[
  |\texttt{smallSubsets}(\texttt{surface},1)|=|\texttt{surface}|+1.
\]
The endpoint theorem
\[
  \texttt{boundedSamplePluginMajorityWithBitFallbackFamilyActualSwitchedLocalInterface\_boolClassifierSpace\_card\_le\_bitCode\_mul\_surfaceCard\_succ\_of\_oneSample\_surjective}
\]
therefore specializes full expressivity at one sample to
\[
  2^{|\texttt{surface}|}
    \le
  2^{\texttt{fallbackBits}}\cdot (|\texttt{surface}|+1).
\]
Conversely,
\[
  \texttt{boundedSamplePluginMajorityWithBitFallbackFamilyActualSwitchedLocalInterface\_not\_surjective\_oneSample\_of\_bitCode\_mul\_surfaceCard\_succ\_lt\_boolClassifierSpace}
\]
blocks the route under the strict opposite inequality.  Thus the cheapest
nonzero sampled override budget buys only a linear multiplicative factor over
the fallback code family; it cannot substitute for the missing exponential
truth-table bits.

This one-sample boundary now also has the sharp additive bit-budget form.  If
the route proves the exact successor envelope
\[
  |\texttt{surface}|+1\le 2^{\texttt{baseBits}},
\]
then
\[
  \texttt{boundedSamplePluginMajorityWithBitFallbackFamilyActualSwitchedLocalInterface\_fallbackBits\_add\_baseBits\_ge\_surfaceCard\_of\_oneSample\_surfaceCard\_succ\_le\_twoPow\_surjective}
\]
proves
\[
  |\texttt{surface}| \le \texttt{fallbackBits}+\texttt{baseBits}.
\]
The matching strict-deficit theorem
\[
  \texttt{boundedSamplePluginMajorityWithBitFallbackFamilyActualSwitchedLocalInterface\_not\_surjective\_oneSample\_of\_surfaceCard\_succ\_le\_twoPow\_and\_fallbackBits\_add\_baseBits\_lt\_surfaceCard}
\]
blocks full-rule expressivity below that budget.  If instead the route states
only the raw manuscript-style bound
\[
  |\texttt{surface}|\le 2^{\texttt{baseBits}},
\]
Lean derives the successor cost explicitly:
\[
  \texttt{boundedSamplePluginMajorityWithBitFallbackFamilyActualSwitchedLocalInterface\_fallbackBits\_add\_baseBits\_succ\_ge\_surfaceCard\_of\_oneSample\_surfaceCard\_le\_twoPow\_surjective}
\]
proves
\[
  |\texttt{surface}| \le \texttt{fallbackBits}+(\texttt{baseBits}+1),
\]
and
\[
  \texttt{boundedSamplePluginMajorityWithBitFallbackFamilyActualSwitchedLocalInterface\_not\_surjective\_oneSample\_of\_surfaceCard\_le\_twoPow\_and\_fallbackBits\_add\_baseBits\_succ\_lt\_surfaceCard}
\]
is the corresponding negative.  This closes a small but important escape hatch:
the one-sample case should use the exact linear support count, but even there
the successor bit on \(|\texttt{surface}|+1\) cannot be rounded away.  These
one-sample routes are now also recorded directly in
\texttt{CruxSynthesis.lean}, so the ledger keeps both the exact linear
successor-count certificate and the weaker raw visible-size variant.

The next sparse radius is also now exact.  Lean proves
\[
  \texttt{PluginSampleMajority.card\_smallSubsets\_two},
\]
namely
\[
  |\texttt{smallSubsets}(\texttt{surface},2)|
    =
  1+|\texttt{surface}|+\binom{|\texttt{surface}|}{2}.
\]
Therefore two sampled overrides give only the empty support, singleton
supports, and unordered pairs.  The endpoint
\[
  \texttt{boundedSamplePluginMajorityWithBitFallbackFamilyActualSwitchedLocalInterface\_boolClassifierSpace\_card\_le\_bitCode\_mul\_one\_add\_surfaceCard\_add\_choose\_two\_of\_twoSample\_surjective}
\]
specializes surjectivity to the concrete quadratic product bound
\[
  2^{|\texttt{surface}|}
    \le
  2^{\texttt{fallbackBits}}\cdot
    \left(1+|\texttt{surface}|+\binom{|\texttt{surface}|}{2}\right),
\]
with strict negative
\[
  \texttt{boundedSamplePluginMajorityWithBitFallbackFamilyActualSwitchedLocalInterface\_not\_surjective\_twoSample\_of\_bitCode\_mul\_one\_add\_surfaceCard\_add\_choose\_two\_lt\_boolClassifierSpace}.
\]
Lean also records the corresponding additive budget certificate: if the exact
quadratic envelope fits into \(\texttt{overheadBits}\) bits, then
\[
  \texttt{boundedSamplePluginMajorityWithBitFallbackFamilyActualSwitchedLocalInterface\_fallbackBits\_add\_overheadBits\_ge\_surfaceCard\_of\_twoSample\_quadraticEnvelope\_le\_twoPow\_surjective}
\]
proves
\[
  |\texttt{surface}| \le \texttt{fallbackBits}+\texttt{overheadBits},
\]
and
\[
  \texttt{boundedSamplePluginMajorityWithBitFallbackFamilyActualSwitchedLocalInterface\_not\_surjective\_twoSample\_of\_quadraticEnvelope\_le\_twoPow\_and\_fallbackBits\_add\_overheadBits\_lt\_surfaceCard}
\]
is the strict-deficit form.  Thus allowing two local sampled repairs does not
change the nature of the obstruction: a quadratic support multiplier is still
far from a free truth-table oracle.  These exact one- and two-sample support-
count routes are now also recorded directly in \texttt{CruxSynthesis.lean}, so
the ledger remembers both the concrete product certificates and the exact
quadratic-envelope bit-budget form.

The canonical exact \((\mathsf{zfeat}(z),a,b)\) route now records this
finite-image burden on the constructive side as well.  In
\texttt{ExactZABTargetInterface.lean}, \texttt{ExactZABERMRoute.lean},
\texttt{CanonicalZABCandidateInterface.lean},
\texttt{CanonicalZABCodeWitness.lean},
\texttt{CanonicalZABRecoveryInterface.lean}, and
\texttt{CanonicalZABERMInterface.lean}, an explicit fixed-order decision-list
code assignment, or equality to the corresponding ERM-selected wrapper, yields
all three finite-image objects: a predictor cover, representative-index cover,
and quotient-code presentation of size at most \(2^{r+2k+1}\).  The identity
and projected bit-vector ERM specializations expose the same objects.
\texttt{CruxSynthesis.lean} now packages this at the synthesis level as
\texttt{CanonicalZABERMRouteCoverage}: the final manuscript-facing exact ERM
route already sits on the stop-grade exact-visible / finite-image /
clocked-payload surface.  This is a positive route for the actual canonical
family, not a product-success shortcut: the manuscript must still prove that
its switched predictors are the family realized by such shared-extractor codes
or by the named ERM wrapper.
The new module \texttt{ExactZABExtractorCollisionObstruction.lean} closes a
structural branch of that same route.  If the shared extractor
\(\mathsf{zfeat} : Z \to \{0,1\}^r\) is noninjective, then every realized exact
\((\mathsf{zfeat}(z),a,b)\) decision-list predictor is constant on the fibers
of \texttt{exactZABVisibleData}.  Lean packages this first at the raw realized
family level via
\[
  \texttt{not\_surjective\_predict\_of\_not\_injective\_of\_realizedByRawExactZABDecisionListFamily},
\]
and then lifts it to the named route surfaces
\[
  \texttt{ExactZABDecisionListTargetData.not\_surjective\_predict\_of\_not\_injective\_zfeat},
\]
\[
  \texttt{CanonicalZABDecisionListCandidateData.not\_surjective\_predict\_of\_not\_injective\_zfeat},
\]
\[
  \texttt{exactZABDecisionListERMFamily\_not\_surjective\_of\_not\_injective\_zfeat},
\]
and the corresponding recovery packages.  When \(Z\) is finite, the
pigeonhole corollary
\[
  \texttt{not\_injective\_of\_card\_gt\_bitVec}
\]
shows that every extractor with \(|Z| > 2^r\) lands in this forbidden branch.
This is genuine crux progress for the named exact \(z{+}a{+}b\) decision-list
route: on the noninjective-extractor branch, the route assumptions do not
merely fail to imply compression, they prevent surjectivity on the full exact
rule class altogether.  What remains live is only the different branch that
first proves the extractor is injective enough for the actual manuscript
surface.
The new module \texttt{ExactZABVisibleCollisionObstruction.lean} shows that
this was not an accident of the raw decision-list combiner.  Lean now proves
that every exact route whose predictors are still functions only of the visible
summary \((\mathsf{zfeat}(z),a,b)\) inherits the same collision blocker.  On
the non-shared side this closes the affine-feature, sparse-threshold, and
affine-decision-list route surfaces via
\[
  \texttt{not\_surjective\_predict\_of\_not\_injective\_of\_realizedByExactZABAffineFeatureFamily},
\]
\[
  \texttt{not\_surjective\_predict\_of\_not\_injective\_of\_realizedByExactZABSparseThresholdAffineFamily},
\]
and
\[
  \texttt{not\_surjective\_predict\_of\_not\_injective\_of\_realizedByExactZABAffineDecisionListFamily}.
\]
On the shared-basis side it closes the corresponding affine-feature,
sparse-threshold, and shared decision-list surfaces via
\[
  \texttt{not\_surjective\_predict\_of\_not\_injective\_of\_realizedBySharedExactZABAffineFeatureFamily},
\]
\[
  \texttt{not\_surjective\_predict\_of\_not\_injective\_of\_realizedBySharedExactZABSparseThresholdAffineFamily},
\]
and
\[
  \texttt{not\_surjective\_predict\_of\_not\_injective\_of\_realizedBySharedExactZABAffineDecisionListFamily},
\]
with the ERM-facing wrapper
\[
  \texttt{sharedExactZABAffineDecisionListERMFamily\_not\_surjective\_of\_not\_injective\_zfeat}
\]
available directly.  So the noninjective-extractor branch is now closed across
the broader exact \((\mathsf{zfeat}(z),a,b)\) route portfolio, not only for the
original raw decision-list presentation.
The new module \texttt{ExactZABWidthCollisionObstruction.lean} specializes
this further to the manuscript-shaped bitvector branch.  If
\(\mathsf{zfeat} : \{0,1\}^n \to \{0,1\}^r\) satisfies \(r<n\), then Lean first
proves
\[
  \texttt{not\_injective\_bitVecExtractor\_of\_width\_lt},
\]
so the extractor is automatically in the forbidden collision branch before any
downstream affine or decision-list reasoning begins.  It then lifts that fact
to the ERM-facing exact route surfaces
\[
  \texttt{exactZABDecisionListERMFamily\_not\_surjective\_of\_width\_lt},
\]
\[
  \texttt{exactZABAffineFeatureERMFamily\_not\_surjective\_of\_width\_lt},
\]
\[
  \texttt{exactZABSparseThresholdAffineERMFamily\_not\_surjective\_of\_width\_lt},
\]
and on the shared-basis side
\[
  \texttt{sharedExactZABAffineFeatureERMFamily\_not\_surjective\_of\_width\_lt},
\]
\[
  \texttt{sharedExactZABSparseThresholdAffineERMFamily\_not\_surjective\_of\_width\_lt},
\]
\[
  \texttt{sharedExactZABAffineDecisionListERMFamily\_not\_surjective\_of\_width\_lt}.
\]
So for bitvector latent data the only live exact-ZAB branch is no longer
``perhaps the extractor is merely lossy but still usable''; the extractor must
at least avoid width compression outright.
The new module \texttt{ExactZABInjectiveBranchObstruction.lean} then turns to
the remaining live branch: injective extractors with \(r\ge n\).  It does not
claim that every such branch is impossible, but it proves that surjectivity on
the current exact/shared exact ERM route surfaces already forces large route
parameters.  For the raw exact decision-list ERM family,
\[
  \texttt{exactZABDecisionListERMFamily\_extractorWidthLowerBound\_of\_surjective\_predict}
\]
shows that surjectivity implies
\[
  r \ge 2^{n+2k}-(2k+1),
\]
with the route-facing negation exposed as
\[
  \texttt{exactZABDecisionListERMFamily\_not\_surjective\_of\_extractorWidth\_lt\_truthTableGap}.
\]
On the shared-basis side,
\[
  \texttt{sharedExactZABAffineFeatureERMFamily\_featureCountLowerBound\_of\_surjective\_predict}
\]
forces at least \(n+2k\) shared affine features, and
\[
  \texttt{sharedExactZABAffineDecisionListERMFamily\_combinerWidthLowerBound\_of\_surjective\_predict}
\]
forces
\[
  p \ge 2^{n+2k}-1.
\]
So the current injective-extractor escape hatch is no longer ``keep
\(\mathsf{zfeat}\) injective but make the downstream witness small'' on these
ERM surfaces.  Any surviving repair would have to change the route class or
accept these large parameter requirements.
The new module
\texttt{ActualSwitchedLocalZABDecisionListWidthObstruction.lean} lifts the same
live-branch arithmetic to the manuscript-facing local-rule interface
\texttt{ActualSwitchedLocalInterface.ZABDecisionListData}.  Abstractly,
\[
  \texttt{ZABDecisionListData.visibleCardGapLowerBound\_of\_surjective\_predict}
\]
forces any surjective actual-local ZAB route to pay extractor width at least
\[
  |\texttt{ExactVisiblePostSwitchSurface}\ Z\ k| - (2k+1).
\]
On the \(n\)-bit input space \(Z=\{0,1\}^n\), the concrete specialization
\[
  \texttt{ZABDecisionListData.extractorWidthLowerBound\_of\_surjective\_predict}
\]
again yields
\[
  r \ge 2^{n+2k}-(2k+1),
\]
and
\[
  \texttt{not\_nonempty\_zabDecisionListData\_of\_surjective\_predict\_of\_extractorWidth\_lt\_truthTableGap}
\]
says there is no actual-local ZAB-data package below that threshold.  So the
``injective but still small'' obstruction is now recorded directly on the
manuscript's local-rule data surface, not only on the ERM wrapper.
The follow-up module \texttt{SameWidthZABObstruction.lean} closes the smallest
remaining injective repair inside that branch.  If the extractor stays
same-width, \(\mathsf{zfeat} : \{0,1\}^n \to \{0,1\}^n\), then once the visible
width \(n+2k\) is at least \(2\), Lean proves
\[
  \texttt{exactZABDecisionListERMFamily\_not\_surjective\_of\_sameWidth}
\]
for the exact ERM surface and
\[
  \texttt{ActualSwitchedLocalInterface.not\_nonempty\_zabDecisionListData\_of\_surjective\_predict\_of\_sameWidth}
\]
for the manuscript-facing actual-local surface.  The concrete full-rule
counterexample appears as
\[
  \texttt{ActualSwitchedLocalInterface.fullRuleActualSwitchedLocalInterface\_not\_zabDecisionListData\_of\_sameWidth}.
\]
So the live injective exact-ZAB branch no longer includes ``same width but
injective'': any surviving repair must widen \(\mathsf{zfeat}\) strictly beyond
the original \(n\) visible bits, not merely prove injectivity.
The next route-facing refinement is
\texttt{ExactZABInjectiveERMRouteCrux.lean}.  It converts the widened-branch
lower bounds into direct contradictions on the named ERM target surface
\texttt{fullExactVisibleRuleFamily}.  Below the verified thresholds, there is no
choice of training samples that makes the raw exact decision-list ERM family,
the shared exact affine-feature ERM family, or the shared exact affine
decision-list ERM family equal the full exact-visible Boolean rule class.  The
three Lean statements are
\[
  \texttt{not\_exists\_exactZABDecisionListERMFamily\_eq\_fullExactVisibleRuleFamily\_of\_extractorWidth\_lt\_truthTableGap},
\]
\[
  \texttt{not\_exists\_sharedExactZABAffineFeatureERMFamily\_eq\_fullExactVisibleRuleFamily\_of\_featureCount\_lt\_visibleWidth},
\]
and
\[
  \texttt{not\_exists\_sharedExactZABAffineDecisionListERMFamily\_eq\_fullExactVisibleRuleFamily\_of\_combinerWidth\_lt\_truthTableGap}.
\]
So the widened injective branch is no longer only a parameter lower-bound
warning: on the present ERM route surface, small widened injective extractors
and small downstream combiners cannot possibly realize the manuscript's full
exact-visible target family at all.
The next strengthening is
\texttt{ExactZABFullRuleObstruction.lean}.  It removes the ERM/sample-choice
wrapper entirely.  Below the same verified thresholds, Lean proves that the
underlying hypothesis classes themselves do not realize the full exact-visible
Boolean family:
\[
  \texttt{fullExactVisibleRuleFamily\_not\_realizedByRawExactZABDecisionListFamily\_of\_extractorWidth\_lt\_truthTableGap},
\]
\[
  \texttt{fullExactVisibleRuleFamily\_not\_realizedBySharedExactZABAffineFeatureFamily\_of\_featureCount\_lt\_visibleWidth},
\]
\[
  \texttt{fullExactVisibleRuleFamily\_not\_realizedBySharedExactZABSparseThresholdAffineFamily\_of\_linearBudget\_lt\_surfaceCard},
\]
and
\[
  \texttt{fullExactVisibleRuleFamily\_not\_realizedBySharedExactZABAffineDecisionListFamily\_of\_combinerWidth\_lt\_truthTableGap}.
\]
So this repair branch is no longer merely blocked on ERM target selection:
below those budgets, changing the sample set or retuning the ERM wrapper inside
the same exact-ZAB classes cannot help, because the classes themselves miss the
full target family.  What remains live is only the genuinely large-parameter
branch with widened extractors or downstream budgets beyond these thresholds.
The next exact-ZAB strengthening is
\texttt{ExactZABDecisionListStructuralObstruction.lean}.  This one is not a
cardinality argument at all.  For the raw exact
\((\mathsf{zfeat}(z), a, b)\) decision-list class, Lean now proves that every
positive-width branch already fails structurally:
\[
  \texttt{not\_realizedByRawExactZABDecisionListFamily\_of\_surjective\_predict\_of\_pos\_widths}.
\]
The proof splits on the shared extractor output.  If some latent state \(z_0\)
already activates a \(\mathsf{zfeat}\) bit before any \(a\)- or \(b\)-bit is
read, then every predictor in the class is constant on the entire
\((a,b)\)-fiber above \(z_0\).  If instead every \(\mathsf{zfeat}(z)\) is
zero, then the class ignores \(z\) entirely.  Either way, once \(n>0\) and
\(k>0\), the route cannot realize a surjective exact family, so in particular
\[
  \texttt{fullExactVisibleRuleFamily\_not\_realizedByRawExactZABDecisionListFamily\_of\_pos\_widths}
\]
and the manuscript-facing actual-local wrapper is impossible:
\[
  \texttt{ActualSwitchedLocalInterface.fullRuleActualSwitchedLocalInterface\_not\_zabDecisionListData\_of\_pos\_widths}.
\]
This closes the raw exact decision-list branch completely for positive latent
and visible widths, independently of the extractor width \(r\).  The still-live
exact-ZAB repair branches are therefore only the richer shared affine-feature,
sparse-threshold, and shared affine decision-list classes.
The next exact-ZAB route analysis,
\texttt{SharedExactZABAffineFeatureCharacterization.lean}, does not close one
of those remaining branches, but it identifies its exact mathematical content.
Lean now proves that the shared exact affine-feature class realizes the full
exact-visible Boolean rule family if and only if the shared affine summary
\[
  u \mapsto \texttt{exactZABAffineFeatureSummary}(\mathsf{zfeat},\mathsf{features},u)
\]
is injective on the exact visible surface.  The key statements are
\[
  \texttt{injective\_summary\_of\_surjective\_predict\_of\_realizedBySharedExactZABAffineFeatureFamily},
\]
\[
  \texttt{fullExactVisibleRuleFamily\_realizedBySharedExactZABAffineFeatureFamily\_of\_injective\_summary},
\]
and
\[
  \texttt{fullExactVisibleRuleFamily\_realizedBySharedExactZABAffineFeatureFamily\_iff\_injective\_summary}.
\]
So this branch is not a small-class theorem by itself.  Once the shared summary
collides, surjectivity fails immediately; but if the summary is injective, the
arbitrary truth-table combiner recovers every exact-visible Boolean rule.  In
other words, the route reduces exactly to an injective codebook condition on
the shared affine summary, not to a standalone compression argument.
The follow-on file \texttt{SharedExactZABAffineFeatureSharpness.lean} shows that
this injective-summary condition is genuinely satisfiable at the lower-bound
threshold on the native bitvector surface.  For \(Z = \texttt{BitVec}\, n\),
\(\mathsf{zfeat} = \mathrm{id}\), and
\(\mathsf{features} = \texttt{canonicalAffineBasis}(n + 2k)\), Lean proves that
\[
  \texttt{exactZABAffineFeatureSummary}(\mathrm{id},
  \texttt{canonicalAffineBasis}(n + 2k), u)
  = \texttt{bitVecZABVisibleData}(u),
\]
so the shared summary is injective and therefore
\[
  \texttt{fullExactVisibleRuleFamily\_realizedBySharedExactZABAffineFeatureFamily\_id\_canonical}.
\]
Hence the previously verified obstruction \(p < n + 2k\) is sharp for this
branch: at \(p = n + 2k\), the shared exact affine-feature route remains fully
expressive.  This is diagnostic rather than obstructive.  It shows that feature
count alone does not block the route once the summary is allowed to be an
injective codebook for the full visible datum \((z,a,b)\).
The next branch analysis,
\texttt{SharedExactZABAffineDecisionListCharacterization.lean}, gives the
parallel exact description for the shared affine decision-list route.  Here the
shared extractor and shared affine basis are fixed, and only the downstream
fixed-order decision-list code varies.  Lean proves that this branch realizes
the full exact-visible Boolean rule family if and only if the first-active
signature
\[
  u \mapsto \texttt{exactZABAffineDecisionListSignature}
    (\mathsf{zfeat},\mathsf{features},u)
\]
is injective on the exact visible surface.  The key statements are
\[
  \texttt{injective\_signature\_of\_surjective\_predict\_of\_realizedBySharedExactZABAffineDecisionListFamily},
\]
\[
  \texttt{fullExactVisibleRuleFamily\_realizedBySharedExactZABAffineDecisionListFamily\_of\_injective\_signature},
\]
and
\[
  \texttt{fullExactVisibleRuleFamily\_realizedBySharedExactZABAffineDecisionListFamily\_iff\_injective\_signature}.
\]
So this branch is not blocked by a small-width argument alone.  It reduces
exactly to an injective first-active codebook condition on the shared summary.
Any future crux for this branch has to attack that injective-signature repair
route directly, not merely the decision-list wrapper as such.
The follow-on file
\texttt{SharedExactZABAffineDecisionListStructuralObstruction.lean} does
exactly that on the native bitvector exact-visible surface.  Lean first proves
the raw structural theorem
\[
  \texttt{not\_injective\_firstActiveFeature\_affineFeatureVector\_of\_one\_lt}:
\]
for any affine feature family on \(\texttt{BitVec}\, d\), once \(d > 1\), the
first-active signature
\[
  x \mapsto
  \texttt{firstActiveFeature?}(\texttt{affineFeatureVector}(\mathsf{features},x))
\]
is never injective.  The reason is structural rather than cardinal: either all
features are empty, so the signature is constantly \texttt{none}, or the
earliest nonempty affine parity test is true at at least two distinct points,
so that earliest signature class already collides.  Pushed back through the
visible-data bijection \(\texttt{bitVecZABVisibleData}\), this yields
\[
  \texttt{not\_injective\_exactZABAffineDecisionListSignature\_id\_of\_one\_lt}
\]
for \(Z = \texttt{BitVec}\, n\), \(\mathsf{zfeat} = \mathrm{id}\), and
arbitrary shared affine bases on the full visible width \(n + 2k\).  Hence Lean
also proves
\[
  \texttt{fullExactVisibleRuleFamily\_not\_realizedBySharedExactZABAffineDecisionListFamily\_id\_of\_one\_lt}.
\]
This is genuine crux work for the identity-extractor shared affine
decision-list route: once \(n + 2k > 1\), no choice of shared affine basis and
no decision-list width \(p\) can recover the full exact-visible Boolean rule
family.  What remains live is only the genuinely widened injective-extractor
branch, where \((z,a,b)\) is first embedded into a larger ambient bitspace
before the shared affine basis is applied.
The next file
\texttt{SharedExactZABAffineDecisionListFiberObstruction.lean} closes that
remaining extractor escape hatch for every nonempty latent type with positive
visible width.  Lean proves
\[
  \texttt{exactZABAffineDecisionListSignature\_fiberPoint},
\]
showing that after freezing one latent point \(z=z_0\), the shared exact affine
decision-list signature is exactly the raw first-active affine signature on the
\(2k\) visible suffix bits \((a,b)\).  Since \(k>0\) implies \(2k>1\), the
earlier raw theorem applies on every such fiber, yielding
\[
  \texttt{not\_injective\_exactZABAffineDecisionListSignature\_of\_pos}
\]
for arbitrary extractors \(\mathsf{zfeat} : Z \to \texttt{BitVec}\, r\) and
arbitrary shared affine bases.  Consequently Lean also proves
\[
  \texttt{fullExactVisibleRuleFamily\_not\_realizedBySharedExactZABAffineDecisionListFamily\_of\_pos}.
\]
The route-facing wrappers
\[
  \texttt{SharedExactZABDecisionListTargetData.not\_surjective\_predict\_of\_pos}
\]
and
\[
  \texttt{SharedExactZABDecisionListRecoveryData.not\_surjective\_predict\_of\_pos}
\]
push the same obstruction onto the exact target and recovery interfaces.  So
the shared exact affine decision-list branch is now closed on every nonempty
exact surface with \(k>0\), independently of extractor width \(r\),
decision-list width \(p\), or the choice of shared extractor.
The recovery side has a separate live quantitative premise.  Lean now tests it
directly in \texttt{BadCodeAgreementObstruction.lean}: if a bad code agrees
with the target on every positive-mass input, its agreement mass is exactly
\(1\).  Hence any strict \(q<1\) recovery package forces every bad code to
disagree with the target somewhere on positive sampling mass:
\[
  \texttt{BitEncodedClassifierFamily.exists\_pos\_mass\_disagreement\_of\_badCode\_agreementBound\_lt\_one}.
\]
The exact ZAB recovery interface exposes the same demand as
\[
  \texttt{ExactZABERMRecoveryData.exists\_pos\_mass\_disagreement\_of\_badCode}.
\]
For an atomic distribution \(\delta_x\), this specializes to a sharper
empty-package obstruction: a bad code agreeing with the target at \(x\) rules
out strict recovery:
\[
  \texttt{not\_ExactZABERMRecoveryData\_of\_pure\_badCode\_agrees},
\]
\[
  \texttt{not\_CanonicalZABERMRecoveryData\_of\_pure\_badCode\_agrees}.
\]
The underlying PMF bridge is now explicit in
\texttt{FiniteIIDAgreement.lean}.  Lean proves that agreement mass is the PMF
outer measure of the agreement set, so agreement mass \(1\) forces agreement on
the entire positive-mass support:
\[
  \texttt{agrees\_on\_support\_of\_agreementMass\_eq\_one}.
\]
Conversely, a single positive-mass disagreement gives the strict quantitative
boundary
\[
  \texttt{agreementMass\_lt\_one\_of\_pos\_mass\_disagreement}.
\]
The latest strengthening fills the gap between pure atoms and full support.  If
a bad code agrees with the target on a finite region \(R\), then the advertised
bad-code threshold must be at least the \(R\)-mass:
\[
  \texttt{finset\_mass\_le\_agreementMass\_of\_agrees\_on},
  \qquad
  \texttt{BitEncodedClassifierFamily.regionMass\_le\_of\_badCode\_agrees\_on}.
\]
Thus any region with mass above \(q\) refutes a uniform agreement bound by \(q\):
\[
  \texttt{BitEncodedClassifierFamily.not\_badCodeAgreementBound\_of\_agrees\_on\_region\_gt}.
\]
The exact ZAB ERM interface exposes the same quantitative obligation:
\[
  \texttt{ExactZABERMRecoveryData.regionMass\_le\_of\_badCode\_agrees\_on},
  \qquad
  \texttt{not\_ExactZABERMRecoveryData\_of\_badCode\_agrees\_on\_region\_gt}.
\]
The current Lean strengthening also localizes the contrapositive.  A uniform
bad-code threshold \(q\) forces every finite region of mass above \(q\) to
contain a positive-mass disagreement point for each bad code:
\[
  \texttt{BitEncodedClassifierFamily.exists\_pos\_mass\_disagreement\_in\_region\_of\_badCodeAgreementBound\_lt\_mass}.
\]
For exact ZAB recovery this becomes
\[
  \texttt{ExactZABERMRecoveryData.exists\_pos\_mass\_disagreement\_in\_region\_of\_badCode}.
\]
So a repair cannot claim a small global bad-code threshold while allowing even a
moderately large all-agreement region for one non-target code under the sampling
PMF.  It must provide localized disagreement, on positive sampling mass, in
every region whose actual mass exceeds the advertised threshold.
The sample-level form is now explicit as well.  In \texttt{FiniteIIDAgreement.lean},
an agreement region of mass \(p\) contributes at least \(p^m\) to the
consistent-sample mass:
\[
  \texttt{regionMass\_pow\_le\_consistentSampleMass\_of\_agrees\_on}.
\]
In \texttt{FiniteIIDFamilyBound.lean}, the samples consistent with any fixed bad
code form a subevent of the deceptive samples:
\[
  \texttt{EncodedFamily.consistentSampleMass\_le\_deceptiveSampleMass\_of\_badCode},
\]
and hence
\[
  \texttt{EncodedFamily.regionMass\_pow\_le\_deceptiveSampleMass\_of\_badCode\_agrees\_on}.
\]
The exact ZAB specialization
\[
  \texttt{exactZAB\_regionMass\_pow\_le\_deceptiveSampleMass\_of\_badCode\_agrees\_on}
\]
turns every large agreeing region for one bad code into a concrete positive
deceptive-sample event under IID sampling.  This does not depend on a finite-probe
or full-surface counting argument; it is a direct sample-distribution burden.
This turns support witnesses into recovery-bound obstructions without adding a
new finite-image or probe assumption.
Equivalently, a pure-atom agreement bound with \(q<1\) forces every code that
matches the target at the atom to be globally equal to the target predictor.
This means the finite-image ERM wrappers do not close the route by themselves:
the manuscript must prove a support/disagreement theorem for the actual
switched bad codes and the actual sampling measure.
The synthesis ledger now records this as route coverage rather than leaving it
only in the local obstruction file.  In \texttt{CruxSynthesis.lean}, the
\texttt{PNPKpolySubrepairClass} list includes product-bound and fielded-switching
finite-image boundary entries, the corresponding clocked-payload boundary
entries, the large-region bad-code disagreement boundary, and the full-visible
and already-surjective Boolean-family obstructions.  These are deliberately
narrow anchors: they show that the product/fielded shortcut has been exhausted
at both the compression and payload interfaces, while the recovery side now has
an explicit positive-mass disagreement obligation under the actual sampling
measure.  The broad
\texttt{kpolyCompressionBridge} no longer remains an uncovered broad escape
class in the synthesis ledger: the later promoted packet packages these
finite-image and counterexample anchors into a theorem-backed blocker for the
entire current broad bridge route.  What still does \emph{not} happen here is
an affirmative small-class theorem for the actual switched predictor family.
The newest strengthening removes an avoidable presentation dependency from
that obstruction.  The counterexample need not be exactly
\texttt{fullExactVisibleRuleFamily}; any exact-visible family whose prediction
map is already surjective onto all Boolean rules on the exact surface refutes
the product-bound-only and fielded-switching-only bridge schemas below the
surface threshold:
\[
  \texttt{not\_productBoundOnlyExactVisibleCompressionBridge\_of\_true\_productBound\_of\_surjective\_predict\_of\_lt\_surfaceCard},
\]
\[
  \texttt{not\_fieldedSwitchingOnlyExactVisibleCompressionBridge\_of\_true\_fieldedSwitching\_of\_surjective\_predict\_of\_lt\_surfaceCard}.
\]
At the finite-image boundary, the same point is recorded as
\[
  \texttt{not\_productBoundSemanticFiniteImageBridge\_of\_true\_productBound\_of\_surjective\_predict\_of\_lt\_predictorCard}
\]
and its fielded-switching analogue.  Thus a repair cannot evade the blocker by
renaming the index type or by presenting the same fully expressive family
through a different exact-visible parameterization.

The new module \texttt{ABVisibleCompressionObstruction.lean} closes the most
charitable reduced-surface workaround for that burden.  Suppose the exact
switched family on \(u=(z,a,b)\) actually factors through the raw visible
projection \((a,b)\).  Lean now proves that any exact visible compression
target, or any clocked \(s\)-bit realization family for the exact surface,
pulls back to the same \(s\)-bit budget on the reduced raw family over
\(\texttt{ABVisibleSurface}\ k\).  Hence, if that reduced family is still
surjective onto all Boolean rules on \((a,b)\), then no such theorem can hold
below the reduced truth-table threshold
\[
  s < |\texttt{ABVisibleSurface}\ k| = 2^{2k}.
\]
So ``the switched decoder ignores \(z\)'' is not yet a same-route rescue.  It
still leaves a separate strict-subclass theorem to prove on the raw
\((a,b)\)-surface itself.  The same module also now propagates this reduced
truth-table lower bound through the charitable invariant-view and fork-view
compression interfaces.  So replacing the exact surface by one of those
intermediate exact views still does not evade the raw \((a,b)\) counting
burden.

The module \texttt{ABDecisionListObstruction.lean} now turns that last sentence
into a concrete strict-subclass theorem for the first raw \((a,b)\) candidate
class.  The fixed-order raw decision-list decoder has only
\[
  2k+1
\]
program bits, while the full Boolean rule class on \((a,b)\) needs
\[
  |\texttt{ABVisibleSurface}\ k| = 2^{2k}
\]
truth-table bits.  Lean proves the sharp positive-width inequality
\[
  2k+1 < 2^{2k} \qquad (k>0).
\]
Therefore a reduced family that is still surjective onto all Boolean rules on
\((a,b)\) cannot be realized by the fixed-order raw decision-list route.  The
same module lifts the result back to the exact surface: if an exact switched
family factors through \((a,b)\) and the reduced family is fully expressive,
then it cannot be realized by the raw exact-surface decision-list class.  This
does not close the full \(Kpoly\) bridge, but it removes one concrete repair
subclass from the live ledger: the author must prove that the actual reduced
switched family is much smaller than all raw Boolean rules, not merely that it
uses the raw \((a,b)\) coordinates.
The new module \texttt{ABDecisionListRouteCrux.lean} pushes this one step
closer to the named positive route interface.  Its theorem
\[
  \texttt{not\_factorsThrough\_ab\_and\_decisionList\_of\_surjective\_predict\_of\_pos}
\]
combines the positive route certificate from
\texttt{ABDecisionListRoute.lean} with the reduced truth-table lower bound and
shows that, for every positive width, the route assumptions themselves are
inconsistent on any branch where the reduced \((a,b)\) family remains fully
expressive.  The invariant-lift form
\[
  \texttt{not\_invariant\_and\_decisionList\_of\_surjective\_lift\_of\_pos}
\]
packages the same contradiction directly at the lifted exact-surface route,
and the concrete corollary
\[
  \texttt{fullExactABRuleFamily\_not\_exactVisibleCompressionTarget\_of\_pos}
\]
shows that the canonical full pullback of all raw \((a,b)\) rules cannot meet
the \(2k+1\) decision-list budget.  This is real crux progress for that named
repair surface: on its full-expressivity branch the route assumptions no
longer merely fail to imply compression, they contradict it.
The synthesis ledger now records these as narrow
\texttt{PNPKpolySubrepairClass} coverage items alongside the clock-wrapper
equivalence, the finite-image budget/transport lemmas, the product/fielded
semantic-boundary and surjective-family anchors, and the reduced \((a,b)\)
finite-image and quotient-code pullback obstructions.  Before the final broad
promotion these sat underneath an uncovered
\texttt{kpolyCompressionBridge} banner; now they feed directly into the
theorem-backed \texttt{KpolyCompressionBridgeCoverage} packet.  The latest
bounded-sample fallback update extends that same narrow ledger with explicit
entries for the canonical \texttt{exactVisibleRuleDecode} endpoint: one saying
the endpoint is already full-rule expressive for every sample bound, and one
saying any sparse-threshold recovery witness on it must satisfy the joint
visible-width and lightest-point barrier.  The next refinement lifts that one
level higher: the same ledger now records the generic surjective-actual-local
joint recovery barrier, the generic surjective-actual-local visible-width
threshold, the generic surjective-actual-local lightest-point threshold, and
the corresponding ``no extractor below threshold'' barrier.  The current
refinement keeps the negative visible-width form as well via the unconditional
surjective-actual-local half-width no-extractor boundary, so the exact-decoder
fallback entries are explicitly exposed as specializations of a broader
route-level obstruction rather than as isolated endpoint facts.

The next reduced-surface module, \texttt{SharedABFeatureRoutes.lean}, pushes
this from abstract counting to the concrete shared-basis families already used
by the route.  On the raw \((a,b)\)-surface it now defines the corresponding
bit families and proves the direct lower bounds:
\begin{itemize}[leftmargin=1.5em]
  \item a shared affine-feature realization forces
  \(|\texttt{ABVisibleSurface}\ k| \le 2^r\),
  \item a shared sparse-threshold realization forces
  \(|\texttt{ABVisibleSurface}\ k| \le 2r\),
  \item a shared decision-list realization forces
  \(|\texttt{ABVisibleSurface}\ k| \le r+1\).
\end{itemize}
Equivalently, if the lifted reduced family on \((a,b)\) is still fully
expressive, then none of these shared raw classes can realize it below the
same reduced truth-table threshold.  So the remaining manuscript-side burden is
not merely to state invariance and choose one shared basis; it is to prove that
the resulting reduced family is genuinely a strict subclass of all Boolean
rules on \((a,b)\).

The companion module \texttt{SharedABFeatureObstruction.lean} packages the same
fact as exact-pullback refuters for the shared-basis route.  If the exact
switched family factors through \((a,b)\), the reduced raw family is
surjective, and the shared combiner budget is below \(2^{2k}\), then the exact
family cannot be realized by the corresponding shared exact class.  Lean proves
this separately for:
\[
  2^r,\qquad 2r,\qquad r+1
\]
shared affine-feature, sparse-threshold, and decision-list budgets.  The
synthesis ledger records these as additional narrow
\texttt{PNPKpolySubrepairClass} items.  Standing alone, these theorems only
rule out full reduced-rule expressivity for named shared classes; they do not
prove a manuscript-faithful small-class theorem for the actual switched family.
But the broad synthesis ledger no longer leaves
\texttt{kpolyCompressionBridge} uncovered: those narrow shared-class blockers
are now explicitly absorbed into the promoted
\texttt{KpolyCompressionBridgeCoverage} route blocker.
The same shared-feature blocker now has the finite-probe form needed after the
new image-budget foundation.  Full reduced-rule surjectivity is unnecessary:
if the reduced family already realizes an injectively indexed finite probe
family \(P\), then a shared affine-feature realization forces
\[
  |P|\le 2^{2^r},
\]
a shared sparse-threshold realization forces
\[
  |P|\le 2^{2r},
\]
and a shared decision-list realization forces
\[
  |P|\le 2^{r+1}.
\]
The exact-pullback versions have the same bounds whenever the exact visible
family factors through \(\texttt{abVisibleData}\).  Lean records the refuters as
\[
  \texttt{not\_realizedBySharedExactABAffineFeatureFamily\_of\_factorsThrough\_ab\_and\_injective\_probe\_of\_lt\_budgetImageCard}
\]
and the matching sparse-threshold and decision-list variants.  Thus a
shared-basis repair cannot merely avoid the full Boolean rule family; it must
prove that the actual reduced predictor image avoids every finite probe family
larger than the named combiner image.
The new module \texttt{SharedABFeatureRouteCrux.lean} pushes this from
``named class obstruction'' to ``named route contradiction.''  It combines the
positive route interfaces from \texttt{SharedABFeatureRoutes.lean} with the
reduced truth-table lower bound and proves that, whenever the reduced
\((a,b)\) family remains fully expressive, the route assumptions themselves are
inconsistent below the corresponding combiner budgets.  Lean now records three
factored contradictions,
\[
  \texttt{not\_factorsThrough\_ab\_and\_sharedAffineFeature\_of\_surjective\_predict\_of\_lt\_cardFormula},
\]
with the matching sparse-threshold and decision-list variants, and three
invariant-lift contradictions for the same named route surfaces.  So on the
full-expressivity branch these shared-feature repairs no longer merely fail to
supply a missing small-class theorem; they directly contradict the exact-visible
compression target forced by their own positive route certificates.  What
remains uncovered is only a different repair shape that first proves the actual
reduced family is a strict subclass of all Boolean rules on \((a,b)\).

The new module \texttt{ABVisibleInvariantSurjectivityObstruction.lean} closes a
more basic loophole on the exact surface itself.  Suppose \(Z\) is nontrivial
and the switched predictor is \texttt{ABVisibleInvariant}, i.e. it depends only
on \((a,b)\) and ignores the exact \(z\)-coordinate.  Then Lean now proves that
the resulting predictor cannot be surjective onto all Boolean rules on
\(\texttt{ExactVisiblePostSwitchSurface}\ Z\ k\).  The reason is structural:
two distinct exact inputs with the same \((a,b)\)-data must receive the same
output, so every \(z\)-fiber is collapsed before any coding budget enters the
picture.  This is strictly stronger than the reduced-surface counting
obstructions above.  It says that an ``ignore \(z\)'' repair is not merely a
compression claim still waiting for a tighter bit bound; it already changes the
target rule class on the exact visible surface.  The same module propagates
this obstruction through the shared affine-feature, sparse-threshold,
decision-list, recovery, and canonical AB candidate interfaces, so the route
cannot recover by relabeling the same \((a,b)\)-only family under a more
charitable exact-surface wrapper.
The newest strengthening removes the remaining ambiguity in that diagnosis.  It
does not only refute surjectivity.  Lean now defines
\[
  \texttt{ABFiberSeparatingRule}
\]
for any exact-surface Boolean rule that distinguishes two inputs with identical
\((a,b)\)-data, and proves
\[
  \texttt{not\_exists\_predict\_eq\_of\_abVisibleInvariant\_of\_abFiberSeparatingRule}.
\]
So an \((a,b)\)-invariant switched family cannot select even one such target
rule.  The concrete hidden-bit target
\[
  \texttt{boolZProjectionRule}(z,a,b)=z
\]
is registered as a separating rule, yielding
\[
  \texttt{not\_exists\_predict\_eq\_boolZProjectionRule\_of\_abVisibleInvariant}.
\]
Thus any same-route repair that discards \(z\) must prove that the target
predictors themselves are constant on every \(z\)-fiber over fixed \((a,b)\);
otherwise the repair misses a named exact-surface target, not merely a large
classifier class.

The support-level version now rules out a weaker recovery story as well.  If a
distribution \(\mu\) assigns positive mass to two exact inputs \(u,v\) with the
same \((a,b)\)-data and a target \(T\) separates them, then every
\((a,b)\)-invariant predictor must disagree with \(T\) at some positive-mass
input.  The generic theorem is
\[
  \texttt{exists\_pos\_mass\_disagreement\_of\_abVisibleInvariant\_predict\_of\_abFiberSeparation},
\]
and the concrete hidden-bit specialization is
\[
  \texttt{exists\_pos\_mass\_disagreement\_of\_abVisibleInvariant\_predict\_boolZProjectionRule}.
\]
Consequently, a repair cannot weaken exact realization to distributional
recovery on any proof-relevant measure that keeps both hidden-\(z\) sides of one
visible fiber in support.  It must either show that the target is fiber-constant
on the exact surface or prove that the proof distribution removes one side of
every target-separating hidden-\(z\) fiber.  The crux ledger records this as the
narrow subrepair
\[
  \texttt{abVisibleInvariantSupportDisagreementObstruction}.
\]
The newest bridge composes this support fact with the PMF agreement-mass
interface.  Under the same positive-mass same-fiber hypotheses, every
\((a,b)\)-invariant predictor has agreement mass strictly below \(1\):
\[
  \texttt{agreementMass\_lt\_one\_of\_abVisibleInvariant\_predict\_of\_abFiberSeparation}.
\]
For the hidden-bit target this becomes
\[
  \texttt{agreementMass\_lt\_one\_of\_abVisibleInvariant\_predict\_boolZProjectionRule}.
\]
The family-level contrapositive is now named as well:
\[
  \texttt{not\_exists\_agreementMass\_eq\_one\_of\_abVisibleInvariant\_of\_abFiberSeparation},
\]
with the hidden-bit specialization
\[
  \texttt{not\_exists\_agreementMass\_eq\_one\_boolZProjectionRule\_of\_abVisibleInvariant}.
\]
Thus an exact \((a,b)\)-only repair cannot be repackaged as perfect
distributional recovery by choosing a different member of the same invariant
family on any measure supporting both hidden-bit alternatives in one visible
fiber.

The new module \texttt{SharedExactZABSparseThresholdSharpness.lean} shows that
the shared exact sparse-threshold \((zfeat(z),a,b)\) branch is genuinely live
once the shared basis is allowed to be large.  On raw \(\texttt{BitVec}\ d\),
Lean constructs the full point-block basis
\(\texttt{allAffinePointBlockFeatures}\ d\), together with an explicit mask and
threshold witness \(\texttt{pointBlockSparseThresholdCode}\), and proves
\[
  \texttt{rawSharedSparseThresholdPredict\_allAffinePointBlockFeatures\_eq},
\]
so every Boolean rule on \(\texttt{BitVec}\ d\) is realized by one shared
sparse-threshold code on that basis.  Pushed back through the exact visible-data
bijection for identity extractors, this yields
\[
  \texttt{fullExactVisibleRuleFamily\_realizedBySharedExactZABSparseThresholdAffineFamily\_id\_allAffinePointBlockFeatures}
\]
for \(Z=\texttt{BitVec}\ n\) whenever \(n+2k>0\).  So the earlier theorem
\[
  \texttt{fullExactVisibleRuleFamily\_not\_realizedBySharedExactZABSparseThresholdAffineFamily\_of\_linearBudget\_lt\_surfaceCard}
\]
is only a small-budget obstruction.  It does not close the sparse-threshold
repair surface outright.  The remaining burden on that branch is to prove that
the manuscript's actual switched family lands in a much smaller sparse-threshold
subclass than this fully expressive point-block realization.

The follow-up module
\texttt{SharedExactZABSparseThresholdInjectiveLiveness.lean} removes the
identity-extractor restriction.  Lean proves
\[
  \texttt{exactZABVisibleData\_injective\_of\_injective\_zfeat}
\]
and then pulls the same point-block witness back along any injective visible
map \((z,a,b)\mapsto(zfeat(z),a,b)\), yielding
\[
  \texttt{fullExactVisibleRuleFamily\_realizedBySharedExactZABSparseThresholdAffineFamily\_of\_injective\_zfeat}.
\]
So on the injective side the shared exact sparse-threshold route is not merely
alive at the native \(\texttt{BitVec}\) surface; it stays fully expressive for
every injective extractor.  Combined with the previously formalized
noninjective-extractor collision obstruction, this means extractor-width
counting alone cannot kill the sparse-threshold branch.  Any remaining crux for
that branch must prove that the manuscript's actual switched family lies in a
strictly smaller sparse-threshold subclass than this ambient injective
realization.

The next module
\texttt{SharedExactZABSparseThresholdInjectiveERMLiveness.lean} upgrades this
from ambient class expressivity to the manuscript-shaped ERM route itself.
Write
\[
  p=\texttt{allAffinePointBlockFeatureCount}(r+2k), \qquad
  U=\texttt{ExactVisiblePostSwitchSurface}\ Z\ k.
\]
Under the standard finite counting gap
\[
  2^{2p}(\lvert U\rvert-1)^m < \lvert U\rvert^m,
\]
Lean proves
\[
  \texttt{exists\_samples\_sharedExactZABSparseThresholdAffineERMFamily\_eq\_fullExactVisibleRuleFamily\_of\_injectiveExactZABVisibleData}
\]
and the extractor-level corollary
\[
  \texttt{exists\_samples\_sharedExactZABSparseThresholdAffineERMFamily\_eq\_fullExactVisibleRuleFamily\_of\_injective\_zfeat}.
\]
So whenever the visible summary \((z,a,b)\mapsto(zfeat(z),a,b)\) is injective,
there exists one labeled sample for each exact visible rule such that ERM over
the shared exact sparse-threshold point-block family returns that rule exactly,
and the resulting indexed ERM family is literally equal to
\(\texttt{fullExactVisibleRuleFamily}\).  This rules out a possible repair
story in which the branch remained expressive only as a raw hypothesis class
but collapsed once the manuscript's ERM wrapper was imposed.  Any remaining
obstruction on the sparse-threshold branch must therefore use extra semantic
restrictions on the sampled subclass, not merely the passage from class
realizability to ERM realizability.

The new consolidation module
\texttt{SharedExactZABSparseThresholdERMCharacterization.lean} packages the two
sides into one route statement.  First Lean proves the converse helper
\[
  \texttt{injective\_zfeat\_of\_injective\_exactZABVisibleData},
\]
so injectivity of the full visible map \((z,a,b)\mapsto(zfeat(z),a,b)\) is
equivalent to injectivity of \(zfeat\) itself.  Then it proves
\[
  \texttt{exists\_samples\_sharedExactZABSparseThresholdAffineERMFamily\_eq\_fullExactVisibleRuleFamily\_iff\_injective\_exactZABVisibleData}
\]
and the extractor-level form
\[
  \texttt{exists\_samples\_sharedExactZABSparseThresholdAffineERMFamily\_eq\_fullExactVisibleRuleFamily\_iff\_injective\_zfeat}.
\]
So, under the same point-block width and finite counting gap, the sparse-threshold
ERM route is neither partly dead nor partly alive: it is exactly an
injective-extractor route.  Noninjective extractors are now formally closed at
the ERM equality surface, while injective extractors remain fully live.  That
sharpens the remaining burden again: any further crux must either prove that
the manuscript's actual extractor is noninjective on the exact surface or prove
that the manuscript's switched family does not in fact land in this point-block
shared sparse-threshold ERM class.

The manuscript-facing bridge
\texttt{ActualSwitchedLocalSharedExactZABSparseThresholdERMData.lean} now
packages exactly that second obligation.  It defines
\texttt{ActualSwitchedLocalInterface.SharedExactZABSparseThresholdERMData},
which says that one actual switched-local predictor family is literally one
shared exact sparse-threshold ERM family on the canonical point-block basis.
Lean then proves the generic obstruction
\[
  \texttt{ActualSwitchedLocalInterface.SharedExactZABSparseThresholdERMData.not\_surjective\_predict\_of\_not\_injective\_zfeat}
\]
and, at the full-rule equality surface,
\[
  \texttt{ActualSwitchedLocalInterface.injective\_zfeat\_of\_eq\_fullExactVisibleRuleFamily}.
\]
So any actual switched-local sparse-threshold ERM realization of the full
exact visible rule family already forces an injective extractor.  On the
concrete \(\texttt{BitVec}\ n\) surface, Lean specializes this to
\[
  \texttt{ActualSwitchedLocalInterface.fullRuleActualSwitchedLocalInterface\_not\_nonempty\_sharedExactZABSparseThresholdERMData\_of\_width\_lt},
\]
which rules out every width-compressing extractor \(\texttt{BitVec}\ n \to
\texttt{BitVec}\ r\) with \(r<n\) for this manuscript-facing route.  This does
not close the injective sparse-threshold branch, but it removes the last
ambiguity about what the actual switched-local packet still has to supply: it
must either provide an injective extractor or leave the point-block shared
sparse-threshold ERM class.

The follow-up characterization file
\texttt{ActualSwitchedLocalSharedExactZABSparseThresholdERMCharacterization.lean}
adds the missing liveness direction at the same manuscript-facing level.  Under
the same point-block width and finite counting gap, Lean first proves the
generic injective-liveness theorem
\[
  \texttt{ActualSwitchedLocalInterface.nonempty\_sharedExactZABSparseThresholdERMData\_of\_injective\_zfeat},
\]
which reindexes the full-rule ERM witness samples along the current predictor
map.  So on the injective branch the manuscript-facing data wrapper is not
special to the full-rule counterexample: once \(zfeat\) is injective, Lean can
realize \emph{any} actual switched-local predictor family as one shared exact
sparse-threshold ERM family on the canonical point-block basis, provided the
same finite counting-gap packet holds.  The file then specializes this generic
construction to the full-rule endpoint and proves
\[
  \texttt{ActualSwitchedLocalInterface.fullRuleActualSwitchedLocalInterface\_nonempty\_sharedExactZABSparseThresholdERMData\_of\_injective\_zfeat},
\]
and then packages the full equivalences
\[
  \texttt{ActualSwitchedLocalInterface.fullRuleActualSwitchedLocalInterface\_nonempty\_sharedExactZABSparseThresholdERMData\_iff\_injective\_exactZABVisibleData}
\]
and
\[
  \texttt{ActualSwitchedLocalInterface.fullRuleActualSwitchedLocalInterface\_nonempty\_sharedExactZABSparseThresholdERMData\_iff\_injective\_zfeat}.
\]
So the full-rule actual switched-local sparse-threshold ERM packet is not just
partially constrained: at this abstraction level it is exactly an
injective-extractor packet.  The noninjective side is closed by the previous
bridge file, and the injective side is genuinely live even after asking for the
manuscript-facing actual-local data wrapper.  Any remaining crux on this branch
must therefore attack either the extractor itself or the claim that the actual
switched family lands in this point-block shared sparse-threshold ERM class; it
can no longer come from a gap between ambient ERM realizability and the
manuscript's actual-local packaging.

The next file
\texttt{ActualSwitchedLocalSharedExactZABSparseThresholdERMWidthCharacterization.lean}
removes the last width-level ambiguity on the native \(\texttt{BitVec}\ n\)
surface.  Lean defines an explicit low-bit-preserving embedding
\(\texttt{BitVec}\ n \to \texttt{BitVec}\ r\) for every \(n \le r\), proves
\[
  \texttt{ActualSwitchedLocalInterface.exists\_injective\_bitVecExtractor\_iff\_width\_le},
\]
and then packages the manuscript-facing existential route theorem
\[
  \texttt{ActualSwitchedLocalInterface.fullRuleActualSwitchedLocalInterface\_exists\_sharedExactZABSparseThresholdERMData\_iff\_width\_le}.
\]
So, under the same point-block width and finite counting gap as before, the
full-rule actual switched-local sparse-threshold ERM packet exists for some
width-\(r\) extractor exactly when \(n \le r\).  This is still not a crux
closure: it proves the surviving branch is precisely the non-compressing
injective-extractor branch.  Any stronger obstruction must therefore use more
than width bookkeeping or the generic actual-local sparse-threshold ERM wrapper.

The next two files sharpen the same packet in a different direction.
\texttt{BitFamilyUniformRecoveryLowerBound.lean} records the generic reverse
counting arithmetic: on any finite input domain with at least two points, the
finite recovery gap
\[
  2^s (|X|-1)^m < |X|^m
\]
already forces \(s < m\).  \texttt{ActualSwitchedLocalSharedExactZABSparseThresholdERMSampleLowerBound.lean}
then specializes that fact to the point-block sparse-threshold route.  Lean
proves that the standard point-block gap hypothesis used by the current
liveness packet implies
\[
  m > 2 \cdot \texttt{allAffinePointBlockFeatureCount}(r + 2k)
    = 2^{2(r+2k)+1},
\]
and therefore, via the earlier width characterization \(n \le r\), any
nontrivial full-rule actual switched-local sparse-threshold ERM packet on
\(\texttt{BitVec}\ n\) forces
\[
  m > 2^{2(n+2k)+1}.
\]
This is still not a crux closure.  It is an honest sample-growth consequence of
the current finite counting-gap liveness theorem, not an unconditional
impossibility theorem for the injective sparse-threshold branch itself.

The next bridge pushes that consequence off the special full-rule counterexample
and onto the bounded-sample plug-in-majority actual-local endpoint.  First,
\texttt{ActualSwitchedLocalSharedExactZABSparseThresholdERMSampleLowerBound.lean}
factors the previous argument through the reusable theorem
\[
  \texttt{ActualSwitchedLocalInterface.visibleWidthBudget\_lt\_sampleSize\_of\_nonempty\_sharedExactZABSparseThresholdERMData\_of\_surjective\_predict},
\]
which applies to any surjective actual switched-local family carrying the same
shared exact sparse-threshold ERM data.  Then
\texttt{ActualSwitchedLocalBoundedSampleSharedExactZABSparseThresholdERMSampleBound.lean}
combines that with the earlier characterization of bounded sample-level plug-in
majority: if the sample bound already reaches the exact-visible alphabet size,
so that the bounded-sample endpoint is full-rule expressive, and if that same
endpoint is also identified with one shared exact sparse-threshold ERM family,
then the same point-block counting-gap packet at \(m = \texttt{sampleBound}\)
forces
\[
  \texttt{sampleBound} > 2^{2(n+2k)+1}.
\]
Equivalently, under that same gap hypothesis, no such bounded-sample
manuscript-facing packet can exist at or below the visible-width budget.  This
is still conditional route-diagnostic work rather than a closure of the
surviving injective sparse-threshold branch: it only rules out bounded-sample
actual-local packets that are both full-rule expressive and demanded to satisfy
the current finite counting-gap liveness inequality at that same bound.

The latest strengthening removes even the need to assume that a bounded-sample
packet witness was supplied ahead of time.  Lean now proves
\[
  \texttt{ActualSwitchedLocalInterface.boundedSamplePluginMajorityActualSwitchedLocalInterface\_visibleWidthBudget\_lt\_sampleBound\_of\_visibleSurfaceCard\_le\_of\_width\_le}.
\]
Here the only extra route-side hypothesis is the noncompressing width condition
\(n \le r\): Lean inserts the explicit low-bit embedding
\(\texttt{bitVecPrefixEmbedding} : \texttt{BitVec}\ n \to \texttt{BitVec}\ r\),
uses the new generic injective-liveness theorem to manufacture the
manuscript-facing sparse-threshold packet automatically on the bounded-sample
endpoint, and then re-runs the same sample-growth argument.  So once bounded
sample plug-in majority is large enough to be full-rule expressive, the
surviving injective-width branch itself already forces
\(\texttt{sampleBound} > 2^{2(n+2k)+1}\) under the current counting-gap packet;
the theorem no longer depends on separately postulating a bounded-sample ERM
equality witness.

The next bounded-sample refinement,
\texttt{ActualSwitchedLocalBoundedSampleSharedExactZABSparseThresholdERMSampleWindowObstruction.lean},
repackages the same counting-gap packet into a genuine forbidden-window
statement.  On \(\texttt{BitVec}\ n\), any bounded-sample plug-in-majority
actual-local packet identified with one shared exact sparse-threshold ERM
family must satisfy either
\[
  \texttt{sampleBound} < 2^{n+2k}
\]
or
\[
  2^{2(n+2k)+1} < \texttt{sampleBound}.
\]
So the whole intermediate regime
\[
  2^{n+2k} \le \texttt{sampleBound} \le 2^{2(n+2k)+1}
\]
is impossible under that same counting-gap hypothesis.  Lean also lifts the
same obstruction to the manuscript-facing recovery packet through the
\texttt{.data} field.  This remains route-facing obstruction work rather than
branch closure: it strengthens the same finite counting-gap consequence into a
forbidden bounded-sample window, but does not touch packets outside that
interval.

The smaller follow-up file
\texttt{ActualSwitchedLocalBoundedSampleSharedExactZABSparseThresholdERMRecoverySampleBoundObstruction.lean}
then lifts the earlier large-sample bridge directly to the manuscript-facing
recovery endpoint itself.  No new counting argument is used: recovery data
already contains the underlying exact-family packet, so once the bounded-sample
plug-in-majority endpoint is large enough for surjectivity, the same point-block
counting-gap hypothesis at \(m = \texttt{sampleBound}\) still forces
\[
  \texttt{sampleBound} > 2^{2(n+2k)+1}.
\]
Equivalently, no bounded-sample manuscript-facing recovery packet can exist at
or below that visible-width budget under the same counting-gap packet.  This is
not a new closure theorem, but it closes a manuscript-facing loophole: the
large-sample burden now hits the recovery packet directly, not only the
underlying data packet or the later forbidden-window repackaging.

The next file,
\texttt{ActualSwitchedLocalBoundedSampleSharedExactZABSparseThresholdERMVisibleBudgetObstruction.lean},
adds the unconditional bounded-sample analogue of the later actual-local
visible-budget obstruction.  Once the bounded sample bound reaches the
exact-visible alphabet size, so that the bounded-sample plug-in-majority
endpoint is surjective, any shared exact sparse-threshold ERM packet on that
endpoint already forces
\[
  \left|\texttt{ExactVisiblePostSwitchSurface}(Z,k)\right|
  \le
  2 \cdot \texttt{allAffinePointBlockFeatureCount}(r+2k).
\]
On \(\texttt{BitVec}\ n\), this again specializes to
\[
  n \le 2r + 2k + 1.
\]
So if \(2r + 2k + 1 < n\), then the bounded-sample manuscript-facing endpoint
cannot carry the shared exact sparse-threshold ERM packet at all, provided only
that its sample bound is large enough for surjectivity.  Unlike the previous
bounded-sample sample-growth theorem, this one uses neither the counting-gap
liveness hypothesis nor any recovery-side agreement estimate.  It is still not
a full closure of the injective branch, but it pushes the unconditional
half-width obstruction all the way to the bounded-sample endpoint.

The next file,
\texttt{ActualSwitchedLocalSharedExactZABSparseThresholdERMRecoveryData.lean},
adds the missing manuscript-facing recovery wrapper for this same point-block
route.  It defines actual-local recovery data to be the earlier exact-family
packet together with one uniform bad-code agreement bound.  Lean then shows that
such data automatically induce the corresponding exact-visible shared
sparse-threshold recovery packet, and therefore close all of the standard
consequences at the actual-local layer:
\begin{itemize}
  \item the exact-visible compression target at budget
  \(2 \cdot \texttt{allAffinePointBlockFeatureCount}(r+2k)\),
  \item a finite predictor cover / representative cover / quotient of size
  \(2^{\,2 \cdot \texttt{allAffinePointBlockFeatureCount}(r+2k)}\),
  \item the bundled clocked \(K\)-poly finite-learning payload at that same
  budget, and
  \item the weighted exact-recovery lower bound
  \[
    1 - 2^{\,2 \cdot \texttt{allAffinePointBlockFeatureCount}(r+2k)} q^m
    \le \texttt{bitExactRecoverySampleMass}(\mu,\dots,m).
  \]
\end{itemize}
The file also inherits the earlier extractor obstructions: noninjective
\(\texttt{zfeat}\), and in particular width-compressing
\(\texttt{BitVec}\ n \to \texttt{BitVec}\ r\) maps with \(r<n\), make even this
stronger recovery packet impossible on surjective actual-local families.  This
is again route-facing infrastructure rather than a new crux closure.  It shows
that once the manuscript can honestly supply the bad-code agreement bound, no
extra finite-image or clocked-payload bookkeeping remains to be proved for this
branch; the remaining burden stays concentrated on the agreement estimate and
the injective extractor regime itself.

The next file,
\texttt{ActualSwitchedLocalSharedExactZABSparseThresholdERMRecoveryAtomObstruction.lean},
turns that recovery wrapper into a genuine obstruction for overly concentrated
recovery laws.  Suppose the actual switched-local family is surjective on the
exact visible surface, the recovery layer carries a bad-code agreement bound
\(q\), and there is one visible point \(x\) whose sampling mass already
exceeds that bound:
\[
  \mu\{x\} > q.
\]
If the exact visible surface also contains some different point \(y \ne x\),
then Lean now shows that no
\texttt{ActualSwitchedLocalInterface.SharedExactZABSparseThresholdERMRecoveryData}
packet can exist.  The proof is direct: surjectivity supplies one local rule
that is constantly false and another that spikes exactly at \(y\); the shared
sparse-threshold target-data packet gives a bad code for the former whose
decoded predictor agrees with it at \(x\); and the earlier region-mass
obstruction then contradicts the bad-code agreement bound because the region
\(\{x\}\) already has mass \(> q\).

This has an immediate pure-atom corollary.  If \(\mu = \delta_x\) is a pure
point mass and \(q<1\), then every surjective nontrivial visible surface is
already impossible for the shared exact sparse-threshold recovery packet.
Lean also lifts the same obstruction to the bounded-sample plug-in-majority
endpoint and to the full-rule \(\texttt{BitVec}\ n\) endpoint whenever
\(0 < n + 2k\), since that guarantees the exact visible surface contains at
least two distinct points.  This is real obstruction work, but it is not full
branch closure: it kills atomic or too-heavy recovery laws, not the remaining
non-atomic injective recovery regime.

The next file,
\texttt{ActualSwitchedLocalSharedExactZABSparseThresholdERMRecoveryConcentrationObstruction.lean},
sharpens that atom obstruction into a surface-wide concentration bound.  Let
\(\mu\) be the recovery law on the exact visible surface, assume the actual
switched-local family is surjective, and suppose the visible surface contains
two distinct points.  Fix any visible point \(y\).  Surjectivity again supplies
the singleton spike at \(y\) and the constant-false target, so the bad-code
agreement packet forces the decoded bad code to disagree with the false target
on at most a \(q\)-mass region.  But outside \(\{y\}\) the spike is false, so
the whole complement already lies inside that agreement region.  Lean therefore
proves
\[
  1 - \mu\{y\} \le q
\]
for every visible point \(y\).  On a nontrivial visible surface this implies
that not all visible points can carry mass \(> 1/2\); choosing one with
\(\mu\{y\} \le 1/2\) yields the unconditional recovery-side lower bound
\[
  \frac12 \le q.
\]
So no surjective actual-local sparse-threshold recovery packet can exist once
\(q < 1/2\) and the exact visible surface has at least two points.  Lean lifts
this obstruction to the bounded-sample plug-in-majority endpoint and to the
full-rule \(\texttt{BitVec}\ n\) endpoint whenever \(0 < n + 2k\), since that
condition guarantees a nontrivial exact visible surface.  The regression file
also records the sharp two-point Bool canary: under the uniform recovery law on
\(\texttt{ExactVisiblePostSwitchSurface Bool 0}\), any full-rule recovery packet
must satisfy \(q \ge 1/2\), and in particular \(q = 0\) is impossible.  This is
again genuine obstruction work, but still not branch closure: it rules out the
entire \(q < 1/2\) recovery regime on nontrivial visible surfaces, while the
remaining \(\,q \ge 1/2\) injective regime stays open.

The next file,
\texttt{ActualSwitchedLocalSharedExactZABSparseThresholdERMRecoveryCardinalityObstruction.lean},
upgrades that two-point obstruction to the full finite surface cardinality.
The helper file \texttt{FinitePMFPointMassBound.lean} proves that every
probability mass function on a finite nonempty type has some point \(y\) with
\[
  \mu\{y\} \le |S|^{-1}.
\]
Combining that averaging fact with the earlier surjective recovery inequality
\(1 - \mu\{y\} \le q\) yields the sharper universal lower bound
\[
  1 - |S|^{-1} \le q
\]
for every surjective actual-local sparse-threshold recovery packet on a finite
exact visible surface \(S\).  Lean lifts that obstruction to the bounded-sample
plug-in-majority endpoint and to the full-rule endpoint.  On the full-rule
\(\texttt{BitVec}\ n\) surface this specializes to
\[
  1 - 2^{-(n+2k)} \le q.
\]
So every \(q\) strictly below the uniform-complement threshold is impossible.
The regression file records the concrete uniform \(\texttt{BitVec}\ 1, k=1\)
canary: any full-rule recovery packet must satisfy
\(q \ge 1 - 2^{-3} = 7/8\), and in particular \(q=0\) is already impossible at
both the full-rule and bounded-sample surjective endpoints.  This is again real
obstruction work, but still not branch closure: it removes the whole
sub-uniform-complement recovery regime while leaving the surviving near-perfect
injective regime open.

The follow-up file,
\texttt{ActualSwitchedLocalSharedExactZABSparseThresholdERMRecoveryThresholdWidthObstruction.lean},
turns that finite-threshold lower bound around.  Suppose some manuscript route
claims a stronger recovery estimate
\[
  q < 1 - 2^{-p}.
\]
Lean then proves that the full-rule \(\texttt{BitVec}\ n\) endpoint must satisfy
\[
  n + 2k < p.
\]
Equivalently, whenever \(p \le n + 2k\), no such full-rule recovery packet can
exist.  The same inverted threshold obstruction lifts to the bounded-sample
plug-in-majority endpoint once the sample bound is large enough for
surjectivity.  So any concrete recovery estimate below the finite
uniform-complement floor automatically becomes a direct visible-width blocker.
The regression file records the uniform \(\texttt{BitVec}\ 1, k=1\) canary in
both directions: hypothetical \(q=1/2\) data already forces \(1+2<4\), and
setting \(p=3\) rules out both the full-rule and bounded-sample endpoints
outright.  This is again real obstruction work, though still not branch
closure: it converts subthreshold recovery claims into width constraints, while
the surviving near-perfect regime with sufficiently large \(p\) remains open.

The boundary-refinement file,
\texttt{ActualSwitchedLocalSharedExactZABSparseThresholdERMRecoveryThresholdVisibleBudgetObstruction.lean},
then combines that threshold-width inversion with the earlier unconditional
surjective ceiling \(n \le 2r + 2k + 1\).  If a manuscript route claims the
strict recovery estimate
\[
  q < 1 - 2^{-(2r + 4k + \mathrm{slack} + 1)},
\]
Lean upgrades the full-rule \(\texttt{BitVec}\ n\) budget to
\[
  n \le 2r + 2k + \mathrm{slack}.
\]
The bounded-sample plug-in-majority endpoint satisfies the same sharpened
budget once the sample bound is large enough for surjectivity.  In particular,
setting \(\mathrm{slack}=0\) kills the old boundary case
\(n = 2r + 2k + 1\) as soon as the route claims
\(q < 1 - 2^{-(2r + 4k + 1)}\).  The regression file records the concrete
uniform canaries: for \(\texttt{BitVec}\ 1, k=1\), hypothetical \(q=1/2\)
data already forces \(1 \le 2\); for \(\texttt{BitVec}\ 3, k=0, r=1\), the
same \(q=1/2\) threshold rules out both the full-rule and bounded-sample
surjective endpoints.  This is still not branch closure, but it removes the
former one-step visible-budget escape hatch whenever the recovery estimate is
strong enough.

The next recovery-side file,
\texttt{ActualSwitchedLocalSharedExactZABSparseThresholdERMRecoveryRegionObstruction.lean},
attacks the surviving injective branch from a different angle.  The exact-visible
wrapper
\texttt{SharedExactZABSparseThresholdERMRecoveryData.regionMass\_le\_of\_badCode\_agrees\_on}
already says that any one bad code can agree with the target on at most a
\(q\)-mass finite region.  On the injective branch, however,
\texttt{SharedExactZABSparseThresholdInjectiveLiveness.lean} proves that the
shared exact sparse-threshold family realizes \emph{every} Boolean rule on the
visible surface.  Lean therefore takes any actual-local target predictor, any
proper finite region \(R\), and any visible point \(x_0 \notin R\), flips the
target only at \(x_0\), and realizes that altered rule by a shared sparse-threshold
code.  The resulting bad code still agrees with the target on all of \(R\), so
the recovery packet must satisfy
\[
  \mu(R) \le q
\]
for every proper finite region \(R\).  The same obstruction lifts not only to
the full-rule endpoint but also to bounded sample plug-in majority, and here
the bounded-sample corollary needs no surjectivity hypothesis at all: the generic
actual-local theorem works with any chosen bounded-sample index.

This is real obstruction work, but still not branch closure.  It does \emph{not}
show that the injective branch is impossible for every finite PMF; it shows that
the branch can survive only when the advertised threshold \(q\) dominates the
mass of every proper finite region.  That is genuinely stronger than the earlier
uniform cardinality floor on nonuniform measures.  The regression file
\texttt{ActualSwitchedLocalSharedExactZABSparseThresholdERMRecoveryRegionObstructionRegression.lean}
uses the two-point \(\texttt{BitVec}\ 1, k=0\) visible surface with a skew
\(4/5\)--\(1/5\) PMF: the heavy singleton region already forces
\[
  4/5 \le q,
\]
and the concrete threshold \(q = 3/4\) rules out both the full-rule endpoint
and the bounded-sample endpoint at sample bound \(2\).  So the new obstruction
is strictly sharper than the old \(1 - 1/|U| = 1/2\) uniform-cardinality floor
on that same surface, while still leaving the near-perfect injective regime
open above the skew region masses.

The next cleanup,
\texttt{FinitePMFHeavyRegionCharacterization.lean} together with
\texttt{ActualSwitchedLocalSharedExactZABSparseThresholdERMRecoveryHeavyRegionCharacterization.lean},
shows that the preceding obstruction is already intrinsic and does not really
depend on searching over arbitrary finite regions.  On any finite PMF, Lean now
proves the exact equivalence
\[
  \exists R \subsetneq U \text{ finite with } q < \mu(R)
  \quad\Longleftrightarrow\quad
  \exists y \in U,\; q < 1 - \mu(y).
\]
So the injective recovery obstruction can be repackaged canonically as a
light-point criterion: a heavy proper region exists exactly when one visible
point has too small a mass.  The new recovery theorem reduces the earlier
arbitrary-region obstruction to this point-complement form, and the regression
file verifies the equivalence concretely on the same skew two-point
\(\texttt{BitVec}\ 1, k=0\) PMF.  In particular, the genuinely non-surjective
bounded-sample endpoint at sample bound \(0\) is already ruled out by the
single inequality
\[
  3/4 < 1 - 1/5 = 4/5.
\]
This is not new branch closure, but it does tighten the earlier narrative:
the surviving injective branch is constrained by point-complements, not merely
by an existential search over regions.

Lean now sharpens that characterization one step further.  The same finite-PMF
file proves that the existential point-complement condition is equivalent to a
single canonical threshold
\[
  q < 1 - \mu(\mathrm{lightestPoint}(\mu)),
\]
where \(\mathrm{lightestPoint}(\mu)\) is a visible point of minimal mass.
So the route obstruction does not just say ``some point is too light''; it is
already controlled by the worst visible atom of \(\mu\).  The recovery file
packages this directly onto bounded-sample plug-in majority as well: the empty
sample index is enough to force
\[
  1 - \mu(\mathrm{lightestPoint}(\mu)) \le q
\]
on the injective branch, with no surjectivity assumption.  On the skew
two-point \(\texttt{BitVec}\ 1, k=0\) PMF this canonical threshold is exactly
\(1 - 1/5 = 4/5\), so the genuinely non-surjective sample-bound-\(0\) endpoint
is ruled out already by the single inequality \(3/4 < 4/5\).

Lean now makes the structural meaning of that threshold fully explicit.  The
same finite-PMF file proves that the \emph{entire family} of proper-region
mass inequalities is equivalent to this one canonical bound:
\[
  \forall R \subsetneq U,\; \mu(R) \le q
  \qquad\Longleftrightarrow\qquad
  1 - \mu(\mathrm{lightestPoint}(\mu)) \le q.
\]
So the strongest proper-region obstruction is not hidden somewhere in a large
search space of finite regions; it is attained by the single extremal region
\(U \setminus \{\mathrm{lightestPoint}(\mu)\}\).  The route-facing recovery
file imports this equivalence directly, so the injective manuscript packet is
now seen to satisfy the whole proper-region family exactly up to the worst-atom
threshold.  On the skew two-point \(\texttt{BitVec}\ 1, k=0\) regression this
means, concretely, that the claim ``every proper region has mass at most
\(3/4\)'' already fails for the same sharp reason \(3/4 < 1 -
\mu(\mathrm{lightestPoint}(\mu))\).

The follow-up file
\texttt{ActualSwitchedLocalSharedExactZABSparseThresholdERMRecoveryInjectiveConcentrationObstruction.lean}
closes the most obvious gap in those recovery thresholds.  The earlier
concentration and cardinality obstructions only applied on the \emph{surjective}
branch, because they obtained bad codes by realizing arbitrary singleton spikes.
The injective-region obstruction is already strong enough to recover the same
numerics without any surjectivity assumption.  Lean simply takes the region
\(U \setminus \{y\}\): if \(zfeat\) is injective and one actual-local index
\(i\) is present, then the point-complement bad code agrees with predictor
\(i\) on all of \(U \setminus \{y\}\), so every injective recovery packet must
satisfy
\[
  1 - \mu(y) \le q.
\]
Combining this with the finite point-mass pigeonhole bound recovers
\[
  1 - |U|^{-1} \le q
\]
for \emph{every} injective manuscript-facing recovery packet, even when the
actual-local family is very far from surjective.  The regression file
\texttt{ActualSwitchedLocalSharedExactZABSparseThresholdERMRecoveryInjectiveConcentrationObstructionRegression.lean}
makes that distinction explicit with a one-index constant-false actual-local
family on \(\texttt{ExactVisiblePostSwitchSurface}(\texttt{BitVec}\ 1,1)\): the
family is not surjective, yet hypothetical injective recovery data already
forces
\[
  1 - 2^{-3} = 7/8 \le q,
\]
so \(q = 3/4\) is impossible there as well.  This is real route-facing
obstruction work, but still not branch closure: it removes the surjectivity
escape hatch for the injective concentration/cardinality floor, while the
surviving near-perfect regime remains open.

The next bounded-sample follow-up removes the last surjectivity assumption from
the manuscript-facing threshold inversion itself.  The older bounded-sample
threshold-width file still needed the sample-size hypothesis
\(|U| \le \texttt{sampleBound}\) to make the endpoint surjective before it
could conclude that
\[
  q < 1 - 2^{-p}
  \quad\Longrightarrow\quad
  n + 2k < p.
\]
But on the injective branch there is no need to realize every Boolean rule via
bounded samples: Lean can simply fix the explicit bounded-sample predictor
index \(\langle [], \texttt{by simp}\rangle\), apply the new injective
concentration theorem to that one actual-local predictor, and recover the same
uniform threshold
\[
  1 - 2^{-(n + 2k)} \le q
\]
with no lower bound on \(\texttt{sampleBound}\).  The file
\texttt{ActualSwitchedLocalBoundedSampleSharedExactZABSparseThresholdERMRecoveryInjectiveThresholdObstruction.lean}
then inverts this bound exactly as before, but now for the genuinely
non-surjective bounded-sample endpoint as well.  Its regression file
\texttt{ActualSwitchedLocalBoundedSampleSharedExactZABSparseThresholdERMRecoveryInjectiveThresholdObstructionRegression.lean}
uses the \(\texttt{BitVec}\ 1, k=1\) surface at sample bound \(0\): the
bounded-sample interface is formally not surjective, yet any hypothetical
injective recovery packet still forces \(7/8 \le q\), and the stricter
threshold \(q = 3/4\) is impossible.  This is again real obstruction work, but
not branch closure: it closes a bounded-sample technical escape hatch on the
injective branch, while the near-perfect regime still survives.

The next file,
\texttt{ActualSwitchedLocalSharedExactZABSparseThresholdERMBitCode.lean},
separates the weaker consequence that does \emph{not} depend on any recovery
estimate.  Lean packages the exact-family identification itself as an
actual-local bounded-code witness:
\[
  \texttt{ActualSwitchedLocalInterface.SharedExactZABSparseThresholdERMData.bitCodeData}.
\]
So once the actual switched-local family is exactly the point-block shared
exact sparse-threshold ERM family, each local rule is already decoded from one
fixed bit family of budget
\[
  2 \cdot \texttt{allAffinePointBlockFeatureCount}(r+2k).
\]
Therefore exact-family identification alone, before any agreement bound is
proved, already yields:
\begin{itemize}
  \item the exact-visible compression target at that budget,
  \item a finite predictor cover / representative cover / quotient of size
  \(2^{\,2 \cdot \texttt{allAffinePointBlockFeatureCount}(r+2k)}\), and
  \item the bundled clocked \(K\)-poly finite-learning payload at that same
  budget.
\end{itemize}
This sharpens the route decomposition.  The recovery-data wrapper is only
needed for the weighted exact-recovery lower bound; it is not needed for the
finite-image or clocked-payload side.  That is still route-facing
infrastructure rather than a new blocker, but it removes one more place where
the manuscript could hide missing bounded-code bookkeeping behind the recovery
argument.

The next file,
\texttt{ActualSwitchedLocalSharedExactZABSparseThresholdERMVisibleBudgetObstruction.lean},
turns that bounded-code packet into an unconditional surjectivity obstruction.
If an actual switched-local family is exactly the point-block shared exact
sparse-threshold ERM family and is still surjective on the exact visible
surface, then Lean now proves
\[
  \left|\texttt{ExactVisiblePostSwitchSurface}(Z,k)\right|
  \le
  2 \cdot \texttt{allAffinePointBlockFeatureCount}(r+2k).
\]
So the exact visible surface must fit inside the point-block sparse-threshold
bit budget before any recovery estimate or finite counting-gap hypothesis is
used.  On \(\texttt{BitVec}\ n\), this specializes to the explicit width bound
\[
  n \le 2r + 2k + 1.
\]
Equivalently, if \(2r + 2k + 1 < n\), then the full-rule actual switched-local
interface cannot carry the shared exact sparse-threshold ERM packet at all.
This is weaker than the earlier conditional \(r \ge n\) characterization of the
surviving injective branch, but it is unconditional: it does not rely on the
counting-gap liveness hypothesis and it does not require any recovery-side
agreement estimate.  So even before the route reaches the harder injective
extractor regime, Lean already rules out all extractors that compress below
roughly half the visible width.  The new route-facing wrapper now makes the
quantifier order explicit as well: below the same threshold, Lean proves
\[
  \neg \exists\, zfeat : \texttt{BitVec}\ n \to \texttt{BitVec}\ r,
    \ \texttt{SharedExactZABSparseThresholdERMData}(\texttt{fullRuleActualSwitchedLocalInterface}, zfeat).
\]
So the visible-budget obstruction is no longer only about a bad fixed extractor;
the full-rule sparse-threshold ERM route cannot rescue itself by choosing a
different extractor later.

The follow-up file
\texttt{ActualSwitchedLocalSharedExactZABSparseThresholdERMRecoveryVisibleBudgetObstruction.lean}
closes the obvious manuscript-facing interface gap in that argument.  Recovery
data already contains the exact-family data, so the same unconditional
half-width obstruction applies directly to the full recovery packet: any
surjective manuscript-facing recovery packet on \(\texttt{BitVec}\ n\) still
forces
\[
  n \le 2r + 2k + 1.
\]
Lean then records the corresponding endpoint corollaries for both the full-rule
actual switched-local interface and the bounded-sample plug-in-majority
interface.  Thus the unconditional visible-budget obstruction is not merely a
data-layer fact; it already kills the manuscript-facing recovery endpoints
below the same half-width threshold, with no \(q\)-estimate and no finite
counting-gap hypothesis.  The updated endpoint wrappers again remove the
extractor choice itself: below the same threshold, Lean proves that there does
not exist any extractor \(zfeat : \texttt{BitVec}\ n \to \texttt{BitVec}\ r\)
supporting the full-rule recovery packet, and likewise none supporting the
bounded-sample plug-in-majority recovery packet once the sample bound is large
enough for surjectivity.

Lean now records a different manuscript-facing consequence of the same recovery
packet: it always yields the bundled clocked \(K\)-poly finite-learning payload
at point-block budget
\[
  s = 2 \cdot \texttt{allAffinePointBlockFeatureCount}(r+2k).
\]
The new file
\texttt{ActualSwitchedLocalSharedExactZABSparseThresholdERMRecoveryPayloadObstruction.lean}
turns that directly into an obstruction theorem.  If the actual switched
predictor family already realizes an injectively indexed finite probe family of
size \(> 2^s\), then no manuscript-facing sparse-threshold recovery packet can
exist.  Lean then sharpens that to a route-native finite-family statement: if
the actual predictor map itself is injective on a finite index type, then
recovery already forces
\[
  |\texttt{Index}| \le 2^s.
\]
So the manuscript cannot hide an oversized actual-local family behind the
finite-learning layer merely by avoiding surjectivity onto the full Boolean
rule space.  The full-rule surjective endpoint is just one corollary; the more
important route-facing statements are the general injective-realization version
and its injective-index specialization.  The latest strengthening now removes
the extractor choice under the same payload obstruction.  Lean proves that once
the actual switched predictor family is too expressive for the payload budget,
there does not exist any extractor \(zfeat\) supporting the manuscript-facing
recovery packet.  This is recorded both for the full-rule surjective endpoint
and for arbitrary finite-index actual-local families whose predictor map is
already injective on too many indices.  The regression file makes this
distinction concrete on the three-point visible surface
\(\texttt{ExactVisiblePostSwitchSurface}(\texttt{Fin}\ 3,0)\): Lean defines an
actual-local family with exactly five indices and five realized predictors,
proves that this family is \emph{not} surjective onto the full eight-rule
Boolean space, and still rules out recovery because the point-block payload
budget there only permits \(2^2 = 4\) predictors.  The strengthened endpoint
now says more: that five-predictor family cannot be rescued by picking some
different \(\texttt{Fin}\ 3 \to \texttt{BitVec}\ 0\) extractor later.  This is
real obstruction work, but not branch closure: it closes a manuscript-facing
finite-learning loophole beyond the pure width/cardinality wrappers, while the
larger injective regime remains open.

Lean now records the same finite-image burden on the quotient-code side
directly.  The new file
\texttt{ActualSwitchedLocalSharedExactZABSparseThresholdERMRecoveryQuotientObstruction.lean}
starts from the fact that any recovery packet already yields a finite
predictor quotient of size at most \(2^s\) at the same point-block budget.
From there Lean proves a route-facing exact consequence: every injectively
realized finite probe family on the exact visible surface must already fit
inside that quotient budget.  Equivalently, if the actual predictor map itself
is injective on a finite index type, then recovery forces
\[
  |\texttt{Index}| \le 2^s
\]
even before one appeals to the clocked finite-learning payload semantics.
Thus a quotient-code repair is not a new escape hatch.  It owes exactly the
same small-image theorem on the actual switched family.  The latest
strengthening now removes the extractor choice here as well: once an injective
realized family or injective predictor index set already exceeds the quotient
budget, Lean proves that there does not exist any extractor \(zfeat\)
supporting the manuscript-facing recovery packet.  The regression file reuses
the same genuinely non-surjective \(\texttt{Fin}\ 5\) family on the three-point
visible surface and shows that recovery is already impossible at the quotient
layer, because five injectively realized predictors cannot be encoded into the
available four-code budget; moreover, that family cannot be rescued by picking
some different \(\texttt{Fin}\ 3 \to \texttt{BitVec}\ 0\) extractor later.
This is not branch closure either, but it closes a different manuscript-facing
loophole: one cannot hide an oversized realized family behind a proposed small
quotient presentation.

Lean now records a different recovery-side obstruction that genuinely uses the
agreement threshold \(q\), not the same \(2^s\) image budget again.  The new
file
\texttt{ActualSwitchedLocalSharedExactZABSparseThresholdERMRecoveryPairwiseAgreementObstruction.lean}
observes that once the actual switched family is identified with one fixed
shared exact sparse-threshold bit family, every other distinct realized
predictor is automatically a bad code against the current target.  Therefore
the recovery hypothesis itself forces a pairwise separation law:
for any two distinct realized predictors \(g_i \neq g_j\),
\[
  \texttt{agreementMass}(\mu,g_i,g_j) \le q.
\]
This is a different kind of burden from the earlier width, payload, or quotient
wrappers.  It says that recovery cannot coexist with even one realized pair
whose overlap under \(\mu\) is too large.  The regression file again reuses the
same genuinely non-surjective \(\texttt{Fin}\ 5\) family on the three-point
visible surface, but now with the pure measure concentrated at one visible
point.  Two distinct realized predictors in that family agree on full
\(\mu\)-mass there, so Lean rules out recovery already for \(q = 3/4\).  Thus
the route cannot escape merely by keeping the realized family small while still
allowing highly overlapping realized predictors under the chosen measure.

Lean now pushes that same \(q\)-dependence to a family-wide obstruction.  The
new file
\texttt{ActualSwitchedLocalSharedExactZABSparseThresholdERMRecoveryHeavyPointCardinalityObstruction.lean}
shows that if one visible point \(x\) already has mass \(\mu(x) > q\), then any
two realized predictors agreeing at \(x\) must in fact be equal globally.  Since
the actual outputs are Boolean, evaluation at that one heavy point can take only
two values, so any injectively indexed actual switched family supporting the
recovery packet must satisfy
\[
  |\mathrm{Index}| \le 2.
\]
This is not another bit-budget wrapper.  It is a measure-side collapse theorem:
one heavy atom above the claimed recovery threshold forces the whole realized
family through a single visible Boolean coordinate.  The regression file again
reuses the genuinely non-surjective \(\texttt{Fin}\ 5\) family with pure mass at
one visible point, and now rules out \(q = 3/4\) for the stronger reason that
an injective five-predictor family cannot survive a heavy-point threshold below
\(1\).

Lean now generalizes that collapse from one coordinate to any finite heavy
region.  The new file
\texttt{ActualSwitchedLocalSharedExactZABSparseThresholdERMRecoveryHeavyRegionCardinalityObstruction.lean}
proves that if a finite visible region \(R\) already has mass
\(\mu(R) > q\), then any two realized predictors agreeing on all of \(R\) must
already be equal globally.  So on any injectively indexed actual switched
family, restriction to \(R\) is itself injective, and the whole realized family
must inject into the Boolean trace space on \(R\):
\[
  |\mathrm{Index}| \le 2^{|R|}.
\]
This is genuinely stronger than the heavy-point theorem when no individual atom
exceeds \(q\) but a small region still does.  The regression file exhibits
exactly that situation on the three-point visible surface with split measure
\((1/2, 1/2, 0)\): no single point has mass above \(3/4\), so the heavy-point
collapse does not apply, yet the two-point region \(\{0,1\}\) already has mass
\(1\).  Lean then reuses the genuinely non-surjective \(\texttt{Fin}\ 5\)
family and rules out recovery at \(q = 3/4\) because an injective five-rule
family cannot fit into the four Boolean traces on a two-point region.

Lean now pushes that heavy-region barrier directly onto the manuscript's
full-rule endpoint, removing the extra injective-branch hypothesis that the
earlier proper-region file still required.  The new file
\texttt{ActualSwitchedLocalSharedExactZABSparseThresholdERMRecoveryFullRuleRegionObstruction.lean}
proves that if a full-rule recovery packet existed and one finite visible
region \(R\) satisfied \(|R| < |\text{surface}|\) together with \(\mu(R) > q\),
then the full Boolean rule family on the whole exact visible surface would have
to inject into the Boolean trace space on \(R\).  Lean formalizes that as the
mass bound
\[
  \mu(R) \le q
\]
for every proper finite region \(R\), with no injectivity or width assumption
on \(zfeat\).  The regression file reuses the skew \(\texttt{BitVec}\ 1\)
measure \((4/5, 1/5)\): Lean now derives \(4/5 \le q\) from any hypothetical
full-rule recovery packet by the new theorem itself, and therefore rules out
\(q = 3/4\) on that endpoint without appealing to the older injective-branch
construction.

Lean now pushes the same region barrier past the full-rule identity case to
arbitrary surjective actual switched families.  The new file
\texttt{ActualSwitchedLocalSharedExactZABSparseThresholdERMRecoverySurjectiveRegionObstruction.lean}
shows that if the actual predictor family is surjective onto all Boolean rules
on the exact visible surface, then every proper finite region still has mass at
most \(q\).  The proof is cleaner than the old injective-branch argument: Lean
chooses two realized predictors directly from surjectivity, namely the
constant-false rule and the same rule flipped at one point outside the region,
and then applies the heavy-region rigidity theorem to force a contradiction.
So the route cannot escape by keeping surjectivity while allowing duplicated
indices or other noninjective presentations of the realized family.  The latest
strengthening now removes the extractor choice on that same branch: once some
proper visible region has mass above \(q\), Lean proves that there does not
exist any extractor \(zfeat\) supporting the manuscript-facing recovery packet
on the surjective actual switched family.

The regression file
\texttt{ActualSwitchedLocalSharedExactZABSparseThresholdERMRecoverySurjectiveRegionObstructionRegression.lean}
makes that distinction explicit.  It defines a duplicated-index actual local
interface whose predictor map ignores one Boolean coordinate of the index, so
the family is still surjective onto the full exact visible rule class but is no
longer injective as an indexed family.  On the same skew \(\texttt{BitVec}\ 1\)
measure, Lean proves that any recovery packet on this duplicated interface
already forces \(4/5 \le q\), and therefore rules out \(q = 3/4\).  This shows
the new theorem is not merely another injective-family cardinality wrapper.
The new existential endpoint says more: the duplicated surjective family cannot
be rescued by switching to some different \(\texttt{BitVec}\ 1 \to
\texttt{BitVec}\ 1\) extractor after the heavy proper region has already
violated the \(q\)-bound.

The follow-up file
\texttt{ActualSwitchedLocalSharedExactZABSparseThresholdERMRecoverySurjectiveHeavyRegionCharacterization.lean}
packages the exact threshold hidden in that surjective-region obstruction.  On
the surjective branch, Lean now proves that the whole proper-region mass family
collapses to one scalar inequality
\[
  1 - \mu(\texttt{lightestPoint}) \le q.
\]
So even without injective indexing, the obstruction is not really ``search over
all proper regions''; it is already the canonical complement mass of the
lightest visible point.  The regression file
\texttt{ActualSwitchedLocalSharedExactZABSparseThresholdERMRecoverySurjectiveHeavyRegionCharacterizationRegression.lean}
keeps this branch nontrivial by reusing the duplicated-index surjective family.
On the same skew \(\texttt{BitVec}\ 1\) witness, Lean proves the concrete
threshold \(4/5 \le q\) through that lightest-point form and again rules out
\(q = 3/4\) on a surjective family whose predictor map is not injective.

The new synthesis upgrade in \texttt{CruxSynthesis.lean} packages those two
surjective actual-local recovery burdens at the route level rather than only as
file-local facts.  First, below the half-width ceiling \(2r+2k+1<n\), Lean now
records that no extractor at all can support the manuscript-facing recovery
packet on any surjective actual-local endpoint on \(\texttt{BitVec}\ n\).
Second, below the intrinsic lightest-point threshold
\(q < 1-\mu(\texttt{lightestPoint})\), Lean likewise records that no extractor
can support that packet on any surjective actual-local endpoint over a finite
visible surface.  The concrete sample-level plug-in-majority route and the
bit-coded exact-visible fallback route are now explicit instances of those
general synthesis theorems, rather than separate endpoint-specific wrappers.

That route-level package has now been tightened in four further directions.
First, \texttt{CruxSynthesis.lean} records the sharpened recovery-threshold
visible-budget barriers
\texttt{PNPKpolySubrepairClass.fullRuleActualLocalRecoveryThresholdVisibleBudgetBoundary},
\texttt{PNPKpolySubrepairClass.fullRuleActualLocalNoExtractorRecoveryThresholdVisibleBudgetObstruction},
\texttt{PNPKpolySubrepairClass.boundedSamplePluginMajorityActualLocalRecoveryThresholdVisibleBudgetBoundary},
and
\texttt{PNPKpolySubrepairClass.boundedSamplePluginMajorityActualLocalNoExtractorRecoveryThresholdVisibleBudgetObstruction}.
So once the route claims
\[
  q < 1 - 2^{-(2r+4k+\mathrm{slack}+1)},
\]
the old one-step escape hatch \(n=2r+2k+1\) disappears at the synthesis level
as well: Lean upgrades it to \(n \le 2r+2k+\mathrm{slack}\) and records the
matching no-extractor obstructions.

Second, the ledger now records the proper-region mass route lock itself:
\texttt{PNPKpolySubrepairClass.surjectiveActualLocalRecoveryProperRegionMassBoundary}
and
\texttt{PNPKpolySubrepairClass.surjectiveActualLocalNoExtractorProperRegionMassBoundary}.
So on every surjective actual-local endpoint, every proper visible region
already has total mass at most \(q\), and any heavier proper region removes the
extractor choice entirely.

Third, the injective-family heavy-region collapse is now promoted too:
\texttt{PNPKpolySubrepairClass.actualLocalRecoveryHeavyRegionTraceCardBoundary}
and
\texttt{PNPKpolySubrepairClass.actualLocalNoExtractorRecoveryHeavyRegionTraceCardObstruction}.
Thus whenever one finite visible region \(R\) already has \(\mu(R)>q\), the
whole injectively indexed realized family must fit into the Boolean trace space
on \(R\),
\[
  |\mathrm{Index}| \le 2^{|R|},
\]
and Lean records the corresponding no-extractor contradiction when that trace
space is too small.  This is strictly stronger than the earlier heavy-point
collapse when no single atom exceeds \(q\) but a small region still does.

Fourth, the non-surjective injective-extractor threshold packet is now
registered directly in the ledger:
\texttt{PNPKpolySubrepairClass.actualLocalInjectiveRecoveryUniformCardinalityThresholdBoundary},
\texttt{PNPKpolySubrepairClass.actualLocalInjectiveRecoveryBitVecCardinalityThresholdBoundary},
and
\texttt{PNPKpolySubrepairClass.actualLocalNoInjectiveExtractorRecoveryBitVecCardinalityThresholdObstruction}.
So even without surjectivity, any injective \(zfeat\) with
\(0<r+(k+k)\) already forces
\[
  q \ge 1-|U|^{-1},
\]
and on \(\texttt{BitVec}\ n\) forces
\[
  q \ge 1 - 2^{-(n+2k)}.
\]
The concrete non-surjective canary is now formal at the synthesis level too:
on the one-index constant-false
\(\texttt{BitVec}\ 1,\ k=1\) endpoint, Lean still derives
\(1-2^{-3}=7/8 \le q\), so \(q=3/4\) already rules out every injective
\(\texttt{BitVec}\ 1 \to \texttt{BitVec}\ 1\) extractor.

If such a theorem cannot be proved, then the proof attempt should be regarded as stalled at a real
mathematical gap, not a missing editorial detail.

\section{Lean Excerpts}

\begin{lstlisting}[caption={Core asymptotic obstruction}]
theorem polylog_isLittleO_logRadiusNeighborhood (k : Nat) {d c : Real}
    (hd : 1 < d) (hc : 0 < c) :
    polylog k =o[atTop] logRadiusNeighborhood d c

theorem logRadiusNeighborhood_not_isBigO_polylog (k : Nat) {d c : Real}
    (hd : 1 < d) (hc : 0 < c) :
    ¬ logRadiusNeighborhood d c =O[atTop] polylog k
\end{lstlisting}

\begin{lstlisting}[caption={Core counting obstruction}]
theorem no_injective_bitCode_of_lt {n s : Nat} (hs : s < 2 ^ n) :
    ¬ ∃ f : LocalRule n → BitCode s, Function.Injective f

theorem no_uniform_code_below_expInput_for_binaryLogNeighborhood
    {d m s : Nat} (hd : 2 ≤ d) (hs : s < 2 ^ m) :
    ¬ ∃ f : LocalRule (d ^ (Nat.log 2 m + 1)) → BitCode s,
      Function.Injective f
\end{lstlisting}

\section{Conclusion}

The present crux-first formal work yields a clean conditional diagnosis.

\begin{enumerate}[leftmargin=1.5em]
  \item Logarithmic-radius locality is \emph{not} the same thing as polylogarithmic complexity.
  \item The full class of local Boolean rules on that visible data is far too large to fit into a
  polylogarithmic coding budget.
  \item Therefore the program needs a separate theorem proving exceptional compressibility
  of the particular switched predictors it generates.
  \item The v13 \(P=NP\) upper-bound decoder also needs a single-message theorem for arbitrary
  satisfying SAT-search outputs.  In Lean this is not just a slogan: correctness for all valid
  SAT-search outputs is equivalent to the fixed-message readout obligation, and with one satisfying
  witness it is equivalent to single-message constancy.  At the manuscript CNF layer, completeness
  of the extraction makes supported arbitrary-output SAT-search correctness equivalent to the
  semantic single-message promise itself; a separate unrepresented-witness countermodel now shows
  this completeness premise cannot be dropped.  In the generalized obstruction, two satisfying witnesses
  with unequal readouts already rule out the existence of any message correct for all valid
  SAT-search outputs.  The weaker selected-output contract is now explicitly separated: ambiguous
  satisfying witnesses can make selected-output correctness true for two different messages while
  arbitrary-output correctness remains impossible.
  \item The v13 product-success step also needs a sequential conditional admissibility theorem;
  unconditional half-success marginals are insufficient, while the checked sequential conditional
  interface is strong enough for the two-step finite product bound.
  \item A generic finite switching theorem now proves the full tower-product bound from sequential
  half-admissibility.  The open crux is whether the manuscript's switched histories satisfy that
  admissibility premise on positive-mass prefixes.  Empty prefixes are automatically harmless by
  finite-history monotonicity, while one failed nonempty prefix half-step is enough to block the
  tower argument.  Independently, a failed final tower-product count blocks sequential
  admissibility and now localizes to an exact failed conditional prefix cut.
  \item The v13 instantiation obligation is now fixed-field-level: both the abstract
  atom-ledger interface and an existential concrete-cell interface are equivalent to the generic
  prefix half-step.  The manuscript must therefore specify the actual finite history field and
  prove that its fixed fibers decompose the relevant histories and satisfy the cellwise
  half-success bounds.  The decomposition part is automatic for any finite cell map; the
  substantive obligation is the half-success bound inside each supplied cell, equivalently a
  proof that every supplied cell has at least as many next-failure prefix points as next-success
  prefix points.  Equivalently, after the field is fixed, each successful prefix point must be
  injectively matched to an unsuccessful prefix point in the same field cell.  This fixed-field
  premise is now also diagnostic in the negative direction: if the matching or prefix certificate
  fails, Lean extracts a supplied cell whose next-failure count is strictly smaller than its
  next-success count.  Product-bound violations therefore localize not only to a fielded position
  but to a concrete bad cell in that position.  This fixed-field premise is strictly stronger than
  global half-success:
  a field that reveals the successful Boolean coordinate already fails cellwise at the first step.
  The ordered fielded-switching interface now gives the positive bridge as well: if every supplied
  field certificate is proved, the tower product bound follows; if one fixed-field prefix
  certificate fails, the remaining fielded suffix is blocked.
  More generally, a positive-mass fixed-field cell that is entirely successful already blocks the
  certificate, so a field that exposes the next success event is admissible only in the vacuous
  zero-success case.  Equivalently, every successful atom in a valid fixed-field certificate must
  be mixed with at least one unsuccessful prefix point.  The same conclusion now holds for any
  supplied field that determines the next success bit cellwise on the current history support: on a
  positive-mass success prefix, such a field blocks the whole fielded suffix, and the theorem
  applies at any later fielded position after the preceding success-event history has been
  accumulated.  This is strictly stronger than the global version: the Boolean-square diagonal
  model has a second-coordinate field that is not globally success-determining but becomes
  success-determining after conditioning.  In the operational matching formulation, any recursive
  matching proof on such a positive-mass prefix directly proves that the supplied field is not
  success-determining on the accumulated history support.  Even more concretely, every successful
  prefix point must have a same-cell failed prefix point; a single successful point whose supplied
  field cell contains only successes is already enough to rule out the matching and certificate.
  \item A repeated exact-success event gives a checked positive-mass obstruction to that atom-ledger
  obligation: after any positive-mass first occurrence, the second repeated step has conditional
  success \(1\), so no abstract, concrete-cell, fixed-field, or recursive same-cell matching v13
  instantiation exists for the repeated two-step claim.  The Boolean-square instance also violates
  the global tower-product inequality, blocking both abstract and concrete v13 certificates by the
  product-bound contrapositive.  The fixed-field version is blocked at the same level: when
  both repeated cuts use the one-cell field, the fielded tower-product bound fails and therefore no
  recursive same-cell matching proof can exist.  This failure is no longer only global: Lean
  localizes any such failure to a concrete fielded position and then to a bad supplied cell whose
  failures are fewer than successes.  In the one-cell repeated model the second conditioned cut is
  explicitly the failed same-cell matching/certificate obligation, with the unique bad cell
  carrying failure count \(0\) and success count \(2\).
  The forced-step generalization blocks the same repair even when the later event is not a literal
  repeat: on a positive accumulated history, any next success event already implied by that history
  leaves the count unchanged and cannot be a half-step.  In the Boolean-square anchor, two valid
  coordinate half-cuts force a third globally half-mass diagonal success event, so the third
  fixed-field switching and same-cell matching obligations fail.
\end{enumerate}

These are currently the sharpest formally checked cruxes we have in Lean.

\begin{thebibliography}{9}
\bibitem{manuscript2025}
Anonymous manuscript under study.
\newblock \emph{P != NP: A Non-Relativizing Proof via Quantale Weakness and Geometric Complexity}.
\newblock October 2025.

\bibitem{mettapedia-pnp}
Mettapedia formal notes.
\newblock \texttt{Mettapedia/Computability/PNP/LocalityObstruction.lean} and
\texttt{Mettapedia/Computability/PNP/CompressionObstruction.lean}.
\end{thebibliography}

\end{document}
